% =====================================================================
%  «Аштавакра-ґіта» — ФОРМАТ-02, робоче видання всієї книги
%
%  Шість шарів на кожен із 298 віршів: деванагарі, IAST, українська
%  транслітерація, переклад, науковий коментар з українськими паралелями,
%  власна схема. Плюс окремо — наш поетичний переклад усієї книги.
%
%  Джерело тексту — власні редакторські картки editorial/cards/*.json
%  і файли books/ashtavakra-gita/format-v2/ та poetychnyi/. Нічого з цього
%  не написано руками у .tex: усе генерує scripts/build_format_v2_chapter.py.
%
%  Складання:
%    python3 scripts/build_verse_schemas.py --pdf
%    python3 scripts/build_format_v2_chapter.py --all
%    cd books/ashtavakra-gita && latexmk -xelatex main.tex
%  або без latexmk:
%    xelatex main.tex && xelatex main.tex
%
%  Потрібні пакети: fontspec, polyglossia, geometry, graphicx; шрифти
%  DejaVu Serif (fonts-dejavu-core) і шрифт із деванагарі
%  (fonts-noto-core або fonts-lohit-deva).
% =====================================================================
\documentclass[10pt,a5paper,openany]{book}

\usepackage{fontspec}
\setmainfont{DejaVu Serif}[Ligatures=TeX]
\setsansfont{DejaVu Sans}[Ligatures=TeX,Scale=MatchLowercase]
\setmonofont{DejaVu Sans Mono}[Scale=MatchLowercase]

% Деванагарі потрібен окремий шрифт: DejaVu його не має, і без цього рядка
% XeLaTeX мовчки лишив би порожні квадратики замість тексту.
\IfFontExistsTF{Noto Serif Devanagari}
  {\newfontfamily\devanagarifont{Noto Serif Devanagari}[Script=Devanagari]}
  {\IfFontExistsTF{Lohit Devanagari}
    {\newfontfamily\devanagarifont{Lohit Devanagari}[Script=Devanagari]}
    {\IfFontExistsTF{Noto Sans Devanagari}
      {\newfontfamily\devanagarifont{Noto Sans Devanagari}[Script=Devanagari]}
      {\newcommand{\devanagarifont}{%
         \errmessage{Немає шрифту з деванагарі: встановіть fonts-noto-core
         або fonts-lohit-deva}}}}}

\usepackage{polyglossia}
\setdefaultlanguage{ukrainian}
\setotherlanguage{english}

\usepackage[a5paper,top=17mm,bottom=19mm,inner=17mm,outer=15mm,headsep=5mm]{geometry}
\usepackage{graphicx}
\usepackage{xcolor}
\IfFileExists{pzdr.tfm}{\usepackage[hidelinks]{hyperref}}{}

\setlength{\parindent}{1.1em}
\setlength{\parskip}{0pt}
\renewcommand{\baselinestretch}{1.04}
\clubpenalty=9000
\widowpenalty=9000
\brokenpenalty=3000
\tolerance=1500
\emergencystretch=3em
\hbadness=10000
\vbadness=10000
\let\cleardoublepage\clearpage

% --- макроси ФОРМАТ-02 ------------------------------------------------
\newcommand{\naukchapter}[3]{%
  \chapter*{#3}%
  \addcontentsline{toc}{chapter}{#3}%
  \markboth{#3}{#3}%
  \noindent{\sffamily\footnotesize Шість шарів на вірш: #2 віршів.\par}%
  \addvspace{0.6em}}

\long\def\naukintro#1{\begin{quote}\small\raggedright #1\end{quote}\addvspace{0.6em}}

\newcommand{\vsec}[1]{%
  \par\addvspace{1.6em}%
  {\sffamily\bfseries\large Вірш~#1\par}%
  \nopagebreak\addvspace{0.3em}}

% Шар 1: деванагарі.
\long\def\vdeva#1{%
  \noindent{\devanagarifont\large #1\par}%
  \nopagebreak\addvspace{0.35em}}

% Шар 2: IAST — латинська транслітерація з діакритикою (ṛ ṣ ṭ ḥ ṃ ñ ṅ ś ā ī ū).
% Окремого шрифту не потребує: рендерить основний DejaVu Serif, у якого є
% всі знаки Latin Extended Additional і в прямому, і в курсивному накресленні
% (перевірено по cmap шрифту). Якщо основного шрифту колись замінять на
% такий, що цих знаків не має, XeLaTeX напише «Missing character» у лог —
% саме це ловить перевірка збірки, тому окремої гілки \IfFontExistsTF тут не треба.
\long\def\viast#1{%
  \noindent{\sffamily\footnotesize\textbf{IAST.}}~%
  {\footnotesize #1\par}%
  \nopagebreak\addvspace{0.35em}}

% Шар 3: українська транслітерація.
\long\def\vtranslit#1{%
  \noindent{\sffamily\footnotesize\textbf{Як звучить.}}~%
  {\footnotesize\itshape #1\par}%
  \nopagebreak\addvspace{0.35em}}

% Шар 4: переклад.
\long\def\vperek#1{\noindent\textbf{Переклад.}~#1\par\addvspace{0.4em}}

% Шар 5: науковий коментар.
\long\def\vcom#1{\noindent\textbf{Коментар.}~#1\par\addvspace{0.4em}}

% Шар 5б: українська паралель.
\long\def\vpar#1#2#3{%
  \addvspace{0.2em}\noindent
  {\sffamily\footnotesize\textbf{#1.}}~{\footnotesize\textit{#2}~#3\par}%
  \addvspace{0.3em}}

\long\def\vquot#1#2#3#4{%
  \begin{quote}\footnotesize\raggedright
  «#1»\par\addvspace{0.15em}%
  — #2. #3. \texttt{\footnotesize #4}\end{quote}}

% Шар 6: ілюстрація.
\long\def\vfig#1#2#3{%
  \addvspace{0.4em}%
  \begin{center}%
    \includegraphics[width=0.92\textwidth]{#1}%
    \par\addvspace{0.2em}%
    {\footnotesize\raggedright #2\par}%
    {\footnotesize\itshape\textcolor{gray}{#3}\par}%
  \end{center}%
  \addvspace{0.2em}}

% Поетичний переклад.
\newcommand{\poemverse}[2]{%
  \par\addvspace{0.9em}%
  {\sffamily\footnotesize #1\par}\nopagebreak
  \begin{quote}\normalsize\raggedright #2\end{quote}}

\begin{document}

% =====================================================================
%  ТИТУЛ
% =====================================================================
\begin{titlepage}
\centering
\vspace*{1.6cm}

{\sffamily\Large АШТАВАКРА-ҐІТА\par}
\vspace{0.5em}
{\sffamily\small Aṣṭāvakragītā\par}

\vspace{1.6cm}
{\sffamily\bfseries\large РОБОЧЕ ПІЛОТНЕ ВИДАННЯ\par}
\vspace{0.5em}
{\sffamily ФОРМАТ-02: шість шарів на вірш\par}
\vspace{0.3em}
{\sffamily Глави 1–20 — 298 віршів\par}

\vspace{1.4cm}
\begin{minipage}{0.86\textwidth}
\small\raggedright
\textbf{Це не завершена книга і не академічне видання.}

\vspace{0.5em}
Це \emph{робоче пілотне видання нового формату}: воно показує, що машина
складає всі шість шарів для кожного вірша книги і що з наявного матеріалу
вже можна скласти PDF.

\vspace{0.5em}
\textbf{Поетичний переклад у цій книзі — наш власний і не канонічний.}
Він не замінює науковий переклад і не претендує на те, щоб його читали
замість санскритського тексту.

\vspace{0.5em}
Усі редакторські картки мають статус \textbf{«чернетка»}. Жодна не позначена
як рецензована, жодна не має підтвердженого аудиту прав. Санскритологічної
вичитки не було.

\vspace{0.5em}
\textbf{Деванагарі тут — механічна транслітерація} поля \texttt{iast} з картки,
а не окреме читання за рукописом. Якщо деванагарі розійдеться з карткою,
правильною є картка.

\vspace{0.5em}
Не цитувати як переклад «Аштавакра-ґіти».
\end{minipage}

\vfill
{\sffamily\footnotesize Електронне робоче видання. Не для продажу.\par}
\end{titlepage}

% =====================================================================
%  ВИХІДНІ ДАНІ
% =====================================================================
\thispagestyle{empty}
\vspace*{\fill}
\begingroup\small
\noindent\textbf{\sffamily Вихідні дані}

\vspace{0.8em}
\noindent
\begin{tabular}{@{}p{0.34\textwidth}p{0.58\textwidth}@{}}
\textsf{Назва} & Аштавакра-ґіта (Aṣṭāvakragītā). Робоче пілотне видання ФОРМАТ-02 \\[0.3em]
\textsf{Зміст} & Глави 1–20 (298 віршів): деванагарі, IAST, українська транслітерація, переклад,
науковий коментар з українськими паралелями, власна схема; окремо — наш поетичний
переклад усієї книги \\[0.3em]
\textsf{Первісний твір} & санскритська поема; за наявними даними — суспільне надбання через вік \\[0.3em]
\textsf{Джерело нашого тексту} & власні картки \texttt{editorial/cards/*.json} \\[0.3em]
\textsf{Статус карток} & «чернетка» — усі, без винятків \\[0.3em]
\textsf{Поетичний переклад} & наш власний, не канонічний і не науковий \\[0.3em]
\textsf{Укладач} & редакція проєкту \texttt{juv4uk/ukrainian-books} \\[0.3em]
\textsf{Ліцензія власного внеску} & «ВОЛЬНІСТЬ» \\[0.3em]
\textsf{Складання} & XeLaTeX, шрифт DejaVu Serif (ним же IAST із діакритикою); деванагарі — Noto Serif Devanagari \\[0.3em]
\end{tabular}

\vspace{1.0em}
\noindent\footnotesize
Правовий статус матеріалів — \texttt{COPYRIGHT\_AND\_SOURCES.md}, паспорт
видання — \texttt{PROVENANCE.md}, правила формату — \texttt{docs/FORMAT\_V2.md},
таблиця транслітерації — \texttt{docs/TRANSLITERATSIIA.md}. Цей PDF не є
юридичним висновком.
\endgroup
\vspace*{\fill}

% =====================================================================
%  ПЕРЕДМОВА
% =====================================================================
\frontmatter
\pagestyle{headings}
\tableofcontents
\clearpage
\chapter*{Передмова}
\addcontentsline{toc}{chapter}{Передмова}
\markboth{Передмова}{Передмова}

Це \textbf{робоче пілотне видання}, а не книга. Мета одна: показати, що
машина справді складає український PDF із кирилицею, діакритикою IAST,
деванагарі, власними схемами й окремою частиною поетичного перекладу — і
показати, скільки власного матеріалу вже придатне до верстки.

\section*{Що тут зверстано}

Усі двадцять глав, 298 віршів, кожен шістьма шарами. Кожен вірш має власну
схему; кожен має українську паралель або прямо позначене її місце. Окремо
йде наш поетичний переклад усієї книги.

\section*{Чого тут немає}

Санскритського читання цифрових свідків (GRETIL, SARIT) як окремого шару.
Деванагарі й транслітерація — наше механічне перетворення того читання, яке
вже звірене і збережене в картці; редакторський апарат транскрипцій не
відтворюється. Це не технічне обмеження, а свідоме редакційне рішення,
описане в \texttt{docs/VYDAVNYCHYI\_VYSNOVOK.md}.

Так само немає рецензії: усі картки — «чернетка», незалежного рецензента не
залучено. Переклади й коментарі робочі, не вивірені.

\section*{Про українські паралелі}

Кожне місце, де згадано українського поета, філософа чи діяча, позначене як
\emph{наша паралель} — тобто наша асоціація, а не твердження про вплив.
Ми не твердимо, що Сковорода читав «Аштавакра-ґіту». Там, де паралель
непевна, у тексті стоїть позначка \textit{потребує перевірки}, а не впевнене
імʼя. Жодної цитати з українських авторів у цьому виданні немає: усі
паралелі — вказівки на твір без цитати. Правила — у
\texttt{books/ashtavakra-gita/paraleli.json} і \texttt{scripts/check\_ukrainian\_parallels.py}.

\section*{Про поетичний переклад}

Друга частина книги — \textbf{наш поетичний переклад}, а не канонічний і не
науковий. Він не замінює переклад у науковій частині: обидва йдуть у книзі,
і читач бачить обидва. Поетична свобода тут — у ритмі й порядку слів, не у
фактах; там, де рима псувала б думку, вжито верлібр.

\mainmatter
% ЗГЕНЕРОВАНО scripts/build_format_v2_chapter.py — не редагувати вручну.
% Джерело: editorial/cards/1.*.json, format-v2/rozdil-1.json, paraleli.json.

\naukchapter{1}{20}{Глава 1. Питання й перші відповіді}
\naukintro{Перша глава не викладає вчення, а робить те, чого не роблять пізніші глави: вона ставить питання і відповідає на нього не доказами, а низкою заперечень. Аштавакра не доводить, що учень вільний; він віднімає від учня все, чим той себе вважав, — стихії, варну, ашраму, діяння, навіть саму практику самадгі (1.15). Те, що лишається після цього віднімання, названо двома словами: sākṣin («свідок») і cit («усвідомлення»).}
\vsec{1.1}
\vdeva{कथं ज्ञानम् अवाप्नोति कथं मुक्तिर् भविष्यति । वैराग्यं च कथं प्राप्तम् एतद् ब्रूहि मम प्रभो ॥}
\vtranslit{катхам джнянам авапноті катхам муктір бгавішяті | вайраґям ча катхам праптам етад брухі мама прабго ||}
\vperek{Як здобувається знання? Як настає визволення? І як досягається безсторонність? Скажи мені це, пане.}
\vcom{Перший вірш книги — не твердження, а питання: три речі, які шукають — jñāna, mukti, vairāgya. Варто тримати в голові, що далі текст заперечить кожну з них саме як щось набуте (пор. 1.15, 2.18). prabho — звертання до Аштавакри; «пане» умовний, але не випадковий відповідник: санскритське звертання підкреслює нерівність того, хто питає, і того, хто відповідає.}
\vpar{Українська паралель (наша інтерпретація)}{Григорій Сковорода, «Алфавіт, чи Буквар миру» (діалог).}{Наша паралель: Сковорода теж будує філософію як низку питань, а не як перелік відповідей, і починає з питання про себе. Це наша інтерпретація, а не твердження про вплив: доказів знайомства Сковороди з «Аштавакра-ґітою» тут не наведено.}
\vfig{ilustratsii/1-1-motion.pdf}{Три питання сходяться в одну відповідь: далі глава заперечить кожне з трьох як щось набуте}{Схема наша, не історична ілюстрація.}

\vsec{1.2}
\vdeva{मुक्तिम् इच्छसि चेत् तात विषयान् विषवत् त्यज । क्षमार्जवदयातोषसत्यं पीयूषवद् भज ॥}
\vtranslit{муктім іччасі чет тата вішаян вішават тяджа | кшамарджавадаятошасатям піюшавад бгаджа ||}
\vperek{Якщо ти прагнеш звільнення, сину, відкинь чуттєві предмети, як отруту; шануй терпимість, щирість, співчуття, вдоволеність і правдивість, як нектар.}
\vcom{Вірш задає умову входу в діалог: звільнення тут не досягають додаванням, а починають із відкидання. Пара «отрута / нектар» протиставляє не приємне й неприємне, а те, що зв’язує, і те, що живить. Другий піввірш — перелік п’яти чеснот у складеному іменнику (dvandva), тобто одна нерозривна сукупність, а не меню. Для українського перекладу принципово, що toṣa — це вдоволеність тим, що є, а не «простота»: «простота» відповідала б іншому слову (sāradya, saralatā) і в цьому переліку відсутня. Прямота ārjava — радше незігнутість, ніж «щирість» у побутовому сенсі; обраний відповідник умовний. Межі термінів див. у примітці до 1.2 в research/ashtavakra/. Твердження про значення toṣa — робоча лексикографічна оцінка, а не рецензований висновок.}
\vpar{Українська паралель (ПОТРЕБУЄ ПЕРЕВІРКИ)}{українська барокова проповідь, Києво-Могилянська гомилетика XVII ст..}{ПОТРЕБУЄ ПЕРЕВІРКИ: найближча українська аналогія до пари отрута / нектар — барокова емблематика, де та сама річ буває ліками або отрутою. Ми не наводимо конкретного твору, бо не перевірили жодного тексту за джерелом, і вигадувати посилання не будемо.}
\vfig{ilustratsii/1-2-balance.pdf}{Отрута й нектар на одній шальці: вага не в речах, а в тому, чим вони стають}{Схема наша, не історична ілюстрація.}

\vsec{1.3}
\vdeva{न पृथ्वी न जलं नाग्निर् न वायुर् द्यौर् न वा भवान् । एषां साक्षिणम् आत्मानं चिद्रूपं विद्धि मुक्तये ॥}
\vtranslit{на прітхві на джалам наґнір на ваюр дяур на ва бгаван | ешам сакшінам атманам чідрупам віддгі муктає ||}
\vperek{Ти не земля, не вода, не вогонь, не вітер, не небо — і не є ними. Заради звільнення пізнай Атман, свідка цих [стихій], чия природа — cit.}
\vcom{Вірш будує заперечення не як список того, чим ти не є, а як зняття ототожнення: п'ять стихій — це складники тіла й світу, і шосте заперечення (na bhavān) закриває вихід через «але ж я — це щось інше». Те, що лишається, названо двома словами: sākṣin («свідок») і cidrūpa («той, чия форма — cit»). Обидва — позитивні визначення після низки заперечень, і саме вони роблять вірш не просто негативним. sākṣin тут не «спостерігач» у побутовому сенсі, а те, що не входить до складу спостережуваного; пор. вживання в 1.7 і далі. cidrūpa — bahuvrīhi, тому «чия природа — cit», а не «свідомість-форма». muktaye (dativus finalis) показує, що пізнання тут не самоціль, а умова звільнення. Переклад «знай» для viddhi — наказовий спосіб, не побажання.}
\vpar{Українська паралель (наша інтерпретація)}{Григорій Сковорода, «Кільце. Дружня розмова про душевний мир».}{Наша паралель: у «Кільце» справжнє «я» теж не збігається з тілом і з речами, які людина собі приписує. Це наша асоціація за будовою думки, а не науковий висновок про запозичення.}
\vfig{ilustratsii/1-3-neti.pdf}{Пʼять стихій закреслено — лишається свідок, чия природа cit}{Схема наша, не історична ілюстрація.}

\vsec{1.4}
\vdeva{यदि देहं पृथक्कृत्य चिति विश्राम्य तिष्ठसि । अधुनैव सुखी शान्तो बन्धमुक्तो भविष्यसि ॥}
\vtranslit{яді дехам прітхаккрітя чіті вішрам'я тіштхасі | адгунайва сукхі шанто бандгамукто бгавішясі ||}
\vperek{Якщо, відокремивши тіло, ти стоїш, спочивши в cit, — то вже зараз ти будеш щасливий, умиротворений, вільний від уз.}
\vcom{Вірш дає не обіцянку на майбутнє, а умову з негайним наслідком: adhunaiva — «вже зараз», і саме це слово є смисловим центром. Майбутній час (bhaviṣyasi) тут не відсуває результат, а означає, що стан настане як прямий наслідок умови, а не через поступовий шлях. Умова потрійна: pṛthakkṛtya (відокремити тіло), viśrāmya citi (спочити в cit), tiṣṭhasi (перебувати). Три дієслова описують не дію, а спосіб стояння. citi — локатив від cit, тобто «у cit», а не «на cit»: це місце, у якому спочивають, а не об'єкт, який досягають. bandhamuktaḥ — bahuvrīhi, «той, у кого узи відпали», а не «той, кого звільнили»: суб'єкт не зазнає дії ззовні.}
\vpar{Українська паралель (наша інтерпретація)}{Леся Українка, «Contra spem spero!».}{Наша паралель: у Лесі Українки дух стоїть, коли тіло вже не тримає, — і саме ця незалежність стану від тіла близька до того, що вірш називає спочити в cit. Ми не цитуємо й не твердимо впливу; це типологічна близькість, яку читач може перевірити сам.}
\vfig{ilustratsii/1-4-motion.pdf}{Умова з негайним наслідком: «вже зараз» — смисловий центр вірша}{Схема наша, не історична ілюстрація.}

\vsec{1.5}
\vdeva{न त्वं विप्रादिको वर्णो नाश्रमी नाक्षगोचरः । असङ्गो ऽसि निराकारो विश्वसाक्षी सुखी भव ॥}
\vtranslit{на твам віпрадіко варно нашрамі накшаґочарах | асанґо ’сі ніракаро вішвасакші сукхі бгава ||}
\vperek{Ти не брахман і не інший стан, не учень і не відречений, і око тебе не бачить. Ти без прив'язаності, без форми, свідок усього — будь щасливий.}
\vcom{Вірш знімає дві ознаки, за якими людину зазвичай визначають, і робить це тими самими словами, якими їх називає традиція. `varṇa` тут — суспільний стан, а не «барва» і не «звук»: `viprādika` означає «той, що починається з брахмана», тобто весь ряд чотирьох станів. `āśramin` — не «ашрама» як місце, а людина в одному з чотирьох життєвих станів. Обидва слова стоять в одному рядку з `na`, і це не пом'якшення: заперечення стосується самої тотожності, а не її гідності. Третє заперечення іншого роду: `akṣagocara` — «не в полі зору ока», тобто не об'єкт чуттєвого сприймання. Перші два знімають суспільні ознаки, третє — можливість бути побаченим узагалі. Після цього лишається тільки ствердне: `asaṅga` (без прив'язаності), `nirākāra` (без форми), `viśvasākṣin` (свідок усього) — і наказ `sukhī bhava`, «будь щасливий», єдиний імператив у вірші. Це відповідь на питання, поставлене в Issues \#3: `varṇa` і `āśrama` в 1.5 не описують мудрого й не є шляхом до нього — вони названі, щоб бути знятими. Порівняй 1.6, де тим самим способом знімаються дгарма й адгарма.}
\vpar{Українська паралель (наша інтерпретація)}{Митрополит Іларіон Київський, «Слово про закон і благодать» (XI ст.).}{Наша паралель: київське «Слово» розрізняє закон, що приписує стан, і благодать, що стану не потребує. Це сучасне типологічне зіставлення, а не доказ впливу; датування «Аштавакра-ґіти» науково не встановлене.}
\vfig{ilustratsii/1-5-neti.pdf}{Варна й ашрама зняті тими самими словами, якими їх називає традиція}{Схема наша, не історична ілюстрація.}

\vsec{1.6}
\vdeva{धर्माधर्मौ सुखं दुःखं मानसानि न ते बिभो । न कर्ताऽसि न भोक्ताऽसि मुक्त एवासि सर्वदा ॥}
\vtranslit{дгармадгармау сукхам духкхам манасані на те бібго | на карта’сі на бгокта’сі мукта евасі сарвада ||}
\vperek{Дхарма й адхарма, щастя й страждання — усе це розумове, не твоє, о всепроникний. Ти не дієць і не споживач — ти вільний завжди.}
\vcom{mānasāni — «розумові»: пари не заперечуються як вигадка, а виводяться з-під авторства того, кому їх приписують. Друга половина містить рідкісну для цього тексту остаточність: mukta eva asi sarvadā — «ти вільний саме завжди». Це варто порівняти з 1.7, де стоїть muktaprāya, «майже вільний». Різниця між сусідніми віршами не стилістична, і переклад не повинен її згладжувати.}
\vpar{Українська паралель (наша інтерпретація)}{Памфіл Юркевич, «Серце і його значення в духовному житті людини, за вченням слова Божого» (1860).}{Наша паралель: коли вірш каже, що дхарма й адхарма розумові, не твої, він відводить моральному судженню місце в думці, а не в речі. Юркевич розводить розум і серце інакше, ніж це робить санскритський текст, але думка про те, що мірило не тотожне людині, у нього є. Ми свідомо не цитуємо Юркевича: доступний цифровий уривок не є бібліотекою першоджерел, а український переклад твору має власні права. Тому тут названо твір, а не рядок.}
\vfig{ilustratsii/1-6-balance.pdf}{Дхарма й адхарма, щастя й страждання важать однаково: усе це розумове}{Схема наша, не історична ілюстрація.}

\vsec{1.7}
\vdeva{एको द्रष्टासि सर्वस्य मुक्तप्रायो ऽसि सर्वदा । अयम् एव हि ते बन्धो द्रष्टारं पश्यसीतरम् ॥}
\vtranslit{еко драштасі сарвася муктапрайо ’сі сарвада | аям ева хі те бандго драштарам пашясітарам ||}
\vperek{Ти єдиний бачитель усього — і завжди майже вільний. Уза твоя ось у чому: ти бачиш іншого бачителя, крім себе.}
\vcom{Вірш перевертає звичне уявлення про неволю. Неволя тут не в тому, що ти робиш погане, а в тому, що ти припускаєш **другого бачителя**. Слово muktaprāyaḥ («майже вільний») ужите навмисно: воно не обіцяє повноти, а показує, наскільки мало лишається. itaram («іншого») — не «щось інше», а саме інший суб'єкт спостереження: щойно з'являється другий, з'являється й уза. Це прямий зв'язок із 1.3, де лишався sākṣin («свідок») після зняття стихій: там свідок один, тут — перевірка, чи він справді один.}
\vpar{Українська паралель (наша інтерпретація)}{Григорій Сковорода, «Ікона Алківіадська» (трактат).}{Наша паралель: Сковорода розрізняє тінь і річ так само, як цей вірш розрізняє побачене й того, хто бачить. Уза вірша — не в тому, що є другий бачитель, а в тому, що ми його шукаємо; Сковорода шукає так само. Це наша асоціація, не твердження про джерело.}
\vfig{ilustratsii/1-7-witness.pdf}{Уза не в тому, що ти бачиш, а в тому, що припускаєш другого бачителя}{Схема наша, не історична ілюстрація.}

\vsec{1.8}
\vdeva{अहं कर्तेत्य् अहंमानम् अहाकृष्णाहि दंशितः । नाहं कर्तेति विश्वासामृतं पीत्वा सुखी भव ॥}
\vtranslit{ахам картетй ахамманам ахакрішнахі дамшітах | нахам картеті вішвасамрітам пітва сукхі бгава ||}
\vperek{«Я діяч» — це самомніння, і воно ж укус великого чорного змія. Випивши нектар упевненості «я не діяч», будь щасливий.}
\vcom{Вірш протиставляє не дію й бездіяльність, а **два самоописи**. Змій — не докір, а образ отрути, що входить разом із думкою «я діяч». Друга пада дає не аргумент, а напій: viśvāsāmṛtam — «нектар упевненості», тобто стан, а не висновок. Звертання до 1.10 важливе: там радджа-сарпа («змія в мотузці») описує помилкове сприйняття світу, а тут змій — сама думка про авторство дії. Це той самий образ у двох ролях.}
\vpar{Українська паралель (ПОТРЕБУЄ ПЕРЕВІРКИ)}{українська поезія XX ст., не визначено.}{ПОТРЕБУЄ ПЕРЕВІРКИ: образ «я діяч» як укус змія близький до української традиції самоспостереження, але конкретного твору, який можна було б назвати з певністю, ми не перевірили. Краще назвати прогалину, ніж підставити перше-ліпше ім'я.}
\vfig{ilustratsii/1-8-illusion.pdf}{Один і той самий образ у двох ролях: «я діяч» — укус, «я не діяч» — нектар}{Схема наша, не історична ілюстрація.}

\vsec{1.9}
\vdeva{एको विशुद्धबोधो ऽहम् इति निश्चयवह्निना । प्रज्वाल्याज्ञानगहनं वीतशोकः सुखी भव ॥}
\vtranslit{еко вішуддгабодго ’хам іті нішчаявахніна | праджваляджнянаґаханам віташоках сукхі бгава ||}
\vperek{«Я — єдине чисте прозріння»: спаливши вогнем упевненості хащу незнання, будь без жалю й щасливий.}
\vcom{Продовження 1.8: там був «нектар упевненості», тут — «вогонь упевненості». Один і той самий niścaya діє двічі, але по-різному: напій змінює стан, вогонь знищує перешкоду. ajñānagahanaṃ — не «незнання» як абстракція, а «хаща», нетрі, у яких застрягають; саме тому потрібен вогонь, а не пояснення. vītaśokaḥ («без жалю») стоїть перед sukhī bhava, тобто жаль названо окремо від щастя — це не те саме, що його відсутність.}
\vpar{Українська паралель (ПОТРЕБУЄ ПЕРЕВІРКИ)}{українська поезія, не визначено.}{ПОТРЕБУЄ ПЕРЕВІРКИ: мотив випалити сумнів трапляється в українській ліриці часто, і саме тому небезпечний: будь-яке ім'я тут буде натяжкою. Потрібна конкретна розвідка за першоджерелом, а не цитата з пам'яті.}
\vfig{ilustratsii/1-9-flame.pdf}{Вогонь упевненості палить не знання, а хащу незнання}{Схема наша, не історична ілюстрація.}

\vsec{1.10}
\vdeva{यत्र विश्वम् इदं भाति कल्पितं रज्जुसर्पवत् । आनन्दपरमानन्दः स बोधस् त्वं सुखं चर ॥}
\vtranslit{ятра вішвам ідам бгаті калпітам раджджусарпават | анандапараманандах са бодгас твам сукхам чара ||}
\vperek{Там, де все це сяє — вигадане, як змія в мотузці, — ти є те прозріння, що є радістю і найвищою радістю. Ходи щасливо.}
\vcom{Образ мотузки й змії — класичний для адвайти: помилка не в тому, що видно, а в тому, що до видимого додано. kalpitam («вигадане») стоїть одразу після viśvam, тобто вигадане саме «все», а не окремий предмет. Друга пада не заперечує світ, а переносить питання: питання не «чи є світ», а «хто той, у кому він сяє». ānandaparamānandaḥ — подвійна радість, і це не тавтологія: перша належить прозрінню, друга — тому, що за ним. Наказ sukhaṃ cara («ходи щасливо») — дієвий, а не описовий: це спосіб перебування, а не обіцянка стану.}
\vpar{Українська паралель (наша інтерпретація)}{Григорій Сковорода, «Потоп Зміїний» (трактат).}{Наша паралель: у Сковороди змій — не побутова тварина, а символ оманного образу світу. Це робить «Потоп Зміїний» найближчим українським текстом до метафори змії в мотузці. Ми називаємо твір, але не твердимо, що Сковорода знав санскритський образ.}
\vfig{ilustratsii/1-10-illusion.pdf}{Змія в мотузці: помилка не в тому, що видно, а в тому, що до видимого додано}{Схема наша, не історична ілюстрація.}

\vsec{1.11}
\vdeva{मुक्ताभिमानी मुक्तो हि बद्धो बद्धाभिमान्य् अपि । किंवदन्तीह सत्येयं या मतिः सा गतिर् भवेत् ॥}
\vtranslit{муктабгімані мукто хі баддго баддгабгіманй апі | кімвадантіха сатєям я матіх са ґатір бгавет ||}
\vperek{Хто вважає себе вільним — той вільний; хто вважає себе зв'язаним — зв'язаний. Що тут ще сказати? Яке переконання, така й доля.}
\vcom{Формула `yā matiḥ sā gatiḥ` — «яке переконання, така й доля» — найпряміше місце книги про те, що стан визначається не фактом, а **самоозначенням**. `abhimānī` — не «той, хто думає» нейтрально, а той, хто **приписує собі**: той самий корінь, що в `ahaṃkāra` (1.6, 15.13). Друга половина (`kiṃvadantīha`) — риторичне «що тут казати», тобто текст сам визнає, що це звучить надто просто.}
\vpar{Українська паралель (наша інтерпретація)}{Григорій Сковорода, «Розмова п'яти подорожніх про істинне щастя в житті».}{Наша паралель: у Сковороди щастя не приходить ззовні — воно вже є в людині, яка його впізнала. Це близько до «хто вважає себе вільним — той вільний», але не те саме: у вірші йдеться про визволення, а не про щастя. Розбіжність ми називаємо, а не згладжуємо.}
\vfig{ilustratsii/1-11-knot.pdf}{Яке переконання, така й доля: вузол затягує й розвʼязує думка}{Схема наша, не історична ілюстрація.}

\vsec{1.12}
\vdeva{असङ्गो निःस्पृहः शान्तो भ्रमात् संसारवान् इव ॥}
\vtranslit{асанґо ніхспріхах шанто бграмат самсараван іва ||}
\vperek{Без прив'язаності, без жадання, втихлий — через оману він видається тим, хто в сансарі.}
\vcom{Вірш короткий: у SARIT він має лише один рядок (два піди), тоді як сусідній 1.13 — два. Чи це справжня коротка форма, чи втрачена половина, за одним свідком не встановити; тому картка лишається чернеткою, а сумнів зафіксовано в `variants`. `saṃsāravān iva` — «**наче** той, хто в сансарі»: `iva` несе всю думку, бо стан у сансарі тут лише видимість, а не факт. Переклад без «наче» знищив би вірш.}
\vpar{Українська паралель (наша інтерпретація)}{Михайло Коцюбинський, «Intermezzo» (1908).}{Наша паралель: у «Intermezzo» людина, відійшовши від залізних рук міста, стає тихою і саме тому вільною. Це наша асоціація за станом, а не твердження про спільне джерело: Коцюбинський пише про втому, вірш — про оману.}
\vfig{ilustratsii/1-12-illusion.pdf}{Серпанок омани: стан у сансарі тут лише видимість, і це несе одне слово — iva}{Схема наша, не історична ілюстрація.}

\vsec{1.13}
\vdeva{कूटस्थं बोधम् अद्वैतम् आत्मानं परिभावय । आभासो ऽहं भ्रमं मुक्त्वा भावं बाह्यम् अथान्तरम् ॥}
\vtranslit{кутастхам бодгам адвайтам атманам парібгавая | абгасо ’хам бграмам муктва бгавам бахям атхантарам ||}
\vperek{Осягай ātman як незмінне, недвоїсте усвідомлення. Відкинувши оману «я — лише відблиск», відкинь і стан зовнішній та внутрішній.}
\vcom{`kūṭastha` — буквально «той, що стоїть на вершині», усталено «незмінний, як ковадло»; це означення ātman, а не опис зусилля. `ābhāso 'ham` — «я — відблиск»: не «я ілюзія», а саме відблиск чогось іншого. Друга половина відкидає не лише внутрішнє, а й **зовнішнє** (`bāhyam atha antaram`) — і саме цей порядок, зовнішнє перед внутрішнім, варто зберегти. Вірші 1.12 і 1.13 у SARIT лежать в одному `lg`; див. `variants` картки 1.12.}
\vpar{Українська паралель (наша інтерпретація)}{Дмитро Чижевський, «Нариси з історії філософії на Україні» (Прага, 1931).}{Наша паралель: Чижевський описує українську традицію як таку, де незмінним началом людини вважають не думку, а глибшу основу. Це бібліографічно перевірюване дослідження, і саме тому воно тут доречне: воно дає читачеві не цитату, а адресу, за якою той самий хід думки можна прочитати в українському матеріалі.}
\vfig{ilustratsii/1-13-witness.pdf}{Незмінне усвідомлення під шаром приписування: відблиск відкидають, а не шукають}{Схема наша, не історична ілюстрація.}

\vsec{1.14}
\vdeva{देहाभिमानपाशेन चिरं बद्धो ऽसि पुत्रक । बोधो ऽहं ज्ञानखड्गेन तन् निष्कृत्य सुखी भव ॥}
\vtranslit{дехабгіманапашена чірам баддго ’сі путрака | бодго ’хам джнянакхадґена тан нішкрітя сукхі бгава ||}
\vperek{Ти довго зв'язаний мотузкою приписування собі тіла, синку. Скажи: «я — усвідомлення» — і, відсікши те мечем знання, будь щасливий.}
\vcom{Два образи в одному вірші: `pāśa` (мотузка, якою в'яжуть) і `khaḍga` (меч, яким ріжуть). Пута не «розв'язують», а **відсікають** — і це та сама різкість, що в 19.1 з кліщами й стрілою. `bodho 'ham` подано як те, що треба **сказати** собі, а не зрозуміти: переклад мусить зберегти цю мовленнєвість, бо вона й робить вірш дієвим. `putraka` («синку») — звертання наставника до учня.}
\vpar{Українська паралель (наша інтерпретація)}{Іван Франко, поема «Мойсей» (1905).}{Наша паралель: у Франка народ лишається в пустелі не через брак дороги, а через рабство, яке він несе в собі. Це найближча українська річ до думки про мотузку приписування, хоч Франко говорить про народ, а вірш — про одну людину. Порівняння наше, не вплив.}
\vfig{ilustratsii/1-14-knot.pdf}{Пута не розвʼязують, а відсікають: мотузка приписування і меч знання}{Схема наша, не історична ілюстрація.}

\vsec{1.15}
\vdeva{निःसङ्गो निष्क्रियो ऽसि त्वं स्वप्रकाशो निरञ्जनः । अयम् एव हि ते बन्धः समाधिम् अनुतिष्ठसि ॥}
\vtranslit{ніхсанґо нішкрійо ’сі твам свапракашо ніраньджанах | аям ева хі те бандгах самадгім анутіштхасі ||}
\vperek{Ти без прив'язаності, без діяння, самосяйний, незаплямлений. Ось саме це твої пута — те, що ти вправляєшся в самадгі.}
\vcom{Вірш перевертає звичне: пута — не світ і не пристрасті, а **сама практика** самадгі. `anutiṣṭhasi` — «вправляєшся, додержуєш», дієслово старанного виконання; саме воно й назване неволею. Це найраніша заява тієї думки, що завершується в 15.20 («відкинь саму медитацію»). Для перекладу важливо не згладити різкість: текст не радить практикувати краще, він каже, що практика — і є пута.}
\vpar{Українська паралель (наша інтерпретація)}{Василь Симоненко, «Ти знаєш, що ти — людина?» (1962).}{Наша паралель: Симоненко ставить гідність людини в неї саму, а не в те, що їй дають ззовні. Це близько до «самосяйний, незаплямлений», але Симоненко говорить про людську гідність у конкретній історії, а вірш — про природу усвідомлення. Ми не стираємо цієї різниці.}
\vfig{ilustratsii/1-15-flame.pdf}{Самосяйне світло не потребує іншого світла — і саме тому практика тут названа путами}{Схема наша, не історична ілюстрація.}

\vsec{1.16}
\vdeva{त्वया व्याप्तम् इदं विश्वं त्वयि प्रोतं यथार्थतः । शुद्धबुद्धस्वरूपस् त्वं मा गमः क्षुद्रचित्तताम् ॥}
\vtranslit{твая в'яптам ідам вішвам тваї протам ятхартхатах | шуддгабуддгасварупас твам ма ґамах кшудрачіттатам ||}
\vperek{Цей всесвіт просякнутий тобою і справді нанизаний на тебе. Ти маєш природу чистого й просвітленого — не йди до дріб'язковості ума.}
\vcom{Два різні образи в першій половині: `vyāpta` («просякнутий, пронизаний») і `prota` («нанизаний», як на нитку). Перше — про заповнення, друге — про наскрізну нитку, на якій усе тримається. `kṣudracittatā` — «дріб'язковість ума», не «дрібні думки»: ідеться про масштаб, а не про кількість. `mā gamaḥ` — заборона, звернена до того, хто вже такий, а не до того, хто має стати.}
\vpar{Українська паралель (наша інтерпретація)}{Богдан-Ігор Антонич, «Зелена Євангелія» (Львів: Ізмарагд, 1938).}{Наша паралель: у Антонича світ — не декорація, а одне живе тіло, і людина в ньому не гість. Це найближча українська річ до «всесвіт просякнутий тобою». Збірка вийшла посмертно 1938 року; жодного рядка ми не цитуємо.}
\vfig{ilustratsii/1-16-thread.pdf}{Два образи в одній половині: просякнути (vyāpta) і нанизати (prota)}{Схема наша, не історична ілюстрація.}

\vsec{1.17}
\vdeva{निरपेक्षो निर्विकारो निर्भरः शीतलाशयः । अगाधबुद्धिर् अक्षुब्धो भव चिन्मात्रवासनः ॥}
\vtranslit{нірапекшо нірвікаро нірбгарах шіталашаях | аґадгабуддгір акшубдго бгава чінматравасанах ||}
\vperek{Без очікування, без видозмін, без тягаря, з прохолодним нутром; із незміренною думкою, незбурений — будь тим, чия схильність є сама лише свідомість.}
\vcom{Єдиний наказ у вірші — `bhava` («будь»), і він стоїть **усередині** переліку означень, а не на початку чи в кінці. `cinmātra-vāsana` — рідкісне поєднання: `vāsanā` зазвичай означає схильність, яку треба усунути (9.8), а тут вона ж стає єдиною, що лишається. Це не суперечність, а звуження: з усіх схильностей дозволена одна.}
\vpar{Українська паралель (наша інтерпретація)}{Олександр Довженко, «Зачарована Десна» (1956).}{Наша паралель: у Довженка спокій — не порожнеча, а повнота, яку людина впізнає, коли перестає квапитись. Це наша асоціація; Довженко пише прозу про дитинство, а вірш — про стан без очікування.}
\vfig{ilustratsii/1-17-space.pdf}{Небо без хмар: єдиний наказ «будь» стоїть усередині переліку означень}{Схема наша, не історична ілюстрація.}

\vsec{1.18}
\vdeva{साकारम् अनृतं विद्धि निराकारं तु निश्चलम् । एतत् तत्त्वोपदेशेन न पुनर्भवसंभवः ॥}
\vtranslit{сакарам анрітам віддгі ніракарам ту нішчалам | етат таттвопадешена на пунарбгавасамбгавах ||}
\vperek{Знай: те, що має образ, — несправжнє; а те, що без образу — непохитне. Цим настановленням про суть немає народження знову.}
\vcom{Останній вірш глави дає найрізкіше з її тверджень: `sākāra` (усе, що має образ) названо `anṛta` — несправжнім. Це не заперечення світу, а заперечення **образності** як ознаки справжнього. Друга половина обіцяє наслідок не моральний, а структурний: `na punarbhavasaṃbhavaḥ` — «немає можливості знову-народження». `upadeśa` — настановлення вчителя, те саме слово, що в 15.1; переклад мусить бути однаковим.}
\vpar{Українська паралель (наша інтерпретація)}{Григорій Сковорода, «Алфавіт, чи Буквар миру» (діалог).}{Наша паралель: Сковорода послідовно розводить образ і те, що за образом стоїть, і саме на цьому розведенні тримається його символіка. Це наша інтерпретація за будовою думки, а не вказівка на джерело вірша.}
\vfig{ilustratsii/1-18-neti.pdf}{Образ і безобразність зняті: заперечується не світ, а образність як ознака справжнього}{Схема наша, не історична ілюстрація.}

\vsec{1.19}
\vdeva{यथैवादर्शमध्यस्थे रूपे ऽन्तः परितस् तु सः । तथैवास्मिञ् छरीरे ऽन्तः परितः परमेश्वरः ॥}
\vtranslit{ятхайвадаршамадгястхе рупе ’нтах парітас ту сах | татхайвасмінь ччаріре ’нтах парітах парамешварах ||}
\vperek{Як образ у дзеркалі — він усередині й навколо нього; так само і в цьому тілі — всередині й навколо верховний володар.}
\vcom{Пара `antaḥ … paritaḥ` («усередині … навколо») стоїть в обох половинах, і саме вона несе думку: не «всередині тіла» як у посудині, а одночасно всередині й поза. Дзеркало — єдиний образ цього вірша; важливо, що образ у дзеркалі не є дзеркалом і не віддільний від нього. `parameśvara` — єдине вживання цього слова в книзі, і воно єдине тут звучить теїстично; переклад не повинен його пом'якшувати.}
\vpar{Українська паралель (наша інтерпретація)}{Григорій Сковорода, «Наркіс. Розмова про те: Пізнай себе» (1769–1771).}{Наша паралель: у «Наркіс» дзеркало води — головний образ, і Сковорода використовує його так само, як цей вірш: те, що видно в дзеркалі, є і всередині, і зовні одночасно. Ми беремо той самий діалог, що й у 1.1, але інший його образ — і кажемо про це прямо, а не вдаємо двох незалежних паралелей.}
\vfig{ilustratsii/1-19-illusion.pdf}{Образ у дзеркалі не є дзеркалом і не віддільний від нього}{Схема наша, не історична ілюстрація.}

\vsec{1.20}
\vdeva{एकं सर्वगतं व्योम बहिर् अन्तर् यथा घटे । नित्यं निरन्तरं ब्रह्म सर्वभूतगणे तथा ॥}
\vtranslit{екам сарваґатам вйома бахір антар ятха ґгате | нітям нірантарам брахма сарвабгутаґане татха ||}
\vperek{Як один усюдисущий простір — зовні й усередині глечика, так повсякчас і без перерви брахман у всій сукупності істот.}
\vcom{Класичний образ глечика: простір не ділиться стінками, хоч і здається «простором у глечику». Пара `bahiḥ antaḥ` повторює `antaḥ … paritaḥ` з 1.19, але тут вона віднесена до простору, а не до образу в дзеркалі. `nirantaram` — «без проміжку», тобто не «вічно» в часовому сенсі, а без розриву в просторі; переклад мусить це розрізняти.}
\vpar{Українська паралель (ПОТРЕБУЄ ПЕРЕВІРКИ)}{українська філософська традиція, не визначено.}{ПОТРЕБУЄ ПЕРЕВІРКИ: образ простору, що не ділиться стінками посудини, в українському матеріалі найчастіше приходить через Сковороду, але конкретного місця, де він говорить саме про посуд і простір, ми за джерелом не перевірили. Тому тут стоїть позначка, а не ім'я.}
\vfig{ilustratsii/1-20-space.pdf}{Простір не ділиться стінками глечика: bahiḥ antaḥ повторює 1.19, але про простір}{Схема наша, не історична ілюстрація.}


% ЗГЕНЕРОВАНО scripts/build_format_v2_chapter.py — не редагувати вручну.
% Джерело: editorial/cards/2.*.json, format-v2/rozdil-2.json, paraleli.json.

\naukchapter{2}{25}{Глава 2. Радість усвідомлення}
\naukintro{Друга глава переходить від заперечення до вигуку. Десять разів у ній звучить aho — «о диво», і саме цей вигук тримає главу разом: він щоразу стоїть там, де учень упізнає в собі те, що раніше приписував іншому. Друга половина глави (2.23–2.25) міняє образ: замість дзеркал і мотузок зʼявляється океан, у якому хвилі світів піднімаються й опадають самі собою. Це не доказ, а спосіб побачити, що рух не належить тому, хто його бачить.}
\vsec{2.1}
\vdeva{अहो निरञ्जनः शान्तो बोधोऽहं प्रकृतेः परः । एतावन्तम् अहं कालं मोहेनैव विडम्बितः ॥}
\vtranslit{ахо ніраньджанах шанто бодго’хам пракрітех парах | етавантам ахам калам мохенайва відамбітах ||}
\vperek{О диво: я — незаплямлене, спокійне усвідомлення, вище за природу. Увесь цей час мене дурила сама лише омана.}
\vcom{Розділ 2 відкривається вигуком подиву (aho) — це не аргумент, а зміна погляду. Три означення (nirañjana, śānta, bodha) ідуть як одне ціле: Я не «має» спокій, а є усвідомленням, яке не плямується. Viḍambitaḥ — не «я помилявся», а «мене дурили»: омана діяла, а не Я. Для українського перекладу важливо не зробити з цього самодокір: текст не звинувачує, а констатує.}
\vpar{Українська паралель (наша інтерпретація)}{Григорій Сковорода, «Кільце. Дружня розмова про душевний мир».}{Наша паралель: у Сковороди справжній мир теж не здобувають ззовні — його впізнають як те, що вже було під нашаруванням думок. Це наша інтерпретація, а не твердження про вплив.}
\vfig{ilustratsii/2-1-flame.pdf}{Полумʼя без диму: усвідомлення, яке не забруднила омана}{Схема наша, не історична ілюстрація.}

\vsec{2.2}
\vdeva{यथा प्रकाशयाम्य् एको देहम् एनं तथा जगत् । अतो मम जगत् सर्वं अथवा न च किञ्चन ॥}
\vtranslit{ятха пракашаямй еко дехам енам татха джаґат | ато мама джаґат сарвам атхава на ча кіньчана ||}
\vperek{Як я один освітлюю це тіло, так освітлюю і світ. Тому весь світ — мій; або ж — аніщо.}
\vcom{Вірш тримає дві відповіді одночасно: «світ мій» і «світ — аніщо». Це не суперечність, а два боки одного: світ не існує окремо від того, хто його освітлює, і саме тому не є чимось набутим. athavā («або ж») — сигнал, що текст свідомо лишає обидві формули, а не обирає одну. Перекладати його як «тобто» було б спрощенням.}
\vpar{Українська паралель (наша інтерпретація)}{Памфіл Юркевич, «Серце і його значення в духовному житті людини» (1860).}{Наша паралель: коли вірш каже, що одне світло освітлює і тіло, і світ, він ставить світло в джерело, а не в речі. Юркевич шукає таке джерело в серці, а не в розумі, — це наша асоціація за будовою думки.}
\vfig{ilustratsii/2-2-witness.pdf}{Одне світло в центрі: воно освітлює і тіло, і світ}{Схема наша, не історична ілюстрація.}

\vsec{2.3}
\vdeva{सशरीरम् अहो विश्वं परित्यज्य मयाधुना । कुतश्चित् कौशलाद् एव परमात्मा विलोक्यते ॥}
\vtranslit{сашарірам ахо вішвам парітяджя маядгуна | куташчіт каушалад ева параматма вілокяте ||}
\vperek{О диво: тепер, відкинувши світ разом із тілом, я — самим лише вмінням, невідомо звідки — бачу вище Я.}
\vcom{Ключове тут — пасив vilokyate: бачення не є зусиллям того, хто бачить. kutaścit («невідомо звідки») підкреслює, що вміння не вироблене вправою. Для української редакції це найтонше місце: активне «я бачу» зраджує сенс, тому в перекладі додано «невідомо звідки» як відповідник kutaścit. saśarīram означає «не залишаючи тіла» — не фізичне відкидання, а відмова від ототожнення.}
\vpar{Українська паралель (ПОТРЕБУЄ ПЕРЕВІРКИ)}{українська містична традиція, не визначено.}{ПОТРЕБУЄ ПЕРЕВІРКИ: мотив уміння, що приходить невідомо звідки, найближчий до українського барокового письма про внутрішню людину, але конкретного твору ми за джерелом не перевірили.}
\vfig{ilustratsii/2-3-neti.pdf}{Світ і тіло відкинуто — лишається вище Я}{Схема наша, не історична ілюстрація.}

\vsec{2.4}
\vdeva{यथा न तोयतो भिन्नास् तरङ्गाः फेनबुद्बुदाः । आत्मनो न तथा भिन्नं विश्वम् आत्मविनिर्गतम् ॥}
\vtranslit{ятха на тоято бгіннас таранґах пхенабудбудах | атмано на татха бгіннам вішвам атмавінірґатам ||}
\vperek{Як хвилі, піна й бульбашки не відмінні від води, так і світ, що вийшов із Я, не відмінний від Я.}
\vcom{Образ води й хвиль працює на одному слові — bhinnam («відмінний»): хвиля не є водою в сенсі «іншою речовиною», але й не тотожна їй як окремий предмет. Це не тотожність і не різниця, а відмова від самої пари. ātmavinirgatam («той, що вийшов із Я») — не «створений», тому в перекладі немає слова «створений».}
\vpar{Українська паралель (наша інтерпретація)}{Григорій Сковорода, «Наркіс. Розмова про те: Пізнай себе».}{Наша паралель: хвиля, що не віддільна від води, — образ, близький до сковородинського розрізнення видимого й того, що в ньому видно. Це наша інтерпретація, не вказівка на джерело вірша.}
\vfig{ilustratsii/2-4-neti.pdf}{Хвиля, піна й бульбашка виявляються водою}{Схема наша, не історична ілюстрація.}

\vsec{2.5}
\vdeva{तन्तु-मात्रो भवेद् एव पटो यद्वद् विचारतः । आत्म-तन्मात्रम् एवेदं तद्वद् विश्वं विचारितम् ॥}
\vtranslit{танту-матро бгавед ева пато ядвад вічаратах | атма-танматрам еведам тадвад вішвам вічарітам ||}
\vperek{Як тканина при дослідженні виявляється лише ниткою, так і цей світ, коли його дослідити, є лише Я.}
\vcom{Це перший вірш книги, де прямо названо метод — vicāra, «дослідження». Тканина не зникає, коли ми бачимо нитку; зникає лише уявлення, що тканина — щось понад нитку. tanmātra («саме те, тільки те») — технічний термін, який у санкх'ї означає тонкі елементи; тут він ужитий у звичайному обмежувальному сенсі, і плутати ці два вживання не слід. Це спостереження про слововживання, а не встановлений факт.}
\vpar{Українська паралель (наша інтерпретація)}{Олександр Довженко, «Зачарована Десна» (1956).}{Наша паралель: у Довженка річ теж розкривається, коли її розібрати до основи, — і основа виявляється простішою й живішою за річ. Наше порівняння стосується способу бачити, а не вчення.}
\vfig{ilustratsii/2-5-thread.pdf}{Тканина розібрана на нитку: так і світ — на Я}{Схема наша, не історична ілюстрація.}

\vsec{2.6}
\vdeva{यथैवेक्षुरसे कॢप्ता तेन व्याप्तैव शर्करा । तथा विश्वं मयि कॢप्तं मया व्याप्तं निरन्तरम् ॥}
\vtranslit{ятхайвекшурасе кліпта тена в'яптайва шаркара | татха вішвам маї кліптам мая в'яптам нірантарам ||}
\vperek{Як цукор уже вкладений у сік цукрової тростини й просякнутий ним, так і світ вкладений у мене й просякнутий мною безперервно.}
\vcom{Построкове зіставлення: редакторська чернетка, не рецензовано.}
\vcom{Рядок а}
\vdeva{यथैवेक्षुरसे कॢप्ता तेन व्याप्तैव शर्करा}
\vtranslit{ятхайвекшурасе кліпта тена в'яптайва шаркара}
\vperek{Як цукор уже вкладений у сік цукрової тростини й просякнутий ним,}
\vcom{yathaiva — підсилена порівняльна вказівка «як»; ikṣurase — «у соці цукрової тростини» (місцевий відмінок). śarkarā — «цукор», жіночий рід; kḷptā та vyāptā (у формі vyāptaiva) описують його як «вкладений» та «просякнутий» соком; tena — «ним». Це порівняння не стверджує, що цукор окремо додали в сік; конкретне членування сандгі слід перевірити за друкованими джерелами.}
\vcom{Рядок б}
\vdeva{तथा विश्वं मयि कॢप्तं मया व्याप्तं निरन्तरम्}
\vtranslit{татха вішвам маї кліптам мая в'яптам нірантарам}
\vperek{так і світ вкладений у мене й просякнутий мною безперервно.}
\vcom{tathā — «так само», друга половина порівняння з yathaiva; viśvaṃ — «світ». mayi — «в мені» (місцевий), mayā — «мною» (орудний), отже санскрит розрізняє місце і спосіб охоплення. kḷptaṃ та vyāptaṃ відповідають формам першого піврядка, але узгоджені з середнім родом viśvaṃ; nirantaram — «безперервно». Цукор і тростинний сік — образ невідокремленості в цьому вірші, а не доказ фізичного творення світу; українські літературні аналогії для 2.6 у чинній книзі ще потребують перевірки.}
\vcom{Образ працює не на «створенні», а на присутності: цукор не доданий до соку — він у ньому вже є, і окремо його не видно. Ключове слово — kḷpta («вкладений, влаштований»), а не sṛṣṭa («створений»); підміна змінила б думку з «світ у мені влаштований» на «світ мною зроблений». Тому в перекладі немає слова «створений».}
\vpar{Українська паралель (ПОТРЕБУЄ ПЕРЕВІРКИ)}{українська етнографічна проза, не визначено.}{ПОТРЕБУЄ ПЕРЕВІРКИ: образ речовини, що вже просякнута своїм смаком, в українській прозі трапляється часто, але назвати твір, де він несе саме цю думку, ми не можемо без звірки.}
\vfig{ilustratsii/2-6-balance.pdf}{Дві шальки: сік тростини й цукор, що вже в ньому}{Схема наша, не історична ілюстрація.}

\vsec{2.7}
\vdeva{आत्माज्ञानाज् जगद् भाति हात्मज्ञानान् न भासते । रज्ज्वज्ञानाद् अहिर् भाति तज्ज्ञानाद् भासते न हि ॥}
\vtranslit{атмаджнянадж джаґад бгаті хатмаджнянан на бгасате | раджджваджнянад ахір бгаті таджджнянад бгасате на хі ||}
\vperek{Від незнання Я світ сяє; від знання Я він не сяє. Від незнання мотузки сяє змія; від знання її вона не сяє.}
\vcom{Построкове зіставлення: редакторська чернетка, не рецензовано.}
\vcom{Рядок а}
\vdeva{आत्माज्ञानाज् जगद् भाति हात्मज्ञानान् न भासते}
\vtranslit{атмаджнянадж джаґад бгаті хатмаджнянан на бгасате}
\vperek{Від незнання Я світ сяє; від знання Я він не сяє.}
\vcom{ātma-ajñānāt («від незнання Я») і ātma-jñānāt («від знання Я») вказують на протиставлення причинних відношень через відкладний відмінок; на стиках слів у чинній IAST стоїть сандгі. jagat — «світ», bhāti — «сяє, являється», na bhāsate — «не сяє». Частку ha, сховану у сполученні hātmajñānān, не можна без текстологічної перевірки перетворювати на окрему богословську тезу. Українське «не сяє» збережено дослівно: вірш не вимагає перекладу «світ фізично зникає».}
\vcom{Рядок б}
\vdeva{रज्ज्वज्ञानाद् अहिर् भाति तज्ज्ञानाद् भासते न हि}
\vtranslit{раджджваджнянад ахір бгаті таджджнянад бгасате на хі}
\vperek{Від незнання мотузки сяє змія; від знання її вона не сяє.}
\vcom{rajju-ajñānāt — «від незнання мотузки», ahiḥ (перед наступним словом у сандгі ahir) — «змія». bhāti повторює дієслово першого рядка, а taj-jñānāt пов'язує друге заперечення зі знанням тієї ж мотузки; bhāsate na hi — «не сяє», з кінцевою часткою hi. Мотузка не перестає існувати: змінюється помилкове впізнавання змії. Це авторський граматично-філософський коментар до образу, не незалежно рецензоване свідчення про всі варіанти санскритського тексту.}
\vcom{Класичний приклад мотузки й змії, але з точністю, яку легко зіпсувати: текст каже na bhāsate — «не сяє», а не «зникає». Зникає не мотузка, а змія; світ не анульовується, а перестає сяяти як щось окреме. Українське «зникає» тут було б зручнішим і хибним.}
\vpar{Українська паралель (наша інтерпретація)}{Григорій Сковорода, «Ікона Алківіадська».}{Наша паралель: Сковорода теж показує, що тінь зникає не тому, що її прибрали, а тому, що побачили річ. Це наша асоціація, а не твердження про спільне джерело.}
\vfig{ilustratsii/2-7-illusion.pdf}{Змія й мотузка: від незнання й від знання}{Схема наша, не історична ілюстрація.}

\vsec{2.8}
\vdeva{प्रकाशो मे निजं रूपं नातिरिक्तोऽस्म्य् अहं ततः । यदा प्रकाशते विश्वं तदाऽहंभास एव हि ॥}
\vtranslit{пракашо ме ніджам рупам натірікто’смй ахам татах | яда пракашате вішвам тада’хамбгаса ева хі ||}
\vperek{Світло — моя власна природа; я не більший за нього. Коли сяє світ, сяє лише сяйво Я.}
\vcom{Построкове зіставлення: редакторська чернетка, не рецензовано.}
\vcom{Рядок а}
\vdeva{प्रकाशो मे निजं रूपं नातिरिक्तोऽस्म्य् अहं ततः}
\vtranslit{пракашо ме ніджам рупам натірікто’смй ахам татах}
\vperek{Світло — моя власна природа; я не більший за нього.}
\vcom{prakāśaḥ («світло, сяйво») у сполуці prakāśo є називним відмінком; me («моє»), nijaṃ rūpaṃ («власна природа») окреслюють самовизначення мовця. nātirikto'smy містить заперечення na та atiriktaḥ asmi — «я не окремий / не надлишковий»; чинний український переклад передає це як «не більший». ahaṃ tataḥ підтримує особову вказівку. Не ототожнюємо дослівність цього перекладу з остаточною санскритологічною інтерпретацією без незалежної звірки сандгі.}
\vcom{Рядок б}
\vdeva{यदा प्रकाशते विश्वं तदाऽहंभास एव हि}
\vtranslit{яда пракашате вішвам тада’хамбгаса ева хі}
\vperek{Коли сяє світ, сяє лише сяйво Я.}
\vcom{yadā … tadā («коли … тоді») поєднує сяйво світу й Я; prakāśate — «сяє, виявляється»; viśvaṃ — «світ». У tadā'haṃbhāsa (із сандгі) складник ahaṃbhāsa описує «сяйво Я» в чинній редакції, однак його точне членування й синтаксичну функцію ще потрібно філологічно перевірити. eva hi — підсилювальне «саме, адже». Не змішувати ahaṃbhāsa із терміном ahaṃkāra («его-творення»). Українська паралель із Павлом Тичиною лишається редакторською асоціацією, не твердженням про вплив.}
\vcom{Тонкість у nātiriktaḥ — «не більший, не надлишковий»: світло не є чимось поза своїм проявом, і того, хто сяє, не можна відділити від сяйва. Складник ahaṃbhāsa рідкісний; перекладено «сяйво Я», а не «его» — его було б ahaṃkāra, інше слово.}
\vpar{Українська паралель (наша інтерпретація)}{Павло Тичина, «Сонячні кларнети» (1918).}{Наша паралель: у Тичини світло теж не позичене ззовні, а є способом, яким світ себе показує. Наше порівняння — типологічне.}
\vfig{ilustratsii/2-8-flame.pdf}{Полумʼя, що не позичає світла: світло і є Я}{Схема наша, не історична ілюстрація.}

\vsec{2.9}
\vdeva{अहो विकल्पितं विश्वम् अज्ञानान् मयि भासते । रूप्यं शुक्तौ फणी रज्जौ वारि सूर्यकरे यथा ॥}
\vtranslit{ахо вікалпітам вішвам аджнянан маї бгасате | руп'ям шуктау пхані раджджау варі сур'якаре ятха ||}
\vperek{О диво: вигаданий світ через незнання сяє в мені — як срібло в перламутрі, як змій у мотузці, як вода в мареві сонячного променя.}
\vcom{Построкове зіставлення: редакторська чернетка, не рецензовано.}
\vcom{Рядок а}
\vdeva{अहो विकल्पितं विश्वम् अज्ञानान् मयि भासते}
\vtranslit{ахо вікалпітам вішвам аджнянан маї бгасате}
\vperek{О диво: вигаданий світ через незнання сяє в мені —}
\vcom{aho — вигук здивування; vikalpitaṃ («помислений, сконструйований») узгоджується з viśvam («світ»). ajñānān в цифровому читанні вказує на незнання як умову омани; mayi — «в мені» (місцевий), bhāsate — «сяє, являється». Форма ajñānān і її сандгі потребують окремої звірки з першодруками; цей коментар не видає електронну транскрипцію за критичне видання.}
\vcom{Рядок б}
\vdeva{रूप्यं शुक्तौ फणी रज्जौ वारि सूर्यकरे यथा}
\vtranslit{руп'ям шуктау пхані раджджау варі сур'якаре ятха}
\vperek{як срібло в перламутрі, як змій у мотузці, як вода в мареві сонячного променя.}
\vcom{rūpyaṃ — «срібло», śuktau — «у перламутрі»; phaṇī — «змія», rajjau — «у мотузці»; vāri — «вода», sūryakare — «у сонячному промені»; yathā — «як». Три порівняння говорять про помилку впізнавання: реальна річ є підставою для гаданої, але саме срібло, змія чи вода з неї не виникають фізично. Тому не плутати цей вірш із 2.10 про причину та форму.}
\vcom{Три приклади підібрані за одним принципом: у кожному є річ, що є, і річ, якої немає, але яка видима. Перламутр є, срібла немає; мотузка є, змія немає; промінь є, води немає. Спільне в них не «світ ілюзорний», а точніше: **помилка має місце в тому, що є, і не додає нічого нового до буття**. Слово vikalpitam («сконструйоване») вказує на роботу розуму, а не на онтологічний статус світу. Порівняй 2.10, де той самий світ описаний уже як такий, що вийшов із мене й повернеться: там акцент на походженні, тут — на способі явлення.}
\vpar{Українська паралель (ПОТРЕБУЄ ПЕРЕВІРКИ)}{українська поезія XX ст., не визначено.}{ПОТРЕБУЄ ПЕРЕВІРКИ: перламутр, у якому сяє срібло, — образ, близький до української лірики про оманливе бачення, але конкретного рядка ми не перевіряли й тому не називаємо.}
\vfig{ilustratsii/2-9-illusion.pdf}{Три омани в одному ряді: срібло, змій, вода}{Схема наша, не історична ілюстрація.}

\vsec{2.10}
\vdeva{मत्तो विनिर्गतं विश्वं मय्य् एव लयम् एष्यति । मृदि कुम्भो जले वीचिः कनके कटकं यथा ॥}
\vtranslit{матто вінірґатам вішвам майй ева лаям ешяті | мріді кумбго джале вічіх канаке катакам ятха ||}
\vperek{Світ, що вийшов із мене, у мені ж і розчиниться — як горщик у глині, як хвиля у воді, як браслет у золоті.}
\vcom{Построкове зіставлення: редакторська чернетка, не рецензовано.}
\vcom{Рядок а}
\vdeva{मत्तो विनिर्गतं विश्वं मय्य् एव लयम् एष्यति}
\vtranslit{матто вінірґатам вішвам майй ева лаям ешяті}
\vperek{Світ, що вийшов із мене, у мені ж і розчиниться —}
\vcom{matto («від мене», вихідний відмінок), vinirgataṃ («те, що вийшло») та viśvaṃ («світ») виражають походження. mayy eva — «саме в мені» (місцевий відмінок з eva); layam eṣyati — «піде до розчинення, повернеться до основи». Ідея тут не про хибне бачення, як у 2.9, а про залежність проявленої форми від її джерела; буквальний смисл не вимагає фізичного знищення світу.}
\vcom{Рядок б}
\vdeva{मृदि कुम्भो जले वीचिः कनके कटकं यथा}
\vtranslit{мріді кумбго джале вічіх канаке катакам ятха}
\vperek{як горщик у глині, як хвиля у воді, як браслет у золоті.}
\vcom{mṛdi — «у глині», kumbho (kumbhaḥ) — «горщик»; jale — «у воді», vīciḥ — «хвиля»; kanake — «у золоті», kaṭakaṃ — «браслет»; yathā — «як». Кожна пара показує, що форма має матеріальну основу: глина залишається глиною, вода — водою, золото — золотом. Це пояснювальна аналогія, не ототожнення трьох речей і не доказ спільного історичного джерела з українською філософією.}
\vcom{Порівняння тут іншого типу, ніж у 2.9. Там ішлося про помилку сприйняття, тут — про **відношення причини й наслідку**. Горщик не є чимось додатковим до глині, браслет — до золота: змінилася форма, не речовина. Це класична для індійської думки модель, де наслідок уже міститься в причині. Слово laya («розчинення») не означає знищення: воно означає повернення до того, з чого виникло. Разом 2.9 і 2.10 дають два різні пояснення одного й того самого: світ і помилковий у способі явлення (2.9), і не відокремлений від джерела (2.10).}
\vpar{Українська паралель (наша інтерпретація)}{Григорій Сковорода, «Діалог, або Розмова про давній світ».}{Наша паралель: думка, що річ розчиняється у своєму началі, у Сковороди теж тримає всю розмову про давній і новий світ. Наша інтерпретація за будовою думки, не за джерелом.}
\vfig{ilustratsii/2-10-neti.pdf}{Горщик, хвиля й браслет розчиняються в своєму началі}{Схема наша, не історична ілюстрація.}

\vsec{2.11}
\vdeva{अहो अहं नमो मह्यं विनाशो यस्य नास्ति मे । ब्रह्मादिस्तम्बपर्यन्तं जगन्नाशे ऽपि तिष्ठतः ॥}
\vtranslit{ахо ахам намо махям вінашо яся насті ме | брахмадістамбапар'янтам джаґаннаше ’пі тіштхатах ||}
\vperek{О диво! Поклін мені самому — тому, в кого немає загибелі, хто стоїть, навіть коли гине світ від Брахми аж до стебла трави.}
\vcom{Зачин «aho ahaṃ namo mahyaṃ» — самопрославлення, і це навмисний парадокс у тексті, що весь будується на запереченні «я». Тут «я» не его, а те, що лишається, коли его знято: тому й поклін спрямований не собі як особі, а тому, що в собі не гине. Діапазон «brahmādistambaparyantam» («від Брахми до стебла трави») — це не поетичне перебільшення, а повний обсяг живих істот у традиційній картині: найвищий бог і найменша рослина. Твердження сильне саме тому, що охоплює обидва краї. Те саме «namo mahyaṃ» повторюється в 2.12 — це рефрен, а не випадковий збіг.}
\vpar{Українська паралель (наша інтерпретація)}{Тарас Шевченко, «Подражаніє Екклезіястові».}{Наша паралель: у Шевченка все, що має кінець, теж минає, і саме на цьому тлі постає те, що стоїть. Наше порівняння — за настроєм, а не твердження про спільну думку.}
\vfig{ilustratsii/2-11-space.pdf}{Небо без хмар: світ минає, а те, що дивиться, стоїть}{Схема наша, не історична ілюстрація.}

\vsec{2.12}
\vdeva{अहो अहं नमो मह्यं एको ऽहं देहवान् अपि । क्वचिन् न गन्ता नागन्ता व्याप्य विश्वम् अवस्थितः ॥}
\vtranslit{ахо ахам намо махям еко ’хам дехаван апі | квачін на ґанта наґанта в'яп'я вішвам авастхітах ||}
\vperek{О диво! Поклін мені самому! Я один — і хоч маю тіло, нікуди не йду й не приходжу: охопивши все, перебуваю.}
\vcom{Парадокс у «dehavān api» («навіть маючи тіло»): тіло є, а руху немає, бо немає куди йти — усе вже охоплене. «na gantā nāgantā» («не той, хто йде, і не той, хто приходить») знімає не дію, а подорож як таку: обидва дієприкметники майбутнього часу, тобто йдеться не про те, що вже сталося, а про те, що не станеться. Разом із 2.11 це пара: там «не гину, коли гине світ» (час), тут «не йду, коли все охоплене» (простір). Два виміри, один висновок.}
\vpar{Українська паралель (наша інтерпретація)}{Василь Стус, «Палімпсести».}{Наша паралель: у Стуса теж є стан, у якому людина не приходить і не йде, а перебуває всупереч усьому, що з нею роблять. Наша асоціація, не вказівка на вплив.}
\vfig{ilustratsii/2-12-space.pdf}{Простір, що нікуди не йде, бо вже охопив усе}{Схема наша, не історична ілюстрація.}

\vsec{2.13}
\vdeva{अहो अहं नमो मह्यं दक्षो नास्तीह मत्समः । असंस्पृश्य शरीरेण येन विश्वं चिरं धृतम् ॥}
\vtranslit{ахо ахам намо махям дакшо настіха матсамах | асамспрішя шарірена єна вішвам чірам дгрітам ||}
\vperek{О диво! Поклін мені самому — тому, хто, не торкаючись тілом, довго втримує все; немає тут нікого, хто дорівнював би мені в цьому.}
\vcom{Парадокс в asaṃspṛśya («не торкнувшись»): світ утримується, і водночас немає дотику. Тіло назване тут не як перешкода, а як те, чим дотик не відбувається — це не заперечення тіла, а заперечення його як інструменту. Слово dakṣa («спритний») ужите з іронією: спритність зазвичай означає вправність у дії, а тут вона полягає в тому, щоб діяти без дотику. Третій рефренний зачин поспіль (2.11, 2.12, 2.13) показує, що це не окремі вигуки, а серія: кожен наступний вірш додає новий вимір до того самого «поклону собі».}
\vpar{Українська паралель (ПОТРЕБУЄ ПЕРЕВІРКИ)}{українська філософська проза, не визначено.}{ПОТРЕБУЄ ПЕРЕВІРКИ: думка про те, що ціле втримується без дотику, в українському матеріалі найчастіше приходить через Сковороду, але конкретного місця ми не звірили.}
\vfig{ilustratsii/2-13-witness.pdf}{Центр без дотику: усе втримується, а він не торкається}{Схема наша, не історична ілюстрація.}

\vsec{2.14}
\vdeva{अहो अहं नमो मह्यं यस्य मे नास्ति किञ्चन । अथवा यस्य मे सर्वं यद् वाङ्मनसगोचरम् ॥}
\vtranslit{ахо ахам намо махям яся ме насті кіньчана | атхава яся ме сарвам яд ванманасаґочарам ||}
\vperek{О диво! Поклін мені самому! У мене немає нічого — або ж у мене є все, що взагалі досяжне слову й розуму.}
\vcom{Вірш ставить дві взаємовиключні відповіді й не обирає між ними: athavā («або ж») залишає обидві. Це не нерішучість, а точний хід: якщо «нічого» і «все» однаково хибні описи, то обидва виражають те саме — те, що не піддається опису ні як володіння, ні як відсутність. Та сама формула athavā na ca kiñcana є і в 2.2, тож це не одноразовий зворот, а повторюваний прийом тексту. Друга половина обмежує коло: sarvaṃ — це не «все взагалі», а лише те, що потрапляє в поле слова й розуму (vāṅmanasagocara). Те, що поза цим полем, не називається ні «всім», ні «нічим».}
\vpar{Українська паралель (наша інтерпретація)}{Григорій Сковорода, «Убогий Жайворонок».}{Наша паралель: у Сковороди вбогість і багатство теж міряються не майном, а тим, чи вміщується в людину те, чого вона хоче. Наше порівняння — типологічне.}
\vfig{ilustratsii/2-14-balance.pdf}{Дві шальки: нічого і все — обидві про мене}{Схема наша, не історична ілюстрація.}

\vsec{2.15}
\vdeva{ज्ञानं ज्ञेयं तथा ज्ञाता त्रितयं नास्ति वास्तवम् । अज्ञानाद् भाति यत्रेदं सो ऽहम् अस्मि निरञ्जनः ॥}
\vtranslit{джнянам джнєям татха джнята трітаям насті ваставам | аджнянад бгаті ятредам со ’хам асмі ніраньджанах ||}
\vperek{Знання, пізнаване й пізнавач — ця трійця не є справжньою. Те, у чому вона сяє через незнання, — це я, незаплямований.}
\vcom{Заперечується не знання як таке, а **трійця** — поділ на того, хто знає, те, що знають, і саме знання. Це найточніше місце глави: доки є трійця, є й поділ, а поділ і є те, що тут названо незнанням. Зверни увагу на порядок: спершу сказано, що трійця не дійсна, і лише потім — що вона сяє. Тобто вона не зникає: вона продовжує являтися, але вже не як справжня. saḥ aham asmi («я є те») — пряма формула ототожнення, і вона стосується не трійці, а того, у чому трійця являється. nirañjana («незаплямований») стоїть останнім словом навмисно: те, у чому сяє незнання, саме ним не забруднюється.}
\vpar{Українська паралель (наша інтерпретація)}{Памфіл Юркевич, «Серце і його значення в духовному житті людини» (1860).}{Наша паралель: коли вірш називає трійцю пізнавача, пізнаваного й знання несправжньою, він питає про те саме, про що Юркевич питає інакше: де в людині те, що не є ані річчю, ані думкою про річ.}
\vfig{ilustratsii/2-15-neti.pdf}{Трійця пізнання відкинута: лишається незаплямований}{Схема наша, не історична ілюстрація.}

\vsec{2.16}
\vdeva{द्वैतमूलम् अहो दुःखं नान्यत् तस्यास्ति भेषजम् । दृश्यम् एतन् मृषा सर्वं एको ऽहं चिद्रसो ऽमलः ॥}
\vtranslit{двайтамулам ахо духкхам нанят тасясті бгешаджам | дрішям етан мріша сарвам еко ’хам чідрасо ’малах ||}
\vperek{О диво! Корінь страждання — двоїстість, і немає від нього інших ліків. Усе це видиме — хибне. Я один: смак свідомості, незаплямований.}
\vcom{Вірш має три частини й будується як діагноз, рецепт і твердження. Діагноз: duḥkha має корінь у двоїстості, а не в обставинах. Рецепт: nānyat tasya asti bheṣajam — «немає інших ліків», тобто лікування не в усуненні окремих страждань, а в усуненні поділу. Твердження: dṛśyam etan mṛṣā — видиме хибне. Останнє найризикованіше, і його не слід читати як «світу немає»: mṛṣā означає «марне, несправжнє» в сенсі несамостійного буття, а не відсутності. cidrasaḥ («чий смак — cit») — рідкісний і точний складник: не «я є свідомість», а «я маю смак свідомості», тобто свідомість не річ, якою я володію, а те, чим я відчувається.}
\vpar{Українська паралель (наша інтерпретація)}{Ліна Костенко, «Сад нетанучих скульптур» (1987).}{Наша паралель: у Костенко двоїстість теж постає джерелом болю, а не умовою мислення. Наше порівняння — за способом ставити питання.}
\vfig{ilustratsii/2-16-knot.pdf}{Вузол двоїстості: корінь болю і єдине, що треба розвʼязати}{Схема наша, не історична ілюстрація.}

\vsec{2.17}
\vdeva{बोधमात्रोऽहम् अज्ञानाद् उपाधिः कल्पितो मया । एवं विमृशतो नित्यं निर्विकल्पे स्थितिर् मम ॥}
\vtranslit{бодгаматро’хам аджнянад упадгіх калпіто мая | евам вімрішато нітям нірвікалпе стхітір мама ||}
\vperek{Я — лише усвідомлення; через незнання я домисив собі обмеження. Хто так постійно розмірковує, той стоїть у не-розрізненні — і це моє стояння.}
\vcom{Тут названо механізм: upādhi («накладка, обмеження») не є чимось зовнішнім — він kalpita, «домислений», і домисил його саме той, хто тепер його знімає (mayā). Це не докір собі: у тексті немає оцінки, лише констатація авторства помилки. nirvikalpa — «без розрізнення», стан, у якому vikalpa (конструювання) не діє; це той самий корінь, що в 2.9 vikalpita.}
\vpar{Українська паралель (ПОТРЕБУЄ ПЕРЕВІРКИ)}{українська релігійна філософія, не визначено.}{ПОТРЕБУЄ ПЕРЕВІРКИ: мотив межі, яку людина домалювала собі сама, близький до української традиції самоспостереження, але назвати твір без звірки ми не можемо.}
\vfig{ilustratsii/2-17-knot.pdf}{Вузол, завʼязаний власною думкою, і стояння поза ним}{Схема наша, не історична ілюстрація.}

\vsec{2.18}
\vdeva{अहो मयि स्थितं विश्वं वस्तुतो न मयि स्थितम् । न मे बन्धोऽस्ति मोक्षो वा भ्रान्तिः शान्ता निराश्रया ॥}
\vtranslit{ахо маї стхітам вішвам вастуто на маї стхітам | на ме бандго’сті мокшо ва бгрантіх шанта нірашрая ||}
\vperek{О диво: світ ніби в мені — а насправді він у мені не вміщується. Немає в мене ні пут, ні визволення; омана вщухла, втративши опору.}
\vcom{Парадокс навмисний: перший піввірш дає звичну формулу, другий її знімає через vastutaḥ («насправді»). Далі — те саме заперечення, що в 15.18: немає ні пут, ні визволення. nirāśrayā («без опори») вказує, чому омана вщухла: у неї не лишилося основи, а не тому, що її перемогли.}
\vpar{Українська паралель (наша інтерпретація)}{Григорій Сковорода, «Алфавіт, чи Буквар миру» (діалог).}{Наша паралель: у Сковороди дух теж описаний як те, що не має краю, у якому його можна було б замкнути. Наша інтерпретація, не твердження про вплив.}
\vfig{ilustratsii/2-18-space.pdf}{Простір, у який не вміщується те, що в ньому ніби є}{Схема наша, не історична ілюстрація.}

\vsec{2.19}
\vdeva{सशरीरम् इदं विश्वं न किञ्चिद् इति निश्चितम् । शुद्धचिन्मात्र आत्मा च तत् कस्मिन् कल्पनाधुना ॥}
\vtranslit{сашарірам ідам вішвам на кіньчід іті нішчітам | шуддгачінматра атма ча тат касмін калпанадгуна ||}
\vperek{Цей світ разом із тілом — ніщо, так установлено; а Я — лише чиста свідомість. У чому ж тоді ще домислювати?}
\vcom{Друга половина — риторичне питання, і воно важливіше за першу: після того як названо «ніщо» і «чиста свідомість», для домислювання не лишається місця (kasmin — «у чому?»). śuddha-cin-mātra зберігає потрійну будову: чисте — свідомість — лише; кожен складник несе вагу.}
\vpar{Українська паралель (наша інтерпретація)}{Іван Вишенський, «Книжка» (кінець XVI ст.).}{Наша паралель: у Вишенського відмова від світу теж робить питання про набутки зайвим. Наше порівняння — типологічне, не про джерело.}
\vfig{ilustratsii/2-19-neti.pdf}{Світ і тіло названі нічим: лишається чиста свідомість}{Схема наша, не історична ілюстрація.}

\vsec{2.20}
\vdeva{शरीरं स्वर्गनरकौ बन्धमोक्षौ भयं तथा । कल्पनामात्रम् एवैतत् किं मे कार्यं चिदात्मनः ॥}
\vtranslit{шарірам сварґанаракау бандгамокшау бгаям татха | калпанаматрам евайтат кім ме кар'ям чідатманах ||}
\vperek{Тіло, небо й пекло, пута й визволення, страх — усе це лише домислювання. Що мені робити, коли моя природа — свідомість?}
\vcom{Два двоїни (dual) у переліку — svarga-narakau і bandha-mokṣau — тримають пари разом: текст не відкидає пекло й не обирає небо, він ставить обидва в один ряд із домислюванням. Питання kiṃ me kāryam не риторичне в сенсі «нічого»: воно справді питає, що лишається робити, і відповіді в цьому вірші немає.}
\vpar{Українська паралель (ПОТРЕБУЄ ПЕРЕВІРКИ)}{українське бароко, не визначено.}{ПОТРЕБУЄ ПЕРЕВІРКИ: перелік того, що є лише домислом, близький до барокової проповіді про марність уявлень, але конкретного тексту ми не перевіряли.}
\vfig{ilustratsii/2-20-neti.pdf}{Шість домислів закреслено: тіло, небо, пекло, пута, воля, страх}{Схема наша, не історична ілюстрація.}

\vsec{2.21}
\vdeva{अहो जनसमूहे ऽपि न द्वैतं पश्यतो मम । अरण्यम् इव संवृत्तं क्व रतिं करवाण्य् अहम् ॥}
\vtranslit{ахо джанасамухе ’пі на двайтам пашято мама | араням іва самвріттам ква ратім караванй ахам ||}
\vperek{О диво! Навіть у натовпі, коли я не бачу двоїстості, все стає як ліс. Де ж мені шукати втіхи?}
\vcom{Парадокс подвійний. Перший: натовп для того, хто не бачить двоїстості, — те саме, що ліс, тобто самотність не залежить від людей навколо. Другий, і він тонший: питання kva ratiṃ karavāṇi («де мені знайти втіху») поставлене там, де втіху вже ніде шукати, бо поділу немає. Це не скарга, а риторичне питання, що закриває можливість відповіді. Форма optativus karavāṇi («мені б робити») додає відтінку роздуму, а не бажання. Зверни увагу: слово rati («втіха») — те саме, що в 15.2 у складі vaiṣayiko rasaḥ («смак до чуттєвого»); тут воно вжите без засудження, але й без місця, куди його покласти.}
\vpar{Українська паралель (наша інтерпретація)}{Михайло Коцюбинський, «Intermezzo» (1908).}{Наша паралель: у Коцюбинського самотність серед людей теж обертається на простір, у якому все стає живим. Наше порівняння — за образом, а не за вченням.}
\vfig{ilustratsii/2-21-witness.pdf}{Натовп без двоїстості стає лісом}{Схема наша, не історична ілюстрація.}

\vsec{2.22}
\vdeva{नाहम् देहो न मे देहो जीवो नाहम् अहं हि चित् । अयम् एव हि मे बन्ध आसीत् या जीविते स्पृहा ॥}
\vtranslit{нахам дехо на ме дехо джіво нахам ахам хі чіт | аям ева хі ме бандга асіт я джівіте спріха ||}
\vperek{Я не тіло, і тіло не моє. Я не живчик — я свідомість. Ось у чому була моя уза: спрага до життя.}
\vcom{Найважливіше тут — **час дієслова**. āsīt — це минулий незавершений час, тобто «була», а не «є». Уза не описується як наявна: вона названа в минулому, бо вже побачена як те, чим вона була. Це не проголошення свободи, а звіт про те, що побачене. Другий точний момент: заперечуються три ототожнення поспіль (тіло, володіння тілом, jīva), і лише четверте — ствердне (aham cit). jīva тут не синонім ātman: це саме «живчик», індивідуальна істота, і її заперечення стоїть окремо від заперечення тіла. Причина узи названа одна — spṛhā, спрага до життя, а не до насолод. Це вужче й точніше, ніж звичне «бажання».}
\vpar{Українська паралель (наша інтерпретація)}{Леся Українка, «Одержима» (1901).}{Наша паралель: у Лесі Українки те, що тримає людину в житті, теж названо прямо і без утіхи. Наша асоціація, не твердження про вплив.}
\vfig{ilustratsii/2-22-knot.pdf}{Вузол спраги жити: ось у чому була уза}{Схема наша, не історична ілюстрація.}

\vsec{2.23}
\vdeva{अहो भुवनकल्लोलैर् विचित्रैर् द्राक् समुत्थितम् । मय्य् अनन्तमहाम्भोधौ चित्तवाते समुद्यते ॥}
\vtranslit{ахо бгуванакаллолайр вічітрайр драк самуттхітам | майй анантамахамбгодгау чіттавате самудяте ||}
\vperek{О диво! Різнобарвні хвилі світів раптом піднялися — бо в мені, безмежному великому океані, здіймається вітер розуму.}
\vcom{Початок триптиха (2.23–2.25), де той самий образ розгортається в три кроки: вітер здіймається, вітер вщухає, хвилі граються самі. Тут важлива **причинність**: хвилі не є океаном і не породжені ним безпосередньо — їх піднімає вітер, тобто citta. Океан лишається океаном і без хвиль. Слово ambhodhi буквально означає «той, що містить води» — це поетичне іменування океану, і воно тримає образ: води багато, хвиля — її форма, а не додаток. Порівняй 2.10, де були горщик у глині й хвиля у воді: там ішлося про невідокремленість наслідку, тут — про те, що рух має причину, і причина не в океані.}
\vpar{Українська паралель (наша інтерпретація)}{Григорій Сковорода, «Сад божественних пісень».}{Наша паралель: у Сковороди буря пристрастей теж постає як рух, що піднімається зсередини, а не приходить ззовні. Наше порівняння — за образом, не за джерелом.}
\vfig{ilustratsii/2-23-boat.pdf}{Бурхливий океан і човен: хвилі світів піднімає вітер розуму}{Схема наша, не історична ілюстрація.}

\vsec{2.24}
\vdeva{मय्य् अनन्तमहाम्भोधौ चित्तवाते प्रशाम्यति । अभाग्याज् जीववणिजो जगत्पोतो विनश्वरः ॥}
\vtranslit{майй анантамахамбгодгау чіттавате прашам'яті | абгаґядж джіваваніджо джаґатпото вінашварах ||}
\vperek{Коли в мені, безмежному океані, вітер розуму вщухає, — на лихо торговцеві-живчику — човен світу гине.}
\vcom{Іронія тут точна й гірка. abhāgyāt («на лихо») сказано з погляду **торговця**: для jīva, який плавав на цьому човні й мав із нього зиск, затишшя — катастрофа. Для океану — спокій. Одне й те саме явище описане двічі, і оцінки протилежні. Метафора подвійна: човен пливе **на** океані, тобто світ не є океаном, а лише тим, що по ньому рухається; і човен гине не від бурі, а від її відсутності. Це відповідає 15.18 («нема ні пут, ні визволення»): те, що зникає, не було справжнім, а лише плавало.}
\vpar{Українська паралель (ПОТРЕБУЄ ПЕРЕВІРКИ)}{українська філософська лірика, не визначено.}{ПОТРЕБУЄ ПЕРЕВІРКИ: образ човна, що тоне разом із торговцем, в українському матеріалі найближчий до філософської лірики, але конкретного твору ми не звірили.}
\vfig{ilustratsii/2-24-boat.pdf}{Човен тоне в ущухлому океані: втрата тільки для торговця}{Схема наша, не історична ілюстрація.}

\vsec{2.25}
\vdeva{मय्य् अनन्तमहाम्भोधाव् आश्चर्यं जीववीचयः । उद्यन्ति घ्नन्ति खेलन्ति प्रविशन्ति स्वभावतः ॥}
\vtranslit{майй анантамахамбгодгав ашчар'ям джівавічаях | удянті ґгнанті кхеланті правішанті свабгаватах ||}
\vperek{У мені, безмежному океані, диво: хвилі живих істот здіймаються, б'ються, граються й входять — самі собою, за власною природою.}
\vcom{Завершення триптиха, і воно знімає те, що було в 2.23–2.24. Там хвилі піднімав вітер, і човен гинув — тобто була причина й був той, кому зле. Тут **причини немає**: svabhāvataḥ («за власною природою») стоїть останнім словом і скасовує питання «хто це робить». Чотири дієслова йдуть без сполучників, підряд — це не опис послідовності, а перелік того, що відбувається водночас. ghnanti («б'ються») — єдине слово з відтінком болю, і воно стоїть у тому ж ряду, що khelanti («граються»), без переходу й без пояснення. Це і є āścaryam («диво»): не те, що все добре, а те, що все це відбувається саме, без господаря.}
\vpar{Українська паралель (наша інтерпретація)}{Богдан-Ігор Антонич, «Зелена Євангелія» (1938).}{Наша паралель: в Антонича світ теж рухається сам собою, за власною природою, і людина лише бачить цей рух. Наше порівняння — типологічне.}
\vfig{ilustratsii/2-25-motion.pdf}{Хвилі живих істот граються й входять самі собою}{Схема наша, не історична ілюстрація.}


% ЗГЕНЕРОВАНО scripts/build_format_v2_chapter.py — не редагувати вручну.
% Джерело: editorial/cards/3.*.json, format-v2/rozdil-3.json, paraleli.json.

\naukchapter{3}{14}{Глава 3. Іспит пізнання}
\naukintro{Третя глава — найнеспокійніша з перших трьох. Вона не описує, як здобути знання, а показує, що знання саме по собі не вирішує нічого: у ній пʼять разів звучить aho там, де мудрець ловить себе на тому, що лишилося. Звідси — три останні вірші, де спокій описано не як здобуток, а як те, після чого питання про здобуток відпадає.}
\vsec{3.1}
\vdeva{अविनाशिनम् आत्मानम् एकं विज्ञाय तत्त्वतः । तवात्मज्ञस्य धीरस्य कथम् अर्थार्जने रतिः ॥}
\vtranslit{авінашінам атманам екам віджняя таттватах | таватмаджняся дгірася катхам артхарджане ратіх ||}
\vperek{Пізнавши істинно Атман — незнищенний і єдиний — як у тебе, атмазнавця й мудрого, може бути втіха в набутті користі?}
\vcom{Питання katham («як?») риторичне й будується на несумісності двох занять: пізнання єдиного й накопичення множинного. Слово artha тут не «сенс», а саме «користь, вигода» — це той самий термін, що в 10.1 серед того, що відкидають. Зверни увагу на порядок означень: спершу те, що пізнано (незнищенне, єдине), потім той, хто пізнав (атмазнавець, мудрий). Втіха в набутті не заборонена — поставлено питання, чи вона взагалі можлива за такого пізнання.}
\vpar{Українська паралель (наша інтерпретація)}{Григорій Сковорода, «Розмова пʼяти подорожніх про істинне щастя в житті».}{Наша паралель: у Сковороди щастя теж не лежить у набутку, і саме тому його не можна ні відняти, ні докупити. Наше порівняння — типологічне.}
\vfig{ilustratsii/3-1-balance.pdf}{Дві шальки: набуток і пізнаний Атман}{Схема наша, не історична ілюстрація.}

\vsec{3.2}
\vdeva{आत्माज्ञानाद् अहो प्रीतिर् विषयभ्रमगोचरे । शुक्तेर् अज्ञानतो लोभो यथा रजतविभ्रमे ॥}
\vtranslit{атмаджнянад ахо прітір вішаябграмаґочаре | шуктер аджнянато лобго ятха раджатавібграме ||}
\vperek{О диво! Через незнання Атмана з'являється прихильність у полі омани чуттєвих предметів — як через незнання перламутру з'являється жадоба в омані срібла.}
\vcom{Вірш будується на точній паралелі з двох частин, і паралель тут не прикраса, а аргумент: prīti (прихильність) співвіднесена з lobha (жадобою), viṣayabhrama (омана чуттєвого) — з rajatavibhrama (омана срібла), ātmājñāna (незнання Атмана) — з śuktyajñāna (незнання перламутру). Причина в обох випадках одна — незнання, а не вада предмета: перламутр не винен, що його беруть за срібло. Це та сама паралель, що в 2.9, але там був перелік трьох прикладів, а тут — один, розгорнутий до кінця. Слово prīti м'якше за spṛhā (2.22) і за lobha: це радше «прихильність», ніж «спрага».}
\vpar{Українська паралель (ПОТРЕБУЄ ПЕРЕВІРКИ)}{українська барокова проповідь, не визначено.}{ПОТРЕБУЄ ПЕРЕВІРКИ: перламутр і срібло як пара омани близькі до барокової емблематики, але конкретного твору ми не звірили й тому не називаємо.}
\vfig{ilustratsii/3-2-illusion.pdf}{Перламутр і срібло: жадоба з незнання}{Схема наша, не історична ілюстрація.}

\vsec{3.3}
\vdeva{विश्वं स्फुरति यत्रेदं तरङ्गा इव सागरे । सो ऽहम् अस्मीति विज्ञाय किं दीन इव धावसि ॥}
\vtranslit{вішвам спхураті ятредам таранґа іва саґаре | со ’хам асміті віджняя кім діна іва дгавасі ||}
\vperek{Пізнавши: «я — те, у чому сяє весь цей світ, як хвилі в морі», — чого ти біжиш, наче вбогий?}
\vcom{Питання kim ... dhāvasi («чого біжиш») спрямоване не до дії взагалі, а до бігу як способу життя: дієслово dhāv означає саме фізичний біг, метушню. Порівняння з dīna («убогий») вказує на мотив, а не на стан: біжать ті, кому бракує. Формула so 'ham asmi («я є те») — та сама, що в 2.15, але там вона завершувала заперечення трійці, а тут відкриває питання. Образ хвиль у морі повторює 2.10 і 2.23–2.25: це наскрізний образ тексту, а не місцева метафора.}
\vpar{Українська паралель (ПОТРЕБУЄ ПЕРЕВІРКИ)}{українська філософська лірика, не визначено.}{ПОТРЕБУЄ ПЕРЕВІРКИ: докір тому, хто шукає те, що вже має, в українській ліриці трапляється, але назвати твір без звірки ми не можемо.}
\vfig{ilustratsii/3-3-motion.pdf}{Хвиля, що шукає моря, у якому вже стоїть}{Схема наша, не історична ілюстрація.}

\vsec{3.4}
\vdeva{श्रुत्वापि शुद्धचैतन्यम् आत्मानम् अतिसुन्दरम् । उपस्थे ऽत्यन्तसंसक्तो मालिन्यम् अधिगच्छति ॥}
\vtranslit{шрутвапі шуддгачайтаням атманам атісундарам | упастхе ’тянтасамсакто маліням адгіґаччаті ||}
\vperek{Навіть почувши про Атман — чисту свідомість, надзвичайно прекрасну, — той, хто надмірно прив'язаний до того, що під рукою, набуває бруду.}
\vcom{Слово upasthe навмисно не уточнене: воно означає і «на лоні», і «поблизу, під рукою», і традиційно вживається з плотським відтінком. Точний зміст залежить від контексту, якого вірш не дає, — і саме тому переклад зберігає двозначність, а не обирає одне значення. Структура вірша — контраст між почутим і робленим: śrutvā api («навіть почувши») підкреслює, що знання про Атман саме собою нічого не змінює. mālinya («бруд») — не моральна оцінка, а стан забрудненості; те саме слово вживане про одяг і про посуд. Тобто йдеться не про гріх, а про те, що річ втрачає чистоту від дотику.}
\vpar{Українська паралель (наша інтерпретація)}{Іван Франко, «Зівʼяле листя» (1896).}{Наша паралель: у Франка знання теж не рятує від привʼязаності — воно лише робить її виднішою. Наше порівняння — за способом бачити людину, а не за вченням.}
\vfig{ilustratsii/3-4-balance.pdf}{Дві шальки: почуте про Атман і те, що під рукою}{Схема наша, не історична ілюстрація.}

\vsec{3.5}
\vdeva{सर्वभूतेषु चात्मानं सर्वभूतानि चात्मनि । मुनेर् जानत आश्चर्यं ममत्वम् अनुवर्तते ॥}
\vtranslit{сарвабгутешу чатманам сарвабгутані чатмані | мунер джаната ашчар'ям маматвам анувартате ||}
\vperek{О диво: у мудреця, який знає Атман у всіх істотах і всі істоти — в Атмані, досі тримається «моє».}
\vcom{Двозмінна конструкція (Атман у всіх, усі в Атмані) — не повтор, а два напрямки одного знання: назовні й усередину. Диво в тому, що навіть за такого знання лишається mamatvam — не «моє» як майно, а мійність як звичка. Дієслово anuvartate («іде слідом») важливе: воно не каже, що мійність панує, а що вона тягнеться позаду. Це найточніше місце глави про межі знання: знання не вичищає звичку автоматично. Порівняй 3.4, де йшлося про того, хто лише почув: там бруд набувають, тут мійність лишається — різні ступені, і текст не згладжує різниці.}
\vpar{Українська паралель (ПОТРЕБУЄ ПЕРЕВІРКИ)}{українська релігійна думка, не визначено.}{ПОТРЕБУЄ ПЕРЕВІРКИ: докір мудрецеві, в якого лишилося своє, близький до української аскетичної літератури, але конкретного місця ми не перевірили.}
\vfig{ilustratsii/3-5-knot.pdf}{Вузол, що лишається навіть у знанні}{Схема наша, не історична ілюстрація.}

\vsec{3.6}
\vdeva{आस्थितः परमाद्वैतं मोक्षार्थे ऽपि व्यवस्थितः । आश्चर्यं कामवशगो विकलः केलिशिक्षया ॥}
\vtranslit{астхітах парамадвайтам мокшартхе ’пі в'явастхітах | ашчар'ям камавашаґо вікалах келішікшая ||}
\vperek{Той, хто утвердився у вищій недвоїстості й стоїть навіть у справі звільнення, — о диво! — під владою бажання: безпорадний, мов у грі-вправлянні.}
\vcom{Вірш починає серію із чотирьох «див» (3.6–3.8, далі 3.10), і кожне будується за однією схемою: високий стан — і те, що в ньому несподівано лишається. Тут протиставлені два утвердження (āsthitaḥ, vyavasthitaḥ — майже синоніми, поставлені поруч навмисно) і один стан, що їх перекреслює. Слово vikala («безпорадний») сильніше за просте «слабкий»: воно означає позбавленість частини, скаліченість. keliśikṣayā («грою-вправлянням») знижує регістр: те, що мало бути шляхом, виявилося забавкою.}
\vpar{Українська паралель (наша інтерпретація)}{Василь Стус, «Веселий цвинтар» (1985).}{Наша паралель: у Стуса людина теж стоїть на своєму — і все одно бажання веде її далі, ніж вона хотіла. Наша асоціація, не твердження про вплив.}
\vfig{ilustratsii/3-6-knot.pdf}{Вузол бажання, затягнутий поверх недвоїстості}{Схема наша, не історична ілюстрація.}

\vsec{3.7}
\vdeva{उद्भूतं ज्ञानदुर्मित्रम् अवधार्यातिदुर्बलः । आश्चर्यं कामम् आकाङ्क्षेत् कालम् अन्तम् अनुश्रितः ॥}
\vtranslit{удбгутам джнянадурмітрам авадгар'ятідурбалах | ашчар'ям камам аканкшет калам антам анушрітах ||}
\vperek{О диво! Зрозумівши, що знання виявилося поганим другом, — надзвичайно слабкий — він усе одно прагнув би бажання, дійшовши до кінця часу.}
\vcom{Складене jñānadurmitram можна прочитати двома способами: «знання, що є поганим другом» і «поганий друг, [який полягає] у знанні». Обидва дають той самий сенс: знання, на яке покладалися, не врятувало. Найсильніше тут — час дієслова й обставина: ākāṅkṣet стоїть в optativus («прагнув би»), а kālam antam anuśritaḥ («дійшовши до кінця часу») означає останню межу, коли вже нічого не змінити. Тобто бажання не просто лишається — воно лишається там, де його вже неможливо задовольнити. Це не докір, а констатація абсурду, і слово āścaryam («диво») вжите саме в цьому сенсі.}
\vpar{Українська паралель (ПОТРЕБУЄ ПЕРЕВІРКИ)}{українська поезія XX ст., не визначено.}{ПОТРЕБУЄ ПЕРЕВІРКИ: мотив знання, що виявилося поганим другом, близький до української лірики про зраду власних переконань, але конкретного твору ми не називаємо без звірки.}
\vfig{ilustratsii/3-7-flame.pdf}{Полумʼя, що не зігріває: знання як поганий друг}{Схема наша, не історична ілюстрація.}

\vsec{3.8}
\vdeva{इहामुत्र विरक्तस्य नित्यानित्यविवेकिनः । आश्चर्यं मोक्षकामस्य मोक्षाद् एव विभीषिका ॥}
\vtranslit{іхамутра вірактася нітянітявівекінах | ашчар'ям мокшакамася мокшад ева вібгішіка ||}
\vperek{О диво: у того, хто відвернувся від цього й того світу, хто розрізняє вічне й невічне, хто бажає звільнення, — страх саме від звільнення.}
\vcom{Найгостріший вірш серії. Три ознаки поставлені поспіль і всі три «правильні»: відречення, розрізнення, бажання звільнення. І саме з третьої виростає страх, бо звільнення означає втрату того, хто звільняється. ablativus mokṣāt («від звільнення») з eva («саме») робить джерелом страху не світ і не смерть, а мету. vibhīṣikā — не «боязкість», а те, що лякає; слово вживане і про опудало. Тобто звільнення постає як те, чим лякають.}
\vpar{Українська паралель (наша інтерпретація)}{Леся Українка, «Кассандра» (1907).}{Наша паралель: у Лесі Українки знання теж не дає спокою, а забирає його, — і страх приходить саме від того, що має визволити. Наше порівняння — типологічне.}
\vfig{ilustratsii/3-8-knot.pdf}{Вузол, завʼязаний на самому звільненні}{Схема наша, не історична ілюстрація.}

\vsec{3.9}
\vdeva{धीरस् तु भोज्यमानो ऽपि पीड्यमानो ऽपि सर्वदा । आत्मानं केवलं पश्यन् न तुष्यति न कुप्यति ॥}
\vtranslit{дгірас ту бгоджямано ’пі підямано ’пі сарвада | атманам кевалам пашян на тушяті на куп'яті ||}
\vperek{Мудрий же — чи його голублять, чи його мучать — завжди, бачачи лише Атман, і не радіє, і не гнівається.}
\vcom{Два дієприкметники пасиву (bhojyamāna, pīḍyamāna) поставлені як рівноцінні умови: похвала й мука описані однією граматичною формою, тобто для мудрого вони вже не різні події. Це не байдужість: na tuṣyati na kupyati — два окремі заперечення, а не одне «не хвилюється». Слово kevalam («лише») робить ключову роботу: бачить не «Атман крізь усе», а лише Атман, і саме тому немає кому радіти. Форма дієприкметника paśyan (теперішній час) вказує на тривалість, а не на одноразове прозріння.}
\vpar{Українська паралель (наша інтерпретація)}{Тарас Шевченко, «Неофіти» (1857).}{Наша паралель: у Шевченка гідність теж не залежить від того, хвалять людину чи мучать. Наше порівняння — за будовою думки, не про джерело.}
\vfig{ilustratsii/3-9-balance.pdf}{Дві шальки: хвала й мука — важать однаково}{Схема наша, не історична ілюстрація.}

\vsec{3.10}
\vdeva{चेष्टमानं शरीरं स्वं पश्यत्य् अन्यशरीरवत् । संस्तवे चापि निन्दायां कथं क्षुभ्येत् महाशयः ॥}
\vtranslit{чештаманам шарірам свам пашятй аняшарірават | самставе чапі ніндаям катхам кшубгєт махашаях ||}
\vperek{Своє тіло, що рухається, він бачить як тіло чуже. То як у хвалі чи в ганьбі збентежився б великодухий?}
\vcom{Вірш закриває серію «див» і робить це тихо, без āścaryam: тут немає вигуку, є питання katham («як?»), на яке відповідь очевидна. ceṣṭamānam («те, що ворушиться») вибране навмисно: тіло описане не як річ, а як процес, і саме тому його легко побачити збоку. Пара saṃstava / nindā («хвала» / «ганьба») — та сама пара, що bhojyamāna / pīḍyamāna у 3.9, але там ішлося про дію на тіло, тут — про слово. mahāśayaḥ («великодухий») буквально означає «той, у кого велике ложе», і це не метафора розуму, а образ простору, у якому вміщається і хвала, і ганьба.}
\vpar{Українська паралель (наша інтерпретація)}{Олександр Довженко, «Щоденник» (1941–1956).}{Наша паралель: у Довженка тіло теж постає як щось, що людині належить менше, ніж їй здається. Наше порівняння — за способом бачити себе.}
\vfig{ilustratsii/3-10-witness.pdf}{Тіло як чуже в колі побаченого, бачитель у центрі}{Схема наша, не історична ілюстрація.}

\vsec{3.11}
\vdeva{मायामात्रम् इदं विश्वं पश्यन् विगतकौतुकः । अपि सन्निहिते मृत्यौ कथं त्रस्यति धीरधीः ॥}
\vtranslit{маяматрам ідам вішвам пашян віґатакаутуках | апі санніхіте мрітяу катхам трасяті дгірадгіх ||}
\vperek{Бачачи, що весь цей світ — тільки мая, і втративши до нього охоту, — навіть коли смерть уже поруч, як злякався б той, чий розум стійкий?}
\vcom{Слово kautuka тут важить більше, ніж здається: це не просто «цікавість», а охота до святкового, потяг до видовища. Тобто йдеться не про втрату знань, а про втрату інтересу до гри. locativus absolutus sannihite mṛtyau («коли смерть присутня») ставить умову максимально жорстко: не «якось колись», а тут і тепер. dhīradhīḥ («чий розум стійкий») — не «хоробрий»: стійкість розуму, а не безстрашність характеру. Питання katham знову риторичне, як у 3.3 і 3.10 — глава послідовно закриває питання питаннями.}
\vpar{Українська паралель (наша інтерпретація)}{Михайло Коцюбинський, «Тіні забутих предків» (1911).}{Наша паралель: у Коцюбинського світ теж постає як те, що показується й зникає, не питаючи людину. Наше порівняння — за образом, а не за вченням.}
\vfig{ilustratsii/3-11-illusion.pdf}{Серпанок маї: світ, що зникає під поглядом}{Схема наша, не історична ілюстрація.}

\vsec{3.12}
\vdeva{निःस्पृहं मानसं यस्य नैराश्ये ऽपि महात्मनः । तस्यात्मज्ञानतृप्तस्य तुलना केन जायते ॥}
\vtranslit{ніхспріхам манасам яся найрашє ’пі махатманах | тасятмаджнянатріптася тулана кена джаяте ||}
\vperek{У кого розум без прагнень, хто навіть у безнадії лишається великим духом, хто вдоволений знанням Атмана, — з ким його й порівнювати?}
\vcom{nairāśye 'pi («навіть у безнадії») — найтонше місце. Безнадія тут не недолік, а умова: відсутність сподівань і є тим, що робить розум вільним від прагнень. Тому й сказано api («навіть») — ніби всупереч очікуванню, що без надії людина меншає. tulanā («порівняння, міра») вжите як річ, що не виникає: не «він не має рівних», а «порівняння не народжується», тобто зникає сама операція зіставлення. Форма jāyate (теперішній, «народжується») показує це як постійний стан, а не як разовий висновок.}
\vpar{Українська паралель (ПОТРЕБУЄ ПЕРЕВІРКИ)}{українська аскетична проза, не визначено.}{ПОТРЕБУЄ ПЕРЕВІРКИ: опис розуму без прагнень, з яким немає з чим порівняти, близький до української аскетичної прози, але конкретного твору ми не перевірили.}
\vfig{ilustratsii/3-12-space.pdf}{Небо без хмар: розум без прагнень}{Схема наша, не історична ілюстрація.}

\vsec{3.13}
\vdeva{स्वभावाद् एव जानानो दृश्यम् एतन् न किञ्चन । इदं ग्राह्यम् इदं त्याज्यं स किं पश्यति धीरधीः ॥}
\vtranslit{свабгавад ева джанано дрішям етан на кіньчана | ідам ґрахям ідам тяджям са кім пашяті дгірадгіх ||}
\vperek{Знаючи самою природою, що все видиме — аніщо, — «це треба взяти, це треба відкинути» — то що ж бачить той, чий розум стійкий?}
\vcom{Дві пари герундивів (grāhyam / tyājyam) подані в лапках як чужа мова: це не позиція мудрого, а те, що він чує навколо або в собі. Заперечується не вибір між «взяти» й «відкинути», а сама наявність пари. svabhāvād eva («самою природою») знімає потребу в доказі: це не висновок, а спосіб бачення. Питання kim paśyati («що бачить») точніше за «чи бачить»: воно питає про предмет, а не про факт, бо коли видиме — аніщо, предмета не лишається.}
\vpar{Українська паралель (наша інтерпретація)}{Григорій Сковорода, «Байки харківські» (1774).}{Наша паралель: у Сковороди вибір між речами теж виявляється зайвим для того, хто вже бачить, що речі не важать. Наше порівняння — типологічне.}
\vfig{ilustratsii/3-13-neti.pdf}{Взяти й відкинути відпадають: усе видиме — аніщо}{Схема наша, не історична ілюстрація.}

\vsec{3.14}
\vdeva{अन्तस्त्यक्तकषायस्य निर्द्वन्द्वस्य निराशिषः । यदृच्छयागतो भोगो न दुःखाय न तुष्टये ॥}
\vtranslit{антастяктакашаяся нірдвандвася нірашішах | ядріччаяґато бгоґо на духкхая на туштає ||}
\vperek{У того, в кого всередині відкинуто каламуть, хто без двоїстості й без сподівань, усе, що випадає само собою, не стає ні стражданням, ні задоволенням.}
\vcom{Давальні duḥkhāya і tuṣṭaye — не «не приносить», а «не йде на»: досвід не зараховується ні туди, ні туди. Це тонше за байдужість: байдужість усе ще ставить оцінку «ніяк», а тут досвід просто не конвертується в жодну з двох категорій. kaṣāya («каламуть, осад») — слово з мови фарб і напоїв: це не гріх і не пристрасть у моральному сенсі, а те, що каламутить прозорість. antaḥ («усередині») стоїть першим членом навмисно: відкинуто не зовнішнє, а осад усередині. yadṛcchayāgataḥ («прийшов сам собою») — те саме yadṛcchā, що в 4.4 і далі; тут воно означає не випадковість, а відсутність замовлення.}
\vpar{Українська паралель (наша інтерпретація)}{Микола Зеров, «Камена» (1924).}{Наша паралель: у Зерова ясність теж приходить після того, як каламуть відстоялася, і тоді все, що випадає, приймають рівно. Наше порівняння — за образом, не за джерелом.}
\vfig{ilustratsii/3-14-space.pdf}{Прозоре небо: каламуть відстоялася}{Схема наша, не історична ілюстрація.}


% ЗГЕНЕРОВАНО scripts/build_format_v2_chapter.py — не редагувати вручну.
% Джерело: editorial/cards/4.*.json, format-v2/rozdil-4.json, paraleli.json.

\naukchapter{4}{6}{Глава 4. Хто пізнав — не боїться}
\naukintro{Четверта глава найкоротша в першій половині книги: шість віршів. Вона не додає нового вчення, а перевіряє попереднє питанням про міру: чим відрізняється той, хто пізнав? Відповідь будується на порівняннях — небо й дим, чотири роди істот, — і жодного разу не на обіцянці нагороди. Останній вірш замикає главу тим, що знімає страх: не як наслідок мужності, а як наслідок знання.}
\vsec{4.1}
\vdeva{हन्तात्मज्ञस्य धीरस्य खेलतो भोगलीलया । न हि संसारवाहीकैर् मूढैः सह समानता ॥}
\vtranslit{хантатмаджняся дгірася кхелато бгоґалілая | на хі самсаравахікайр мудгайх саха саманата ||}
\vperek{О! Атмазнавцеві, мудрому, що грається грою досвідів, немає рівні з дурнями, які тягнуть сансару.}
\vcom{Построкове зіставлення: редакторська чернетка, не рецензовано.}
\vcom{Рядок а}
\vdeva{हन्तात्मज्ञस्य धीरस्य खेलतो भोगलीलया}
\vtranslit{хантатмаджняся дгірася кхелато бгоґалілая}
\vperek{О! Атмазнавцеві, мудрому, що грається грою досвідів,}
\vcom{Вигук hanta відкриває подив. Форми ātmajñasya й dhīrasya означують того, хто знає Атман; khelato та bhogalīlayā вводять образ гри досвідів. Це перший смисловий фрагмент одного речення, а не самостійна теза.}
\vcom{Рядок б}
\vdeva{न हि संसारवाहीकैर् मूढैः सह समानता}
\vtranslit{на хі самсаравахікайр мудгайх саха саманата}
\vperek{немає рівні з дурнями, які тягнуть сансару.}
\vcom{Заперечення na hi ... samānatā завершує думку про відсутність рівності; saṃsāravāhīkair mūḍhaiḥ saha називає тих, із ким здійснено порівняння. Образ тягаря сансари взято з чинної картки; атрибуція мовця між цифровими свідками лишається спірною.}
\vcom{Слово vāhīka («погонич, той, хто везе») вибране точно: сансара тут не стан, у якому перебувають, а вантаж, який тягнуть. Дурість полягає не в незнанні, а в тому, що людина впряглася. Протиставлення будується на двох способах дії: khelataḥ («того, хто грається») і vāhīkaiḥ («тих, хто тягне»). Те саме діяння, різна якість. līlā («гра») — слово з релігійної мови, де гра не протистоїть серйозності, а означає дію без потреби. Вигук hanta на початку задає тон подиву, а не докору.}
\vpar{Українська паралель (наша інтерпретація)}{Леся Українка, «Лісова пісня» (1911).}{Наша паралель: у Лесі Українки гра теж не протистоїть досвіду, а є способом жити серед нього, не застигаючи. Наше порівняння — типологічне, не твердження про джерело.}
\vfig{ilustratsii/4-1-motion.pdf}{Дві дороги: гра досвідів і тягар сансари}{Схема наша, не історична ілюстрація.}

\vsec{4.2}
\vdeva{यत् पदं प्रेप्सवो दीनाः शक्राद्याः सर्वदेवताः । अहो तत्र स्थितो योगी न हर्षम् उपगच्छति ॥}
\vtranslit{ят падам препсаво дінах шакрадях сарвадеватах | ахо татра стхіто йоґі на харшам упаґаччаті ||}
\vperek{О диво: те, до чого прагнуть убогі всі боги, починаючи з Шакри, — йог, що вже стоїть у цьому, не відчуває захоплення.}
\vcom{Построкове зіставлення: редакторська чернетка, не рецензовано.}
\vcom{Рядок а}
\vdeva{यत् पदं प्रेप्सवो दीनाः शक्राद्याः सर्वदेवताः}
\vtranslit{ят падам препсаво дінах шакрадях сарвадеватах}
\vperek{О диво: те, до чого прагнуть убогі всі боги, починаючи з Шакри, —}
\vcom{Форми yat padaṃ і prepsavaḥ описують стан, до якого прагнуть; dīnāḥ застосовано навіть до śakrādyāḥ sarvadevatāḥ. Український вигук «О диво» перенесено на початок усього речення з наступного санскритського піврядка — це переставлення для плавного читання, не дослівний порядок.}
\vcom{Рядок б}
\vdeva{अहो तत्र स्थितो योगी न हर्षम् उपगच्छति}
\vtranslit{ахо татра стхіто йоґі на харшам упаґаччаті}
\vperek{йог, що вже стоїть у цьому, не відчуває захоплення.}
\vcom{У другому піврядку aho виражає подив, tatra sthito yogī називає йога, який уже перебуває в цьому стані, na harṣam upagacchati заперечує захоплення. Порівняння прагнення та незворушності — змістова структура вірша, а не оцінка реальних релігійних громад.}
\vcom{Слово dīnāḥ («убогі») вжите про богів навмисно: убогість тут не в майні, а в прагненні — хто прагне, той має брак. Дезидератив prepsavaḥ («ті, що хочуть досягти») підкреслює саме бажання досягти, а не досягнення. Йог описаний дієприкметником sthitaḥ («той, хто стоїть») — це стан, а не подія; і саме тому harṣa («захоплення») не виникає: захоплення потребує новизни, а тут нічого нового не сталося. Вигук aho той самий, що в серії «див» глави 3, і тут він позначає межу між прагненням і перебуванням.}
\vpar{Українська паралель (наша інтерпретація)}{Григорій Сковорода, «Вдячний Еродій» (трактат).}{Наша паралель: у Сковороди прагнення саме по собі теж викриває брак, і тому досягнуте не викликає захоплення. Наша інтерпретація за будовою думки, а не вказівка на вплив.}
\vfig{ilustratsii/4-2-balance.pdf}{Дві шальки: прагнення богів і спокій того, хто вже стоїть}{Схема наша, не історична ілюстрація.}

\vsec{4.3}
\vdeva{तज्ज्ञस्य पुण्यपापाभ्यां स्पर्शो ह्य् अन्तर् न जायते । न ह्य् आकाशस्य धूमेन दृश्यमाना ऽपि संगतिः ॥}
\vtranslit{таджджняся пуняпапабгям спаршо хй антар на джаяте | на хй акашася дгумена дрішямана ’пі самґатіх ||}
\vperek{Того, хто осягнув Істину, не зачіпають ані заслуги, ані провини. Так само й небо не зливається з димом, хоча на око вони ніби поєднані.}
\vcom{Построкове зіставлення: редакторська чернетка, не рецензовано.}
\vcom{Рядок а}
\vdeva{तज्ज्ञस्य पुण्यपापाभ्यां स्पर्शो ह्य् अन्तर् न जायते}
\vtranslit{таджджняся пуняпапабгям спаршо хй антар на джаяте}
\vperek{Того, хто осягнув Істину, не зачіпають ані заслуги, ані провини.}
\vcom{Родовий tajjñasya («того, хто знає Те») й подвійна форма puṇyapāpābhyāṃ у конструкції sparśo hy antar na jāyate пояснюють заперечений дотик заслуг і провин. «Істина» — чинний перекладацький відповідник, а не доведене буквальне значення займенника tat поза контекстом.}
\vcom{Рядок б}
\vdeva{न ह्य् आकाशस्य धूमेन दृश्यमाना ऽपि संगतिः}
\vtranslit{на хй акашася дгумена дрішямана ’пі самґатіх}
\vperek{Так само й небо не зливається з димом, хоча на око вони ніби поєднані.}
\vcom{Порівняння na hy ākāśasya dhūmena та dṛśyamānā 'pi saṃgatiḥ розрізняє видиму суміжність диму й неба та їхню сутнісну єдність. Це аналогія самого вірша; вона не є твердженням про фізичну природу атмосфери.}
\vcom{У вірші 4.3 використано порівняння неба і диму, щоби розрізнити видимий зв'язок та сутнісне ототожнення. Зовнішнє спостереження диму на тлі неба ще не означає, що небо набуває властивостей диму. Аналогічно фігура «того, хто знає Те» в недвоїстому вченні не визначається заслугами чи провинами як власною сутністю. З цього не випливає юридичний або етичний дозвіл чинити шкоду: вірш формулює філософську тезу про природу свідка, а не правила поведінки. У робочому перекладі «осягнув Істину» — тлумачна експлікація займенника tat, можлива буквальніша редакція — «той, хто знає Те»; на рецензії слід обрати один варіант. Потрібно окремо звірити різницю «з'єднання» / «дотик» і мовну форму фінального іменника saṅgati / saṃgatiḥ. **Атрибуція мовця не усталена за звіреними електронними свідками:** GRETIL/TextGrid перед 4.1 подає «aṣṭāvakra uvāca», а SARIT в заголовку глави 4 — «janaka uvāca». Поки не залучено незалежне друковане видання чи рукописні свідчення, не слід беззастережно приписувати твердження Аштавакрі або Джанаці. Два електронні свідки не є критичним виданням; коментарі сучасних авторів тут не передруковано.}
\vpar{Українська паралель (наша інтерпретація)}{Памфіл Юркевич, «Серце і його значення в духовному житті людини» (1860).}{Наша паралель: коли вірш каже, що небо не зливається з димом, він відділяє те, що проходить, від того, на чому воно проходить. Юркевич розводить серце й привхідне в людині так само рішуче. Наше порівняння — за способом розрізняти.}
\vfig{ilustratsii/4-3-space.pdf}{Небо й дим: видно разом, а не зливаються}{Схема наша, не історична ілюстрація.}
\vfig{ilustratsii-naratyvni/4-3-dym-i-nebo.pdf}{Дим від вогнища проходить крізь небо, але не змішується з ним}{Наша ілюстрація до вірша 4.3 — не історичне зображення.}

\vsec{4.4}
\vdeva{आत्मैवेदं जगत् सर्वं ज्ञातं येन महात्मना । यदृच्छया वर्तमानं तं निषेद्धुं क्षमेत कः ॥}
\vtranslit{атмайведам джаґат сарвам джнятам єна махатмана | ядріччая вартаманам там нішеддгум кшамета ках ||}
\vperek{Хто здатен стримати того, хто пізнав увесь світ як єдиний Атман і живе так, як складаються обставини?}
\vcom{Построкове зіставлення: редакторська чернетка, не рецензовано.}
\vcom{Рядок а}
\vdeva{आत्मैवेदं जगत् सर्वं ज्ञातं येन महात्मना}
\vtranslit{атмайведам джаґат сарвам джнятам єна махатмана}
\vperek{Хто здатен стримати того, хто пізнав увесь світ як єдиний Атман}
\vcom{Конструкція ātmaivedaṃ jagat sarvaṃ jñātaṃ yena mahātmanā називає пізнання всього світу як Атмана. Чинний український переклад починається з питання «Хто здатен стримати...», яке в санскриті завершальне: український поділ передає смислові половини, але не дзеркальний порядок слів.}
\vcom{Рядок б}
\vdeva{यदृच्छया वर्तमानं तं निषेद्धुं क्षमेत कः}
\vtranslit{ядріччая вартаманам там нішеддгум кшамета ках}
\vperek{і живе так, як складаються обставини?}
\vcom{yadṛcchayā vartamānaṃ описує життя відповідно до обставин, а taṃ niṣeddhuṃ kṣameta kaḥ формулює риторичне питання, хто може стримати цю людину. Тут не виводиться дозвіл нехтувати правами інших чи відповідальністю.}
\vcom{Вірш поєднує дві теми: пізнання недвоїстої природи світу та свободу від нав'язаної поведінкової схеми. Питання «хто здатен стримати?» є риторичним і характеризує свободу того, хто усвідомив світ як Атман. Слово yadṛcchayā не слід автоматично прирівнювати до примхливого бажання або безвідповідальності: у цьому контексті воно допускає читання «за збігом обставин, природним чином». Українське «єдиний Атман» — смислове відтворення підсилювального eva, а не окремий прикметник оригіналу; перед остаточним виданням перевірити, чи краще писати «сам Атман». Важливо, що vartamānaṃ taṃ граматично підтримує тлумачення про людину, яка живе спонтанно, а не тезу, ніби «світ робить усе, що йому хочеться». Цей коментар — редакторська інтерпретація, не незалежна наукова рецензія. Додатковий текстологічний блокер: GRETIL і SARIT не збігаються в атрибуції мовця всієї четвертої глави, тому цей вірш не приписано беззастережно жодному співрозмовникові.}
\vpar{Українська паралель (наша інтерпретація)}{Михайло Коцюбинський, «Intermezzo» (1908).}{Наша паралель: у Коцюбинського той, хто бачить ціле, теж не піддається керуванню ззовні. Наше порівняння — типологічне.}
\vfig{ilustratsii/4-4-witness.pdf}{Центр, якого не втримати: він бачить усе як єдине}{Схема наша, не історична ілюстрація.}

\vsec{4.5}
\vdeva{आब्रह्मस्तम्बपर्यन्ते भूतग्रामे चतुर्विधे । विज्ञस्यैव हि सामर्थ्यं इच्छा ऽनिच्छाविवर्जने ॥}
\vtranslit{абрахмастамбапар'янте бгутаґраме чатурвідге | віджнясяйва хі самартхям ічча ’ніччавіварджане ||}
\vperek{Серед чотирьох родів істот — від Брахми до стеблини трави — лише той, хто пізнав, спроможний полишити і бажання, і небажання.}
\vcom{Построкове зіставлення: редакторська чернетка, не рецензовано.}
\vcom{Рядок а}
\vdeva{आब्रह्मस्तम्बपर्यन्ते भूतग्रामे चतुर्विधे}
\vtranslit{абрахмастамбапар'янте бгутаґраме чатурвідге}
\vperek{Серед чотирьох родів істот — від Брахми до стеблини трави —}
\vcom{ābrahmastambaparyante bhūtagrāme caturvidhe задає межі сукупності істот та згадує чотири роди. Образ «від Брахми до стеблини» передає контраст міфологічного найвищого й звичайної рослини; докладна класифікація потребує окремого текстологічного підтвердження.}
\vcom{Рядок б}
\vdeva{विज्ञस्यैव हि सामर्थ्यं इच्छा ऽनिच्छाविवर्जने}
\vtranslit{віджнясяйва хі самартхям ічча ’ніччавіварджане}
\vperek{лише той, хто пізнав, спроможний полишити і бажання, і небажання.}
\vcom{vijñasyaiva hi sāmarthyaṃ пов'язує спроможність із тим, хто знає; icchā 'nicchāvivarjane позначає полишення бажання й небажання. Український «лише» зберігає eva, але не проголошує знання соціальним привілеєм.}
\vcom{Вірш протиставляє всеосяжний образ світу істот — від найвищого міфологічного персонажа до найскромнішої рослини — стану знання. Ключовий вислів vijñasyaiva hi sāmarthyam указує не на соціальний ранг, а на спроможність «того, хто знає». Пара icchā / anicchā охоплює внутрішні схильності до прийняття й відкидання; переклад «бажання й небажання» ближчий до буквального заперечного словотвору, тоді як «бажання й відраза» є сильнішою філософською інтерпретацією, яку варто перевірити в незалежному коментарі. «Чотири роди» не слід подавати як сучасну біологічну систематику: вірш використовує традиційну космологічну категоризацію, але не перелічує її членів. У GRETIL/TextGrid перед 4.1 позначено «aṣṭāvakra uvāca», тоді як у SARIT заголовок четвертої глави має «janaka uvāca». Через невирішену суперечність заголовків коментар не приписує ці слова певному мовцеві. Це власне робоче пояснення за двома електронними транскрипціями, не критичне видання й не незалежна фахова рецензія.}
\vpar{Українська паралель (наша інтерпретація)}{Іван Франко, «Мойсей» (1905).}{Наша паралель: у Франка вибір, що не зводиться до бажання й небажання, теж постає як те, на що здатний не кожен. Наше порівняння — за будовою думки, не про вплив.}
\vfig{ilustratsii/4-5-neti.pdf}{Бажання й небажання відкинуто разом}{Схема наша, не історична ілюстрація.}

\vsec{4.6}
\vdeva{आत्मानम् अद्वयं कश्चिज् जानाति जगद् ईश्वरं । यद् वेत्ति तत् स कुरुते न भयं तस्य कुत्रचित् ॥}
\vtranslit{атманам адваям кашчідж джанаті джаґад ішварам | яд ветті тат са куруте на бгаям тася кутрачіт ||}
\vperek{Хтось знає Атман — недвоїстий, володаря світу. Що він знає, те й робить, — і страху в нього немає ніде.}
\vcom{Построкове зіставлення: редакторська чернетка, не рецензовано.}
\vcom{Рядок а}
\vdeva{आत्मानम् अद्वयं कश्चिज् जानाति जगद् ईश्वरं}
\vtranslit{атманам адваям кашчідж джанаті джаґад ішварам}
\vperek{Хтось знає Атман — недвоїстий, володаря світу.}
\vcom{ātmānam advayaṃ ... jagad īśvaraṃ є об'єктом пізнання; kaścij jānāti називає неназваного «когось». У словосполученні про володаря світу збережено чинну редакцію без привнесення нових богословських атрибуцій.}
\vcom{Рядок б}
\vdeva{यद् वेत्ति तत् स कुरुते न भयं तस्य कुत्रचित्}
\vtranslit{яд ветті тат са куруте на бгаям тася кутрачіт}
\vperek{Що він знає, те й робить, — і страху в нього немає ніде.}
\vcom{yad vetti tat sa kurute зіставляє знання і дію; na bhayaṃ tasya kutracit заперечує страх будь-де. Це формула філософського тексту, не перевірене твердження про психологічну безстрашність конкретної людини.}
\vcom{kaścit («хтось») на початку — не випадкове слово: вірш не обіцяє, що це стосується кожного, і не називає його. Друга половина дає формулу yad vetti tat kurute («що знає, те й робить»), і вона сильніша за звичне «чини за знанням»: тут знання й дія не узгоджуються, а ототожнюються. kutracit з na («ніде») закриває простір повністю — страх не має місця, а не просто відсутній тут. Це останній вірш глави, і він повертає до її початку: 4.1 протиставляв того, хто грається, тому, хто тягне; тут сказано, чому — бо немає страху, а отже немає й потреби тягнути.}
\vpar{Українська паралель (наша інтерпретація)}{Василь Стус, «Час творчості».}{Наша паралель: у Стуса теж є стан, у якому людина робить те, що знає, і не має страху ніде. Наша асоціація, а не твердження про джерело.}
\vfig{ilustratsii/4-6-witness.pdf}{Той, хто знає недвоїсте, робить те, що знає}{Схема наша, не історична ілюстрація.}


% ЗГЕНЕРОВАНО scripts/build_format_v2_chapter.py — не редагувати вручну.
% Джерело: editorial/cards/5.*.json, format-v2/rozdil-5.json, paraleli.json.

\naukchapter{5}{4}{Глава 5. Розчинення}
\naukintro{Пʼята глава — чотири вірші й один рефрен. Рефрен тут не прикраса: він стоїть у кінці кожної строфи тими самими словами, і саме повтор і є думкою. Глава не описує розчинення, а показує, що робити з ним, коли відкидати вже нічого: у першому вірші це сказано прямо — у чистого немає звʼязку ні з чим, тож і відкидати йому нічого.}
\vsec{5.1}
\vdeva{न ते सङ्गो ऽस्ति केनापि किं शुद्धस् त्यक्तुम् इच्छसि । सङ्घातविलयं कुर्वन्न् एवम् एव लयं व्रज ॥}
\vtranslit{на те санґо ’сті кенапі кім шуддгас тяктум іччасі | санґгатавілаям курванн евам ева лаям враджа ||}
\vperek{У тебе немає зв'язку ні з чим — то чого ж ти, чистий, хочеш відкинути? Здійснюючи розчинення скупчення, так само йди до розчинення.}
\vcom{Построкове зіставлення: редакторська чернетка, не рецензовано.}
\vcom{Рядок а}
\vdeva{न ते सङ्गो ऽस्ति केनापि किं शुद्धस् त्यक्तुम् इच्छसि}
\vtranslit{на те санґо ’сті кенапі кім шуддгас тяктум іччасі}
\vperek{У тебе немає зв'язку ні з чим — то чого ж ти, чистий, хочеш відкинути?}
\vcom{na ... saṅgo 'sti заперечує прив'язаність; kenāpi вказує на відсутність зв'язку з будь-чим, kiṃ ... icchasi формує питання, tyaktum — інфінітив «відкинути». Тлумачення стосується внутрішньої свободи, а не наказу полишити родину чи працю.}
\vcom{Рядок б}
\vdeva{सङ्घातविलयं कुर्वन्न् एवम् एव लयं व्रज}
\vtranslit{санґгатавілаям курванн евам ева лаям враджа}
\vperek{Здійснюючи розчинення скупчення, так само йди до розчинення.}
\vcom{saṅghātavilayaṃ — розчинення сукупності, kurvann — дієприкметник «здійснюючи», vraja — наказова форма «йди». evam eva layaṃ vraja — рефрен усіх чотирьох віршів глави 5; розрізнення vilaya / laya як філософських термінів потребує незалежної фахової звірки.}
\vcom{Питання kim ... icchasi («чого хочеш») викриває безглуздість самого наміру відректися: відкидають те, до чого прив'язані, а зв'язку немає (na te saṅgaḥ asti kenāpi). Слово saṅghāta («скупчення») — технічне: тіло описане не як річ, а як агрегат, суму. Тому vilaya («розчинення») і laya («розчинення») стоять в одному вірші двічі: перше — те, що відбувається з тілом само собою, друге — те, до чого слід іти. Різниця в тому, що перше трапляється, а друге обирають. РЕФРЕН ГЛАВИ. Усі чотири вірші глави 5 закінчуються однаково: evam eva layaṃ vraja («так само йди до розчинення»). Це не чотири окремі повчання, а одна літанія з незмінним приспівом: змінюється тільки перша половина — обґрунтування, а висновок лишається той самий. Тому главу слід читати як одне ціле.}
\vpar{Українська паралель (наша інтерпретація)}{Григорій Сковорода, «Асхань. Симфонія, наречена книга Асхань, про пізнання самого себе».}{Наша паралель: у Сковороди пізнання себе теж нічого не додає й нічого не відкидає — воно знімає перешкоду, і тільки. Наша інтерпретація, не твердження про вплив.}
\vfig{ilustratsii/5-1-neti.pdf}{Відкидати нічого: у чистого немає звʼязку}{Схема наша, не історична ілюстрація.}

\vsec{5.2}
\vdeva{उदेति भवतो विश्वं वारिधेर् इव बुद्बुदः । इति ज्ञात्वा एकम् आत्मानम् एवम् एव लयं व्रज ॥}
\vtranslit{удеті бгавато вішвам варідгер іва будбудах | іті джнятва екам атманам евам ева лаям враджа ||}
\vperek{Світ постає з тебе, як бульбашка з океану. Знаючи, що Атман єдиний, — так само йди до розчинення.}
\vcom{Построкове зіставлення: редакторська чернетка, не рецензовано.}
\vcom{Рядок а}
\vdeva{उदेति भवतो विश्वं वारिधेर् इव बुद्बुदः}
\vtranslit{удеті бгавато вішвам варідгер іва будбудах}
\vperek{Світ постає з тебе, як бульбашка з океану.}
\vcom{udeti означає «постає»; bhavataḥ допускає значення «з тебе» або «твого», vāridher iva budbudaḥ порівнює світ із бульбашкою в океані. Це образ залежного виникнення, а не самостійний доказ неіснування світу.}
\vcom{Рядок б}
\vdeva{इति ज्ञात्वा एकम् आत्मानम् एवम् एव लयं व्रज}
\vtranslit{іті джнятва екам атманам евам ева лаям враджа}
\vperek{Знаючи, що Атман єдиний, — так само йди до розчинення.}
\vcom{iti jñātvā — «так знаючи», ekam ātmānam стосується єдності Атмана. evam eva layaṃ vraja — наказовий рефрен усіх чотирьох віршів цієї глави; український відповідник збережено без змін.}
\vcom{Порівняння з бульбашкою точніше за звичне «хвиля в морі» (2.10, 2.23–2.25): хвиля є форма тієї ж води, а бульбашка — порожнеча, обтягнута плівкою з тієї ж води. Вона не має власного вмісту. bhavataḥ стоїть в ablativus/genitivus, і обидва прочитання працюють: «з тебе» (джерело) і «твого» (належність). Висновок вірша не «світ нереальний», а «Атман єдиний» — і саме з єдності випливає розчинення, а не навпаки. РЕФРЕН ГЛАВИ. Усі чотири вірші глави 5 закінчуються однаково: evam eva layaṃ vraja («так само йди до розчинення»). Це не чотири окремі повчання, а одна літанія з незмінним приспівом: змінюється тільки перша половина — обґрунтування, а висновок лишається той самий. Тому главу слід читати як одне ціле.}
\vpar{Українська паралель (наша інтерпретація)}{Богдан-Ігор Антонич, «Ротації» (1938).}{Наша паралель: в Антонича світ теж постає з того, хто на нього дивиться, як бульбашка з океану. Наше порівняння — за образом.}
\vfig{ilustratsii/5-2-space.pdf}{Бульбашка й океан: світ постає з тебе}{Схема наша, не історична ілюстрація.}
\vfig{ilustratsii-naratyvni/5-2-bulbashka-z-okeanu.pdf}{Людина сидить на березі, а з океану одна за одною піднімаються бульбашки}{Наша ілюстрація до вірша 5.2 — не історичне зображення.}

\vsec{5.3}
\vdeva{प्रत्यक्षम् अप्य् अवस्तुत्वाद् विश्वं नास्त्य् अमले त्वयि । रज्जुसर्प इव व्यक्तम् एवम् एव लयं व्रज ॥}
\vtranslit{пратякшам апй авастутвад вішвам настй амале тваї | раджджусарпа іва в'яктам евам ева лаям враджа ||}
\vperek{Хоч світ і перед очима, він не існує в тобі, незаплямованому, — бо він не річ, як змій у мотузці, явно. Так само йди до розчинення.}
\vcom{Построкове зіставлення: редакторська чернетка, не рецензовано.}
\vcom{Рядок а}
\vdeva{प्रत्यक्षम् अप्य् अवस्तुत्वाद् विश्वं नास्त्य् अमले त्वयि}
\vtranslit{пратякшам апй авастутвад вішвам настй амале тваї}
\vperek{Хоч світ і перед очима, він не існує в тобі, незаплямованому, — бо він не річ,}
\vcom{pratyakṣam api означає «хоч і видимий перед очима»; avastutvād — «через відсутність самостійної речовості/сутності», viśvaṃ nāsti — «світ не існує»; amale tvayi звертається до незаплямованого адресата. Це відтворення редакторської картки, а не нове критичне читання.}
\vcom{Рядок б}
\vdeva{रज्जुसर्प इव व्यक्तम् एवम् एव लयं व्रज}
\vtranslit{раджджусарпа іва в'яктам евам ева лаям враджа}
\vperek{як змій у мотузці, явно. Так само йди до розчинення.}
\vcom{rajjusarpa iva — порівняння з уявним змієм у мотузці, vyaktam — «виразно, явно». evam eva layaṃ vraja зберігає незмінний рефрен глави; філософське пояснення образу слід відрізняти від лексичного значення санскриту.}
\vcom{Найсильніше тут — avastutvāt («через несуттєвість»). Це не «світ оманливий», а точніше: світ не є vastu, річчю, самостійною сутністю. pratyakṣam api («хоч і перед очима») визнає видимість повністю — вірш не заперечує, що ми бачимо. Заперечується лише те, що видиме має власне буття. Змій у мотузці (той самий приклад, що в 2.9) поданий тут з vyaktam («явно»): помилка не тьмяна, вона гранично чітка — і саме тому переконлива. РЕФРЕН ГЛАВИ. Усі чотири вірші глави 5 закінчуються однаково: evam eva layaṃ vraja («так само йди до розчинення»). Це не чотири окремі повчання, а одна літанія з незмінним приспівом: змінюється тільки перша половина — обґрунтування, а висновок лишається той самий. Тому главу слід читати як одне ціле.}
\vpar{Українська паралель (наша інтерпретація)}{Іван Франко, «Сойчине крило» (1905).}{Наша паралель: у Франка видиме й справжнє теж не збігаються, і саме на цьому розходженні тримається оповідь. Наше порівняння — типологічне, не про джерело.}
\vfig{ilustratsii/5-3-illusion.pdf}{Світ перед очима й змій у мотузці}{Схема наша, не історична ілюстрація.}

\vsec{5.4}
\vdeva{समदुःखसुखः पूर्ण आशानैराश्ययोः समः । समजीवितमृत्युः सन्न् एवम् एव लयं व्रज ॥}
\vtranslit{самадухкхасукхах пурна ашанайрашяйох самах | самаджівітамрітюх санн евам ева лаям враджа ||}
\vperek{Рівний у стражданні й радості, повний, рівний у сподіванні й безнадії, рівний у житті й смерті, — так само йди до розчинення.}
\vcom{Построкове зіставлення: редакторська чернетка, не рецензовано.}
\vcom{Рядок а}
\vdeva{समदुःखसुखः पूर्ण आशानैराश्ययोः समः}
\vtranslit{самадухкхасукхах пурна ашанайрашяйох самах}
\vperek{Рівний у стражданні й радості, повний, рівний у сподіванні й безнадії,}
\vcom{samaduḥkhasukhaḥ — рівність до страждання й радості; pūrṇaḥ — «повний»; āśānairāśyayoḥ samaḥ — рівний у парі сподівання/безнадія. Два протиставлення містяться в одному друкованому піврядку; це не два самостійні переклади.}
\vcom{Рядок б}
\vdeva{समजीवितमृत्युः सन्न् एवम् एव लयं व्रज}
\vtranslit{самаджівітамрітюх санн евам ева лаям враджа}
\vperek{рівний у житті й смерті, — так само йди до розчинення.}
\vcom{samajīvitamṛtyuḥ — рівність у житті й смерті, san — дієприкметник «бувши». evam eva layaṃ vraja — той самий наказовий рефрен, що в 5.1–5.3; український порядок слів залишено без змін.}
\vcom{Вірш складається з трьох пар, і всі три побудовані на sama («рівний»): страждання/радість, сподівання/безнадія, життя/смерть. Пари охоплюють три виміри — відчуття, очікування, буття. Слово pūrṇaḥ («повний») стоїть між першою і другою парою не випадково: рівність не порожнеча, а повнота, якій не бракує нічого з обох боків. Це єдиний вірш глави, де немає ні прикладу, ні заперечення — лише перелік ознак. Тим самим літанія закінчується не аргументом, а описом того, хто йде. РЕФРЕН ГЛАВИ. Усі чотири вірші глави 5 закінчуються однаково: evam eva layaṃ vraja («так само йди до розчинення»). Це не чотири окремі повчання, а одна літанія з незмінним приспівом: змінюється тільки перша половина — обґрунтування, а висновок лишається той самий. Тому главу слід читати як одне ціле.}
\vpar{Українська паралель (наша інтерпретація)}{Леся Українка, «Камінний господар» (1912).}{Наша паралель: у Лесі Українки рівність перед останнім теж не залежить від надії чи її відсутності. Наше порівняння — за будовою думки.}
\vfig{ilustratsii/5-4-balance.pdf}{Дві шальки, що стоять рівно: життя й смерть}{Схема наша, не історична ілюстрація.}


% ЗГЕНЕРОВАНО scripts/build_format_v2_chapter.py — не редагувати вручну.
% Джерело: editorial/cards/6.*.json, format-v2/rozdil-6.json, paraleli.json.

\naukchapter{6}{4}{Глава 6. Горщик і простір}
\naukintro{Шоста глава тримається на рефрені, довшому за рефрен пʼятої: чотири вірші закінчуються тими самими словами про відкидання, взяття й розчинення. Кожна строфа дає новий образ — простір і горщик, океан і хвиля, перламутр і срібло, — і після кожного повторює, що з цим знанням нічого не треба робити. Це глава про те, що знання не є дією.}
\vsec{6.1}
\vdeva{आकाशवद् अनन्तो ऽहं घटवत् प्रकृतं जगत् । इति ज्ञानं तथैतस्य न त्यागो न ग्रहो लयः ॥}
\vtranslit{акашавад ананто ’хам ґгатават пракрітам джаґат | іті джнянам татхайтася на тяґо на ґрахо лаях ||}
\vperek{Я безмежний, як простір, а світ — як горщик. Таке знання — і в нього немає ні відкидання, ні взяття, ні розчинення.}
\vcom{Построкове зіставлення: редакторська чернетка, не рецензовано.}
\vcom{Рядок а}
\vdeva{आकाशवद् अनन्तो ऽहं घटवत् प्रकृतं जगत्}
\vtranslit{акашавад ананто ’хам ґгатават пракрітам джаґат}
\vperek{Я безмежний, як простір, а світ — як горщик.}
\vcom{ākāśavad і ghaṭavat — порівняння «як простір» та «як горщик». ananto 'ham стосується безмежного «я», тоді як prakṛtaṃ jagat — створеного світу. Це метафоричне зіставлення, а не буквальний опис простору.}
\vcom{Рядок б}
\vdeva{इति ज्ञानं तथैतस्य न त्यागो न ग्रहो लयः}
\vtranslit{іті джнянам татхайтася на тяґо на ґрахо лаях}
\vperek{Таке знання — і в нього немає ні відкидання, ні взяття, ні розчинення.}
\vcom{iti jñānaṃ — «таке знання»; na tyāgo na graho layaḥ заперечує відкидання, взяття і розчинення. Сенс рефрену — відсутність потреби щось робити з таким знанням, а не вимога знищити світ.}
\vcom{Образ простору й горщика працює точніше, ніж здається: простір не вміщується в горщик і не обмежується його стінками — горщик лише окреслює частину простору, не додаючи й не забираючи нічого. Тому рефрен і каже: немає ні tyāga (відкидання), ні graha (взяття), ні навіть laya (розчинення) — усі три дії припускають, що щось справді додалося або зникло. Слово prakṛtam («зроблене») підкреслює: горщик — виріб, а простір — ні. РЕФРЕН ГЛАВИ. Усі чотири вірші глави 6 закінчуються однаково: iti jñānaṃ tathaitasya na tyāgo na graho layaḥ («таке знання — і в нього немає ні відкидання, ні взяття, ні розчинення»). Зверни увагу: рефрен глави 5 вимагав laya («йди до розчинення»), а рефрен глави 6 заперечує навіть laya. Це не суперечність, а наступний крок: спершу розчинення як те, до чого йдуть, потім — знання, у якому немає кому йти й нічого розчиняти. Главу слід читати як одне ціле.}
\vpar{Українська паралель (ПОТРЕБУЄ ПЕРЕВІРКИ)}{українська філософська традиція, не визначено.}{ПОТРЕБУЄ ПЕРЕВІРКИ: образ простору, що не ділиться стінками посуду, найчастіше приходить до українського читача через Сковороду; саме місце, де він говорить про посуд і простір, ми за джерелом не звірили й тому не називаємо.}
\vfig{ilustratsii/6-1-space.pdf}{Простір і горщик: стінки ділять посуд, а не простір}{Схема наша, не історична ілюстрація.}

\vsec{6.2}
\vdeva{महोदधिर् इवाहं स प्रपञ्चो वीचिसंनिभः । इति ज्ञानं तथैतस्य न त्यागो न ग्रहो लयः ॥}
\vtranslit{маходадгір івахам са прапаньчо вічісамнібгах | іті джнянам татхайтася на тяґо на ґрахо лаях ||}
\vperek{Я — як великий океан, а цей плин світу — як хвилі. Таке знання — і в нього немає ні відкидання, ні взяття, ні розчинення.}
\vcom{Построкове зіставлення: редакторська чернетка, не рецензовано.}
\vcom{Рядок а}
\vdeva{महोदधिर् इवाहं स प्रपञ्चो वीचिसंनिभः}
\vtranslit{маходадгір івахам са прапаньчо вічісамнібгах}
\vperek{Я — як великий океан, а цей плин світу — як хвилі.}
\vcom{mahodadhir ivāhaṃ — «я, мов великий океан»; prapañcaḥ — розгортання світу; vīcisaṃnibhaḥ — «подібний до хвиль». Хвиля тут є образом явищ, а не незалежною від океану сутністю.}
\vcom{Рядок б}
\vdeva{इति ज्ञानं तथैतस्य न त्यागो न ग्रहो लयः}
\vtranslit{іті джнянам татхайтася на тяґо на ґрахо лаях}
\vperek{Таке знання — і в нього немає ні відкидання, ні взяття, ні розчинення.}
\vcom{Спільний рефрен iti jñānaṃ ... na tyāgo na graho layaḥ заперечує потребу брати, відкидати або розчиняти явища. Порівняння океану з хвилями не є твердженням про фізичну будову світу.}
\vcom{prapañca — важке слово: буквально «розгортання, розпросторення», а в ужитку «різноманіття світу» з відтінком заплутаності. Тут воно вжите без засудження: хвилі не є вадою океану. Порівняй 2.23–2.25, де той самий образ розгорнутий у три кроки; тут він стиснутий до одного порівняння, бо мета інша — не показати, як хвилі виникають, а показати, що з ними не треба нічого робити. РЕФРЕН ГЛАВИ. Усі чотири вірші глави 6 закінчуються однаково: iti jñānaṃ tathaitasya na tyāgo na graho layaḥ («таке знання — і в нього немає ні відкидання, ні взяття, ні розчинення»). Зверни увагу: рефрен глави 5 вимагав laya («йди до розчинення»), а рефрен глави 6 заперечує навіть laya. Це не суперечність, а наступний крок: спершу розчинення як те, до чого йдуть, потім — знання, у якому немає кому йти й нічого розчиняти. Главу слід читати як одне ціле.}
\vpar{Українська паралель (наша інтерпретація)}{Григорій Сковорода, «Вдячний Еродій» (трактат).}{Наша паралель: у Сковороді море як образ цілого, у якому частка не має власного окремого життя, близьке до цього вірша. Наше порівняння — типологічне, не твердження про джерело.}
\vfig{ilustratsii/6-2-motion.pdf}{Хвилі й океан: у знанні немає ні взяття, ні відкидання}{Схема наша, не історична ілюстрація.}

\vsec{6.3}
\vdeva{अहं स शुक्तिसङ्काशो रूप्यवद् विश्वकल्पना । इति ज्ञानं तथैतस्य न त्यागो न ग्रहो लयः ॥}
\vtranslit{ахам са шуктісанкашо руп'явад вішвакалпана | іті джнянам татхайтася на тяґо на ґрахо лаях ||}
\vperek{Я — як перламутр, а уявлення світу — як срібло в ньому. Таке знання — і в нього немає ні відкидання, ні взяття, ні розчинення.}
\vcom{Построкове зіставлення: редакторська чернетка, не рецензовано.}
\vcom{Рядок а}
\vdeva{अहं स शुक्तिसङ्काशो रूप्यवद् विश्वकल्पना}
\vtranslit{ахам са шуктісанкашо руп'явад вішвакалпана}
\vperek{Я — як перламутр, а уявлення світу — як срібло в ньому.}
\vcom{śuktisaṅkāśaḥ — «подібний до перламутру»; rūpyavat — «як срібло»; viśvakalpanā — уявлення світу. Порівняння показує різницю між тим, що бачать, і тим, що помилково приписують побаченому.}
\vcom{Рядок б}
\vdeva{इति ज्ञानं तथैतस्य न त्यागो न ग्रहो लयः}
\vtranslit{іті джнянам татхайтася на тяґо на ґрахо лаях}
\vperek{Таке знання — і в нього немає ні відкидання, ні взяття, ні розчинення.}
\vcom{Рефрен na tyāgo na graho layaḥ заперечує відкидання, взяття й розчинення: у метафорі не потрібно знищувати неіснуюче «срібло». Це пояснювальна аналогія хибного сприйняття, не рукописний доказ.}
\vcom{Третє й останнє з трьох порівнянь глави (простір/горщик, океан/хвилі, перламутр/срібло), і воно єдине, де друга річ не існує зовсім: горщик і хвиля є, срібла немає. Тому тут і стоїть kalpanā («уявлення, конструкція») — не «світ», а саме уявлення світу. Зверни увагу: перламутр не заперечується і не змінюється — зникає тільки срібло, а разом із ним і питання, що з ним робити. Це і є сенс рефрену. РЕФРЕН ГЛАВИ. Усі чотири вірші глави 6 закінчуються однаково: iti jñānaṃ tathaitasya na tyāgo na graho layaḥ («таке знання — і в нього немає ні відкидання, ні взяття, ні розчинення»). Зверни увагу: рефрен глави 5 вимагав laya («йди до розчинення»), а рефрен глави 6 заперечує навіть laya. Це не суперечність, а наступний крок: спершу розчинення як те, до чого йдуть, потім — знання, у якому немає кому йти й нічого розчиняти. Главу слід читати як одне ціле.}
\vpar{Українська паралель (ПОТРЕБУЄ ПЕРЕВІРКИ)}{українська барокова емблематика, не визначено.}{ПОТРЕБУЄ ПЕРЕВІРКИ: пара перламутр / срібло як образ оманливого уявлення знайома бароковій емблематиці, але конкретного твору ми не перевірили.}
\vfig{ilustratsii/6-3-illusion.pdf}{Перламутр і срібло: уявлення світу}{Схема наша, не історична ілюстрація.}

\vsec{6.4}
\vdeva{अहं वा सर्वभूतेषु सर्वभूतान्य् अथो मयि । इति ज्ञानं तथैतस्य न त्यागो न ग्रहो लयः ॥}
\vtranslit{ахам ва сарвабгутешу сарвабгутанй атхо маї | іті джнянам татхайтася на тяґо на ґрахо лаях ||}
\vperek{Я — у всіх істотах, а всі істоти — у мені. Таке знання — і в нього немає ні відкидання, ні взяття, ні розчинення.}
\vcom{Построкове зіставлення: редакторська чернетка, не рецензовано.}
\vcom{Рядок а}
\vdeva{अहं वा सर्वभूतेषु सर्वभूतान्य् अथो मयि}
\vtranslit{ахам ва сарвабгутешу сарвабгутанй атхо маї}
\vperek{Я — у всіх істотах, а всі істоти — у мені.}
\vcom{sarvabhūteṣu — «у всіх істотах» (місцевий множини), sarvabhūtāni — «усі істоти», mayi — «у мені». atho поєднує два напрями взаємного перебування.}
\vcom{Рядок б}
\vdeva{इति ज्ञानं तथैतस्य न त्यागो न ग्रहो लयः}
\vtranslit{іті джнянам татхайтася на тяґо на ґрахо лаях}
\vperek{Таке знання — і в нього немає ні відкидання, ні взяття, ні розчинення.}
\vcom{Підсумкове iti jñānaṃ ... na tyāgo na graho layaḥ повторює заперечення відкидання, взяття й розчинення. Цей рефрен завершує главу після трьох образних порівнянь.}
\vcom{Двозмінна формула та сама, що в 3.5, і тут вона завершує главу. Після трьох порівнянь (простір, океан, перламутр) висновок поданий без образу взагалі — лише відношення. atho («а також») з'єднує два напрямки так, що вони не зливаються в один: я в істотах і істоти в мені — не те саме твердження, прочитане двічі, а два боки одного відношення. Саме тому рефрен знову заперечує tyāga і graha: у такому відношенні немає нічого, що можна було б узяти або відкинути. РЕФРЕН ГЛАВИ. Усі чотири вірші глави 6 закінчуються однаково: iti jñānaṃ tathaitasya na tyāgo na graho layaḥ («таке знання — і в нього немає ні відкидання, ні взяття, ні розчинення»). Зверни увагу: рефрен глави 5 вимагав laya («йди до розчинення»), а рефрен глави 6 заперечує навіть laya. Це не суперечність, а наступний крок: спершу розчинення як те, до чого йдуть, потім — знання, у якому немає кому йти й нічого розчиняти. Главу слід читати як одне ціле.}
\vpar{Українська паралель (наша інтерпретація)}{Павло Тичина, «Плуг» (1920).}{Наша паралель: у Тичини взаємне перебування всього в усьому теж стає способом бачити світ, а не твердженням про будову речей. Наше порівняння — за образом.}
\vfig{ilustratsii/6-4-witness.pdf}{Взаємне перебування: я в усіх істотах, усі істоти в мені}{Схема наша, не історична ілюстрація.}


% ЗГЕНЕРОВАНО scripts/build_format_v2_chapter.py — не редагувати вручну.
% Джерело: editorial/cards/7.*.json, format-v2/rozdil-7.json, paraleli.json.

\naukchapter{7}{5}{Глава 7. Океан, у якому пливе світ}
\naukintro{Сьома глава коротка й тримається на двох повторах. Перші три вірші починаються тими самими словами — у мені, безмежному океані, — і щоразу додають новий рух: човен, хвилю, домислювання. Останні два закінчуються словами саме в цьому я стою. Отже, глава не про те, що світ ілюзорний, а про те, де саме стоїть той, хто це бачить.}
\vsec{7.1}
\vdeva{मय्य् अनन्तमहाम्भोधौ विश्वपोत इतस्ततः । भ्रमति स्वान्तवातेन न ममास्त्य् असहिष्णुता ॥}
\vtranslit{майй анантамахамбгодгау вішвапота ітастатах | бграматі свантаватена на мамастй асахішнута ||}
\vperek{У мені, безмежному океані, човен світу блукає туди-сюди вітром власного нутра — і в мене немає нетерпимості.}
\vcom{`pota` — «човен, судно», а не «хвиля»: світ тут не хвилювання океану, а рухома річ **на** ньому. Це відрізняє 7.1 від 7.2, де світ уже `vīci` (хвиля). `svānta-vāta` («вітер власного нутра») — рушій усередині, не зовні. `asahiṣṇutā` — «нетерпимість, нездатність терпіти»: океан не страждає від човна, і саме це сказано, а не «спокій».}
\vpar{Українська паралель (наша інтерпретація)}{Олександр Довженко, «Поема про море» (1956).}{Наша паралель: у Довженка море теж стає мірою того, що з людиною діється, а сама людина лишається тим, у чому цей рух відбувається. Наше порівняння — за образом, не про джерело.}
\vfig{ilustratsii/7-1-motion.pdf}{Човен світу в безмежному океані, гнаний вітром нутра}{Схема наша, не історична ілюстрація.}

\vsec{7.2}
\vdeva{मय्य् अनन्तमहाम्भोधौ जगद्वीचिः स्वभावतः । उदेतु वास्तम् आयातु न मे वृद्धिर् न च क्षतिः ॥}
\vtranslit{майй анантамахамбгодгау джаґадвічіх свабгаватах | удету вастам аяту на ме вріддгір на ча кшатіх ||}
\vperek{У мені, безмежному океані, хвиля світу — за власною природою. Нехай піднімається або спадає: у мене немає ні приросту, ні втрати.}
\vcom{Два імперативи (`udetu`, `āyātu`) — це не побажання, а **дозвіл**: текст не просить хвилю поводитися інакше. `svabhāvataḥ` («за власною природою») знімає питання причини. Пара `vṛddhi` / `kṣati` — приріст і втрата: океан не стає більшим від хвилі й не меншим без неї. Разом із 7.1 це друга метафора того самого, і різниця між човном і хвилею не випадкова.}
\vpar{Українська паралель (наша інтерпретація)}{Богдан-Ігор Антонич, «Книга Лева» (1936).}{Наша паралель: в Антонича хвиля теж піднімається й спадає за власною природою, не питаючи, чи це комусь на приріст. Наше порівняння — типологічне.}
\vfig{ilustratsii/7-2-motion.pdf}{Хвиля світу: піднімається й спадає без приросту й втрати}{Схема наша, не історична ілюстрація.}

\vsec{7.3}
\vdeva{मय्य् अनन्तमहाम्भोधौ विश्वं नाम विकल्पना । अतिशान्तो निराकार एतद् एवाहम् आस्थितः ॥}
\vtranslit{майй анантамахамбгодгау вішвам нама вікалпана | атішанто ніракара етад евахам астхітах ||}
\vperek{У мені, безмежному океані те, що зветься світом, є домислюванням. Вельми втихлий, без форми — саме в цьому я стою.}
\vcom{`nāma` тут не «ім'я», а знецінювальна частка: «те, що **зветься** світом». `vikalpanā` — та сама «конструкція розуму», що в 2.9 і 2.17; словник карток має лишатися узгодженим. Формула `etad evāham āsthitaḥ` («саме в цьому я стою») повторюється в 7.3 і 7.4 — це не прикраса, а спосіб зчепити два вірші в одне твердження.}
\vpar{Українська паралель (ПОТРЕБУЄ ПЕРЕВІРКИ)}{українська містична проза, не визначено.}{ПОТРЕБУЄ ПЕРЕВІРКИ: думка, що світ є домислом, у українському матеріалі найближча до містичної прози, але назвати твір без звірки ми не можемо.}
\vfig{ilustratsii/7-3-neti.pdf}{Світ і форма відкинуті: лишається стояння в тиші}{Схема наша, не історична ілюстрація.}

\vsec{7.4}
\vdeva{नात्मा भावेषु नो भावस् तत्रानन्ते निरञ्जने । इत्य् असक्तो ऽस्पृहः शान्त एतद् एवाहम् आस्थितः ॥}
\vtranslit{натма бгавешу но бгавас татрананте ніраньджане | ітй асакто ’спріхах шанта етад евахам астхітах ||}
\vperek{Я не в станах буття, і стан — не там, у безмежному, незаплямленому. Так — неприв'язаний, без спраги, втихлий — саме в цьому я стою.}
\vcom{Вірш навмисно важкий: заперечуються обидві сторони — і `ātman` як один зі `bhāva`, і `bhāva` як те, що перебуває в безмежному. Це не заклик, а **усунення альтернативи**, і спрощувати його в перекладі не можна. `nirañjana` — той самий епітет, що в 1.15 і 2.15: незаплямленість тут не метафора чистоти, а вказівка на відсутність того, до чого можна приліпитися.}
\vpar{Українська паралель (наша інтерпретація)}{Іван Вишенський, «Писання» (кінець XVI — початок XVII ст.).}{Наша паралель: у Вишенського стан непривʼязаності й відсутності спраги теж описано як стояння, а не як дію. Наше порівняння — типологічне, не твердження про вплив.}
\vfig{ilustratsii/7-4-neti.pdf}{Стани буття не мої: непривʼязаний, без спраги, втихлий}{Схема наша, не історична ілюстрація.}

\vsec{7.5}
\vdeva{अहो चिन्मात्रम् एवाहम् इन्द्रजालोपमं जगत् । अतो मम कथं कुत्र हेयोपादेयकल्पना ॥}
\vtranslit{ахо чінматрам евахам індраджалопамам джаґат | ато мама катхам кутра хейопадеякалпана ||}
\vperek{О диво: я — саме лише свідомість, а світ подібний до чару. Тож звідки й де в мені домислювання про те, що відкинути і що прийняти?}
\vcom{`cinmātra` («саме лише свідомість») — ключове слово розділу; переклад не повинен додавати «чиста», бо `mātra` означає «не більше ніж». `indrajāla` — той самий образ чару, що в 10.2, і словник має бути спільним. Питання `kathaṃ kutra` подвійне — «як» і «де» — і саме подвійність робить його риторичним: для вибору не лишається ні способу, ні місця. Глава 7 на цьому закінчується.}
\vpar{Українська паралель (ПОТРЕБУЄ ПЕРЕВІРКИ)}{українське бароко, не визначено.}{ПОТРЕБУЄ ПЕРЕВІРКИ: питання, звідки береться вибір між прийняти й відкинути, близьке до барокової проповіді про марність вибору, але конкретного тексту ми не перевірили.}
\vfig{ilustratsii/7-5-illusion.pdf}{Світ як чар і свідомість, що його бачить}{Схема наша, не історична ілюстрація.}


% ЗГЕНЕРОВАНО scripts/build_format_v2_chapter.py — не редагувати вручну.
% Джерело: editorial/cards/8.*.json, format-v2/rozdil-8.json, paraleli.json.

\naukchapter{8}{4}{Глава 8. Полон і визволення — в розумі}
\naukintro{Восьма глава — чотири вірші, дві пари. Перша пара (8.1 і 8.2) побудована як дослівне заперечення: ті самі шість дієслів, у другому вірші кожне з часткою na. Друга пара міняє мірило: спершу привʼязаність до поглядів (8.3), далі саме слово я (8.4). Визволення тут не додає нічого — воно знімає рівно те, що було названо полоном.}
\vsec{8.1}
\vdeva{तदा बन्धो यदा चित्तं किञ्चिद् वाञ्छति शोचति । किञ्चिन् मुञ्चति गृह्णाति किञ्चिद् धृष्यति कुप्यति ॥}
\viast{tadā bandho yadā cittaṃ kiñcid vāñchati śocati | kiñcin muñcati gṛhṇāti kiñcid dhṛṣyati kupyati ||}
\vtranslit{тада бандго яда чіттам кіньчід ваньччаті шочаті | кіньчін муньчаті ґріхнаті кіньчід дгрішяті куп'яті ||}
\vperek{Полон — тоді, коли розум чогось бажає, сумує, щось відпускає або бере, чомусь радіє чи гнівається.}
\vcom{Построкове зіставлення: редакторська чернетка, не рецензовано.}
\vcom{Рядок а}
\vdeva{तदा बन्धो यदा चित्तं किञ्चिद् वाञ्छति शोचति}
\viast{tadā bandho yadā cittaṃ kiñcid vāñchati śocati}
\vtranslit{тада бандго яда чіттам кіньчід ваньччаті шочаті}
\vperek{Полон — тоді, коли розум чогось бажає, сумує,}
\vcom{Кореляція tadā ... yadā пояснює, коли постає полон. У першому піврядку розум бажає й сумує; це два з шести дієслів переліку. Переклад залишає речення відкритим до наступного рядка.}
\vcom{Рядок б}
\vdeva{किञ्चिन् मुञ्चति गृह्णाति किञ्चिद् धृष्यति कुप्यति}
\viast{kiñcin muñcati gṛhṇāti kiñcid dhṛṣyati kupyati}
\vtranslit{кіньчін муньчаті ґріхнаті кіньчід дгрішяті куп'яті}
\vperek{щось відпускає або бере, чомусь радіє чи гнівається.}
\vcom{Другий піврядок завершує перелік: розум відпускає, бере, радіє та гнівається. Форма dhṛṣyati в джерельній картці не має остаточного лексикографічного розв'язання; відповідник «радіє» залишено без самовільного виправлення санскриту.}
\vcom{Визначення полону тут не метафізичне, а **описове**: шість дієслів, і кожне — звичайний рух розуму. Зверніть увагу, що «бере» і «відпускає» стоять в одному ряду з гнівом: вибір сам по собі є ознакою полону. Перелік не оцінює, а перелічує — і саме тому діє.}
\vpar{Українська паралель (наша інтерпретація)}{Григорій Сковорода, «Сад божественних пісень».}{Наша паралель: у Сковороди буря пристрастей теж описана як рух усередині людини, а не як те, що приходить ззовні. Наше порівняння — за образом, не за джерелом.}
\vfig{ilustratsii/8-1-knot.pdf}{Вузол полону: шість рухів розуму}{Схема наша, не історична ілюстрація.}
\vfig{ilustratsii-naratyvni/8-1-polon.pdf}{Людина за ґратами тягнеться рукою до птаха, що сидить на волі}{Наша ілюстрація до вірша 8.1 — не історичне зображення.}

\vsec{8.2}
\vdeva{तदा मुक्तिर् यदा चित्तं न वाञ्छति न शोचति । न मुञ्चति न गृह्णाति न हृष्यति न कुप्यति ॥}
\viast{tadā muktir yadā cittaṃ na vāñchati na śocati | na muñcati na gṛhṇāti na hṛṣyati na kupyati ||}
\vtranslit{тада муктір яда чіттам на ваньччаті на шочаті | на муньчаті на ґріхнаті на хрішяті на куп'яті ||}
\vperek{Визволення — тоді, коли розум не бажає, не сумує, не відпускає й не бере, не радіє й не гнівається.}
\vcom{Построкове зіставлення: редакторська чернетка, не рецензовано.}
\vcom{Рядок а}
\vdeva{तदा मुक्तिर् यदा चित्तं न वाञ्छति न शोचति}
\viast{tadā muktir yadā cittaṃ na vāñchati na śocati}
\vtranslit{тада муктір яда чіттам на ваньччаті на шочаті}
\vperek{Визволення — тоді, коли розум не бажає, не сумує,}
\vcom{Кореляція tadā ... yadā пов'язує визволення зі станом розуму. Форми na vāñchati й na śocati називають перші два заперечені рухи — «не бажає» і «не сумує». Це початок одного речення, яке продовжується в другому піврядку.}
\vcom{Рядок б}
\vdeva{न मुञ्चति न गृह्णाति न हृष्यति न कुप्यति}
\viast{na muñcati na gṛhṇāti na hṛṣyati na kupyati}
\vtranslit{на муньчаті на ґріхнаті на хрішяті на куп'яті}
\vperek{не відпускає й не бере, не радіє й не гнівається.}
\vcom{Чотири повтори na завершують симетричний перелік: не відпускає, не бере, не радіє, не гнівається. Порівняння з 8.1 є редакторським зіставленням двох чинних карток; 8.2 має hṛṣyati, що не дає підстав автоматично замінювати dhṛṣyati у 8.1.}
\vcom{Вірш є **дослівним запереченням** попереднього: ті самі шість дієслів, кожне з часткою na. Це важливо для перекладу: будь-яка стилістична варіація тут зруйнує паралелізм, який і є аргументом. Визволення не додає нічого — воно знімає рівно те, що було названо полоном.}
\vpar{Українська паралель (наша інтерпретація)}{Леся Українка, «Contra spem spero!» (1890).}{Наша паралель: у Лесі Українки рівновага теж не залежить від того, чи є надія. Наше порівняння — типологічне: обидва тексти описують стан, а не здобуток.}
\vfig{ilustratsii/8-2-knot.pdf}{Той самий вузол, але кожен рух заперечено}{Схема наша, не історична ілюстрація.}

\vsec{8.3}
\vdeva{तदा बन्धो यदा चित्तं सक्तं कास्वपि दृष्टिषु । तदा मोक्षो यदा चित्तम् असक्तं सर्वदृष्टिषु ॥}
\viast{tadā bandho yadā cittaṃ saktaṃ kāsvapi dṛṣṭiṣu | tadā mokṣo yadā cittam asaktaṃ sarvadṛṣṭiṣu ||}
\vtranslit{тада бандго яда чіттам сактам касвапі дріштішу | тада мокшо яда чіттам асактам сарвадріштішу ||}
\vperek{Полон — тоді, коли розум прив'язаний до якихось поглядів; визволення — коли розум неприв'язаний до всіх поглядів.}
\vcom{Построкове зіставлення: редакторська чернетка, не рецензовано.}
\vcom{Рядок а}
\vdeva{तदा बन्धो यदा चित्तं सक्तं कास्वपि दृष्टिषु}
\viast{tadā bandho yadā cittaṃ saktaṃ kāsvapi dṛṣṭiṣu}
\vtranslit{тада бандго яда чіттам сактам касвапі дріштішу}
\vperek{Полон — тоді, коли розум прив'язаний до якихось поглядів;}
\vcom{Фраза tadā bandho yadā формує умову полону; saktaṃ означає прив'язаність розуму, kāsvapi dṛṣṭiṣu — «до якихось поглядів». Це перша половина протиставлення, а не твердження про історичний вплив на українських авторів.}
\vcom{Рядок б}
\vdeva{तदा मोक्षो यदा चित्तम् असक्तं सर्वदृष्टिषु}
\viast{tadā mokṣo yadā cittam asaktaṃ sarvadṛṣṭiṣu}
\vtranslit{тада мокшо яда чіттам асактам сарвадріштішу}
\vperek{визволення — коли розум неприв'язаний до всіх поглядів.}
\vcom{Друга половина паралельна першій: tadā mokṣo yadā вводить визволення, asaktaṃ — неприв'язаність, sarvadṛṣṭiṣu — «до всіх поглядів». Контраст kāsvapi / sarva передає різний обсяг прив'язаності; переклад збережено дослівно.}
\vcom{`dṛṣṭi` — «погляд» у подвійному сенсі: і зір, і переконання. Асиметрія навмисна: `kāsvapi` («до якихось») проти `sarva` («до всіх»). Тобто достатньо **однієї** прив'язаності, щоб бути в полоні, але потрібна неприв'язаність до **всіх** — і це не той самий обсяг вимог, що звучить суворіше, ніж здається з першого читання.}
\vpar{Українська паралель (наша інтерпретація)}{Микола Хвильовий, «Я (Романтика)» (1924).}{Наша паралель: у Хвильового привʼязаність до власного погляду теж постає як те, що тримає міцніше за будь-яку зовнішню силу. Наше порівняння — за будовою думки.}
\vfig{ilustratsii/8-3-balance.pdf}{Дві шальки: привʼязаність до поглядів і свобода від них}{Схема наша, не історична ілюстрація.}

\vsec{8.4}
\vdeva{यदा नाहं तदा मोक्षो यदाहं बन्धनं तदा । मत्वेति हेलया किञ्चिन् मा गृहाण विमुञ्च मा ॥}
\viast{yadā nāhaṃ tadā mokṣo yadāhaṃ bandhanaṃ tadā | matveti helayā kiñcin mā gṛhāṇa vimuñca mā ||}
\vtranslit{яда нахам тада мокшо ядахам бандганам тада | матветі хелая кіньчін ма ґріхана вімуньча ма ||}
\vperek{Коли немає «я» — визволення; коли є «я» — полон. Так розуміючи, легко нічого не бери й нічого не відпускай.}
\vcom{Построкове зіставлення: редакторська чернетка, не рецензовано.}
\vcom{Рядок а}
\vdeva{यदा नाहं तदा मोक्षो यदाहं बन्धनं तदा}
\viast{yadā nāhaṃ tadā mokṣo yadāhaṃ bandhanaṃ tadā}
\vtranslit{яда нахам тада мокшо ядахам бандганам тада}
\vperek{Коли немає «я» — визволення; коли є «я» — полон.}
\vcom{Протиставлення yadā nāhaṃ / yadāhaṃ пов'язує відсутність чи присутність «я» з визволенням та полоном. Це філософське твердження тексту, а не психологічний діагноз; чинний український відповідник лишається без змін.}
\vcom{Рядок б}
\vdeva{मत्वेति हेलया किञ्चिन् मा गृहाण विमुञ्च मा}
\viast{matveti helayā kiñcin mā gṛhāṇa vimuñca mā}
\vtranslit{матветі хелая кіньчін ма ґріхана вімуньча ма}
\vperek{Так розуміючи, легко нічого не бери й нічого не відпускай.}
\vcom{matveti вводить розуміння попередньої тези, helayā — відтінок легкості. Подвійне mā при наказових формах gṛhāṇa та vimuñca передає дві заборони: не брати й не відпускати; жодних нових санскритських читань тут не встановлюємо.}
\vcom{Тут уперше названо причину, а не симптом: попередні три вірші описували рухи розуму, цей називає їхнє джерело — `aham`. `helayā` («легко, жартома») різко змінює тон: те, що виглядало як велика вимога, подане як щось майже незначне. І `mā … mā` — подвійна заборона: не бери й не відпускай, тобто не роби й самого жесту вибору.}
\vpar{Українська паралель (наша інтерпретація)}{Іван Багряний, «Сад Гетсиманський» (1950).}{Наша паралель: у Багряного теж є думка, що остання вʼязниця людини — вона сама. Наше порівняння — типологічне, не твердження про вплив.}
\vfig{ilustratsii/8-4-neti.pdf}{Я і його відсутність: легко не брати й не відпускати}{Схема наша, не історична ілюстрація.}


% ЗГЕНЕРОВАНО scripts/build_format_v2_chapter.py — не редагувати вручну.
% Джерело: editorial/cards/9.*.json, format-v2/rozdil-9.json, paraleli.json.

\naukchapter{9}{8}{Глава 9. Відречення}
\naukintro{Девʼята глава — вісім віршів про відречення, і жодного разу воно не постає як зусилля. Спершу пари протилежностей утихають самі (9.1), далі втихає бажання жити (9.2), далі річ виявляється без суті (9.3). Найгостріший вірш — 9.8: сансара названа просто усталеними схильностями, і тоді відречення зводиться до одного — відпустити їх.}
\vsec{9.1}
\vdeva{कृताकृते च द्वन्द्वानि कदा शान्तानि कस्य वा । एवं ज्ञात्वेह निर्वेदाद् भव त्यागपरो ऽव्रती ॥}
\viast{kṛtākṛte ca dvandvāni kadā śāntāni kasya vā | evaṃ jñātveha nirvedād bhava tyāgaparo 'vratī ||}
\vtranslit{крітакріте ча двандвані када шантані кася ва | евам джнятвеха нірведад бгава тяґапаро ’враті ||}
\vperek{Зроблене й незроблене, пари протилежностей — коли вони втихли і чиї вони? Так знаючи, будь тут, у відреченні, схильним до відкидання — і без обітниць.}
\vcom{Построкове зіставлення: редакторська чернетка, не рецензовано.}
\vcom{Рядок а}
\vdeva{कृताकृते च द्वन्द्वानि कदा शान्तानि कस्य वा}
\viast{kṛtākṛte ca dvandvāni kadā śāntāni kasya vā}
\vtranslit{крітакріте ча двандвані када шантані кася ва}
\vperek{Зроблене й незроблене, пари протилежностей — коли вони втихли і чиї вони?}
\vcom{Питальні слова kadā («коли?») і kasya vā («чиї?») завершують перший друкований піврядок. dvandvāni — «пари протилежностей»; kṛtākṛte потребує окремої морфологічної звірки, тому українське «зроблене й незроблене» лишається чинним робочим відповідником, а не остаточно встановленим граматичним рішенням.}
\vcom{Рядок б}
\vdeva{एवं ज्ञात्वेह निर्वेदाद् भव त्यागपरो ऽव्रती}
\viast{evaṃ jñātveha nirvedād bhava tyāgaparo 'vratī}
\vtranslit{евам джнятвеха нірведад бгава тяґапаро ’враті}
\vperek{Так знаючи, будь тут, у відреченні, схильним до відкидання — і без обітниць.}
\vcom{evaṃ jñātvā («так пізнавши») передує наказовому bhava («будь»); nirvedād вказує на підставу відречення. tyāgaparaḥ («схильний до відкидання») поєднано з avratī («без обітниць»): внутрішнє відречення не перетворюється на припис нової обітниці. GRETIL і SARIT по-різному позначають сегментацію nirvedād bhava / nirvedādbhava та 'vratī / +avratī; це не доказ незалежного рукописного різночитання.}
\vcom{`avratī` («без обітниць») — не дрібниця: текст водночас радить `tyāga` (відкидання) і знімає `vrata` (обітницю). Тобто відкидання не повинно стати новою дисципліною — інакше воно перетвориться на те, від чого відкидають. Питання `kasya vā` («і чиї ж вони?») залишене без відповіді навмисно.}
\vpar{Українська паралель (наша інтерпретація)}{Олександр Потебня, «Думка й мова» (1862).}{Наша паралель: Потебня теж питає, чиї насправді ті пари, якими ми міряємо світ, — і відповідь у нього не в речах, а в тому, хто міряє. Наше порівняння — типологічне.}
\vfig{ilustratsii/9-1-neti.pdf}{Пари протилежностей втихли: лишається відречення без обітниць}{Схема наша, не історична ілюстрація.}

\vsec{9.2}
\vdeva{कस्यापि तात धन्यस्य लोकचेष्टावलोकनात् । जीवितेच्छा बुभुक्षा च बुभुत्सोपशमं गताः ॥}
\viast{kasyāpi tāta dhanyasya lokaceṣṭāvalokanāt | jīvitecchā bubhukṣā ca bubhutsopaśamaṃ gatāḥ ||}
\vtranslit{касяпі тата дганяся локачештавалоканат | джівітечча бубгукша ча бубгутсопашамам ґатах ||}
\vperek{У якогось щасливого, любий, від споглядання того, як поводиться світ, бажання жити, голод і жага пізнавати дійшли до втишення.}
\vcom{Построкове зіставлення: редакторська чернетка, не рецензовано.}
\vcom{Рядок а}
\vdeva{कस्यापि तात धन्यस्य लोकचेष्टावलोकनात्}
\viast{kasyāpi tāta dhanyasya lokaceṣṭāvalokanāt}
\vtranslit{касяпі тата дганяся локачештавалоканат}
\vperek{У якогось щасливого, любий, від споглядання того, як поводиться світ,}
\vcom{Родовий зворот kasyāpi ... dhanyasya («якогось щасливого») окреслює особу, а tāta — звертання до співрозмовника («любий»). Аблатив lokaceṣṭāvalokanāt пояснює підставу: від споглядання руху й учинків світу. Це повідомлення про досвід, а не наказ споглядати за певним правилом.}
\vcom{Рядок б}
\vdeva{जीवितेच्छा बुभुक्षा च बुभुत्सोपशमं गताः}
\viast{jīvitecchā bubhukṣā ca bubhutsopaśamaṃ gatāḥ}
\vtranslit{джівітечча бубгукша ча бубгутсопашамам ґатах}
\vperek{бажання жити, голод і жага пізнавати дійшли до втишення.}
\vcom{jīvitecchā («бажання жити») та bubhukṣā («голод») названі разом; у третьому понятті йдеться про жагу пізнавати. SARIT має bubhutsopaśamaṃ gatāḥ, а GRETIL/TextGrid, за задокументованою звіркою Issue \#227, — bubhutsopaśamaḥ gatāḥ: різниця -ṃ / -ḥ відкрита й змінює можливий синтаксичний розбір. Не вважати електронні читання рукописною колацією. Втишення бажань у цьому контексті не дорівнює заклику до бездіяльності.}
\vcom{Три подібні іменники стоять поруч: `jīvitecchā` (бажання жити), `bubhukṣā` (голод) і `bubhutsā` (жага пізнавати) — останнє те саме слово, що в 10.5, і словник має бути спільним. Причина названа одна й зовнішня: **споглядання** того, як поводиться світ, а не зусилля над собою. `tāta` — ласкаве звертання, не «батько».}
\vpar{Українська паралель (наша інтерпретація)}{Леся Українка, «Оргія» (1913).}{Наша паралель: у Лесі Українки бажання жити теж утихає не від зусилля, а від побаченого. Наше порівняння — за способом бачити.}
\vfig{ilustratsii/9-2-flame.pdf}{Полумʼя, що втихло: бажання жити й жага пізнавати}{Схема наша, не історична ілюстрація.}

\vsec{9.3}
\vdeva{अनित्यं सर्वम् एवेदं तापत्रितयदूषितम् । असारं निन्दितं हेयम् इति निश्चित्य शाम्यति ॥}
\viast{anityaṃ sarvam evedaṃ tāpatritayadūṣitam | asāraṃ ninditaṃ heyam iti niścitya śāmyati ||}
\vtranslit{анітям сарвам еведам тапатрітаядушітам | асарам ніндітам хеям іті нішчітя шам'яті ||}
\vperek{Усе це нестале, зіпсоване троїстою мукою, без суті, осуджене, гідне відкинення — так вирішивши, людина втихає.}
\vcom{Построкове зіставлення: редакторська чернетка, не рецензовано.}
\vcom{Рядок а}
\vdeva{अनित्यं सर्वम् एवेदं तापत्रितयदूषितम्}
\viast{anityaṃ sarvam evedaṃ tāpatritayadūṣitam}
\vtranslit{анітям сарвам еведам тапатрітаядушітам}
\vperek{Усе це нестале, зіпсоване троїстою мукою,}
\vcom{anityaṃ («нестале») характеризує sarvam eva idam («усе це»); tāpa-tritaya-dūṣitam — «уражене троїстою мукою». Сам цей піврядок не називає окремо трьох її різновидів: конкретний перелік належить до пояснення традиції, а не до буквального санскритського рядка. Українське «усе це» зберігає вказівний наголос idam.}
\vcom{Рядок б}
\vdeva{असारं निन्दितं हेयम् इति निश्चित्य शाम्यति}
\viast{asāraṃ ninditaṃ heyam iti niścitya śāmyati}
\vtranslit{асарам ніндітам хеям іті нішчітя шам'яті}
\vperek{без суті, осуджене, гідне відкинення — так вирішивши, людина втихає.}
\vcom{asāraṃ («без суті»), ninditaṃ («осуджене»), heyam («те, що належить відкинути») продовжують низку оцінок мінливого. iti niścitya — «так визначивши / дійшовши висновку», після чого śāmyati означає «втихає». Ця синтаксична послідовність описує втишення після розпізнання, а не наказ нехтувати щоденними обов'язками чи здоров'ям.}
\vcom{`tāpa-tritaya` («троїста мука») — усталена формула: мука від себе, від інших і від природи; переклад має лишити її як формулу, а не розтлумачувати в тексті вірша. П'ять означень ідуть підряд, і кожне сильніше за попереднє: від `anitya` (нестале) до `heya` (гідне відкинення). Втишення тут — **наслідок рішення**, а не почуття.}
\vpar{Українська паралель (наша інтерпретація)}{Ольга Кобилянська, «Земля» (1902).}{Наша паралель: у Кобилянської речі, за які тримаються, теж виявляються несталими й без суті. Наше порівняння — типологічне, не твердження про джерело.}
\vfig{ilustratsii/9-3-neti.pdf}{Нестале й троїста мука відкинуті — людина втихає}{Схема наша, не історична ілюстрація.}

\vsec{9.4}
\vdeva{को ऽसौ कालो वयः किं वा यत्र द्वन्द्वानि नो नृणाम् । तान्य् उपेक्ष्य यथाप्राप्तवर्ती सिद्धिम् अवाप्नुयात् ॥}
\viast{ko 'sau kālo vayaḥ kiṃ vā yatra dvandvāni no nṛṇām | tāny upekṣya yathāprāptavartī siddhim avāpnuyāt ||}
\vtranslit{ко ’сау кало ваях кім ва ятра двандвані но нрінам | танй упекшя ятхапраптаварті сіддгім авапнуят ||}
\vperek{Який же той час, вік або що — коли в людей немає пар протилежностей? Знехтувавши їх, хто живе як випало, той досяг би здійснення.}
\vcom{Построкове зіставлення: редакторська чернетка, не рецензовано.}
\vcom{Рядок а}
\vdeva{को ऽसौ कालो वयः किं वा यत्र द्वन्द्वानि नो नृणाम्}
\viast{ko 'sau kālo vayaḥ kiṃ vā yatra dvandvāni no nṛṇām}
\vtranslit{ко ’сау кало ваях кім ва ятра двандвані но нрінам}
\vperek{Який же той час, вік або що — коли в людей немає пар протилежностей?}
\vcom{Початкове ko 'sau kālo («який же той час») та vayaḥ kiṃ vā («вік або що») утворюють риторичне питання. У сполученні yatra dvandvāni no nṛṇām слово no передає заперечення наявності «пар протилежностей» (dvandvāni) серед людей (nṛṇām), а не присвійне «наших». Питання не обіцяє, що певний вік звільнить людину від протилежностей.}
\vcom{Рядок б}
\vdeva{तान्य् उपेक्ष्य यथाप्राप्तवर्ती सिद्धिम् अवाप्नुयात्}
\viast{tāny upekṣya yathāprāptavartī siddhim avāpnuyāt}
\vtranslit{танй упекшя ятхапраптаварті сіддгім авапнуят}
\vperek{Знехтувавши їх, хто живе як випало, той досяг би здійснення.}
\vcom{tāny upekṣya («не зважаючи на них / відкинувши їх») відсилає до попередніх dvandvāni. yathāprāptavartī буквально окреслює людину, яка живе відповідно до того, що трапляється; це не твердження, ніби обов'язки зникають. siddhim avāpnuyāt містить потенційно-умовну форму avāpnuyāt («досяг би»): переклад не повинен перетворювати її на безумовну обіцянку.}
\vcom{Питання риторичне й знецінювальне: не існує такого віку, коли пар немає, — отже, чекати на них марно. `yathāprāpta-vartī` («той, хто живе як випало») — той самий принцип, що `yadṛcchā` у 3.14 і 4.4, і в перекладі він має звучати так само. `avāpnuyāt` — умовний спосіб: «досяг би», а не «досягне»; це не обіцянка.}
\vpar{Українська паралель (ПОТРЕБУЄ ПЕРЕВІРКИ)}{українська філософська лірика, не визначено.}{ПОТРЕБУЄ ПЕРЕВІРКИ: думка, що жити як випало — це і є здійснення, близька до української філософської лірики, але конкретного твору ми не звірили.}
\vfig{ilustratsii/9-4-balance.pdf}{Дві шальки: пари протилежностей і життя як випало}{Схема наша, не історична ілюстрація.}

\vsec{9.5}
\vdeva{नाना मतं महर्षीणां साधूनां योगिनां तथा । दृष्ट्वा निर्वेदम् आपन्नः को न शाम्यति मानवः ॥}
\viast{nānā mataṃ maharṣīṇāṃ sādhūnāṃ yogināṃ tathā | dṛṣṭvā nirvedam āpannaḥ ko na śāmyati mānavaḥ ||}
\vtranslit{нана матам махаршінам садгунам йоґінам татха | дріштва нірведам апаннах ко на шам'яті манавах ||}
\vperek{Побачивши, скільки різних поглядів мають великі мудреці, праведні й йоги, — хто, дійшовши відречення, не втихне?}
\vcom{Построкове зіставлення: редакторська чернетка, не рецензовано.}
\vcom{Рядок а}
\vdeva{नाना मतं महर्षीणां साधूनां योगिनां तथा}
\viast{nānā mataṃ maharṣīṇāṃ sādhūnāṃ yogināṃ tathā}
\vtranslit{нана матам махаршінам садгунам йоґінам татха}
\vperek{Побачивши, скільки різних поглядів мають великі мудреці, праведні й йоги, —}
\vcom{nānā mataṃ («різноманітні погляди») стосується maharṣi, sādhu та yogin — трьох позначень духовних авторитетів, наведених у родовому множини. Текстологічна примітка: SARIT 9.5 передає maharṣīṇāṃ, тоді як GRETIL Avg\_9.5 має maharṣīṣāṃ; це відмінність цифрових читань, а не доведене різночитання рукописів. Тут відтворено IAST чинної редакторської картки.}
\vcom{Рядок б}
\vdeva{दृष्ट्वा निर्वेदम् आपन्नः को न शाम्यति मानवः}
\viast{dṛṣṭvā nirvedam āpannaḥ ko na śāmyati mānavaḥ}
\vtranslit{дріштва нірведам апаннах ко на шам'яті манавах}
\vperek{хто, дійшовши відречення, не втихне?}
\vcom{dṛṣṭvā («побачивши») позначає попередню дію, nirvedam āpannaḥ — людину, яка дійшла стану розчарування чи неприв'язаності, а ko na śāmyati mānavaḥ формує риторичне питання. Вірш не є доказом непридатності кожного вчителя; це інтерпретація зіткнення різних поглядів, а не наказ відмовитися від знання.}
\vcom{Риторичне питання з відповіддю «ніхто». Важливий перелік: `maharṣi`, `sādhu`, `yogin` — три різні авторитети, і текст показує, що **вони не згодні між собою**. Відречення тут походить не від безладдя світу, а від розбіжності авторитетів — це сильніше й конкретніше.}
\vpar{Українська паралель (наша інтерпретація)}{Євген Плужник, «Рання осінь» (1927).}{Наша паралель: у Плужника розгубленість перед безліччю людських поглядів теж приводить до тиші, а не до нового погляду. Наше порівняння — за будовою думки.}
\vfig{ilustratsii/9-5-motion.pdf}{Багато поглядів сходяться до тиші}{Схема наша, не історична ілюстрація.}

\vsec{9.6}
\vdeva{कृत्वा मूर्तिपरिज्ञानं चैतन्यस्य न किं गुरुः । निर्वेदसमतायुक्त्या यस् तारयति संसृतेः ॥}
\viast{kṛtvā mūrtiparijñānaṃ caitanyasya na kiṃ guruḥ | nirvedasamatāyuktyā yas tārayati saṃsṛteḥ ||}
\vtranslit{крітва муртіпаріджнянам чайтаняся на кім ґурух | нірведасаматаюктя яс тараяті самсрітех ||}
\vperek{Хто, пізнавши образ свідомості, — чи не вчитель? — той через відречення, рівність і з'єднаність переправляє з сансари.}
\vcom{Построкове зіставлення: редакторська чернетка, не рецензовано.}
\vcom{Рядок а}
\vdeva{कृत्वा मूर्तिपरिज्ञानं चैतन्यस्य न किं गुरुः}
\viast{kṛtvā mūrtiparijñānaṃ caitanyasya na kiṃ guruḥ}
\vtranslit{крітва муртіпаріджнянам чайтаняся на кім ґурух}
\vperek{Хто, пізнавши образ свідомості, — чи не вчитель? —}
\vcom{kṛtvā mūrtiparijñānaṃ та caitanyasya пов'язують пізнання з образом або виявом свідомості; na kiṃ guruḥ має питальну силу «чи не вчитель?». Синтаксис нетривіальний: перша частина не містить наказу безумовно підкорятися наставникові.}
\vcom{Рядок б}
\vdeva{निर्वेदसमतायुक्त्या यस् तारयति संसृतेः}
\viast{nirvedasamatāyuktyā yas tārayati saṃsṛteḥ}
\vtranslit{нірведасаматаюктя яс тараяті самсрітех}
\vperek{той через відречення, рівність і з'єднаність переправляє з сансари.}
\vcom{nirveda-samatā-yuktyā — форма орудного відмінка: засіб, пов'язаний із неприв'язаністю, рівновагою та способом їх поєднання; yas tārayati saṃsṛteḥ — «хто переправляє із сансари». Образ переправи не вказує конкретної фізичної річки чи історичної події; додавати їх до наукового перекладу не слід.}
\vcom{Єдине місце в книзі, де `guru` названо **позитивно** і з означенням: учитель — той, хто пізнав `caitanya`. Але й тут текст не вимагає послуху: `tārayati` — «переправляє», дієслово руху, а не влади. Три складники (`nirveda`, `samatā`, `yukti`) подані як спосіб дії, а не як чесноти вчителя.}
\vpar{Українська паралель (наша інтерпретація)}{Григорій Сковорода, «Кільце. Дружня розмова про душевний мир».}{Наша паралель: у Сковороди рівність теж є тим, через що людина виходить із кола, у якому її тримали. Наше порівняння — типологічне.}
\vfig{ilustratsii/9-6-witness.pdf}{Переправа з сансари: центр — пізнаний образ свідомості}{Схема наша, не історична ілюстрація.}

\vsec{9.7}
\vdeva{पश्य भूतविकारांस् त्वं भूतमात्रान् यथार्थतः । तत्क्षणाद् बन्धनिर्मुक्तः स्वरूपस्थो भविष्यसि ॥}
\viast{paśya bhūtavikārāṃs tvaṃ bhūtamātrān yathārthataḥ | tatkṣaṇād bandhanirmuktaḥ svarūpastho bhaviṣyasi ||}
\vtranslit{пашя бгутавікарамс твам бгутаматран ятхартхатах | таткшанад бандганірмуктах сварупастхо бгавішясі ||}
\vperek{Дивись на видозміни сущого як на самі лише елементи, як є насправді — і тієї ж миті, вільний від пут, станеш тим, хто стоїть у власній формі.}
\vcom{Построкове зіставлення: редакторська чернетка, не рецензовано.}
\vcom{Рядок а}
\vdeva{पश्य भूतविकारांस् त्वं भूतमात्रान् यथार्थतः}
\viast{paśya bhūtavikārāṃs tvaṃ bhūtamātrān yathārthataḥ}
\vtranslit{пашя бгутавікарамс твам бгутаматран ятхартхатах}
\vperek{Дивись на видозміни сущого як на самі лише елементи, як є насправді —}
\vcom{paśya — наказовий спосіб «дивись»; bhūta-vikārān і bhūta-mātrān співвідносять видозміни й елементи. yathārthataḥ уточнює спосіб сприйняття — «як є насправді». Це заклик до розрізнення в межах філософського тексту, не твердження про сучасну фізику чи медицину.}
\vcom{Рядок б}
\vdeva{तत्क्षणाद् बन्धनिर्मुक्तः स्वरूपस्थो भविष्यसि}
\viast{tatkṣaṇād bandhanirmuktaḥ svarūpastho bhaviṣyasi}
\vtranslit{таткшанад бандганірмуктах сварупастхо бгавішясі}
\vperek{і тієї ж миті, вільний від пут, станеш тим, хто стоїть у власній формі.}
\vcom{tatkṣaṇāt («тієї ж миті») визначає часовий перехід; bandha-nirmuktaḥ («звільнений від пут») і svarūpa-sthaḥ («той, хто перебуває у власній формі») описують очікуваний стан. bhaviṣyasi є майбутнім часом «станеш», а не минулим описом уже досягнутого.}
\vcom{`tatkṣaṇāt` («тієї ж миті») — найсильніше слово вірша: не поступово, не через практику, а негайно. `bhūtamātra` («самі лише елементи») знімає оцінку з речей: вони не погані й не добрі, вони — складники. `svarūpastha` («той, хто стоїть у власній формі») — стан, а не дія; і `bhaviṣyasi` — майбутній час, тому це не опис того, що вже є, а вказівка на те, що станеться.}
\vpar{Українська паралель (наша інтерпретація)}{Володимир Свідзінський, «Медобір» (1975, посмертно).}{Наша паралель: у Свідзінського річ теж розкладається на те, з чого вона є, і саме так стає видимою. Наше порівняння — за образом.}
\vfig{ilustratsii/9-7-thread.pdf}{Видозміни розібрані на елементи, як намистини на нитці}{Схема наша, не історична ілюстрація.}

\vsec{9.8}
\vdeva{वासना एव संसार इति सर्वा विमुञ्च ताः । तत्त्यागो वासनात्यागात् स्थितिरद्य यथा तथा ॥}
\viast{vāsanā eva saṃsāra iti sarvā vimuñca tāḥ | tattyāgo vāsanātyāgāt sthitiradya yathā tathā ||}
\vtranslit{васана ева самсара іті сарва вімуньча тах | таттяґо васанатяґат стхітірадя ятха татха ||}
\vperek{Сансара — це саме усталені схильності, так; відпусти їх усі. Їх відкидання — через відкидання схильностей; стояння в собі — нині, так чи інакше.}
\vcom{Построкове зіставлення: редакторська чернетка, не рецензовано.}
\vcom{Рядок а}
\vdeva{वासना एव संसार इति सर्वा विमुञ्च ताः}
\viast{vāsanā eva saṃsāra iti sarvā vimuñca tāḥ}
\vtranslit{васана ева самсара іті сарва вімуньча тах}
\vperek{Сансара — це саме усталені схильності, так; відпусти їх усі.}
\vcom{vāsanā пов'язано із saṃsāra, а sarvā ... tāḥ указує на множину схильностей, які закликають відпустити; vimuñca — наказова форма. Це філософський контекст звички й прив'язаності, не порада відмовлятися від медичної допомоги, соціальних зв'язків чи відповідальності.}
\vcom{Рядок б}
\vdeva{तत्त्यागो वासनात्यागात् स्थितिरद्य यथा तथा}
\viast{tattyāgo vāsanātyāgāt sthitiradya yathā tathā}
\vtranslit{таттяґо васанатяґат стхітірадя ятха татха}
\vperek{Їх відкидання — через відкидання схильностей; стояння в собі — нині, так чи інакше.}
\vcom{tat-tyāgaḥ та vāsanā-tyāgāt повторюють корінь tyāga («відкидання») у різних граматичних відношеннях; sthitir adya «перебування нині», yathā tathā «так чи інакше». Переклад не повинен вигадувати конкретного ритуалу або гарантувати миттєвий зовнішній результат.}
\vcom{Вірш закінчує главу визначенням, а не обіцянкою: сансара — не світ, а `vāsanā`, усталені схильності. `yathā tathā` («так чи інакше») знімає вимогу до способу: результат не залежить від того, як саме це станеться. Друга половина повторює те саме слово двічі (`tyāga`), і повтор тут не недогляд — він показує, що відкидання схильностей і є відкиданням, а не засобом до нього.}
\vpar{Українська паралель (наша інтерпретація)}{Леся Українка, «Бояриня» (1910).}{Наша паралель: у Лесі Українки схильності, засвоєні звичкою, теж тримають людину міцніше за будь-який наказ. Наше порівняння — типологічне.}
\vfig{ilustratsii/9-8-knot.pdf}{Вузол схильностей і стояння в собі}{Схема наша, не історична ілюстрація.}


% ЗГЕНЕРОВАНО scripts/build_format_v2_chapter.py — не редагувати вручну.
% Джерело: editorial/cards/10.*.json, format-v2/rozdil-10.json, paraleli.json.

\naukchapter{10}{8}{Глава 10. Досить уже}
\naukintro{Десята глава — найрізкіша за тоном. Вона не пояснює й не доводить: вона перелічує те, що вже було, і каже досить. Спершу відкидаються бажання й користь (10.1), далі побутові речі названо сном і чарами (10.2), далі спрага названа самою сансарою (10.3–10.4). Останні три вірші звернені до минулого: усе це вже гинуло не раз, і саме тому питання стоїть не про нове, а про припинення.}
\vsec{10.1}
\vdeva{विहाय वैरिणं कामम् अर्थं चानर्थसंकुलं । धर्मम् अप्य् एतयोर् हेतुं सर्वत्रानादरं कुरु ॥}
\vtranslit{віхая вайрінам камам артхам чанартхасамкулам | дгармам апй етайор хетум сарватранадарам куру ||}
\vperek{Покинувши вороже бажання і сповнену лиха користь — навіть дхарму, що є причиною обох, — усюди зневажай.}
\vcom{Построкове зіставлення: редакторська чернетка, не рецензовано.}
\vcom{Рядок а}
\vdeva{विहाय वैरिणं कामम् अर्थं चानर्थसंकुलं}
\vtranslit{віхая вайрінам камам артхам чанартхасамкулам}
\vperek{Покинувши вороже бажання і сповнену лиха користь —}
\vcom{vihāya — дієприслівникова форма «покинувши»; vairiṇaṃ kāmam і arthaṃ cānarthasaṃkulaṃ називають бажання та користь, сповнену лиха. Цей епітет стосується artha, а не автоматично всіх наступних понять.}
\vcom{Рядок б}
\vdeva{धर्मम् अप्य् एतयोर् हेतुं सर्वत्रानादरं कुरु}
\vtranslit{дгармам апь етайор хетум сарватранадарам куру}
\vperek{навіть дхарму, що є причиною обох, — усюди зневажай.}
\vcom{dharmam api вводить навіть дхарму; etayor hetuṃ відносить «причину обох» до попередніх понять, sarvatrānādaraṃ kuru має наказову силу. Порівняння з українською етикою потребує контексту: це теза вірша, а не заклик до порушення закону.}
\vcom{Дхарма стоїть у переліку **останньою** і з підсиленням api («навіть»): спершу відкидають бажання й користь, а потім те, що їх породжує. anartha-saṃkula («сповнене лиха») стосується саме artha, і переклад не повинен переносити цю оцінку на дхарму. «Зневажай» (anādara) — сильне слово; м'якше «не зважай» змінило б ступінь.}
\vpar{Українська паралель (наша інтерпретація)}{Пантелеймон Куліш, «Хуторна поезія» (1882).}{Наша паралель: у Куліша відмова від користі теж подана як умова спокою, а не як втрата. Наше порівняння — типологічне, не твердження про джерело.}
\vfig{ilustratsii/10-1-neti.pdf}{Бажання, користь і дхарма відкинуті разом}{Схема наша, не історична ілюстрація.}

\vsec{10.2}
\vdeva{स्वप्नेन्द्रजालवत् पश्य दिनानि त्रीणि पञ्च वा । मित्रक्षेत्रधनागारदारदायादिसंपदः ॥}
\vtranslit{свапнендраджалават пашя дінані тріні паньча ва | мітракшетрадганаґарадарадаядісампадах ||}
\vperek{Дивись на друзів, поле, гроші, дім, дружину, спадок і решту багатств як на сон і чари — на три дні або п'ять.}
\vcom{Построкове зіставлення: редакторська чернетка, не рецензовано.}
\vcom{Рядок а}
\vdeva{स्वप्नेन्द्रजालवत् पश्य दिनानि त्रीणि पञ्च वा}
\vtranslit{свапнендраджалават пашья дінані тріні панча ва}
\vperek{Дивись на друзів, поле, гроші, дім, дружину, спадок і решту багатств}
\vcom{svapnendrajālavat paśya пропонує погляд «як на сон і чари», dināni trīṇi pañca vā згадує три чи п'ять днів. Український повний переклад переносить перелік предметів із наступного піврядка перед порівнянням; рядковий фрагмент тому є смисловим, а не буквально послідовним підрядником.}
\vcom{Рядок б}
\vdeva{मित्रक्षेत्रधनागारदारदायादिसंपदः}
\vtranslit{мітракшетрадганаґарадарадаядісампадах}
\vperek{як на сон і чари — на три дні або п'ять.}
\vcom{mitrakṣetradhanāgāradāradāyādisaṃpadaḥ складає конкретний перелік дружби, поля, грошей, житла, родини та спадку. В українській цілісній редакції цей перелік стоїть на початку речення; це прозоре редакторське переставлення, а не текстова розбіжність санскритських свідків.}
\vcom{Перелік навмисно побутовий: не «світ» загалом, а конкретні речі, за які тримаються. «Три дні або п'ять» — не точний строк, а спосіб сказати «коротко»; переклад зберігає числа, бо вони й творять ефект. indrajāla — не просто «ілюзія», а фокус, свідома підміна видимого.}
\vpar{Українська паралель (ПОТРЕБУЄ ПЕРЕВІРКИ)}{українська барокова проповідь, не визначено.}{ПОТРЕБУЄ ПЕРЕВІРКИ: перелік побутових речей як сну й чару близький до барокової проповіді про марність, але конкретного твору ми не перевірили.}
\vfig{ilustratsii/10-2-illusion.pdf}{Побутові речі під серпанком: сон і чари}{Схема наша, не історична ілюстрація.}

\vsec{10.3}
\vdeva{यत्र यत्र भवेत् तृष्णा संसारं विद्धि तत्र वै । प्रौढवैराग्यम् आश्रित्य वीततृष्णः सुखी भव ॥}
\vtranslit{ятра ятра бгавет трішна самсарам віддгі татра вай | праудгавайраґям ашрітя вітатрішнах сукхі бгава ||}
\vperek{Де б не виникла спрага — знай: саме там і сансара. Опираючись на зрілу безсторонність, стань тим, у кого спрага минула, — і будь щасливий.}
\vcom{Построкове зіставлення: редакторська чернетка, не рецензовано.}
\vcom{Рядок а}
\vdeva{यत्र यत्र भवेत् तृष्णा संसारं विद्धि तत्र वै}
\vtranslit{ятра ятра бгавет трішна самсарам віддгі татра вай}
\vperek{Де б не виникла спрага — знай: саме там і сансара.}
\vcom{Пара yatra yatra ... tatra vai утворює відповідність «де б ... саме там»; bhavet tṛṣṇā і saṃsāraṃ viddhi поєднують спрагу та сансару через наказ «знай». Тут сансара названа у зв'язку зі спрагою, а не як географічне місце.}
\vcom{Рядок б}
\vdeva{प्रौढवैराग्यम् आश्रित्य वीततृष्णः सुखी भव}
\vtranslit{праудгавайраґьям ашрітья вітатрішнах сукхі бгава}
\vperek{Опираючись на зрілу безсторонність, стань тим, у кого спрага минула, — і будь щасливий.}
\vcom{pra uḍha-vairāgyam āśritya описує опертя на зрілу неприв'язаність; vītatṛṣṇaḥ sukhī bhava завершується наказом «будь щасливий». Складене prauḍhavairāgyam збережено в IAST картки; коментар не додає нового варіанта читання.}
\vcom{Тут уперше в розділі дано не заперечення, а **вказівку**: `yatra yatra … tatra` («де б … саме там»). Сансара — не місце й не стан світу, а точка, де виникла спрага. prauḍha («зріла») важливе: безсторонність має бути витриманою, а не поривчастою. Висновок «будь щасливий» — наказ, а не побажання.}
\vpar{Українська паралель (наша інтерпретація)}{Володимир Винниченко, «Сонячна машина» (1928).}{Наша паралель: у Винниченка спрага теж постає тим, що тримає людину в колі, незалежно від того, чи вона цього хоче. Наше порівняння — типологічне.}
\vfig{ilustratsii/10-3-knot.pdf}{Вузол спраги: де вона виникла, там і сансара}{Схема наша, не історична ілюстрація.}

\vsec{10.4}
\vdeva{तृष्णामात्रात्मको बन्धस् तन्नाशो मोक्ष उच्यते । भवासंसक्तिमात्रेण प्राप्तितुष्टिर् मुहुर् मुहुः ॥}
\vtranslit{трішнаматратмако бандгас таннашо мокша учяте | бгавасамсактіматрена праптітуштір мухур мухух ||}
\vperek{Полон має природу самої лише спраги; її знищення називають визволенням. Від самої лише неприв'язаності до буття приходить вдоволеність досягнутим — знову й знову.}
\vcom{Построкове зіставлення: редакторська чернетка, не рецензовано.}
\vcom{Рядок а}
\vdeva{तृष्णामात्रात्मको बन्धस् तन्नाशो मोक्ष उच्यते}
\vtranslit{трішнаматратмако бандгас таннашо мокша учьяте}
\vperek{Полон має природу самої лише спраги; її знищення називають визволенням.}
\vcom{tṛṣṇāmātrātmako bandhas визначає полон через саму спрагу; tannāśo mokṣa ucyate називає її зникнення визволенням. Це визначення всередині вірша, не універсальна психологічна діагностика.}
\vcom{Рядок б}
\vdeva{भवासंसक्तिमात्रेण प्राप्तितुष्टिर् मुहुर् मुहुः}
\vtranslit{бгавасамсактіматрена праптітуштір мухур мухух}
\vperek{Від самої лише неприв'язаності до буття приходить вдоволеність досягнутим — знову й знову.}
\vcom{bhavāsaṃsaktimātreṇa указує на неприв'язаність, prāptituṣṭir — на вдоволеність досягнутим, muhur muhuḥ — на повторюваність. Українське «знову й знову» зберігає повтор без додавання нового джерела.}
\vcom{Тут полон означено **однією** ознакою — спрага, і визволення теж однією — її знищення. Це найточніше місце глави, і воно коротке; розширювати переклад не варто. `toṣa` входить сюди як складник `prāpti-tuṣṭi` («вдоволеність досягнутим»), а не як самостійна чеснота — це варто пам'ятати при звірці з 1.2, де `toṣa` стоїть окремо. `muhur muhuḥ` — не «завжди», а «раз у раз»: вдоволеність повторюється, а не тримається одного разу назавжди.}
\vpar{Українська паралель (наша інтерпретація)}{Іван Франко, «Смерть Каїна» (1889).}{Наша паралель: у Франка полон теж має природу самої лише спраги, а не зовнішньої перешкоди. Наше порівняння — за будовою думки.}
\vfig{ilustratsii/10-4-knot.pdf}{Один вузол у двох назвах: полон і визволення}{Схема наша, не історична ілюстрація.}

\vsec{10.5}
\vdeva{त्वम् एकश् चेतनः शुद्धो जडं विश्वम् असत् तथा । अविद्यापि न किञ्चित् सा का बुभुत्सा तथापि ते ॥}
\vtranslit{твам екаш четанах шуддго джадам вішвам асат татха | авідяпі на кіньчіт са ка бубгутса татхапі те ||}
\vperek{Ти один — свідомий, чистий; світ же — неживий і несправжній. Навіть незнання — аніщо. То яке в тебе, при всьому тому, бажання пізнавати?}
\vcom{Построкове зіставлення: редакторська чернетка, не рецензовано.}
\vcom{Рядок а}
\vdeva{त्वम् एकश् चेतनः शुद्धो जडं विश्वम् असत् तथा}
\vtranslit{твам екаш четанах шуддго джадам вішвам асат татха}
\vperek{Ти один — свідомий, чистий; світ же — неживий і несправжній.}
\vcom{tvam ekaś cetanaḥ śuddho характеризує адресата як свідомого й чистого, jaḍaṃ viśvam asat tathā — світ як неживий та несправжній у межах позиції твору. asat не дорівнює моральній категорії «злий».}
\vcom{Рядок б}
\vdeva{अविद्यापि न किञ्चित् सा का बुभुत्सा तथापि ते}
\vtranslit{авідьяпі на кінчіт са ка бубгутса татхапі те}
\vperek{Навіть незнання — аніщо. То яке в тебе, при всьому тому, бажання пізнавати?}
\vcom{avidyāpi na kiñcit sā — твердження про незнання, kā bubhutsā tathāpi te — риторичне питання про бажання пізнавати. Це філософська постановка питання, не заклик відмовлятися від навчання чи наукової перевірки.}
\vcom{`asat` тут — «несправжнє», не «зло»; це різні речі, і переклад не повинен їх зливати. Питання `kā bubhutsā` замикає коло: якщо навіть незнання ніщо, то пізнавати нічого — і саме тому питання адресоване тому, хто вже слухає пояснення. Це найгостріше місце глави: воно ставить під сумнів користь самої розмови.}
\vpar{Українська паралель (ПОТРЕБУЄ ПЕРЕВІРКИ)}{українська релігійна філософія, не визначено.}{ПОТРЕБУЄ ПЕРЕВІРКИ: питання, навіщо пізнавати те, що й так не є річчю, близьке до української релігійної філософії, але конкретного місця ми не звірили.}
\vfig{ilustratsii/10-5-witness.pdf}{Свідомий у центрі, несправжній світ навколо}{Схема наша, не історична ілюстрація.}

\vsec{10.6}
\vdeva{राज्यं सुताः कलत्राणि शरीराणि सुखानि च । संसक्तस्यापि नष्टानि तव जन्मनि जन्मनि ॥}
\vtranslit{раджям сутах калатрані шарірані сукхані ча | самсактасяпі наштані тава джанмані джанмані ||}
\vperek{Царство, сини, дружини, тіла й утіхи — усе це вже гинуло для тебе, прив'язаного, — з народження в народження.}
\vcom{Построкове зіставлення: редакторська чернетка, не рецензовано.}
\vcom{Рядок а}
\vdeva{राज्यं सुताः कलत्राणि शरीराणि सुखानि च}
\vtranslit{раджьям сутах калатрані шарірані сукхані ча}
\vperek{Царство, сини, дружини, тіла й утіхи —}
\vcom{rājyaṃ sutāḥ kalatrāṇi śarīrāṇi sukhāni ca перелічує царство, дітей, дружин, тіла й утіхи. Початок із царства враховує становище співрозмовника, але не встановлює зовнішніх історичних дат.}
\vcom{Рядок б}
\vdeva{संसक्तस्यापि नष्टानि तव जन्मनि जन्मनि}
\vtranslit{самсактасьяпі наштані тава джанмані джанмані}
\vperek{усе це вже гинуло для тебе, прив'язаного, — з народження в народження.}
\vcom{saṃsaktasyāpi naṣṭāni tava janmani janmani підкреслює втрати для прив'язаного в повторюваних народженнях. Українське «вже гинуло» передає минулу форму naṣṭāni, не проголошуючи переродження емпірично доведеним фактом.}
\vcom{`naṣṭāni` — доконаний дієприкметник: не «загине», а «гинуло вже не раз». Переклад має зберегти цю повторюваність, інакше вірш звучить як погроза, а він — як нагадування. Звертання адресоване цареві, і царство названо **першим** у переліку: порядок не випадковий.}
\vpar{Українська паралель (ПОТРЕБУЄ ПЕРЕВІРКИ)}{українська поезія XIX ст., не визначено.}{ПОТРЕБУЄ ПЕРЕВІРКИ: думка, що все набуте вже гинуло не раз, трапляється в українській поезії, але називати твір без звірки ми не будемо.}
\vfig{ilustratsii/10-6-motion.pdf}{Низка набутків, що вже гинули з народження в народження}{Схема наша, не історична ілюстрація.}

\vsec{10.7}
\vdeva{अलम् अर्थेन कामेन सुकृतेनापि कर्मणा । एभ्यः संसारकान्तारे न विश्रान्तम् अभून् मनः ॥}
\vtranslit{алам артхена камена сукрітенапі кармана | ебгях самсаракантаре на вішрантам абгун манах ||}
\vperek{Досить уже й багатства, і бажання, і навіть добрих учинків: у цьому лісі сансари твій розум так і не знав відпочинку.}
\vcom{Построкове зіставлення: редакторська чернетка, не рецензовано.}
\vcom{Рядок а}
\vdeva{अलम् अर्थेन कामेन सुकृतेनापि कर्मणा}
\vtranslit{алам артхена камена сукрітенапі кармана}
\vperek{Досить уже й багатства, і бажання, і навіть добрих учинків:}
\vcom{alam arthena kāmena sukṛtenāpi karmaṇā об'єднує багатство, бажання і навіть заслуги/добрі вчинки в одному заклику «досить». sukṛta не слід ототожнювати із запереченням людяності чи допомоги іншим.}
\vcom{Рядок б}
\vdeva{एभ्यः संसारकान्तारे न विश्रान्तम् अभून् मनः}
\vtranslit{ебгьях самсаракантаре на вішрантам абгун манах}
\vperek{у цьому лісі сансари твій розум так і не знав відпочинку.}
\vcom{ebhyaḥ saṃsārakāntāre na viśrāntam abhūn manaḥ містить образ лісу сансари та розуму без відпочинку. Для українського читача ліс — впізнавана образна ситуація, але ми не приписуємо її запозиченню з українського фольклору.}
\vcom{`sukṛta` тут — «заслуга, добре зроблене», а не «добро» взагалі; і текст відкидає саме заслугу, ставлячи її в один ряд із багатством і бажанням. Це найнезручніше місце глави для читача, який звик, що добрі вчинки беззастережно рахуються. `kāntāra` — «ліс», образ блукання, а не місце. `alam` — «годі вже», з інтонацією втоми, а не заборони.}
\vpar{Українська паралель (ПОТРЕБУЄ ПЕРЕВІРКИ)}{українська проза XX ст., не визначено.}{ПОТРЕБУЄ ПЕРЕВІРКИ: образ лісу сансари, у якому розум не знав відпочинку, близький до української прози про виснаження, але конкретного твору ми не перевірили.}
\vfig{ilustratsii/10-7-flame.pdf}{Полумʼя без відпочинку: ліс сансари}{Схема наша, не історична ілюстрація.}

\vsec{10.8}
\vdeva{कृतं न कति जन्मानि कायेन मनसा गिरा । दुःखम् आयासदं कर्म तद् अद्याप्य् उपरम्यताम् ॥}
\vtranslit{крітам на каті джанмані каєна манаса ґіра | духкхам аясадам карма тад адяпй упарам'ятам ||}
\vperek{Скільки ж за ці народження зроблено тілом, розумом і мовою діл, що дають біль і втому! Тож нехай це припиниться хоч тепер.}
\vcom{Построкове зіставлення: редакторська чернетка, не рецензовано.}
\vcom{Рядок а}
\vdeva{कृतं न कति जन्मानि कायेन मनसा गिरा}
\vtranslit{крітам на каті джанмані кайена манаса ґіра}
\vperek{Скільки ж за ці народження зроблено тілом, розумом і мовою діл, що дають біль і втому!}
\vcom{kṛtaṃ na kati janmāni kāyena manasā girā поєднує повторювані народження з трьома інструментальними формами «тілом, розумом, мовою». Український переклад переносить згадку про біль та втому з другого санскритського піврядка в перше речення; сегментація смислова, не побуквена.}
\vcom{Рядок б}
\vdeva{दुःखम् आयासदं कर्म तद् अद्याप्य् उपरम्यताम्}
\vtranslit{духкхам аясадам карма тад адьяпь упарамьятам}
\vperek{Тож нехай це припиниться хоч тепер.}
\vcom{duḥkham āyāsadaṃ karma називає дію, що приносить біль і втому, tad adyāpy uparamyatām подає побажання її припинення. Даний наказ не вимагає стверджувати історичну відмову від праці конкретної особи.}
\vcom{Потрійний інструментал — тіло, розум, мова — це класична формула трьох дверей дії, і переклад має зберегти всі три, а не згортати їх у «вчинки». `uparamyatām` — пасивний імператив: не «припини», а «нехай припиниться»; дія знімається не зусиллям того, хто діє. `adya api` — «хоч тепер», з відтінком докору, а не нейтрального «зараз». Цим віршем глава 10 закінчується.}
\vpar{Українська паралель (наша інтерпретація)}{Василь Стефаник, «Камінний хрест» (1900).}{Наша паралель: у Стефаника втома від зробленого теж стає останнім аргументом людини. Наше порівняння — типологічне, не про джерело.}
\vfig{ilustratsii/10-8-neti.pdf}{Діла тілом, розумом і мовою відкинуті}{Схема наша, не історична ілюстрація.}


% ЗГЕНЕРОВАНО scripts/build_format_v2_chapter.py — не редагувати вручну.
% Джерело: editorial/cards/11.*.json, format-v2/rozdil-11.json, paraleli.json.

\naukchapter{11}{8}{Глава 11. Переконання}
\naukintro{Одинадцята глава має одну будову на всі вісім віршів: цитата переконання — і наслідок для того, хто його має. niścayī, той, хто має переконання, повторюється на початку кожного вірша, і саме цей повтор і є формою глави. Тут важливо, що переконання не треба доводити: текст описує не аргумент, а стан того, хто вже так думає.}
\vsec{11.1}
\vdeva{भावाभावविकारश् च स्वभावाद् इति निश्चयी । निर्विकारो गतक्लेशः सुखेनैवोपशाम्यति ॥}
\viast{bhāvābhāvavikāraś ca svabhāvād iti niścayī | nirvikāro gatakleśaḥ sukhenaivopaśāmyati ||}
\vtranslit{бгавабгававікараш ча свабгавад іті нішчаї | нірвікаро ґатаклешах сукхенайвопашам'яті ||}
\vperek{«Видозміна буття й небуття — за власною природою» — хто має таке переконання, той без видозмін, з минулими турбаціями, втихає легко.}
\vcom{Построкове зіставлення: редакторська чернетка, не рецензовано.}
\vcom{Рядок а}
\vdeva{भावाभावविकारश् च स्वभावाद् इति निश्चयी}
\viast{bhāvābhāvavikāraś ca svabhāvād iti niścayī}
\vtranslit{бгавабгававікараш ча свабгавад іті нішчаї}
\vperek{«Видозміна буття й небуття — за власною природою» — хто має таке переконання,}
\vcom{bhāva й abhāva повʼязані з vikāra; внутрішнє членування складного слова потребує фахової звірки. svabhāvād — «за власною природою», iti niścayī — «той, хто має таке переконання».}
\vcom{Рядок б}
\vdeva{निर्विकारो गतक्लेशः सुखेनैवोपशाम्यति}
\viast{nirvikāro gatakleśaḥ sukhenaivopaśāmyati}
\vtranslit{нірвікаро ґатаклешах сукхенайвопашам'яті}
\vperek{той без видозмін, з минулими турбаціями, втихає легко.}
\vcom{nirvikāraḥ означає відсутність власної видозміни; gata-kleśaḥ — згаслі турбації; sukhena eva upaśāmyati — втишення саме з легкістю. Наслідок переконання в другому рядку граматично продовжує niścayī з першого.}
\vcom{`niścayī` («той, хто має переконання») повторюється на початку 11.1–11.4 — це структурна формула всієї глави, і переклад має тримати її однаково. `nirvikāraḥ` уточнює: не «без змін» у сенсі нерухомості, а без **власних** видозмін від того, що змінюється назовні. `sukhena eva` («саме з легкістю») — втишення не коштує зусиль, і це важливо.}
\vpar{Українська паралель (наша інтерпретація)}{Григорій Сковорода, «Потоп зміїний» (трактат).}{Наша паралель: у Сковороди видозміна речей теж відбувається за власною природою, і саме тому переконання про це дає спокій. Наше порівняння — типологічне, не твердження про джерело.}
\vfig{ilustratsii/11-1-motion.pdf}{Буття й небуття як видозміна за власною природою}{Схема наша, не історична ілюстрація.}

\vsec{11.2}
\vdeva{ईश्वरः सर्वनिर्माता नेहान्य इति निश्चयी । अन्तर्गलितसर्वाशः शान्तः क्वापि न सज्जते ॥}
\viast{īśvaraḥ sarvanirmātā nehānya iti niścayī | antargalitasarvāśaḥ śāntaḥ kvāpi na sajjate ||}
\vtranslit{ішварах сарванірмата неханя іті нішчаї | антарґалітасарвашах шантах квапі на саджджате ||}
\vperek{«Творець усього — Ішвара, і немає тут іншого» — хто має таке переконання, той із розчиненими всередині сподіваннями, втихлий, ніде не застрягає.}
\vcom{Построкове зіставлення: редакторська чернетка, не рецензовано.}
\vcom{Рядок а}
\vdeva{ईश्वरः सर्वनिर्माता नेहान्य इति निश्चयी}
\viast{īśvaraḥ sarvanirmātā nehānya iti niścayī}
\vtranslit{ішварах сарванірмата неханя іті нішчаї}
\vperek{«Творець усього — Ішвара, і немає тут іншого» — хто має таке переконання,}
\vcom{īśvaraḥ («Ішвара»), sarvanirmātā («творець усього») та nehānya («тут немає іншого») формують перше твердження. iti niścayī означає «той, хто має таке переконання» — дослівно повторювану граматичну рамку 11.1–11.8. Рядок описує позицію мовця, а не доводить конкретне історичне богослов'я.}
\vcom{Рядок б}
\vdeva{अन्तर्गलितसर्वाशः शान्तः क्वापि न सज्जते}
\viast{antargalitasarvāśaḥ śāntaḥ kvāpi na sajjate}
\vtranslit{антарґалітасарвашах шантах квапі на саджджате}
\vperek{той із розчиненими всередині сподіваннями, втихлий, ніде не застрягає.}
\vcom{antargalitasarvāśaḥ — складний прикметник про внутрішньо згаслі сподівання, śāntaḥ — «втихлий», kvāpi na sajjate — «ніде не чіпляється / не застрягає». Не ототожнювати образ внутрішньої свободи з байдужістю до обов'язків інших людей.}
\vcom{У 11.2 Ішвару (`īśvaraḥ`) прямо названо `sarvanirmātā` («творець усього») й заперечено іншого (`na iha anyaḥ`); сама назва Ішвари вже трапляється в 4.6 як `jagad īśvaram` (GRETIL) / `jagadīśvaram` (SARIT). Це не теологія, а спосіб зняти питання причини. `sajjate` — «застрягає, чіпляється»: наслідок не в тому, що людина стала всемогутньою, а в тому, що вона більше ні за що не тримається.}
\vpar{Українська паралель (наша інтерпретація)}{Митрополит Іларіон Київський, «Слово про закон і благодать» (XI ст.).}{Наша паралель: у київському Слові теж є думка, що за всім стоїть один початок, і тому сподівання людини не мають де застрягти. Наше порівняння — типологічне; датування цієї книги науково не встановлене.}
\vfig{ilustratsii/11-2-witness.pdf}{Один творець у центрі, і немає тут іншого}{Схема наша, не історична ілюстрація.}

\vsec{11.3}
\vdeva{आपदः संपदः काले दैवादेवेति निश्चयी । तृप्तः स्वस्थेन्द्रियो नित्यं न वाञ्छति न शोचति ॥}
\viast{āpadaḥ saṃpadaḥ kāle daivādeveti niścayī | tṛptaḥ svasthendriyo nityaṃ na vāñchati na śocati ||}
\vtranslit{ападах сампадах кале дайвадеветі нішчаї | тріптах свастхендрійо нітям на ваньччаті на шочаті ||}
\vperek{«Нещастя й достаток у свій час — саме від долі» — хто має таке переконання, той вдоволений, із чуттями в порядку, повсякчас не бажає й не сумує.}
\vcom{Построкове зіставлення: редакторська чернетка, не рецензовано.}
\vcom{Рядок а}
\vdeva{आपदः संपदः काले दैवादेवेति निश्चयी}
\viast{āpadaḥ saṃpadaḥ kāle daivādeveti niścayī}
\vtranslit{ападах сампадах кале дайвадеветі нішчаї}
\vperek{«Нещастя й достаток у свій час — саме від долі» — хто має таке переконання,}
\vcom{āpadaḥ та saṃpadaḥ позначають нещастя і достаток; kāle — «у свій час», daivād eva — «саме від долі». iti niścayī позначає переконання, як у 11.1–11.2. Це філософська позиція персонажа, не емпіричне твердження, що всі життєві події спричиняє доля.}
\vcom{Рядок б}
\vdeva{तृप्तः स्वस्थेन्द्रियो नित्यं न वाञ्छति न शोचति}
\viast{tṛptaḥ svasthendriyo nityaṃ na vāñchati na śocati}
\vtranslit{тріптах свастхендрійо нітям на ваньччаті на шочаті}
\vperek{той вдоволений, із чуттями в порядку, повсякчас не бажає й не сумує.}
\vcom{tṛptaḥ — «вдоволений», svasthendriyaḥ — «із чуттями у стані рівноваги», nityam — повсякчас. Пара na vāñchati / na śocati заперечує і прагнення, і сумування; український повтор «не ... й не ...» збережено без нової обіцянки абсолютної нечутливості.}
\vcom{`daiva` — «доля» як те, що розподіляє, а не як примха. Формула «не бажає й не сумує» — та сама пара, що в 8.1–8.2, і переклад має її впізнавати. `svasthendriya` («із чуттями в порядку») — рідкісний складник: ідеться не про приборкання чуттів, а про їхню справність.}
\vpar{Українська паралель (наша інтерпретація)}{Пантелеймон Куліш, «Чорна рада» (1857).}{Наша паралель: у Куліша доля теж постає як те, що роздає нещастя й достаток не за заслугою. Наше порівняння — за будовою думки.}
\vfig{ilustratsii/11-3-balance.pdf}{Дві шальки: нещастя й достаток у свій час}{Схема наша, не історична ілюстрація.}

\vsec{11.4}
\vdeva{सुखदुःखे जन्ममृत्यू दैवादेवेति निश्चयी । साध्यादर्शी निरायासः कुर्वन्नपि न लिप्यते ॥}
\viast{sukhaduḥkhe janmamṛtyū daivādeveti niścayī | sādhyādarśī nirāyāsaḥ kurvannapi na lipyate ||}
\vtranslit{сукхадухкхе джанмамрітю дайвадеветі нішчаї | садгядарші ніраясах курваннапі на ліп'яте ||}
\vperek{«Щастя й страждання, народження й смерть — саме від долі» — хто має таке переконання, той, не бачачи нічого досяжного, без натуги, навіть діючи, не забруднюється.}
\vcom{Построкове зіставлення: редакторська чернетка, не рецензовано.}
\vcom{Рядок а}
\vdeva{सुखदुःखे जन्ममृत्यू दैवादेवेति निश्चयी}
\viast{sukhaduḥkhe janmamṛtyū daivādeveti niścayī}
\vtranslit{сукхадухкхе джанмамрітю дайвадеветі нішчаї}
\vperek{«Щастя й страждання, народження й смерть — саме від долі» — хто має таке переконання,}
\vcom{sukhaduḥkhe — двоїна «щастя і страждання», janmamṛtyū — пара «народження і смерть», daivād eva — «саме від долі». Рефрен iti niścayī повторює формулу 11.3; перелік доповнює попереднє зіставлення протилежностей.}
\vcom{Рядок б}
\vdeva{साध्यादर्शी निरायासः कुर्वन्नपि न लिप्यते}
\viast{sādhyādarśī nirāyāsaḥ kurvannapi na lipyate}
\vtranslit{садгядарші ніраясах курваннапі на ліп'яте}
\vperek{той, не бачачи нічого досяжного, без натуги, навіть діючи, не забруднюється.}
\vcom{sādhyādarśī потребує обережної передачі як «той, хто не бачить досяжної мети»; nirāyāsaḥ — «без натуги», kurvann api na lipyate — «навіть діючи, не забруднюється». Тут розрізнено відсутність чіпляння за результат і фізичну відмову від будь-якої дії.}
\vcom{`sādhyādarśī` («той, хто не бачить досяжного») — найтонше слово глави: дія не припиняється, але зникає те, заради чого вона нібито робиться. `kurvann api na lipyate` — «навіть діючи, не забруднюється»: це не відмова від дії, а її розпредметнення. Перелік розширюється від 11.3 до 11.4 (додано народження й смерть), і це наростання навмисне.}
\vpar{Українська паралель (наша інтерпретація)}{Тарас Шевченко, «Давидові псалми» (1845).}{Наша паралель: у Шевченкових переспівах псалмів щастя й страждання теж не міряються людською міркою. Наше порівняння — типологічне.}
\vfig{ilustratsii/11-4-balance.pdf}{Дві шальки: народження й смерть, щастя й страждання}{Схема наша, не історична ілюстрація.}

\vsec{11.5}
\vdeva{चिन्तया जायते दुःखं नान्यथेहेति निश्चयी । तया हीनः सुखी शान्तः सर्वत्र गलितस्पृहः ॥}
\viast{cintayā jāyate duḥkhaṃ nānyatheheti niścayī | tayā hīnaḥ sukhī śāntaḥ sarvatra galitaspṛhaḥ ||}
\vtranslit{чінтая джаяте духкхам нанятхехеті нішчаї | тая хінах сукхі шантах сарватра ґалітаспріхах ||}
\vperek{«Страждання народжується з думання, не інакше» — хто має таке переконання, той, позбавлений його, щасливий, втихлий, скрізь із опалою спрагою.}
\vcom{Построкове зіставлення: редакторська чернетка, не рецензовано.}
\vcom{Рядок а}
\vdeva{चिन्तया जायते दुःखं नान्यथेहेति निश्चयी}
\viast{cintayā jāyate duḥkhaṃ nānyatheheti niścayī}
\vtranslit{чінтая джаяте духкхам нанятхехеті нішчаї}
\vperek{«Страждання народжується з думання, не інакше» — хто має таке переконання,}
\vcom{cintayā — орудний від cintā: турботливе мисленнєве обертання; jāyate duḥkhaṃ — «народжується страждання», nānyathā — «не інакше». Рефрен iti niścayī подає це як висловлене переконання, а не універсальне медичне пояснення страждання.}
\vcom{Рядок б}
\vdeva{तया हीनः सुखी शान्तः सर्वत्र गलितस्पृहः}
\viast{tayā hīnaḥ sukhī śāntaḥ sarvatra galitaspṛhaḥ}
\vtranslit{тая хінах сукхі шантах сарватра ґалітаспріхах}
\vperek{той, позбавлений його, щасливий, втихлий, скрізь із опалою спрагою.}
\vcom{tayā hīnaḥ — «позбавлений її» (попередньої cintā), sukhī śāntaḥ — «щасливий, втихлий», sarvatra galitaspṛhaḥ — «скрізь із згаслим прагненням». Це не твердження про шкідливість інтелектуальної праці; в контексті йдеться про прив'язливе хвилювання.}
\vcom{`cintā` — не «думка» взагалі, а **турботливе обертання думки**, і саме воно назване джерелом страждання. Переклад «думання» тут кращий за «думка», бо остання звучить нейтрально. `galitaspṛha` («із опалою спрагою») — образ, який текст уже вживав у 11.2 (`antargalita`); словник лишається спільним.}
\vpar{Українська паралель (наша інтерпретація)}{Памфіл Юркевич, «Серце і його значення в духовному житті людини» (1860).}{Наша паралель: коли вірш каже, що страждання народжується з думання, він ставить джерело болю в думку, а не в річ. Юркевич шукає це джерело в серці — наше порівняння за способом ставити питання.}
\vfig{ilustratsii/11-5-flame.pdf}{Полумʼя, що живиться думанням}{Схема наша, не історична ілюстрація.}

\vsec{11.6}
\vdeva{नाहं देहो न मे देहो बोधो ऽहम् इति निश्चयी । कैवल्यम् इव संप्राप्तो न स्मरत्य् अकृतं कृतम् ॥}
\viast{nāhaṃ deho na me deho bodho 'ham iti niścayī | kaivalyam iva saṃprāpto na smaraty akṛtaṃ kṛtam ||}
\vtranslit{нахам дехо на ме дехо бодго ’хам іті нішчаї | кайвалям іва сампрапто на смаратй акрітам крітам ||}
\vperek{«Я не тіло, і тіло не моє; я — усвідомлення» — хто має таке переконання, той, наче досягши окремішності, не пам'ятає незробленого як зробленого.}
\vcom{Построкове зіставлення: редакторська чернетка, не рецензовано.}
\vcom{Рядок а}
\vdeva{नाहं देहो न मे देहो बोधो ऽहम् इति निश्चयी}
\viast{nāhaṃ deho na me deho bodho 'ham iti niścayī}
\vtranslit{нахам дехо на ме дехо бодго ’хам іті нішчаї}
\vperek{«Я не тіло, і тіло не моє; я — усвідомлення» — хто має таке переконання,}
\vcom{nāhaṃ deho, na me deho — два заперечення ототожнення із тілом та володіння ним; bodho 'ham — «я — усвідомлення». iti niścayī повторює рамку глави; твердження стосується самопізнання в адвайті, а не фактичної відсутності тіла.}
\vcom{Рядок б}
\vdeva{कैवल्यम् इव संप्राप्तो न स्मरत्य् अकृतं कृतम्}
\viast{kaivalyam iva saṃprāpto na smaraty akṛtaṃ kṛtam}
\vtranslit{кайвалям іва сампрапто на смаратй акрітам крітам}
\vperek{той, наче досягши окремішності, не пам'ятає незробленого як зробленого.}
\vcom{kaivalyam iva saṃprāptaḥ буквально містить iva («наче»), тому стан не проголошено беззастережно досягнутим. na smaraty akṛtaṃ kṛtam — стислий і неоднозначний вислів про зроблене й незроблене; чинний український переклад дослівно збережено, але він потребує фахового зіставлення, а не вигаданої однозначної психологічної інтерпретації.}
\vcom{Формула «не тіло, і тіло не моє» та сама, що в 8.4 і 2.22 — тут вона подана як переконання, а не як висновок. `kaivalya` («окремішність») з'являється один раз і з `iva` («наче»): текст не каже, що стан досягнуто. `akṛtaṃ kṛtam` — «незроблене як зроблене»: ідеться не про забудькуватість, а про те, що минуле більше не рахується.}
\vpar{Українська паралель (наша інтерпретація)}{Григорій Сковорода, «Наркіс. Розмова про те: Пізнай себе».}{Наша паралель: у Сковороди справжня людина теж не збігається з тілом і не належить йому. Наше порівняння — типологічне, не про вплив.}
\vfig{ilustratsii/11-6-neti.pdf}{Тіло й володіння тілом відкинуті: лишається усвідомлення}{Схема наша, не історична ілюстрація.}

\vsec{11.7}
\vdeva{आब्रह्मस्तम्बपर्यन्तम् अहम् एवेति निश्चयी । निर्विकल्पः शुचिः शान्तः प्राप्ताप्राप्तविनिर्वृतः ॥}
\viast{ābrahmastambaparyantam aham eveti niścayī | nirvikalpaḥ śuciḥ śāntaḥ prāptāprāptavinirvṛtaḥ ||}
\vtranslit{абрахмастамбапар'янтам ахам еветі нішчаї | нірвікалпах шучіх шантах праптапраптавінірврітах ||}
\vperek{«Від Брахми до стеблини — це саме я» — хто має таке переконання, той без домислювання, чистий, втихлий, вдоволений і досягнутим, і недосягнутим.}
\vcom{Построкове зіставлення: редакторська чернетка, не рецензовано.}
\vcom{Рядок а}
\vdeva{आब्रह्मस्तम्बपर्यन्तम् अहम् एवेति निश्चयी}
\viast{ābrahmastambaparyantam aham eveti niścayī}
\vtranslit{абрахмастамбапар'янтам ахам еветі нішчаї}
\vperek{«Від Брахми до стеблини — це саме я» — хто має таке переконання,}
\vcom{ābrahmastambaparyantam окреслює межу «від Брахми до стеблини», aham eva — «це саме я», iti niścayī — «хто має таке переконання». Це метафізична формула всеохопності, а не біологічне ототожнення окремої людини з рослиною.}
\vcom{Рядок б}
\vdeva{निर्विकल्पः शुचिः शान्तः प्राप्ताप्राप्तविनिर्वृतः}
\viast{nirvikalpaḥ śuciḥ śāntaḥ prāptāprāptavinirvṛtaḥ}
\vtranslit{нірвікалпах шучіх шантах праптапраптавінірврітах}
\vperek{той без домислювання, чистий, втихлий, вдоволений і досягнутим, і недосягнутим.}
\vcom{nirvikalpaḥ — «без домислювання або розрізнювальних конструкцій», śuciḥ і śāntaḥ — «чистий і втихлий»; prāptāprāptavinirvṛtaḥ об'єднує «досягнуте» й «недосягнуте». Чинну українську пару збережено як етичний образ неприв'язаності до набуття.}
\vcom{`ābrahmastambaparyantam` — та сама формула межі, що в 2.11 (`brahmādistambaparyantam`), і переклад має бути однаковий. `prāptāprāpta-vinirvṛta` — «вдоволений і досягнутим, і недосягнутим»: це не байдужість, а відсутність різниці між двома станами. `śuci` тут — не ритуальна чистота, а відсутність домішки, бо стоїть поруч із `nirvikalpa`.}
\vpar{Українська паралель (наша інтерпретація)}{Іван Вишенський, «Книжка» (кінець XVI ст.).}{Наша паралель: у Вишенського межа між високим і низьким у світі теж знімається перед тим, що одне. Наше порівняння — типологічне.}
\vfig{ilustratsii/11-7-neti.pdf}{Від Брахми до стеблини — це саме я}{Схема наша, не історична ілюстрація.}

\vsec{11.8}
\vdeva{नानाश्चर्यम् इदं विश्वं न किञ्चिद् इति निश्चयी । निर्वासनः स्फूर्तिमात्रो न किञ्चिद् इव शाम्यति ॥}
\viast{nānāścaryam idaṃ viśvaṃ na kiñcid iti niścayī | nirvāsanaḥ sphūrtimātro na kiñcid iva śāmyati ||}
\vtranslit{нанашчар'ям ідам вішвам на кіньчід іті нішчаї | нірвасанах спхуртіматро на кіньчід іва шам'яті ||}
\vperek{«Цей світ, різноманітно дивовижний, — аніщо» — хто має таке переконання, той без усталених схильностей, саме лише сяйнування, втихає, наче аніщо.}
\vcom{Построкове зіставлення: редакторська чернетка, не рецензовано.}
\vcom{Рядок а}
\vdeva{नानाश्चर्यम् इदं विश्वं न किञ्चिद् इति निश्चयी}
\viast{nānāścaryam idaṃ viśvaṃ na kiñcid iti niścayī}
\vtranslit{нанашчар'ям ідам вішвам на кіньчід іті нішчаї}
\vperek{«Цей світ, різноманітно дивовижний, — аніщо» — хто має таке переконання,}
\vcom{nānāścaryam idaṃ viśvaṃ — «цей різноманітно дивовижний світ», na kiñcid — «аніщо»; рамка iti niścayī завершує восьмикратний повтор глави 11. Українське «аніщо» є перекладацьким переданням філософського судження, не запереченням емпіричних об'єктів.}
\vcom{Рядок б}
\vdeva{निर्वासनः स्फूर्तिमात्रो न किञ्चिद् इव शाम्यति}
\viast{nirvāsanaḥ sphūrtimātro na kiñcid iva śāmyati}
\vtranslit{нірвасанах спхуртіматро на кіньчід іва шам'яті}
\vperek{той без усталених схильностей, саме лише сяйнування, втихає, наче аніщо.}
\vcom{nirvāsanaḥ — «без закріплених схильностей», sphūrtimātro — «саме лише сяйнування / виявлення», na kiñcid iva śāmyati — «втихає, наче аніщо». Частка iva («наче») має залишитися помітною: вірш не оголошує людину буквально неіснуючою.}
\vcom{Глава закінчується тим самим словом, що й почалася в 11.1: `niścayī`, але тепер воно стоїть при «аніщо». `nirvāsana` («без усталених схильностей») — те саме `vāsanā`, що в 9.8; словник спільний. `sphūrtimātra` («саме лише сяйнування») і `na kiñcid iva` («наче аніщо») разом дають обережність, яку легко втратити: текст не каже, що людина зникає, — він каже «наче».}
\vpar{Українська паралель (наша інтерпретація)}{Богдан-Ігор Антонич, «Зелена Євангелія» (1938).}{Наша паралель: в Антонича світ теж постає дивовижним і водночас не вартим того, щоб за нього триматися. Наше порівняння — за образом.}
\vfig{ilustratsii/11-8-illusion.pdf}{Дивовижний світ під серпанком: аніщо}{Схема наша, не історична ілюстрація.}


% ЗГЕНЕРОВАНО scripts/build_format_v2_chapter.py — не редагувати вручну.
% Джерело: editorial/cards/12.*.json, format-v2/rozdil-12.json, paraleli.json.

\naukchapter{12}{8}{Глава 12. Саме в цьому я стою}
\naukintro{Дванадцята глава вся закінчується тими самими словами — саме в цьому я стою, — і цей рефрен тут не приспів, а мірило: кожен вірш спершу перелічує, чого людина вже не робить, і лише потім каже, де вона стоїть. Перші три вірші йдуть від тіла до мови й далі до думки; наступні три знімають пари — приймати й відкидати, обряд і його припинення. Останній вірш не додає нічого нового: він просто називає того, хто так стоїть.}
\vsec{12.1}
\vdeva{कायकृत्यासहः पूर्वं ततो वाग्विस्तरासहः । अथ चिन्तासहस्तस्माद् एवम् एवाहम् आस्थितः ॥}
\vtranslit{каякрітясахах пурвам тато ваґвістарасахах | атха чінтасахастасмад евам евахам астхітах ||}
\vperek{Спершу нездатний до тілесних справ, потім нездатний до поширення мови, далі нездатний до думання — тому саме в цьому я стою.}
\vcom{Тут відмова йде **не за бажанням, а через неспроможність**: `asahaḥ` — «нездатний витримати», а не «той, хто відмовився». Це рідкісний для тексту випадок, коли стан описано знизу — через те, чого людина вже не може. Формула `evam evāham āsthitaḥ` повторюється в 12.1–12.4, як `niścayī` у главі 11: переклад мусить тримати її однаково.}
\vpar{Українська паралель (наша інтерпретація)}{Василь Стус, «Час творчості».}{Наша паралель: у Стуса нездатність до звичних справ теж постає не як втрата, а як те, у чому людина нарешті стоїть. Наше порівняння — типологічне.}
\vfig{ilustratsii/12-1-neti.pdf}{Нездатність до тіла, мови й думки — і стояння в цьому}{Схема наша, не історична ілюстрація.}

\vsec{12.2}
\vdeva{प्रीत्यभावेन शब्दादेर् अदृश्यत्वेन चात्मनः । विक्षेपैकाग्र्यहृदय एवम् एवाहम् आस्थितः ॥}
\vtranslit{прітябгавена шабдадер адрішятвена чатманах | вікшепайкаґр'яхрідая евам евахам астхітах ||}
\vperek{Через відсутність прихильності до звуку й подібного та через невидимість його для самого себе — серце, для якого розсіяність і зосередженість одне — саме в цьому я стою.}
\vcom{`vikṣepa-aikāgrya` в одному складнику — найважливіше місце вірша: розсіяність і зосередженість названі **одним** словом, і це не поетична вільність, а твердження, що для того, хто стоїть у собі, різниці між ними немає. `adṛśyatva` («невидимість») стосується `ātman` — себе не видно собі як об'єкт.}
\vpar{Українська паралель (наша інтерпретація)}{Леся Українка, «У пущі» (1909).}{Наша паралель: у Лесі Українки зосередженість і розсіяність теж виявляються одним станом, а не двома. Наше порівняння — за будовою думки.}
\vfig{ilustratsii/12-2-witness.pdf}{Серце в центрі: розсіяність і зосередженість — одне}{Схема наша, не історична ілюстрація.}

\vsec{12.3}
\vdeva{समाध्यासादिविक्षिप्तौ व्यवहारः समाधये । एवं विलोक्य नियमम् एवम् एवाहम् आस्थितः ॥}
\vtranslit{самадгясадівікшіптау в'явахарах самадгає | евам вілокя ніямам евам евахам астхітах ||}
\vperek{Коли розсіяність походить від прагнення самадгі та подібного, то й звичайна діяльність — задля самадгі. Так побачивши це правило, саме в цьому я стою.}
\vcom{Найдошкульніший вірш глави: **прагнення самадгі саме стає розсіяністю** (`vikṣipti`). І далі — іронія, яку легко проґавити: якщо навіть буденна діяльність робиться «задля самадгі», то правило (`niyama`) виявляється всеохопним, і виходу з нього немає. Переклад мусить зберегти цю замкненість, а не зробити з неї пораду.}
\vpar{Українська паралель (ПОТРЕБУЄ ПЕРЕВІРКИ)}{українська релігійна проза, не визначено.}{ПОТРЕБУЄ ПЕРЕВІРКИ: думка, що звичайна діяльність і самадгі — одне, близька до української релігійної прози, але конкретного твору ми не звірили.}
\vfig{ilustratsii/12-3-balance.pdf}{Дві шальки: прагнення самадгі й звичайна діяльність}{Схема наша, не історична ілюстрація.}

\vsec{12.4}
\vdeva{हेयोपादेयविरहाद् एवं हर्षविषादयोः । अभावाद् अद्य हे ब्रह्मन् नेवम् एवाहम् आस्थितः ॥}
\vtranslit{хейопадеявірахад евам харшавішадайох | абгавад адя хе брахман невам евахам астхітах ||}
\vperek{Через відсутність того, що відкидати й приймати, а отже — через відсутність радості й смутку, нині, о Брахмане, саме в цьому я стою.}
\vcom{`he brahman` — звертання до Джанаки як до брахмана, єдине таке в розділі. Ланцюг причинний: немає пари «відкидати/приймати» → немає радості й смутку → нема чого досягати. **Відкрите питання:** у SARIT наприкінці стоїть `nevamevāhamāsthitaḥ` (sandhi `na + evameva`), і я прочитав це як ствердне `evameva` за паралелізмом. Перевірка решти глави це підтверджує лише частково: 12.5, 12.6 і 12.7 усі закінчуються на `evamevāhamāsthitaḥ` **без** `na`, тобто 12.4 — єдиний виняток у главі. Отже `na` може бути або справжнім (і тоді сенс вірша протилежний), або помилкою транскрипції. **GRETIL недоступний, другого свідка немає** — картка лишається чернеткою, а не «перевірено».}
\vpar{Українська паралель (наша інтерпретація)}{Ліна Костенко, «Маруся Чурай» (1979).}{Наша паралель: у Костенко теж є стан, у якому немає що відкидати й немає чого приймати, і саме в ньому людина стоїть. Наше порівняння — типологічне.}
\vfig{ilustratsii/12-4-neti.pdf}{Відкидання й приймання, радість і смуток відкинуті}{Схема наша, не історична ілюстрація.}

\vsec{12.5}
\vdeva{आश्रमानाश्रमं ध्यानं चित्तस्वीकृतवर्जनम् । विकल्पं मम वीक्ष्य एतैर् एवम् एवाहम् आस्थितः ॥}
\vtranslit{ашраманашрамам дгянам чіттасвікрітаварджанам | вікалпам мама вікшя етайр евам евахам астхітах ||}
\vperek{Медитація зі стадією життя і без неї, прийняття й відкидання, зроблені умом, — усе це домислювання. Побачивши це в собі, через саме це я так стою.}
\vcom{Єдине місце в книзі, де `āśrama` (стадія життя) названо **поряд з її відсутністю** — і обидва варіанти відкинуті як `vikalpa`. Це важливо для нашого словника: текст не обирає між «життям у стадії» і «поза нею», він знімає саму пару. `citta-svīkṛta-varjana` — прийняття й відкидання не речей, а **вироблені умом**; переклад мусить зберегти, що йдеться про акт ума, а не про оцінку світу.}
\vpar{Українська паралель (ПОТРЕБУЄ ПЕРЕВІРКИ)}{українська філософська проза, не визначено.}{ПОТРЕБУЄ ПЕРЕВІРКИ: думка, що медитація й обряди є домислом, найближча до української філософської прози, але назвати твір без звірки ми не можемо.}
\vfig{ilustratsii/12-5-knot.pdf}{Вузол домислу: медитація й обряд}{Схема наша, не історична ілюстрація.}

\vsec{12.6}
\vdeva{कर्मानुष्ठानम् अज्ञानाद् यथैव उपरमस् तथा । बुद्ध्वा सम्यग् इदं तत्त्वम् एवम् एवाहम् आस्थितः ॥}
\vtranslit{кармануштханам аджнянад ятхайва упарамас татха | буддгва сам'яґ ідам таттвам евам евахам астхітах ||}
\vperek{Виконання обрядів через незнання і припинення їх — однаково. Зрозумівши правильно цю суть, я так стою.}
\vcom{`yathaiva … tathā` («так само») ставить дію й відмову від дії на один рівень — і обидві названі `ajñāna`, коли виходять із незнання. Це не заклик припинити обряди: текст каже, що припинення не краще за виконання. Переклад мусить звучати рівно, без інтонації «краще не робити».}
\vpar{Українська паралель (ПОТРЕБУЄ ПЕРЕВІРКИ)}{українське бароко, не визначено.}{ПОТРЕБУЄ ПЕРЕВІРКИ: рівність обряду й припинення обряду — думка, близька до барокової полеміки, але конкретного тексту ми не перевірили.}
\vfig{ilustratsii/12-6-balance.pdf}{Дві шальки: обряд і його припинення — однаково}{Схема наша, не історична ілюстрація.}

\vsec{12.7}
\vdeva{अचिन्त्यं चिन्त्यमानो ऽपि चिन्तारूपं भजत्य् असौ । त्यक्त्वा तद्भावनं तस्माद् एवम् एवाहम् आस्थितः ॥}
\vtranslit{ачінтям чінтямано ’пі чінтарупам бгаджатй асау | тяктва тадбгаванам тасмад евам евахам астхітах ||}
\vperek{Незбагненне, навіть коли його осмислюють, набирає форми думки. Тому, відкинувши те плекання, я так стою.}
\vcom{Механізм названий точно: незбагненне не залишається собою, коли його осмислюють — воно **набирає форми думки** (`cintārūpa`). Тому відкидається не предмет, а `bhāvana` — плекання, повторюване обертання. Це та сама `cintā`, що в 11.5, і словник має бути спільним.}
\vpar{Українська паралель (наша інтерпретація)}{Григорій Сковорода, «Алфавіт, чи Буквар миру» (діалог).}{Наша паралель: у Сковороди незбагненне теж набирає форми думки, і саме тому його не можна втримати плеканням. Наше порівняння — типологічне.}
\vfig{ilustratsii/12-7-illusion.pdf}{Незбагненне під серпанком набирає форми думки}{Схема наша, не історична ілюстрація.}

\vsec{12.8}
\vdeva{एवम् एव कृतं येन स कृतार्थो भवेद् असौ । एवम् एव स्वभावो यः स कृतार्थो भवेद् असौ ॥}
\vtranslit{евам ева крітам єна са крітартхо бгавед асау | евам ева свабгаво ях са крітартхо бгавед асау ||}
\vperek{У кого це так зроблено — той стає тим, чию мету здійснено; у кого така природа — той стає тим, чию мету здійснено.}
\vcom{Вірш повторює рядок двічі й міняє одне слово: `kṛtam` (зроблено) → `svabhāvaḥ` (природа). Це не тавтологія, а розрізнення: спершу те, що людина зробила, потім те, чим вона стала. `kṛtārtha` — «той, чию мету здійснено», усталене слово; переклад тримає його однаково в обох половинах. Глава закінчується не настановою, а констатацією.}
\vpar{Українська паралель (наша інтерпретація)}{Микола Зеров, «Камена» (1924).}{Наша паралель: у Зерова завершеність теж постає як наслідок способу бачити, а не як наслідок зусилля. Наше порівняння — за будовою думки.}
\vfig{ilustratsii/12-8-flame.pdf}{Полумʼя здійсненої мети}{Схема наша, не історична ілюстрація.}


% ЗГЕНЕРОВАНО scripts/build_format_v2_chapter.py — не редагувати вручну.
% Джерело: editorial/cards/13.*.json, format-v2/rozdil-13.json, paraleli.json.

\naukchapter{13}{7}{Глава 13. Я сиджу, як мені завгодно}
\naukintro{Тринадцята глава — сім віршів і один рефрен, який стоїть у кінці кожного: я сиджу, як мені завгодно. Слово сиджу тут не про позу: глава саме й показує, що поза нічого не додає (13.5). Рефрен тримає іншу думку — що стан того, хто пізнав, не залежить ні від дії, ні від недіяння, і саме тому його не можна ані здобути, ані втратити.}
\vsec{13.1}
\vdeva{अकिंचनभवं स्वास्थ्यं कौपीनत्वे ऽपि दुर्लभम् । त्यागादाने विहायास्माद् अहम् आसे यथासुखम् ॥}
\viast{akiṃcanabhavaṃ svāsthyaṃ kaupīnatve 'pi durlabham | tyāgādāne vihāyāsmād aham āse yathāsukham ||}
\vtranslit{акімчанабгавам свастхям каупінатве ’пі дурлабгам | тяґадане віхаясмад ахам асе ятхасукхам ||}
\vperek{Стан неволодіння — це бути в собі; навіть у жебрацькому поясі він рідкісний. Відкинувши віддавання й приймання, я сиджу, як мені завгодно.}
\vcom{Построкове зіставлення: редакторська чернетка, не рецензовано.}
\vcom{Рядок а}
\vdeva{अकिंचनभवं स्वास्थ्यं कौपीनत्वे ऽपि दुर्लभम्}
\viast{akiṃcanabhavaṃ svāsthyaṃ kaupīnatve 'pi durlabham}
\vtranslit{акімчанабгавам свастхям каупінатве ’пі дурлабгам}
\vperek{Стан неволодіння — це бути в собі; навіть у жебрацькому поясі він рідкісний.}
\vcom{akiṃcana-bhavaṃ називає стан неволодіння, а svāsthyaṃ — перебування у власному стані; це не медичне твердження про здоров’я. kaupīnatve 'pi durlabham підкреслює, що навіть зовнішня ознака аскези (стегнова пов’язка) не гарантує внутрішньої свободи.}
\vcom{Рядок б}
\vdeva{त्यागादाने विहायास्माद् अहम् आसे यथासुखम्}
\viast{tyāgādāne vihāyāsmād aham āse yathāsukham}
\vtranslit{тяґадане віхаясмад ахам асе ятхасукхам}
\vperek{Відкинувши віддавання й приймання, я сиджу, як мені завгодно.}
\vcom{tyāgādāne — пара «віддавання й приймання»; vihāya означає «відкинувши». Після неї aham āse yathāsukham — повторюваний рефрен глави («я сиджу, як мені завгодно»). Думка стосується ставлення до володіння, а не буквальної вимоги сидіти.}
\vcom{Найважливіше тут — svāsthya. Слово означає буквально «перебування в собі» і в сучасній мові значить «здоров'я»; у вірші воно вжите в першому сенсі, і це не метафора, а визначення: бути здоровим — значить бути в собі. Друга точність — kaupīna: жебрацький пояс названий не як приклад аскези, а як те, у чому svāsthya все одно рідкісне. Тобто зовнішня ознака відречення нічого не гарантує. Пара tyāga / ādāna («віддавання» / «приймання») — та сама, що heya / upādeya у 7.5, і тут її теж відкидають цілою, а не по частинах. РЕФРЕН ГЛАВИ. Шість із семи віршів (13.1, 13.3–13.7) закінчуються однаково: aham āse yathāsukham («я сиджу — як мені завгодно»). ās — не «бути», а саме СИДІТИ, перебувати на місці; yathāsukham — «як зручно, за бажанням». Разом це формула, яка скасовує будь-яку поставу як вимогу: не треба ні стояти, ні лежати правильно. 13.2 — єдиний вірш глави без цього рефрену, і його треба читати окремо.}
\vpar{Українська паралель (наша інтерпретація)}{Григорій Сковорода, «Асхань. Симфонія, наречена книга Асхань, про пізнання самого себе».}{Наша паралель: у Сковороди бути в собі теж означає не мати нічого свого — і саме тому цей стан не залежить від достатку. Наше порівняння — типологічне.}
\vfig{ilustratsii/13-1-neti.pdf}{Віддавання й приймання відкинуті: сиджу, як завгодно}{Схема наша, не історична ілюстрація.}
\vfig{ilustratsii-naratyvni/13-1-sydyt-iak-zavhodno.pdf}{Людина сидить під деревом, а жебрацький пояс і чаша лежать осторонь}{Наша ілюстрація до вірша 13.1 — не історичне зображення.}

\vsec{13.2}
\vdeva{कुत्रापि खेदः कायस्य जिह्वा कुत्रापि खिद्यते । मनः कुत्रापि तत् त्यक्त्वा पुरुषार्थे स्थितः सुखम् ॥}
\viast{kutrāpi khedaḥ kāyasya jihvā kutrāpi khidyate | manaḥ kutrāpi tat tyaktvā puruṣārthe sthitaḥ sukham ||}
\vtranslit{кутрапі кхедах каяся джіхва кутрапі кхідяте | манах кутрапі тат тяктва пурушартхе стхітах сукхам ||}
\vperek{Десь втомлюється тіло, десь втомлюється язик, десь розум — відкинувши це, я стою в меті людини, щасливо.}
\vcom{Построкове зіставлення: редакторська чернетка, не рецензовано.}
\vcom{Рядок а}
\vdeva{कुत्रापि खेदः कायस्य जिह्वा कुत्रापि खिद्यते}
\viast{kutrāpi khedaḥ kāyasya jihvā kutrāpi khidyate}
\vtranslit{кутрапі кхедах каяся джіхва кутрапі кхідяте}
\vperek{Десь втомлюється тіло, десь втомлюється язик,}
\vcom{kutrāpi повторено двічі («десь»), зіставляючи khedaḥ kāyasya («втому тіла») та jihvā ... khidyate («втомлюється язик», тобто орган мовлення). Перший піврядок не містить ще висновку щодо того, що слід відкинути.}
\vcom{Рядок б}
\vdeva{मनः कुत्रापि तत् त्यक्त्वा पुरुषार्थे स्थितः सुखम्}
\viast{manaḥ kutrāpi tat tyaktvā puruṣārthe sthitaḥ sukham}
\vtranslit{манах кутрапі тат тяктва пурушартхе стхітах сукхам}
\vperek{десь розум — відкинувши це, я стою в меті людини, щасливо.}
\vcom{manaḥ kutrāpi додає третю сферу — розум; tat tyaktvā («відкинувши це») відсилає до названих утом. puruṣārthe sthitaḥ sukham — перебування в людській меті або її здійсненні, щасливо; це окрема думка, і вірш 13.2 не закінчується звичайним рефреном aham āse yathāsukham.}
\vcom{ЄДИНИЙ вірш глави без рефрену aham āse yathāsukham, і це не випадково: тут ідеться не про те, як сидіти, а про те, ЩО втомлюється. Три рівні названі поспіль — тіло, язик, розум — і кожен втомлюється «десь», тобто не завжди й не однаково. Слово kheda («втома») м'якше за duḥkha («страждання») з глави 8: це не біль, а зношування. puruṣārtha («мета людини») вжите в однині, хоч традиційно їх чотири, — тобто йдеться не про вибір між ними, а про те, що всі вони вже не потрібні.}
\vpar{Українська паралель (наша інтерпретація)}{Олександр Довженко, «Щоденник» (1941–1956).}{Наша паралель: у Довженка втома тіла, мови й думки теж описується окремо, і саме розрізнення цих трьох втом показує, що людина не зводиться до жодної з них. Наше порівняння — за будовою думки.}
\vfig{ilustratsii/13-2-neti.pdf}{Три втоми відкинуті: лишається мета людини}{Схема наша, не історична ілюстрація.}

\vsec{13.3}
\vdeva{कृतं किम् अपि नैव स्याद् इति संचिन्त्य तत्त्वतः । यदा यत् कर्तुम् आयाति तत् कृत्वासे यथासुखम् ॥}
\viast{kṛtaṃ kim api naiva syād iti saṃcintya tattvataḥ | yadā yat kartum āyāti tat kṛtvāse yathāsukham ||}
\vtranslit{крітам кім апі найва сяд іті самчінтя таттватах | яда ят картум аяті тат крітвасе ятхасукхам ||}
\vperek{Подумавши по-справжньому: «зробленого не буде нічого» — що б не випало зробити, зробивши те, я сиджу, як мені завгодно.}
\vcom{Построкове зіставлення: редакторська чернетка, не рецензовано.}
\vcom{Рядок а}
\vdeva{कृतं किम् अपि नैव स्याद् इति संचिन्त्य तत्त्वतः}
\viast{kṛtaṃ kim api naiva syād iti saṃcintya tattvataḥ}
\vtranslit{крітам кім апі найва сяд іті самчінтя таттватах}
\vperek{Подумавши по-справжньому: «зробленого не буде нічого» —}
\vcom{kṛtaṃ kim api naiva syād задає заперечення результату дії; iti saṃcintya tattvataḥ означає «так розваживши по суті». Оптатив syāt не варто без пояснення підмінювати твердженням, що жодна дія фізично не відбувається.}
\vcom{Рядок б}
\vdeva{यदा यत् कर्तुम् आयाति तत् कृत्वासे यथासुखम्}
\viast{yadā yat kartum āyāti tat kṛtvāse yathāsukham}
\vtranslit{яда ят картум аяті тат крітвасе ятхасукхам}
\vperek{що б не випало зробити, зробивши те, я сиджу, як мені завгодно.}
\vcom{yadā yat kartum āyāti описує те, що випадає зробити; tat kṛtvāse поєднує «зробивши те» з «сиджу». Рефрен yathāsukham збережено у чинному перекладі: відсутність привласнення результату не заперечує самого виконання дії.}
\vcom{Парадокс у тому, що дія не скасовується: yadā yat kartum āyāti («коли що випадає зробити») — дієслово āyāti означає саме те, що приходить само, без запрошення. Отже, діють, але без того, хто міг би зарахувати її собі. kṛtvāse — злитна форма kṛtvā + āse, і в ній дві дії стоять впритул: зробивши — сиджу. Умова naiva syāt («нехай не буде») стоїть у optativus, тобто це не констатація, а спосіб бачення, прийнятий заздалегідь. РЕФРЕН ГЛАВИ. Шість із семи віршів (13.1, 13.3–13.7) закінчуються однаково: aham āse yathāsukham («я сиджу — як мені завгодно»). ās — не «бути», а саме СИДІТИ, перебувати на місці; yathāsukham — «як зручно, за бажанням». Разом це формула, яка скасовує будь-яку поставу як вимогу: не треба ні стояти, ні лежати правильно. 13.2 — єдиний вірш глави без цього рефрену, і його треба читати окремо.}
\vpar{Українська паралель (ПОТРЕБУЄ ПЕРЕВІРКИ)}{українська філософська лірика, не визначено.}{ПОТРЕБУЄ ПЕРЕВІРКИ: думка, що зробленого не буде нічого, близька до української філософської лірики, але конкретного твору ми не звірили.}
\vfig{ilustratsii/13-3-space.pdf}{Небо без хмар: зробленого не буде нічого}{Схема наша, не історична ілюстрація.}

\vsec{13.4}
\vdeva{कर्मनैष्कर्म्यनिर्बन्धभावा देहस्थयोगिनः । संयोगायोगविरहाद् अहम् आसे यथासुखम् ॥}
\viast{karmanaiṣkarmyanirbandhabhāvā dehasthayoginaḥ | saṃyogāyogavirahād aham āse yathāsukham ||}
\vtranslit{карманайшкарм'янірбандгабгава дехастхайоґінах | самйоґайоґавірахад ахам асе ятхасукхам ||}
\vperek{У йога, що перебуває в тілі, є стани, до яких його змушують дія й недіяння; та через відсутність з'єднання й роз'єднання я сиджу, як мені завгодно.}
\vcom{Построкове зіставлення: редакторська чернетка, не рецензовано.}
\vcom{Рядок а}
\vdeva{कर्मनैष्कर्म्यनिर्बन्धभावा देहस्थयोगिनः}
\viast{karmanaiṣkarmyanirbandhabhāvā dehasthayoginaḥ}
\vtranslit{карманайшкарм'янірбандгабгава дехастхайоґінах}
\vperek{У йога, що перебуває в тілі, є стани, до яких його змушують дія й недіяння;}
\vcom{Складне слово karmanaiṣkarmyanirbandhabhāvā зіставляє дію (karma), недіяння (naiṣkarmya) й наполягання або примус (nirbandha); dehasthayoginaḥ вказує на йога, який перебуває в тілі. Це твердження про напругу між двома способами поведінки, а не доказ певної історичної школи.}
\vcom{Рядок б}
\vdeva{संयोगायोगविरहाद् अहम् आसे यथासुखम्}
\viast{saṃyogāyogavirahād aham āse yathāsukham}
\vtranslit{самйоґайоґавірахад ахам асе ятхасукхам}
\vperek{та через відсутність з'єднання й роз'єднання я сиджу, як мені завгодно.}
\vcom{saṃyogāyogavirahād — аблативна конструкція з ідеєю відсутності з'єднання й нез'єднання. Далі aham āse yathāsukham повертає рефрен особистого перебування; не слід перетворювати це на припис відмовитися від щоденних обов’язків.}
\vcom{Слово nirbandha («змушування») — ключове: доки є тіло, є й примус то до дії, то до недіяння, і жоден вибір не вільний. saṃyoga / ayoga («з'єднання» / «роз'єднання») — пара з мови йоги, де з'єднання зазвичай означає зв'язок з об'єктом; тут відкидається сама можливість такого зв'язку. Зверни увагу на порядок: спершу визнано примус (він є), потім названо те, чого немає (з'єднання). Вірш не обіцяє свободи від тіла — лише від рахунку. РЕФРЕН ГЛАВИ. Шість із семи віршів (13.1, 13.3–13.7) закінчуються однаково: aham āse yathāsukham («я сиджу — як мені завгодно»). ās — не «бути», а саме СИДІТИ, перебувати на місці; yathāsukham — «як зручно, за бажанням». Разом це формула, яка скасовує будь-яку поставу як вимогу: не треба ні стояти, ні лежати правильно. 13.2 — єдиний вірш глави без цього рефрену, і його треба читати окремо.}
\vpar{Українська паралель (ПОТРЕБУЄ ПЕРЕВІРКИ)}{українська релігійна проза, не визначено.}{ПОТРЕБУЄ ПЕРЕВІРКИ: розрізнення з'єднання й роз'єднання як того, що змушує людину до станів, найближче до української релігійної прози, але назвати твір без звірки ми не можемо.}
\vfig{ilustratsii/13-4-knot.pdf}{Вузол дії й недіяння розвʼязано}{Схема наша, не історична ілюстрація.}

\vsec{13.5}
\vdeva{अर्थानर्थौ न मे स्थित्या गत्या न शयनेन वा । तिष्ठन् गच्छन् स्वपन् तस्माद् अहम् आसे यथासुखम् ॥}
\viast{arthānarthau na me sthityā gatyā na śayanena vā | tiṣṭhan gacchan svapan tasmād aham āse yathāsukham ||}
\vtranslit{артханартхау на ме стхітя ґатя на шаянена ва | тіштхан ґаччан свапан тасмад ахам асе ятхасукхам ||}
\vperek{Ні користі, ні шкоди немає мені від стояння, ходіння чи лежання — тому, стоячи, ходячи чи сплячи, я сиджу, як мені завгодно.}
\vcom{Построкове зіставлення: редакторська чернетка, не рецензовано.}
\vcom{Рядок а}
\vdeva{अर्थानर्थौ न मे स्थित्या गत्या न शयनेन वा}
\viast{arthānarthau na me sthityā gatyā na śayanena vā}
\vtranslit{артханартхау на ме стхітя ґатя на шаянена ва}
\vperek{Ні користі, ні шкоди немає мені від стояння, ходіння чи лежання —}
\vcom{arthānarthau — двоїна «користь і шкода»; sthityā, gatyā, śayanena — орудні форми «стоянням, рухом, лежанням». Це оцінка способів тілесного перебування, а не заборона ходити чи лежати.}
\vcom{Рядок б}
\vdeva{तिष्ठन् गच्छन् स्वपन् तस्माद् अहम् आसे यथासुखम्}
\viast{tiṣṭhan gacchan svapan tasmād aham āse yathāsukham}
\vtranslit{тіштхан ґаччан свапан тасмад ахам асе ятхасукхам}
\vperek{тому, стоячи, ходячи чи сплячи, я сиджу, як мені завгодно.}
\vcom{tiṣṭhan, gacchan, svapan — три дієприкметники «стоячи, йдучи, сплячи»; tasmād — «тому». Слідом aham āse yathāsukham зберігає словесний парадокс «сиджу» попри зміну постави; рефрен уособлює сталість внутрішнього ставлення, а не конкретну фізичну позу.}
\vcom{Найгостріший парадокс глави, і він у самій граматиці: три дієприкметники (tiṣṭhan, gacchan, svapan — «стоячи, ходячи, сплячи») стоять поруч із дієсловом āse («сиджу»). Формально це суперечність, і вона навмисна: ās тут уже не про поставу тіла, а про перебування взагалі. Пара artha / anartha («користь» / «шкода») — та сама, що в 7.2 vṛddhi / kṣati, і знову відкидається цілою: не «немає шкоди», а немає ні того, ні того. РЕФРЕН ГЛАВИ. Шість із семи віршів (13.1, 13.3–13.7) закінчуються однаково: aham āse yathāsukham («я сиджу — як мені завгодно»). ās — не «бути», а саме СИДІТИ, перебувати на місці; yathāsukham — «як зручно, за бажанням». Разом це формула, яка скасовує будь-яку поставу як вимогу: не треба ні стояти, ні лежати правильно. 13.2 — єдиний вірш глави без цього рефрену, і його треба читати окремо.}
\vpar{Українська паралель (наша інтерпретація)}{Василь Стефаник, «Камінний хрест» (1900).}{Наша паралель: у Стефаника тілесне положення людини теж не міняє того, що з нею сталося всередині. Наше порівняння — типологічне.}
\vfig{ilustratsii/13-5-rest.pdf}{Людина під зайшлим сонцем: поза не додає ні користі, ні шкоди}{Схема наша, не історична ілюстрація.}

\vsec{13.6}
\vdeva{स्वपतो नास्ति मे हानिः सिद्धिर् यत्नवतो न वा । नाशोल्लासौ विहायास्माद् अहम् आसे यथासुखम् ॥}
\viast{svapato nāsti me hāniḥ siddhir yatnavato na vā | nāśollāsau vihāyāsmād aham āse yathāsukham ||}
\vtranslit{свапато насті ме ханіх сіддгір ятнавато на ва | нашолласау віхаясмад ахам асе ятхасукхам ||}
\vperek{У сплячого немає мені втрати, як і немає досягнення в того, хто докладає зусиль; відкинувши загибель і піднесення, я сиджу, як мені завгодно.}
\vcom{Построкове зіставлення: редакторська чернетка, не рецензовано.}
\vcom{Рядок а}
\vdeva{स्वपतो नास्ति मे हानिः सिद्धिर् यत्नवतो न वा}
\viast{svapato nāsti me hāniḥ siddhir yatnavato na vā}
\vtranslit{свапато насті ме ханіх сіддгір ятнавато на ва}
\vperek{У сплячого немає мені втрати, як і немає досягнення в того, хто докладає зусиль;}
\vcom{svapataḥ (у сандгі svapato) — родовий однини дієприкметника «той, хто спить»; na asti — «немає»; me — «мені / в мене»; hāniḥ — «втрата». siddhiḥ (у сандгі siddhir) — «успіх, досягнення»; yatnavataḥ (у сандгі yatnavato) — «того, хто докладає зусиль»; na vā завершує заперечення щодо досягнення. Це не наказ спати або відмовитися від праці: рядок зіставляє дві оцінки — втрату й здобуток.}
\vcom{Рядок б}
\vdeva{नाशोल्लासौ विहायास्माद् अहम् आसे यथासुखम्}
\viast{nāśollāsau vihāyāsmād aham āse yathāsukham}
\vtranslit{нашолласау віхаясмад ахам асе ятхасукхам}
\vperek{відкинувши загибель і піднесення, я сиджу, як мені завгодно.}
\vcom{nāśa-ullāsau — двоїна, «загибель і піднесення»; vihāya — дієприслівникова форма «відкинувши». Форма vihāyāsmād передає злиття vihāya + asmāt («від цього / тому»); точний синтаксичний зв'язок варто окремо перевірити за друкованими джерелами. aham āse yathāsukham — «я сиджу, як мені завгодно». Рефрен глави не скасовує потреби уважно розглядати межу слів і сандгі.}
\vcom{Вірш будується на запереченні двох протилежних чеснот: сон не шкодить, зусилля не дає. Тим самим знімається сама вісь, на якій тримається будь-яка практика, — «старанність». ullāsa («піднесення») — слово з відтінком радості, не просто зростання; у парі з nāśa («загибель») воно описує не події, а те, як людина їх переживає. Порівняй 3.9, де похвала й мука були описані однією граматичною формою: тут той самий хід, але на рівні оцінки. РЕФРЕН ГЛАВИ. Шість із семи віршів (13.1, 13.3–13.7) закінчуються однаково: aham āse yathāsukham («я сиджу — як мені завгодно»). ās — не «бути», а саме СИДІТИ, перебувати на місці; yathāsukham — «як зручно, за бажанням». Разом це формула, яка скасовує будь-яку поставу як вимогу: не треба ні стояти, ні лежати правильно. 13.2 — єдиний вірш глави без цього рефрену, і його треба читати окремо.}
\vpar{Українська паралель (наша інтерпретація)}{Ліна Костенко, «Берестечко» (1999).}{Наша паралель: у Костенко теж є стан, у якому ні сон, ні зусилля не дають ні втрати, ні досягнення. Наше порівняння — за будовою думки.}
\vfig{ilustratsii/13-6-balance.pdf}{Дві шальки: втрата й досягнення — обидві відкинуті}{Схема наша, не історична ілюстрація.}

\vsec{13.7}
\vdeva{सुखादिरूपा नियमं भावेष्व् आलोक्य भूरिशः । शुभाशुभे विहायास्माद् अहम् आसे यथासुखम् ॥}
\viast{sukhādirūpā niyamaṃ bhāveṣv ālokya bhūriśaḥ | śubhāśubhe vihāyāsmād aham āse yathāsukham ||}
\vtranslit{сукхадірупа ніямам бгавешв алокя бгурішах | шубгашубге віхаясмад ахам асе ятхасукхам ||}
\vperek{Побачивши не раз у всіх станах те правило, що має вигляд задоволення, — відкинувши добре й погане, я сиджу, як мені завгодно.}
\vcom{Построкове зіставлення: редакторська чернетка, не рецензовано.}
\vcom{Рядок а}
\vdeva{सुखादिरूपा नियमं भावेष्व् आलोक्य भूरिशः}
\viast{sukhādirūpā niyamaṃ bhāveṣv ālokya bhūriśaḥ}
\vtranslit{сукхадірупа ніямам бгавешв алокя бгурішах}
\vperek{Побачивши не раз у всіх станах те правило, що має вигляд задоволення, —}
\vcom{bhāveṣv — місцевий множини «у станах»; ālokya — дієприслівникова форма «побачивши»; bhūriśaḥ — «багаторазово». Чинна IAST-картка має niyamaṃ («правило»), що відображено в її українському перекладі; SARIT має aniyamaṃ («відсутність правила»). Це змістове різночитання не розв'язано, а sukhādirūpā також має проблемне граматичне узгодження. Построковість фіксує читання картки, не видає його за доведений оригінал.}
\vcom{Рядок б}
\vdeva{शुभाशुभे विहायास्माद् अहम् आसे यथासुखम्}
\viast{śubhāśubhe vihāyāsmād aham āse yathāsukham}
\vtranslit{шубгашубге віхаясмад ахам асе ятхасукхам}
\vperek{відкинувши добре й погане, я сиджу, як мені завгодно.}
\vcom{śubhāśubhe — двоїна «добре й погане»; vihāya — «відкинувши»; у vihāyāsmād видно сандгі з asmāt («із цього / тому»). aham āse — «я сиджу», а yathāsukham — «як мені зручно». Це повторюваний фінал глави 13, але він не вирішує різночитання niyamaṃ / aniyamaṃ у першому піврядку.}
\vcom{Найважче місце глави, і не через сенс, а через текст. Свідки розходяться в одному слові: SARIT подає aniyamaṃ («безладдя, відсутність правила»), GRETIL — niyamaṃ («правило»). Сенс при цьому стає ПРОТИЛЕЖНИМ: побачити, що все, що має вигляд задоволення, насправді безладдя — або що воно підлягає правилу. Крім того, форма sukhādirūpā (жіночий рід, однина) не узгоджується ні з niyamaṃ/aniyamaṃ (середній рід), ні з bhāveṣu (множина, чоловічий/середній). Найімовірніше, обидва свідки успадкували зіпсоване читання, і правильним було sukhādirūpāṇi (середній рід, множина) — «речі, чия подоба — задоволення тощо». У перекладі взято читання SARIT, але це ВИБІР, а не розв'язання: потрібен третій свідок. РЕФРЕН ГЛАВИ. Шість із семи віршів (13.1, 13.3–13.7) закінчуються однаково: aham āse yathāsukham («я сиджу — як мені завгодно»). ās — не «бути», а саме СИДІТИ, перебувати на місці; yathāsukham — «як зручно, за бажанням». Разом це формула, яка скасовує будь-яку поставу як вимогу: не треба ні стояти, ні лежати правильно. 13.2 — єдиний вірш глави без цього рефрену, і його треба читати окремо.}
\vpar{Українська паралель (ПОТРЕБУЄ ПЕРЕВІРКИ)}{українська барокова поезія, не визначено.}{ПОТРЕБУЄ ПЕРЕВІРКИ: правило, що має вигляд задоволення, близьке до барокового мотиву оманливої втіхи, але конкретного тексту ми не перевірили.}
\vfig{ilustratsii/13-7-illusion.pdf}{Серпанок задоволення: під ним — добре й погане}{Схема наша, не історична ілюстрація.}


% ЗГЕНЕРОВАНО scripts/build_format_v2_chapter.py — не редагувати вручну.
% Джерело: editorial/cards/14.*.json, format-v2/rozdil-14.json, paraleli.json.

\naukchapter{14}{4}{Глава 14. Порожній розум}
\naukintro{Чотирнадцята глава — чотири вірші, і всі вони про те саме з різних боків: коли спрага згасла, питання про набутки стають безпредметними. Перший вірш описує такого, хто наче заснув і прокинувся; другий перелічує те, чого вже нема про що питати; третій знімає саму турботу про визволення; четвертий визнає, що ззовні таку людину годі відрізнити від божевільної.}
\vsec{14.1}
\vdeva{प्रकृत्या शून्यचित्तो यः प्रमादाद् भावभावनः । निद्रितो बोधित इव क्षीणसंसरणो हि सः ॥}
\viast{prakṛtyā śūnyacitto yaḥ pramādād bhāvabhāvanaḥ | nidrito bodhita iva kṣīṇasaṃsaraṇo hi saḥ ||}
\vtranslit{пракрітя шунячітто ях прамадад бгавабгаванах | нідріто бодгіта іва кшінасамсарано хі сах ||}
\vperek{Хто за природою має порожній розум — навіть якщо з недбальства плекає [в собі] стани, — той наче заснув і прокинувся: блукання його вичерпане.}
\vcom{Построкове зіставлення: редакторська чернетка, не рецензовано.}
\vcom{Рядок а}
\vdeva{प्रकृत्या शून्यचित्तो यः प्रमादाद् भावभावनः}
\viast{prakṛtyā śūnyacitto yaḥ pramādād bhāvabhāvanaḥ}
\vtranslit{пракрітя шунячітто ях прамадад бгавабгаванах}
\vperek{Хто за природою має порожній розум — навіть якщо з недбальства плекає [в собі] стани,}
\vcom{Конструкція prakṛtyā śūnyacitto yaḥ характеризує того, чий розум за природою порожній; pramādād bhāvabhāvanaḥ додає застереження про недбалість та плекання станів. Це початок вислову, не встановлення окремої школи тлумачення.}
\vcom{Рядок б}
\vdeva{निद्रितो बोधित इव क्षीणसंसरणो हि सः}
\viast{nidrito bodhita iva kṣīṇasaṃsaraṇo hi saḥ}
\vtranslit{нідріто бодгіта іва кшінасамсарано хі сах}
\vperek{— той наче заснув і прокинувся: блукання його вичерпане.}
\vcom{Друга половина подає порівняння nidrito bodhita iva й завершення kṣīṇasaṃsaraṇo hi saḥ. Форма kṣīṇasaṃsaraṇo лишається читанням чинної картки; цифрове kṣīṇasaṃsmaraṇo іншого свідка не визнано тут помилкою чи рукописно вирішеним варіантом.}
\vcom{Вірш навмисно знімає вимогу старанності. pramādāt («через недбальство») — це не похвала, але й не докір: недбальство нічого не змінює, бо порожнеча розуму тут не досягнута, а природна (prakṛtyā). Складене bhāvabhāvana («плекання станів») описує звичну роботу розуму — і вона названа марною. nidritaḥ / bodhitaḥ iva («наче заснулий і розбуджений») — не два стани, а один, побачений із двох боків: людина прокинулась, але ще не ворушиться. БУДОВА ГЛАВИ. Глава 14 не має спільного рефрену — на відміну від 5, 6 і 13. Натомість усі чотири вірші тримаються однієї думки, поданої з чотирьох боків: хто такий цей чоловік зсередини (14.1), що для нього зникає (14.2), про що він уже не турбується (14.3) і як його бачать ззовні (14.4).}
\vpar{Українська паралель (наша інтерпретація)}{Євген Плужник, «Рання осінь» (1927).}{Наша паралель: у Плужника людина, що наче заснула й прокинулася, теж опиняється поза власним блуканням. Наше порівняння — типологічне.}
\vfig{ilustratsii/14-1-flame.pdf}{Полумʼя, що прокинулося: блукання вичерпане}{Схема наша, не історична ілюстрація.}

\vsec{14.2}
\vdeva{क्व धनानि क्व मित्राणि क्व मे विषयदस्यवः । क्व शास्त्रं क्व च विज्ञानं यदा मे गलिता स्पृहा ॥}
\viast{kva dhanāni kva mitrāṇi kva me viṣayadasyavaḥ | kva śāstraṃ kva ca vijñānaṃ yadā me galitā spṛhā ||}
\vtranslit{ква дганані ква мітрані ква ме вішаядасявах | ква шастрам ква ча віджнянам яда ме ґаліта спріха ||}
\vperek{Де багатства, де друзі, де мої слуги — чуттєві предмети? Де писання, де знання — коли в мене розтанула спрага?}
\vcom{Построкове зіставлення: редакторська чернетка, не рецензовано.}
\vcom{Рядок а}
\vdeva{क्व धनानि क्व मित्राणि क्व मे विषयदस्यवः}
\viast{kva dhanāni kva mitrāṇi kva me viṣayadasyavaḥ}
\vtranslit{ква дганані ква мітрані ква ме вішаядасявах}
\vperek{Де багатства, де друзі, де мої слуги — чуттєві предмети?}
\vcom{Повтор kva («де?») створює низку риторичних питань: про багатства, друзів і viṣayadasyavaḥ — «слуг», пов'язаних із чуттєвими предметами. Відповідник взято з чинного перекладу, без нового тлумачення складу.}
\vcom{Рядок б}
\vdeva{क्व शास्त्रं क्व च विज्ञानं यदा मे गलिता स्पृहा}
\viast{kva śāstraṃ kva ca vijñānaṃ yadā me galitā spṛhā}
\vtranslit{ква шастрам ква ча віджнянам яда ме ґаліта спріха}
\vperek{Де писання, де знання — коли в мене розтанула спрага?}
\vcom{Питання kva śāstraṃ kva ca vijñānaṃ продовжує перелік, а yadā me galitā spṛhā вводить умову: зникнення жаги. Не слід автоматично сприймати цей образ як історичну критику української науки чи освіти.}
\vcom{П'ять разів kva («де») — і жодної відповіді: питання не шукає місця, а показує, що місця немає. Зверни увагу на viṣayadasyavaḥ («раби предметів»): предмети названі не панами, а слугами, і це перевертає звичну картину — не вони тримають людину, а вона їх утримує. galitā («розтанула, опала») — слово про те, що сходить само, як пух або як лушпиння; спрага не переможена, вона зійшла. БУДОВА ГЛАВИ. Глава 14 не має спільного рефрену — на відміну від 5, 6 і 13. Натомість усі чотири вірші тримаються однієї думки, поданої з чотирьох боків: хто такий цей чоловік зсередини (14.1), що для нього зникає (14.2), про що він уже не турбується (14.3) і як його бачать ззовні (14.4).}
\vpar{Українська паралель (наша інтерпретація)}{Богдан-Ігор Антонич, «Зелена Євангелія» (1938).}{Наша паралель: в Антонича питання де багатства, де друзі теж обертається на порожній перелік, коли згасає спрага. Наше порівняння — за образом.}
\vfig{ilustratsii/14-2-neti.pdf}{Багатства, друзі, писання й знання відпадають}{Схема наша, не історична ілюстрація.}
\vfig{ilustratsii-naratyvni/14-2-porozhnii-dvir.pdf}{Порожній двір: обриси речей, друзів і слуг лишилися, а їх самих немає}{Наша ілюстрація до вірша 14.2 — не історичне зображення.}

\vsec{14.3}
\vdeva{विज्ञाते साक्षिपुरुषे परमात्मनि चेश्वरे । नैराश्ये बन्धमोक्षे च न चिन्ता मुक्तये मम ॥}
\viast{vijñāte sākṣipuruṣe paramātmani ceśvare | nairāśye bandhamokṣe ca na cintā muktaye mama ||}
\vtranslit{віджняте сакшіпуруше параматмані чешваре | найрашє бандгамокше ча на чінта муктає мама ||}
\vperek{Коли пізнано пурушу-свідка, вищого Атмана й Ішвару — а також безнадію, узу й визволення — у мене немає турботи про визволення.}
\vcom{Построкове зіставлення: редакторська чернетка, не рецензовано.}
\vcom{Рядок а}
\vdeva{विज्ञाते साक्षिपुरुषे परमात्मनि चेश्वरे}
\viast{vijñāte sākṣipuruṣe paramātmani ceśvare}
\vtranslit{віджняте сакшіпуруше параматмані чешваре}
\vperek{Коли пізнано пурушу-свідка, вищого Атмана й Ішвару —}
\vcom{Перший піврядок подає locativus vijñāte та низку означень: sākṣipuruṣe, paramātmani, ceśvare. Український текст картки передає свідка, вищого Атмана й Ішвару; питання їхнього доктринального співвідношення потребує окремого коментаря.}
\vcom{Рядок б}
\vdeva{नैराश्ये बन्धमोक्षे च न चिन्ता मुक्तये मम}
\viast{nairāśye bandhamokṣe ca na cintā muktaye mama}
\vtranslit{найрашє бандгамокше ча на чінта муктає мама}
\vperek{а також безнадію, узу й визволення — у мене немає турботи про визволення.}
\vcom{Другий піврядок згадує nairāśye, bandhamokṣe й заперечення турботи na cintā muktaye mama. Це редакторське граматичне пояснення наявного перекладу, не самостійна незалежна колація цифрових свідків.}
\vcom{Тут уперше в книзі з'являється складене sākṣipuruṣa — «пуруша-свідок». Це не те саме, що просто ātman: sākṣin додає саме те, що свідок не бере участі, а лише присутній при. Найтонше — перелік: у ньому стоять поруч Ішвара і nairāśya («безнадія»), уза й визволення, тобто протилежності названі в одному ряді й пізнані разом. Тому й na cintā muktaye («немає турботи про визволення»): коли визволення пізнане як одне з пари, турбуватися про нього нема про що. БУДОВА ГЛАВИ. Глава 14 не має спільного рефрену — на відміну від 5, 6 і 13. Натомість усі чотири вірші тримаються однієї думки, поданої з чотирьох боків: хто такий цей чоловік зсередини (14.1), що для нього зникає (14.2), про що він уже не турбується (14.3) і як його бачать ззовні (14.4).}
\vpar{Українська паралель (ПОТРЕБУЄ ПЕРЕВІРКИ)}{українська релігійна філософія, не визначено.}{ПОТРЕБУЄ ПЕРЕВІРКИ: перелік пізнаного, після якого зникає турбота про визволення, близький до української релігійної філософії, але конкретного місця ми не звірили.}
\vfig{ilustratsii/14-3-witness.pdf}{Свідок у центрі: уза й визволення втрачають вагу}{Схема наша, не історична ілюстрація.}

\vsec{14.4}
\vdeva{अन्तर्विकल्पशून्यस्य बहिः स्वच्छन्दचारिणः । भ्रान्तस्येव दशास् तास् तास् तादृशा एव जानते ॥}
\viast{antarvikalpaśūnyasya bahiḥ svacchandacāriṇaḥ | bhrāntasyeva daśās tās tās tādṛśā eva jānate ||}
\vtranslit{антарвікалпашуняся бахіх сваччандачарінах | бгрантасєва дашас тас тас тадріша ева джанате ||}
\vperek{У того, хто всередині порожній від розрізнень, а зовні ходить, як сам захоче, стани бувають усякі, наче в божевільного; і знають його тільки такі самі, як він.}
\vcom{Построкове зіставлення: редакторська чернетка, не рецензовано.}
\vcom{Рядок а}
\vdeva{अन्तर्विकल्पशून्यस्य बहिः स्वच्छन्दचारिणः}
\viast{antarvikalpaśūnyasya bahiḥ svacchandacāriṇaḥ}
\vtranslit{антарвікалпашуняся бахіх сваччандачарінах}
\vperek{У того, хто всередині порожній від розрізнень, а зовні ходить, як сам захоче,}
\vcom{Вираз antarvikalpaśūnyasya / bahiḥ svacchandacāriṇaḥ зіставляє внутрішню свободу від розрізнень та зовнішню довільність руху. Обидва означення стосуються тієї самої постаті; порівняння не є медичним описом.}
\vcom{Рядок б}
\vdeva{भ्रान्तस्येव दशास् तास् तास् तादृशा एव जानते}
\viast{bhrāntasyeva daśās tās tās tādṛśā eva jānate}
\vtranslit{бгрантасєва дашас тас тас тадріша ева джанате}
\vperek{стани бувають усякі, наче в божевільного; і знають його тільки такі самі, як він.}
\vcom{Порівняння bhrāntasyeva та форми daśās tās tās змальовують різні видимі стани; tādṛśā eva jānate обмежує коло тих, хто їх розуміє. Український відповідник лишено чинним, не заявляючи історичного впливу на українських авторів.}
\vcom{bhrānta означає і «заблукалий», і «божевільний» — двозначність тут не прикраса, а точність: ззовні така людина не відрізняється від божевільної, бо різниця всередині, а всередину не видно. tādṛśāḥ eva («лише такі ж») закриває тему пізнання: зрозуміти його може тільки той, хто в такому самому стані, тобто свідчення ззовні неможливе. Це прямий виклик будь-якому зовнішньому описові, і він поставлений останнім віршем глави. БУДОВА ГЛАВИ. Глава 14 не має спільного рефрену — на відміну від 5, 6 і 13. Натомість усі чотири вірші тримаються однієї думки, поданої з чотирьох боків: хто такий цей чоловік зсередини (14.1), що для нього зникає (14.2), про що він уже не турбується (14.3) і як його бачать ззовні (14.4).}
\vpar{Українська паралель (наша інтерпретація)}{Василь Стус, «Веселий цвинтар» (1985).}{Наша паралель: у Стуса людину теж знають тільки такі самі, як вона, — і саме це, а не осуд, стає мірою. Наше порівняння — типологічне.}
\vfig{ilustratsii/14-4-witness.pdf}{Порожній усередині, різний зовні}{Схема наша, не історична ілюстрація.}


% ЗГЕНЕРОВАНО scripts/build_format_v2_chapter.py — не редагувати вручну.
% Джерело: editorial/cards/15.*.json, format-v2/rozdil-15.json, paraleli.json.

\naukchapter{15}{20}{Глава 15. Живи легко}
\naukintro{Пʼятнадцята глава — двадцять віршів, і в ній кілька разів звучить те саме, що стало її назвою: живи легко. Це не заклик до безтурботності, а висновок із попереднього: якщо ти не тіло й не діяч (15.4–15.5), то приріст і втрата до тебе не доходять (15.10–15.11). Глава довга, бо повторює ту саму думку з різних боків, і в кінці доходить до найрадикальнішого: відкинути саму медитацію (15.20).}
\vsec{15.1}
\vdeva{यथातथोपदेशेन कृतार्थः सत्त्वबुद्धिमान् । आजीवम् अपि जिज्ञासुः परस् तत्र विमुह्यति ॥}
\vtranslit{ятхататхопадешена крітартхах саттвабуддгіман | аджівам апі джіджнясух парас татра вімухяті ||}
\vperek{Хто має ясний розум, той здійснює мету настановою «як є»; а допитливий, хоч би все життя шукав, — інший у цьому збивається з пуття.}
\vcom{Построкове зіставлення: редакторська чернетка, не рецензовано.}
\vcom{Рядок а}
\vdeva{यथातथोपदेशेन कृतार्थः सत्त्वबुद्धिमान्}
\vtranslit{ятхататхопадешена крітартхах саттвабуддгіман}
\vperek{Хто має ясний розум, той здійснює мету настановою «як є»;}
\vcom{yathātathopadeśena — орудний відмінок «через настанову» з неоднозначним yathā-tathā: у чинному українському перекладі «як є», проте можливе й інше тлумачення способу навчання; kṛtārthaḥ — «той, хто досяг мети», sattvabuddhimān — «той, чий розум ясний». Граматично це описаний стан учня, а не емпіричне твердження, ніби освіта непотрібна.}
\vcom{Рядок б}
\vdeva{आजीवम् अपि जिज्ञासुः परस् तत्र विमुह्यति}
\vtranslit{аджівам апі джіджнясух парас татра вімухяті}
\vperek{а допитливий, хоч би все життя шукав, — інший у цьому збивається з пуття.}
\vcom{ājīvam api — «хоч би протягом життя», jijñāsuḥ — «той, хто прагне знати», paraḥ (у сандгі paras) — «інший», tatra — «тут/у цьому», vimuhyati — «плутається, збивається з пуття». Протиставлення ясності й нескінченного пошуку належить голосу вірша; воно саме собою не засуджує всіх дослідників чи навчання.}
\vcom{`sattvabuddhimān` — «той, чий розум ясний» — стоїть як умова, а не як наслідок: спершу ясність, потім настанова. `yathātathā` («як є») я читаю як прислівник способу при `upadeśa`; інше можливе прочитання — «настанова „так-не-так“», тобто безладний навчальний текст, і тоді сенс інший. Друга половина протиставляє не знання незнанню, а **ясність допитливості**: `jijñāsu` тут не похвала. Це узгоджується з 9.2, де втихає саме `bubhutsā`.}
\vpar{Українська паралель (наша інтерпретація)}{Памфіл Юркевич, «Серце і його значення в духовному житті людини» (1860).}{Наша паралель: коли вірш протиставляє ясний розум допитливому, він ставить настанову вище за пошук. Юркевич теж відводить пізнанню місце, підпорядковане цілому, а не навпаки. Наше порівняння — за будовою думки.}
\vfig{ilustratsii/15-1-motion.pdf}{Дві дороги: настанова як є і довгий пошук}{Схема наша, не історична ілюстрація.}

\vsec{15.2}
\vdeva{मोक्षो विषयवैरस्यं बन्धो वैषयिको रसः । एतावद् एव विज्ञानं यथेच्छसि तथा कुरु ॥}
\vtranslit{мокшо вішаявайрасям бандго вайшаїко расах | етавад ева віджнянам ятхеччасі татха куру ||}
\vperek{Визволення — коли чуттєве втрачає свій смак; неволя — коли воно й далі вабить. Ось і все розуміння. Чини, як бажаєш.}
\vcom{Построкове зіставлення: редакторська чернетка, не рецензовано.}
\vcom{Рядок а}
\vdeva{मोक्षो विषयवैरस्यं बन्धो वैषयिको रसः}
\vtranslit{мокшо вішаявайрасям бандго вайшаїко расах}
\vperek{Визволення — коли чуттєве втрачає свій смак; неволя — коли воно й далі вабить.}
\vcom{mokṣo (mokṣaḥ) — «визволення», viṣayavairasyaṃ — «втрата смаку до чуттєвих предметів»; bandho (bandhaḥ) — «зв'язаність, неволя», vaiṣayiko rasaḥ — «смак до чуттєвого». Рима понять працює через rasa — смак, а не заборону самих матеріальних речей. Це філософська дефініція у творі, а не медичний опис втрати смакових відчуттів.}
\vcom{Рядок б}
\vdeva{एतावद् एव विज्ञानं यथेच्छसि तथा कुरु}
\vtranslit{етавад ева віджнянам ятхеччасі татха куру}
\vperek{Ось і все розуміння. Чини, як бажаєш.}
\vcom{etāvad eva vijñānaṃ — «ось і все розуміння/розрізнення»; yathecchasi (yathā icchasi) — «як бажаєш», tathā kuru — «так чини», наказовий спосіб від kṛ. Формула залишає простір дії після пояснення визволення, але в санскриті немає окремого слова «анархія» чи безумовної політичної програми; будь-яке таке зіставлення було б сучасною інтерпретацією.}
\vcom{Протиставлення viṣaya-vairasyaṃ і vaiṣayika rasaḥ побудовано на образі «смаку»: прив'язаність визначається не самим існуванням предметів, а ставленням до них. Фінальне yathecchasi tathā kuru стоїть після etāvad eva vijñānaṃ; тому «чини, як бажаєш» не слід відривати від попереднього визначення визволення. Формула не містить слова «анархія» і сама собою не обґрунтовує твердження про політичне вчення; порівняння з анархізмом належить до окремого сучасного інтерпретаційного есе, а не до перекладу. Джерельна база цієї чернетки: SARIT verse\_15.2 і GRETIL Avg\_15.2; незалежна санскритологічна рецензія ще потрібна.}
\vpar{Українська паралель (наша інтерпретація)}{Григорій Сковорода, «Убогий Жайворонок».}{Наша паралель: у Сковороди смак до речей теж є мірилом неволі, а не самі речі. Наше порівняння — типологічне, не твердження про джерело.}
\vfig{ilustratsii/15-2-balance.pdf}{Дві шальки: смак чуттєвого і його відсутність}{Схема наша, не історична ілюстрація.}

\vsec{15.3}
\vdeva{वाग्मिप्राज्ञमहोद्योगं जनं मूकजडालसम् । करोति तत्त्वबोधो ऽयम् अतस् त्यक्तो बुभुक्षुभिः ॥}
\vtranslit{ваґміпраджнямаходйоґам джанам мукаджадаласам | кароті таттвабодго ’ям атас тякто бубгукшубгіх ||}
\vperek{Це пізнання суті робить красномовних, тямущих і вельми завзятих людей німими, тупими й лінивими — тому голодні його відкидають.}
\vcom{Найдошкульніше місце глави, і воно **сатиричне**: пізнання суті не підносить людину, а відбирає в неї красномовство, тямущість і завзятість. `bubhukṣubhiḥ` — «голодними»: тими, хто заробляє словом і зусиллям. Тобто текст сам пояснює, чому його вчення не користується попитом, і робить це без жалю. Переклад мусить зберегти різкість трьох пар, не пом'якшуючи їх до «мовчазних, спокійних, бездіяльних».}
\vpar{Українська паралель (ПОТРЕБУЄ ПЕРЕВІРКИ)}{українська барокова проповідь, не визначено.}{ПОТРЕБУЄ ПЕРЕВІРКИ: думка, що знання суті робить красномовних німими, близька до барокової проповіді про марність учення, але конкретного тексту ми не перевірили.}
\vfig{ilustratsii/15-3-neti.pdf}{Красномовні й завзяті стають німими}{Схема наша, не історична ілюстрація.}

\vsec{15.4}
\vdeva{न त्वं देहो न ते देहो भोक्ता कर्ता न वा भवान् । चिद्रूपो ऽसि सदा साक्षी निरपेक्षः सुखं चर ॥}
\vtranslit{на твам дехо на те дехо бгокта карта на ва бгаван | чідрупо ’сі сада сакші нірапекшах сукхам чара ||}
\vperek{Ти не тіло, і тіло не твоє; і не ти — той, хто споживає чи чинить. Ти маєш природу свідомості, повсякчас свідок, без залежності — живи легко.}
\vcom{`sākṣin` («свідок») тут єдиний раз у главі й стоїть у трійці з `cidrūpa` і `nirapekṣa`; переклад мусить тримати його як термін, а не як метафору. `sukhaṃ cara` — наказовий спосіб, «живи легко»: це не опис стану, а запрошення. Формула «не тіло, і тіло не твоє» та сама, що в 8.4, 11.6 і 2.22 — тут до неї додано заперечення **діяча й споживача**, чого в тих місцях не було.}
\vpar{Українська паралель (наша інтерпретація)}{Григорій Сковорода, «Наркіс. Розмова про те: Пізнай себе».}{Наша паралель: у Сковороди справжня людина теж не є ані тілом, ані тим, хто споживає. Наше порівняння — типологічне.}
\vfig{ilustratsii/15-4-neti.pdf}{Тіло й споживач відкинуті: лишається свідок}{Схема наша, не історична ілюстрація.}

\vsec{15.5}
\vdeva{रागद्वेषौ मनोधर्मौ न मनस् ते कदाचन । निर्विकल्पो ऽसि बोधात्मा निर्विकारः सुखं चर ॥}
\vtranslit{раґадвешау манодгармау на манас те кадачана | нірвікалпо ’сі бодгатма нірвікарах сукхам чара ||}
\vperek{Пристрасть і відраза — властивості ума, і ум ніколи не твій. Ти без домислювання, твоя суть — усвідомлення, ти без видозмін — живи легко.}
\vcom{Ланцюг тут двоступеневий: спершу `rāga`/`dveṣa` віднесено до `manas`, а потім заперечено сам `manas` як власність. `nirvikāra` («без видозмін») — те саме слово, що в 11.1; словник спільний. `sukhaṃ cara` повторює 15.4 дослівно, і в перекладі це має бути той самий вислів — інакше зникне наскрізний рефрен глави.}
\vpar{Українська паралель (наша інтерпретація)}{Микола Хвильовий, «Я (Романтика)» (1924).}{Наша паралель: у Хвильового пристрасть і відраза теж показані як властивості ума, а не як голос самої речі. Наше порівняння — за способом бачити.}
\vfig{ilustratsii/15-5-balance.pdf}{Дві шальки: пристрасть і відраза — обидві не твої}{Схема наша, не історична ілюстрація.}

\vsec{15.6}
\vdeva{सर्वभूतेषु चात्मानं सर्वभूतानि चात्मनि । विज्ञाय निरहंकारो निर्ममस् त्वं सुखी भव ॥}
\vtranslit{сарвабгутешу чатманам сарвабгутані чатмані | віджняя нірахамкаро нірмамас твам сукхі бгава ||}
\vperek{Пізнавши себе у всіх істотах і всі істоти — у собі, ти без «я-роблю», без «моє» — будь щасливий.}
\vcom{Двобічність тут обов'язкова: `ātmānam` у всіх істотах **і** всі істоти в `ātman`. Переклад, що зберіг лише один напрямок, втрачає суть. `nirahaṃkāra`/`nirmama` — не «скромний» і не «безкорисливий», а буквально «без `ahaṃkāra`» і «без `mama`»: два займенники, від яких відмовляються. `sukhī bhava` — наказовий спосіб, як і `sukhaṃ cara` у 15.4–15.5.}
\vpar{Українська паралель (наша інтерпретація)}{Митрополит Іларіон Київський, «Слово про закон і благодать» (XI ст.).}{Наша паралель: у київському Слові теж є думка, що людина впізнає себе не в собі одній. Наше порівняння — типологічне; датування цієї книги науково не встановлене.}
\vfig{ilustratsii/15-6-witness.pdf}{Взаємне перебування: ти в усіх, усі в тобі}{Схема наша, не історична ілюстрація.}

\vsec{15.7}
\vdeva{विश्वं स्फुरति यत्रेदं तरङ्गा इव सागरे । तत्त्वम् एव न सन्देहश् चिन्मूर्ते विज्वरो भव ॥}
\vtranslit{вішвам спхураті ятредам таранґа іва саґаре | таттвам ева на сандехаш чінмурте віджваро бгава ||}
\vperek{Те, в чому цей всесвіт постає, як хвилі в морі, — це саме ти, немає сумніву, о втілення свідомості; будь без трепету.}
\vcom{Образ хвиль і моря — єдиний такий у книзі, і він робить роботу, яку не робить жодне означення: хвиля не **відмінна** від моря й не **тотожна** йому як річ. `cinmūrte` — звертання «о втілення свідомості», споріднене з `cidrūpa` у 15.4; переклад мусить тримати спільний корінь. `vijvara` («без трепету») — від `jvara`, гарячка; це не спокій, а відсутність пропасниці.}
\vpar{Українська паралель (наша інтерпретація)}{Богдан-Ігор Антонич, «Ротації» (1938).}{Наша паралель: в Антонича всесвіт теж постає як хвиля в морі, а не як споруда. Наше порівняння — за образом.}
\vfig{ilustratsii/15-7-motion.pdf}{Всесвіт як хвилі в морі: це саме ти}{Схема наша, не історична ілюстрація.}

\vsec{15.8}
\vdeva{श्रद्धत्स्व तात श्रद्धत्स्व नात्र मोहं कुरुष्व भोः । ज्ञानस्वरूपो भगवान् आत्मा त्वं प्रकृतेः परः ॥}
\vtranslit{шраддгатсва тата шраддгатсва натра мохам курушва бгох | джнянасварупо бгаґаван атма твам пракрітех парах ||}
\vperek{Повір, любий, повір; не роби тут затьмарення, друже. Твоя суть — знання; бгаґаван ātman — ти, понад природою.}
\vcom{Єдине місце в книзі, де вчення вимагає **віри** — і вимагає двічі, повтором `śraddhatsva`. Це суперечить загальній настанові тексту «перевір сам», і я не згладжую цю суперечність: вона є. `bhagavān ātmā` — рідкісне поєднання, де `ātman` названо `bhagavant`; переклад мусить зберегти обидва слова, не зводячи до «Бог у тобі». `prakṛteḥ paraḥ` — «понад природою», тобто не складник її.}
\vpar{Українська паралель (ПОТРЕБУЄ ПЕРЕВІРКИ)}{українська релігійна поезія, не визначено.}{ПОТРЕБУЄ ПЕРЕВІРКИ: пряме звертання повір, любий із повторенням близьке до української релігійної поезії, але конкретного твору ми не звірили.}
\vfig{ilustratsii/15-8-flame.pdf}{Полумʼя знання: повір — і не роби затьмарення}{Схема наша, не історична ілюстрація.}

\vsec{15.9}
\vdeva{गुणैः संवेष्टितो देहस् तिष्ठत्य् आयाति याति च । आत्मा न गन्ता नागन्ता किम् एनम् अनुशोचसि ॥}
\vtranslit{ґунайх самвештіто дехас тіштхатй аяті яті ча | атма на ґанта наґанта кім енам анушочасі ||}
\vperek{Тіло, обгорнуте гунами, стоїть, приходить і йде; ātman не йде і не приходить — чого ти за ним жалкуєш?}
\vcom{Заперечення подвійне й симетричне: `na gantā na agantā` — не йде **і не приходить**, тобто не підпадає під жоден бік пари. Питання `kim enam anuśocasi` — пряме звертання до того, хто сумує; у главі це вже друге таке звертання після 15.6, і разом вони задають тон: не пояснення, а докір. `anuśocasi` — «жалкуєш за ним», не «сумуєш» взагалі.}
\vpar{Українська паралель (наша інтерпретація)}{Тарас Шевченко, «Марія» (1859).}{Наша паралель: у Шевченка те, що не йде й не приходить, теж протиставлене тілу, яке приходить і відходить. Наше порівняння — типологічне, не про джерело.}
\vfig{ilustratsii/15-9-motion.pdf}{Тіло приходить і йде, а те, що бачить, стоїть}{Схема наша, не історична ілюстрація.}
\vfig{ilustratsii-naratyvni/15-9-tilo-ide.pdf}{Тіло йде дорогою, лишаючи сліди, а той, хто сидить, не йде}{Наша ілюстрація до вірша 15.9 — не історичне зображення.}

\vsec{15.10}
\vdeva{देहस् तिष्ठतु कल्पान्तं गच्छत्य् अद्यैव वा पुनः । क्व वृद्धिः क्व च वा हानिस् तव चिन्मात्ररूपिणः ॥}
\vtranslit{дехас тіштхату калпантам ґаччатй адяйва ва пунах | ква вріддгіх ква ча ва ханіс тава чінматрарупінах ||}
\vperek{Нехай тіло стоїть до кінця світового періоду — або нехай іде хоч сьогодні. Де приріст і де втрата для тебе, чия природа — сама лише свідомість?}
\vcom{Пара наказових способів (`tiṣṭhatu` … `gacchatu`) задає байдужість до тривалості життя: обидва варіанти рівноцінні. `kalpānta` («кінець світового періоду») — не поетична гіпербола, а максимально можливий строк, і саме тому друга частина звучить різко: `adya eva` («хоч сьогодні»). `cinmātrarūpa` («чия природа — сама лише свідомість») — те саме, що `cinmūrte` у 15.7 і `cidrūpa` у 15.4; словник лишається спільним.}
\vpar{Українська паралель (наша інтерпретація)}{Ліна Костенко, «Сад нетанучих скульптур» (1987).}{Наша паралель: у Костенко питання про приріст і втрату теж виявляється поставленим не до тієї міри. Наше порівняння — за будовою думки.}
\vfig{ilustratsii/15-10-balance.pdf}{Дві шальки: тіло надовго і тіло на день}{Схема наша, не історична ілюстрація.}

\vsec{15.11}
\vdeva{त्वय्य् अनन्तमहाम्भोधौ विश्ववीचिः स्वभावतः । उदेतु वास्तम् आयातु न ते वृद्धिर् न वा क्षतिः ॥}
\vtranslit{твайй анантамахамбгодгау вішвавічіх свабгаватах | удету вастам аяту на те вріддгір на ва кшатіх ||}
\vperek{У тобі, безмежному великому океані, хвиля всесвіту за власною природою нехай піднімається або спадає — немає в тебе ні приросту, ні втрати.}
\vcom{Той самий образ, що в 15.7, але розгорнутий: там хвилі були в морі, тут ти **сам** — океан, і всесвіт — одна хвиля в ньому. `svabhāvataḥ` («за власною природою») важливе: хвиля не робиться нею ззовні, вона така за природою. Пара `vṛddhi`/`kṣati` повторює 15.10 дослівно — це рефрен двох віршів, і в перекладі він мусить звучати однаково.}
\vpar{Українська паралель (наша інтерпретація)}{Олександр Довженко, «Поема про море» (1956).}{Наша паралель: у Довженка хвиля теж не додає й не забирає нічого в моря, у якому вона піднялася. Наше порівняння — за образом.}
\vfig{ilustratsii/15-11-motion.pdf}{Хвиля всесвіту в безмежному океані}{Схема наша, не історична ілюстрація.}

\vsec{15.12}
\vdeva{तात चिन्मात्ररूपो ऽसि न ते भिन्नम् इदं जगत् । अतः कस्य कथं कुत्र हेयोपादेयकल्पना ॥}
\vtranslit{тата чінматрарупо ’сі на те бгіннам ідам джаґат | атах кася катхам кутра хейопадеякалпана ||}
\vperek{Любий, твоя природа — сама лише свідомість; цей світ не відмінний від тебе. Тому — для кого, як і де вигадувати те, що відкидати й приймати?}
\vcom{Три питання підряд (`kasya katham kutra`) — не риторичний повтор, а розбір по трьох осях: для кого, яким чином, де. `kalpanā` — не «вибір», а **вигадування**: текст не каже, що вибір між добрим і поганим складний, він каже, що його нема кому робити. `tāta` — те саме ласкаве звертання, що в 9.2 і 15.8.}
\vpar{Українська паралель (наша інтерпретація)}{Іван Вишенський, «Книжка» (кінець XVI ст.).}{Наша паралель: у Вишенського світ теж не відділений від того, хто його бачить, і саме тому вибір між прийняти й відкинути втрачає предмет. Наше порівняння — типологічне.}
\vfig{ilustratsii/15-12-neti.pdf}{Відкидати й приймати нічого: світ не відмінний від тебе}{Схема наша, не історична ілюстрація.}

\vsec{15.13}
\vdeva{एकस्मिन् अव्यये शान्ते चिदाकाशे ऽमले त्वयि । कुतो जन्म कुतो कर्म कुतो ऽहंकार एव च ॥}
\vtranslit{екасмін ав'яє шанте чідакаше ’мале тваї | куто джанма куто карма куто ’хамкара ева ча ||}
\vperek{У тобі — одному, незмінному, втихлому, незаплямованому просторі свідомості — звідки народження, звідки діяння, звідки ж «я-роблю»?}
\vcom{`cidākāśa` («простір свідомості») — єдине таке складне слово в книзі; переклад мусить лишити його як термін, а не як метафору. Чотири означення (`eka`, `avyaya`, `śānta`, `amala`) звужують значення до одного, і кожне наступне сильніше. Три `kutaḥ` ідуть у порядку виникнення: народження → діяння → `ahaṃkāra`, і останнє поставлене як те, що робить перші два можливими.}
\vpar{Українська паралель (ПОТРЕБУЄ ПЕРЕВІРКИ)}{українська містична проза, не визначено.}{ПОТРЕБУЄ ПЕРЕВІРКИ: питання, звідки народження й діяння в тому, що незмінне, найближче до української містичної прози, але назвати твір без звірки ми не можемо.}
\vfig{ilustratsii/15-13-space.pdf}{Порожній простір свідомості: звідки народження й діяння?}{Схема наша, не історична ілюстрація.}

\vsec{15.14}
\vdeva{यद् त्वं पश्यसि तत्रैकस् त्वम् एव प्रतिभाससे । किं पृथग् भासते स्वर्णात् कटकाङ्गदनूपुरम् ॥}
\vtranslit{яд твам пашясі татрайкас твам ева пратібгасасе | кім прітхаґ бгасате сварнат катаканґаданупурам ||}
\vperek{Те, що ти бачиш, — там сяєш один саме ти. Чи сяє окремо від золота браслет, наруч або ножна прикраса?}
\vcom{Риторичне питання з відповіддю «ні»: три прикраси названі поіменно, щоб показати, що різниця між ними — лише форма, а речовина одна. Це той самий хід, що з хвилями в 15.7 і 15.11, але тут він зроблений на **речі**, а не на стихії. `pratibhāsase` — «сяєш, постаєш», споріднене з `sphurati` в 15.7; переклад мусить тримати цю спорідненість.}
\vpar{Українська паралель (наша інтерпретація)}{Григорій Сковорода, «Ікона Алківіадська».}{Наша паралель: у Сковороди прикраса й золото теж співвідносяться як видиме й те, що в ньому сяє. Наше порівняння — за образом.}
\vfig{ilustratsii/15-14-thread.pdf}{Прикраси на нитці: усі вони — золото}{Схема наша, не історична ілюстрація.}
\vfig{ilustratsii-naratyvni/15-14-brastelet-i-zoloto.pdf}{Злиток золота на плиті й браслет поруч: окремо від золота браслет не сяє}{Наша ілюстрація до вірша 15.14 — не історичне зображення.}

\vsec{15.15}
\vdeva{अयं सो ऽहम् अयं नाहम् विभागम् इति संत्यज । सर्वम् आत्मेति निश्चित्य निःसंकल्पः सुखी भव ॥}
\vtranslit{аям со ’хам аям нахам вібгаґам іті самтяджа | сарвам атметі нішчітя ніхсамкалпах сукхі бгава ||}
\vperek{«Це — він, це я; це — не я» — відкинь такий поділ. Вирішивши, що все — ātman, будь без побудов і щасливий.}
\vcom{Тут відкидається не одна зі сторін поділу, а **сама навичка ділити**: `vibhāgam … saṃtyaja`. `niḥsaṃkalpa` — не «без думок», а без `saṃkalpa`, наміру-побудови, тобто без попереднього складання того, що має бути. Це те саме слово, що `kalpanā` в 15.12, і переклад мусить бути спорідненим.}
\vpar{Українська паралель (наша інтерпретація)}{Микола Зеров, «Камена» (1924).}{Наша паралель: у Зерова поділ на це і не це теж знімається ясністю, а не зусиллям. Наше порівняння — типологічне.}
\vfig{ilustratsii/15-15-neti.pdf}{Поділ на це і не це відкинуто: усе — ātman}{Схема наша, не історична ілюстрація.}

\vsec{15.16}
\vdeva{तवैवाज्ञानतो विश्वं त्वम् एकः परमार्थतः । त्वत्तो ऽअन्यो नास्ति संसारी नासंसारी च कश्चन ॥}
\vtranslit{тавайваджнянато вішвам твам еках парамартхатах | тватто ’анйо насті самсарі насамсарі ча кашчана ||}
\vperek{Всесвіт — саме через твоє незнання; у справжньому сенсі ти один. Немає іншого, ніж ти — ні того, хто в сансарі, ні того, хто поза нею.}
\vcom{Заперечення знову симетричне: немає ні `saṃsārī`, ні `asaṃsārī` — тобто не існує ні того, хто в сансарі, ні того, хто з неї вийшов. Це найрадикальніше з тверджень глави, і його легко зіпсувати перекладом «немає ні сансари, ні визволення»: у тексті йдеться про **осіб**, а не про стани. `paramārthataḥ` («у справжньому сенсі») — усталений термін, не «насправді» в побутовому значенні.}
\vpar{Українська паралель (ПОТРЕБУЄ ПЕРЕВІРКИ)}{українська філософська проза, не визначено.}{ПОТРЕБУЄ ПЕРЕВІРКИ: думка, що всесвіт постає через незнання, найближча до української філософської прози, але конкретного твору ми не перевірили.}
\vfig{ilustratsii/15-16-space.pdf}{Всесвіт як горщик, а простір — один}{Схема наша, не історична ілюстрація.}

\vsec{15.17}
\vdeva{भ्रान्तिमात्रम् इदं विश्वं न किञ्चिद् इति निश्चयी । निर्वासनः स्फूर्तिमात्रो न किञ्चिद् इव शाम्यति ॥}
\vtranslit{бгрантіматрам ідам вішвам на кіньчід іті нішчаї | нірвасанах спхуртіматро на кіньчід іва шам'яті ||}
\vperek{«Цей всесвіт — саме лише блукання, аніщо» — хто має таке переконання, той без усталених схильностей, саме лише сяйнування, втихає, наче аніщо.}
\vcom{**Цей вірш майже дослівно збігається з 11.8.** Різниця лише в першому складнику: тут `bhrāntimātram` («саме лише блукання»), там `nānāścaryam` («різноманітно дивовижне»). Друга половина — `nirvāsanaḥ sphūrtimātro na kiñcid iva śāmyati` — ідентична в обох місцях. Це не помилка картки й не недогляд: текст справді повторює формулу, змінюючи лише означення світу. Я зберіг переклад другої половини **дослівно той самий**, що в 11.8, щоб повтор було видно. `bhrānti` — «блукання» в сенсі омани-ходіння, не «ілюзія» як технічний термін.}
\vpar{Українська паралель (наша інтерпретація)}{Іван Франко, «Мойсей» (1905).}{Наша паралель: у Франка блукання теж може виявитися тим, що не має власної суті. Наше порівняння — типологічне.}
\vfig{ilustratsii/15-17-illusion.pdf}{Всесвіт під серпанком: блукання, аніщо}{Схема наша, не історична ілюстрація.}

\vsec{15.18}
\vdeva{एक एव भवाम्भोधाव् आसीद् अस्ति भविष्यति । न ते बन्धो ऽस्ति मोक्षो वा कृतकृत्यः सुखं चर ॥}
\vtranslit{ека ева бгавамбгодгав асід асті бгавішяті | на те бандго ’сті мокшо ва крітакрітях сукхам чара ||}
\vperek{Один-єдиний в океані буття — був, є, буде. Немає в тебе ні пут, ні визволення; той, чию справу зроблено, — живи легко.}
\vcom{`na te bandho 'sti mokṣo vā` — та сама симетрична пара, що в 15.16: немає ні пут, ні визволення. Переклад мусить тримати їх **як пару**, бо звичне «ти вільний» знищило б думку. Три часи (`āsīt asti bhaviṣyati`) стоять у ряд, як у 19.3, але там вони скасовувалися, а тут стверджуються. `kṛtakṛtya` — те саме, що `kṛtārtha` в 12.8, і переклад має бути спорідненим.}
\vpar{Українська паралель (наша інтерпретація)}{Василь Стус, «Час творчості».}{Наша паралель: у Стуса той, чию справу зроблено, теж не має ні пут, ні визволення — він просто живе. Наше порівняння — за будовою думки.}
\vfig{ilustratsii/15-18-space.pdf}{Один в океані буття: був, є, буде}{Схема наша, не історична ілюстрація.}

\vsec{15.19}
\vdeva{मा संकल्पविकल्पाभ्यां चित्तं क्षोभय चिन्मय । उपशाम्य सुखं तिष्ठ स्वात्मन्य् आनन्दविग्रहे ॥}
\vtranslit{ма самкалпавікалпабгям чіттам кшобгая чінмая | упашам'я сукхам тіштха сватманй анандавіґрахе ||}
\vperek{Не збурюй ум наміром і домислюванням, о сповнений свідомості. Втихнувши, стій легко у собі — втіленні блаженства.}
\vcom{`saṃkalpa` і `vikalpa` названі парою як те, що збурює: перше — складання наміру, друге — роздвоєння-домислювання. Це ті самі два корені, що в 15.12 (`kalpanā`) і 15.15 (`niḥsaṃkalpa`), і разом вони утворюють послідовний словник глави. `ānandavigraha` («втілення блаженства») — єдине таке поєднання в книзі; це не «блаженна людина», а означення **себе**.}
\vpar{Українська паралель (наша інтерпретація)}{Павло Тичина, «Сонячні кларнети» (1918).}{Наша паралель: у Тичини втихнути й стояти легко теж означає перестати збурювати те, що й так сяє. Наше порівняння — за образом.}
\vfig{ilustratsii/15-19-flame.pdf}{Полумʼя без вітру: втихнувши, стій легко}{Схема наша, не історична ілюстрація.}

\vsec{15.20}
\vdeva{त्यजैव ध्यानं सर्वत्र मा किञ्चिद् धृदि धारय । आत्मा त्वं मुक्त एवासि किं विमृश्य करिष्यसि ॥}
\vtranslit{тяджайва дгянам сарватра ма кіньчід дгріді дгарая | атма твам мукта евасі кім вімрішя карішясі ||}
\vperek{Відкинь саму медитацію, скрізь не тримай нічого в серці. Ти — ātman, ти вже вільний: що ти робитимеш, далі розмірковуючи?}
\vcom{Останній вірш глави відкидає **медитацію** — те, з чого починалася майже вся релігійна практика. `tyaja eva` з підсиленням: не «зменш», а «відкинь саме». `mukta eva asi` — «ти вже вільний», і `eva` тут не прикраса: без нього фраза звучала б як обіцянка. Завершальне питання `kiṃ … kariṣyasi` не закликає до дії — воно показує, що робити нічого. Це узгоджується з 15.3, де пізнання суті відбирає завзятість.}
\vpar{Українська паралель (наша інтерпретація)}{Леся Українка, «Оргія» (1913).}{Наша паралель: у Лесі Українки відкинути саму медитацію — це теж крок, а не втрата, бо мета вже не попереду. Наше порівняння — типологічне.}
\vfig{ilustratsii/15-20-neti.pdf}{Медитація відкинута: ти вже вільний}{Схема наша, не історична ілюстрація.}


% ЗГЕНЕРОВАНО scripts/build_format_v2_chapter.py — не редагувати вручну.
% Джерело: editorial/cards/16.*.json, format-v2/rozdil-16.json, paraleli.json.

\naukchapter{16}{11}{Глава 16. Поки не забудеш усе}
\naukintro{Шістнадцята глава починається й закінчується тим самим: спокою не буде, поки не забудеш усе. Між цими двома віршами — одинадцять спроб сказати те саме інакше: через натугу (16.3), через чотири цілі людини (16.5), через пару потяг і відраза (16.8). Останній вірш навмисно повертає читача до початку й додає, що навіть боги-наставники тут не помічники.}
\vsec{16.1}
\vdeva{आचक्ष्व शृणु वा तात । नानाशास्त्राण्यनेकशः । तथापि न तव स्वास्थ्यं । सर्वविस्मरणाद् ऋते ॥}
\vtranslit{ачакшва шріну ва тата | нанашастранянекашах | татхапі на тава свастхям | сарвавісмаранад ріте ||}
\vperek{Розповідай чи слухай, любий, хоч би скільки різних наук, — спокою тобі не буде, поки не забудеш усе.}
\vcom{Глава відкривається не запереченням знання, а запереченням його достатності. Дія в парі «говорити / слухати» охоплює обидві ролі вченого — того, хто навчає, і того, хто вчиться, — і обидві виявляються однаково неспроможними. Ключове слово — svāsthya, «перебування в собі»: воно описує стан, а не досягнення, і саме тому його не можна здобути додаванням знань. Умова поставлена одна й негативна: sarvavismaraṇa, забуття всього. Текст не пояснює, що саме забути й як; традиційні коментарі читають це як забуття розрізнень і набутих уявлень, але сам вірш такого тлумачення не містить. Не сказано тут і про ātman, cit чи sākṣin — жодного з цих термінів у вірші немає.}
\vpar{Українська паралель (наша інтерпретація)}{Григорій Сковорода, «Байки харківські» (1774).}{Наша паралель: у Сковороди багато знаного теж не дає спокою, поки не зміниться сам спосіб тримати знане. Наше порівняння — типологічне.}
\vfig{ilustratsii/16-1-motion.pdf}{Багато наук, що не дають спокою}{Схема наша, не історична ілюстрація.}

\vsec{16.2}
\vdeva{भोगं कर्म समाधिं वा । कुरु विज्ञ तथापि ते । चित्तं निरस्तसर्वाशम् । अत्यर्थं रोचयिष्यति ॥}
\vtranslit{бгоґам карма самадгім ва | куру віджня татхапі те | чіттам нірастасарвашам | атяртхам рочаїшяті ||}
\vperek{Насолоджуйся, дій або зосереджуйся, обізнаний, — та по-справжньому смакуватиме тобі лише тоді, коли розум відкине всі надії.}
\vcom{Вірш перебирає три типові стратегії: чуттєву насолоду, дію (обряд, обов'язок) і зосередження — тобто весь набір засобів, які зазвичай пропонують як шлях. Заперечення тут не в тому, що ці заняття погані: їх дозволено робити далі. Змінено умову смаку — він з'являється не від заняття, а від відсутності надії (nirasta-sarvāśa). Це тонша думка, ніж «відмовся від усього»: не дію скасовано, а припинено сподівання, вкладене в дію. Через каузативну форму rocayiṣyati речення можна зрозуміти й так, що смак дає сама дія, коли її не супроводжує надія; переклад вище тримає саме цю, слабшу вимогу. Традиційні коментарі тут бачать указівку на нішкаху (недіяння-зречення), але текст слова «недіяння» не містить.}
\vpar{Українська паралель (наша інтерпретація)}{Леся Українка, «У пущі» (1909).}{Наша паралель: у Лесі Українки смак до дії теж з'являється лише тоді, коли розум відпускає надії. Наше порівняння — за будовою думки.}
\vfig{ilustratsii/16-2-neti.pdf}{Надії відкинуті — і тоді смак повертається}{Схема наша, не історична ілюстрація.}

\vsec{16.3}
\vdeva{आयासात् सकलो दुःखी । नैनं जानाति कश्चन । अनेनैवोपदेशेन । धन्यः प्राप्नोति निर्वृतिम् ॥}
\vtranslit{аясат сакало духкхі | найнам джанаті кашчана | аненайвопадешена | дганях прапноті нірврітім ||}
\vperek{Від натуги всі нещасні, і ніхто з них цього не бачить. Блаженний же досягає вгамування самим лише цим настановленням.}
\vcom{Вірш ставить діагноз: причина страждання — āyāsa, натуга, а не брак чогось. Друга половина вводить протиставлення між тими, хто не бачить, і тим, хто чує настановлення; але текст не пояснює, що саме робить це настановлення дієвим — чи зміст, чи готовність слухача. Порівняння з 16.1 показує, що «настановлення» тут не нове знання, а вказівка на припинення натуги: обидва вірші зводять причину до дії, яку треба припинити, а не до відомостей, які треба додати. Термінів ātman, cit, sākṣin, vairāgya, toṣa, āśrama у вірші немає.}
\vpar{Українська паралель (ПОТРЕБУЄ ПЕРЕВІРКИ)}{українська філософська лірика, не визначено.}{ПОТРЕБУЄ ПЕРЕВІРКИ: думка, що від натуги всі нещасні, близька до української філософської лірики, але конкретного твору ми не звірили.}
\vfig{ilustratsii/16-3-knot.pdf}{Вузол натуги: всі нещасні, і ніхто не бачить}{Схема наша, не історична ілюстрація.}

\vsec{16.4}
\vdeva{व्यापारे खिद्यते यस् तु । निमेषोन्मेषयोर् अपि । तस्यालस्यधुरीणस्य । सुखं नान्यस्य कस्यचित् ॥}
\vtranslit{в'япаре кхідяте яс ту | німешонмешайор апі | тасяласядгурінася | сукхам наняся касячіт ||}
\vperek{Кому важко навіть повіки звести — той перший серед лінивих: немає щастя ні йому, ні комусь іншому [такому].}
\vcom{Вірш різкий за тоном: він не співчуває втомі, а виставляє її на посміх. Найменша дія — повіки звести — уже вимірює, скільки сили лишилося. Разом із 16.3 це дає цікаву пару: там натуга робить нещасним, а тут відсутність сили на найменшу дію теж не веде до щастя. Отже текст не хвалить і неробства: він описує глухий кут, у якому будь-яке ставлення до дії — і надмірне, і нульове — однаково не дає спокою. Друга половина двозначна (див. граматику), і від вибору залежить, чи вірш картає лінивого, чи взагалі відмовляє в щасті всім, крім нього. Традиційні коментарі читають це як осуд лінощів; сам текст такого висновку не формулює.}
\vpar{Українська паралель (наша інтерпретація)}{Іван Франко, «Зівʼяле листя» (1896).}{Наша паралель: у Франка безсилля теж описане без жалю й без поблажливості, як стан, а не як провина. Наше порівняння — за способом бачити.}
\vfig{ilustratsii/16-4-balance.pdf}{Дві шальки: лінощі й відсутність щастя}{Схема наша, не історична ілюстрація.}

\vsec{16.5}
\vdeva{इदं कृतम् इदं नेति । द्वन्द्वैर् मुक्तं यदा मनः । धर्मार्थकाममोक्षेषु । निरपेक्षं तदा भवेत् ॥}
\vtranslit{ідам крітам ідам неті | двандвайр муктам яда манах | дгармартхакамамокшешу | нірапекшам тада бгавет ||}
\vperek{Коли розум звільниться від пари «це я зробив — а це ні», він стане байдужим і до дгарми, і до артги, і до ками, і до мокші.}
\vcom{Тут уперше в главі названо чотири цілі людини — і всі чотири поставлено в один ряд як те, до чого розум стає байдужим. Це сильніше, ніж заперечення однієї з них: дгарма й мокша втрачають привілей бути «правильними» цілями. Механізм указано точно: не самі цілі пошкоджують розум, а пара оцінок «зробив / не зробив», у якій вони постають. Слово mokṣa тут ужите як назва четвертої цілі, а не як стан; тому переклад передає його як «мокша», не «звільнення», щоб не змішати з 20.12, де mukti стоїть у низці заперечень. Оптатив bhavet робить речення умовним: це не обіцянка, а опис умови.}
\vpar{Українська паралель (наша інтерпретація)}{Олександр Потебня, «Думка й мова» (1862).}{Наша паралель: Потебня теж показує, що пара «я зробив — а це ні» є справою думки, а не речей. Наше порівняння — типологічне.}
\vfig{ilustratsii/16-5-neti.pdf}{Чотири цілі людини відпадають разом}{Схема наша, не історична ілюстрація.}

\vsec{16.6}
\vdeva{विरक्तो विषयद्वेष्टा । रागी विषयलोलुपः । ग्रहमोक्षविहीनस् तु । न विरक्तो न रागवान् ॥}
\vtranslit{віракто вішаядвешта | раґі вішаялолупах | ґрахамокшавіхінас ту | на віракто на раґаван ||}
\vperek{Відвернений — це той, хто ненавидить чуттєве; пристрасний — той, хто до нього жадібний. А хто не хапає і не відкидає — той уже ні відвернений, ні пристрасний.}
\vcom{Вірш руйнує звичну шкалу «пристрасть погано, відверненість добре». Ненависть до чуттєвого поставлено в один ряд із жадобою: обидві — рух навколо того самого предмета, лише з різним знаком. Тому virakta тут не чеснота, а діагноз. Справжня межа проведена за іншим критерієм — graha і mokṣa, хапання й відкидання, тобто сама реакція, а не її напрям. Важливо, що vairāgya — термін, якого в цьому вірші немає: є virakta, дієприкметник, і його якраз поставлено під сумнів. Хто шукає тут похвалу відреченню, той прочитає вірш навпаки.}
\vpar{Українська паралель (наша інтерпретація)}{Михайло Коцюбинський, «Intermezzo» (1908).}{Наша паралель: у Коцюбинського той, хто не хапає й не відкидає, теж стоїть поза парою відвернений і пристрасний. Наше порівняння — типологічне.}
\vfig{ilustratsii/16-6-balance.pdf}{Дві шальки: відвернений і пристрасний}{Схема наша, не історична ілюстрація.}

\vsec{16.7}
\vdeva{हेयोपादेयता तावत् । संसारविटपाङ्कुरः । स्पृहा जीवति यावद् वै । निर्विचारदशास्पदम् ॥}
\vtranslit{хейопадеята тават | самсаравітапанкурах | спріха джіваті явад вай | нірвічарадашаспадам ||}
\vperek{Розрізнення «відкинути це, взяти те» — лише пагін на дереві сансари; він росте, доки жива жадоба, а жадоба — осідок життя без роздуму.}
\vcom{Метафора дерева працює в двох напрямках: пагін — це те, що видно, і водночас те, що тримає корінь. Розрізнення «відкинути / взяти» тут не мудрість, а симптом. Друга половина вказує на причину: жадоба (spṛhā) і стан без роздуму (nirvicāra-daśā) живлять одне одного — без роздуму легко бажати, а бажання не потребує роздуму. Варто зауважити, що vijñāna, viveka чи vicāra в позитивному значенні тут не з'являються: текст не пропонує «правильного розрізнення» на противагу хибному, він указує на розрізнення як таке. Це не заклик до нерозумності; сам вірш — акт розрізнення, і ця суперечність у тексті лишається нерозв'язаною.}
\vpar{Українська паралель (наша інтерпретація)}{Іван Франко, «Мойсей» (1905).}{Наша паралель: у Франка розрізнення відкинути це, взяти те теж постає як пагін, що росте з живого кореня. Наше порівняння — за образом.}
\vfig{ilustratsii/16-7-thread.pdf}{Пагін на дереві сансари: відкинути й узяти}{Схема наша, не історична ілюстрація.}

\vsec{16.8}
\vdeva{प्रवृत्तौ जायते रागो । निवृत्तौ द्वेष एव हि । निर्द्वन्द्वो बालवद् धीमान् । एवम् एव व्यवस्थितः ॥}
\vtranslit{правріттау джаяте раґо | нівріттау двеша ева хі | нірдвандво балавад дгіман | евам ева в'явастхітах ||}
\vperek{Коли дієш — народжується потяг, коли спиняєшся — відраза. Мудрий, вільний від цієї пари, стоїть так, як стоїть дитина.}
\vcom{Вірш дає механіку пари: потяг і відраза не два окремі почуття, а два боки одного маятника, і кожен викликає інший. Тому вихід не в тому, щоб обрати бік (діяти або спинитися), а в тому, щоб вийти з коливання. Порівняння з дитиною — єдиний образ, який текст пропонує замість опису стану; воно нічого не доводить, а показує. Це місце часто читають як «природність» (akṛtrima) у дусі пізніших інтерпретацій, але слово «природний» у вірші відсутнє; є лише bālavat, «як дитина». Не названо тут і ātman чи cit: вірш говорить про механіку реагування, не про природу свідомості.}
\vpar{Українська паралель (наша інтерпретація)}{Григорій Сковорода, «Алфавіт, чи Буквар миру» (діалог).}{Наша паралель: у Сковороди дитина теж стоїть як міра простоти, недосяжної для того, хто вже знає. Наше порівняння — типологічне.}
\vfig{ilustratsii/16-8-motion.pdf}{Потяг і відраза, а поза ними — дитина}{Схема наша, не історична ілюстрація.}

\vsec{16.9}
\vdeva{हातुम् इच्छति संसारं । रागी दुःखजिहासया । वीतरागो हि निर्दुःखस् । तस्मिन्न् अपि न खिद्यति ॥}
\vtranslit{хатум іччаті самсарам | раґі духкхаджіхасая | вітараґо хі нірдухкхас | тасмінн апі на кхідяті ||}
\vperek{Пристрасний хоче втекти з сансари, бо хоче позбутися страждання. А в кого пристрасті вже немає, той не страждає — і навіть у сансарі не мучиться.}
\vcom{Тут показано пастку: бажання вийти з сансари саме є пристрастю, і тому воно тримає в сансарі. Межа проходить не між «у сансарі» і «поза сансарі», а між наявністю й відсутністю пристрасті. Формула nirduḥkhaḥ ... na khidyati подвійна: спершу стан (немає страждання), потім поведінка (не мучиться навіть там). Це не обіцянка, що все стане приємним, а вказівка, що сама присутність сансари перестає бути мірилом. Слово vairāgya, яке тут напрошується, у вірші не вжите: є vītarāga, «той, у кого пристрасть пішла», — опис, а не назва чесноти.}
\vpar{Українська паралель (наша інтерпретація)}{Василь Стус, «Палімпсести».}{Наша паралель: у Стуса втеча від страждання теж не приносить того, чого шукали. Наше порівняння — за будовою думки.}
\vfig{ilustratsii/16-9-knot.pdf}{Вузол втечі: страх болю тримає міцніше}{Схема наша, не історична ілюстрація.}

\vsec{16.10}
\vdeva{यस्याभिमानो मोक्षे ऽपि । देहे ऽपि ममता तथा । न च ज्ञानी न वा योगी । केवलं दुःखभाग् असौ ॥}
\vtranslit{ясябгімано мокше ’пі | дехе ’пі мамата татха | на ча джняні на ва йоґі | кевалам духкхабгаґ асау ||}
\vperek{Якщо людина претендує навіть на звільнення і водночас відчуває тіло як «своє», вона не знавець і не йог — вона просто має свою частку страждання.}
\vcom{Найгостріше місце глави: під підозру потрапляє не лише прив'язаність до тіла, а й зацікавленість у звільненні. abhimāna — слово з присмаком претензії, і саме воно стоїть поруч із mokṣa. Друга частина вірша — тверезий висновок: із двох претензій не виростає ані знання (jñānī), ані практика (yogī), лишається тільки частка страждання (duḥkha-bhāj). Тут mukti не названо — є mokṣa; ані ātman, ані cit у вірші немає. Традиційні коментарі часто перетворюють це на заклик відмовитися від бажання звільнення; текст робить слабше й точніше твердження — він заперечує претензію, а не саму мету.}
\vpar{Українська паралель (наша інтерпретація)}{Памфіл Юркевич, «Серце і його значення в духовному житті людини» (1860).}{Наша паралель: коли вірш каже, що претензія на звільнення при відчутті тіла як свого нічого не дає, він розводить слово й стан. Юркевич розводить їх так само рішуче. Наше порівняння — за способом розрізняти.}
\vfig{ilustratsii/16-10-witness.pdf}{Претензія на волю й відчуття тіла як свого}{Схема наша, не історична ілюстрація.}

\vsec{16.11}
\vdeva{हरो यद्य् उपदेष्टा ते । हरिः कमलजो ऽपि वा । तथापि न तव स्वास्थ्यं । सर्वविस्मरणाद् ऋते ॥}
\vtranslit{харо ядй упадешта те | харіх камаладжо ’пі ва | татхапі на тава свастхям | сарвавісмаранад ріте ||}
\vperek{Навіть якби твоїм наставником був сам Шіва, Вішну чи Брахма, — спокою тобі не буде, поки не забудеш усе.}
\vcom{Глава закінчується тим самим реченням, що й почалася (пор. 16.1), з однією заміною: замість «наук» — імена найвищих богів. Це обрамлення робить висновок сильнішим: не допомогло знання — не допоможе й найвищий авторитет, бо обидва додають, а умова вимагає віднімання. Ототожнення Гара з Шівою, Гарі з Вішну і Камаладжі з Брахмою — загальноприйняте в традиції, і текст його не пояснює; тут воно подане як тлумачення, а не як сказане у вірші. Слово svāsthya повторене з 16.1 дослівно, і саме воно, а не mukti, є метою в цій главі.}
\vpar{Українська паралель (ПОТРЕБУЄ ПЕРЕВІРКИ)}{українська релігійна проза, не визначено.}{ПОТРЕБУЄ ПЕРЕВІРКИ: думка, що й самі боги не дадуть спокою, поки людина не забуде все, близька до української релігійної прози, але конкретного твору ми не перевірили.}
\vfig{ilustratsii/16-11-neti.pdf}{Навіть боги-наставники не дають спокою}{Схема наша, не історична ілюстрація.}


% ЗГЕНЕРОВАНО scripts/build_format_v2_chapter.py — не редагувати вручну.
% Джерело: editorial/cards/17.*.json, format-v2/rozdil-17.json, paraleli.json.

\naukchapter{17}{20}{Глава 17. Звільнений}
\naukintro{Сімнадцята глава — портрет, а не настанова. У ній майже немає наказів: вірші описують, що робить і чого не робить звільнений, і щоразу повертаються до одного — він у собі. Найдивніші місця — 17.10 і 17.20: текст прямо відмовляється назвати його стан, і ця відмова тут не поетична фігура, а точність.}
\vsec{17.1}
\vdeva{तेन ज्ञानफलं प्राप्तं । योगाभ्यासफलं तथा । तृप्तः स्वच्छेन्द्रियो नित्यम् । एकाकी रमते तु यः ॥}
\vtranslit{тена джнянапхалам праптам | йоґабгясапхалам татха | тріптах сваччендрійо нітям | екакі рамате ту ях ||}
\vperek{Хто вдоволений, чиї чуття прозорі й хто завжди сам — той і справді радіє: він уже має плід знання і плід вправ у йозі.}
\vcom{Глава 17 описує того, хто вже пройшов, і починається з переліку, а не з доказу. Показово, що «плід знання» і «плід йоги» поставлено поруч як одне й те саме досягнення: текст не протиставляє пізнання й практику, він каже, що обидві дали те саме. Слово tṛpta, «вдоволений, насичений», тут ключове: воно описує не досягнення, а відсутність голоду. Поруч стоїть ekākī, «самотній», — не як сум, а як відсутність потреби в іншому. Термін toṣa, який напрошується як відповідник, у вірші не вжитий: є tṛpta від √tṛp; це різні, хоч і близькі слова, і змішувати їх не слід.}
\vpar{Українська паралель (наша інтерпретація)}{Григорій Сковорода, «Розмова пʼяти подорожніх про істинне щастя в житті».}{Наша паралель: у Сковороди вдоволеність теж постає як плід, а не як вихідна умова. Наше порівняння — типологічне.}
\vfig{ilustratsii/17-1-balance.pdf}{Дві шальки: плід знання і плід вправ}{Схема наша, не історична ілюстрація.}

\vsec{17.2}
\vdeva{न कदाचिज् जगत्य् अस्मिंस् । तत्त्वज्ञो हन्त खिद्यति । यत एकेन तेनेदं । पूर्णं ब्रह्माण्डमण्डलम् ॥}
\vtranslit{на кадачідж джаґатй асмімс | таттваджньо ханта кхідяті | ята екена тенедам | пурнам брахмандамандалам ||}
\vperek{Ніколи в цьому світі не мучиться той, хто знає істину: о, тим одним сповнене все — усе коло світів.}
\vcom{Аргумент вірша — від повноти до відсутності страждання: мучитися може лише той, кому чогось бракує, а бракувати нічого, якщо ціле вже сповнене. «Коло Брахми» (brahmāṇḍa-maṇḍala) — не філософський термін, а образ усесвіту як замкненої кулі; цим підкреслено, що поза ним нічого немає й тому додати нічого. Слово tattva тут означає «істинне, дійсне», а не «категорію» школи санкх'я. Вигук hanta зберігає розмовну інтонацію, і в перекладі він вартий збереження: це не доказ, а подив. Про ātman, cit чи sākṣin вірш не говорить.}
\vpar{Українська паралель (наша інтерпретація)}{Памфіл Юркевич, «Серце і його значення в духовному житті людини» (1860).}{Наша паралель: коли вірш каже, що той, хто знає істину, не мучиться, він ставить знання в джерело стану, а не в наслідок. Юркевич шукає це джерело інакше — наше порівняння за будовою думки.}
\vfig{ilustratsii/17-2-witness.pdf}{Усе коло світів у центрі — і той, хто не мучиться}{Схема наша, не історична ілюстрація.}

\vsec{17.3}
\vdeva{न जातु विषयाः के ऽपि । स्वारामं हर्षयन्त्य् अमी । सल्लकीपल्लवप्रीतम् । इवेभं निम्बपल्लवाः ॥}
\vtranslit{на джату вішаях ке ’пі | сварамам харшаянтй амі | саллакіпаллавапрітам | івебгам німбапаллавах ||}
\vperek{Жодні чуттєві речі не тішать того, кому втіха — у ньому самому, як не тішить слона, що любить пагони салакі, гірке листя німби.}
\vcom{Порівняння побудоване на смаку: слон має улюблену їжу й тому не бере іншої, хоч вона й лежить поруч. Так і тут — відмова не від болю й не з аскези, а від того, що не смакує. Це важливо для розуміння глави: йдеться не про придушення бажання, а про зміну мірила. Обидві рослини названі конкретно, і це робить образ побутовим, а не абстрактним; жодного філософського терміна в цьому вірші немає — ні ātman, ні cit, ні mukti.}
\vpar{Українська паралель (наша інтерпретація)}{Іван Вишенський, «Книжка» (кінець XVI ст.).}{Наша паралель: у Вишенського втіха ззовні теж не замінює втіхи всередині. Наше порівняння — типологічне.}
\vfig{ilustratsii/17-3-balance.pdf}{Дві шальки: чуттєві речі й втіха всередині}{Схема наша, не історична ілюстрація.}

\vsec{17.4}
\vdeva{यस् तु भोगेषु भुक्तेषु । न भवत्य् अधिवासितः । अभुक्तेषु निराकाङ्क्षी । तादृशो भवदुर्लभः ॥}
\vtranslit{яс ту бгоґешу бгуктешу | на бгаватй адгівасітах | абгуктешу ніраканкші | тадрішо бгавадурлабгах ||}
\vperek{Хто вже зазнаних утіх не тягне за собою, як просякнутий їхнім запахом, і не прагне незазнаних — така людина в цьому світі рідкість.}
\vcom{Вірш розрізняє дві небезпеки, які зазвичай зливаються: пам'ять про зазнану втіху (adhivāsita, «просоченість») і потяг до ще не зазнаної. Перша тонша за другу, бо не вимагає жодного бажання — досить, щоб минуле лишило смак. Саме тому такий стан названо рідкісним: відмовитися від того, чого ще немає, легко; перестати носити в собі те, що було, — важко. Слово vairāgya тут не вжите, хоч саме його зазвичай ставлять на це місце; текст обходиться описом.}
\vpar{Українська паралель (наша інтерпретація)}{Тарас Шевченко, «І мертвим, і живим…» (1845).}{Наша паралель: у Шевченка рідкість людини, яка не тягне за собою зазнаного й не прагне незазнаного, теж названа прямо. Наше порівняння — типологічне.}
\vfig{ilustratsii/17-4-neti.pdf}{Зазнане й незазнане не тягне за собою}{Схема наша, не історична ілюстрація.}

\vsec{17.5}
\vdeva{बुभुक्षुर् इह संसारे । मुमुक्षुर् अपि दृश्यते । भोगमोक्षनिराकाङ्क्षी । विरलो हि महाशयः ॥}
\vtranslit{бубгукшур іха самсаре | мумукшур апі дрішяте | бгоґамокшаніраканкші | вірало хі махашаях ||}
\vperek{У цьому світі видно і тих, хто голодний до нього, і тих, хто прагне звільнення. Але хто не прагне ні втіх, ні звільнення — такий великий і рідкісний.}
\vcom{Два дезидеративи (bubhukṣu, mumukṣu) поставлено в один граматичний ряд, і цим сказано головне: бажання звільнитися має ту саму форму, що й бажання насолоди. Різниця лише в предметі, не в структурі. Тому третій член ряду — той, хто не прагне ані того, ані того, — і є рідкісний. Це не пропозиція байдужості: вірш не каже, що така людина чогось не робить, він каже, що вона не прагне. Традиційні коментарі часто додають, що такий уже вільний; сам текст цього слова не вживає — ні mukta, ні mukti тут немає.}
\vpar{Українська паралель (наша інтерпретація)}{Ліна Костенко, «Маруся Чурай» (1979).}{Наша паралель: у Костенко той, хто не прагне ні втіх, ні волі, теж постає як найрідкісніший. Наше порівняння — за будовою думки.}
\vfig{ilustratsii/17-5-motion.pdf}{Дві дороги прагнення — і третя, порожня}{Схема наша, не історична ілюстрація.}

\vsec{17.6}
\vdeva{धर्मार्थकाममोक्षेषु । जीविते मरणे तथा । कस्याप्य् उदारचित्तस्य । हेयोपादेयता न हि ॥}
\vtranslit{дгармартхакамамокшешу | джівіте маране татха | касяпй ударачіттася | хейопадеята на хі ||}
\vperek{Для великодушної людини немає «відкинути» й «узяти» — ні щодо дгарми, артги, ками й мокші, ні щодо життя й смерті.}
\vcom{Той самий перелік, що в 16.5, але розширений життям і смертю: текст послідовно відмовляється визнавати будь-яку пару за ієрархію. Слово udāracitta, «великодушний», — єдина оцінка в вірші, і вона стосується не поведінки, а обсягу: вузьке серце мусить вибирати, широке не має потреби. Зверніть увагу, що мокша знову стоїть у переліку цілей, а не як стан — той самий ужиток, що в 16.5. Термінів ātman, cit, sākṣin, vairāgya тут немає.}
\vpar{Українська паралель (наша інтерпретація)}{Микола Зеров, «Камена» (1924).}{Наша паралель: у Зерова відсутність відкинути й узяти теж є ознакою великодушності, а не байдужості. Наше порівняння — типологічне.}
\vfig{ilustratsii/17-6-neti.pdf}{Шість пар відкинуті: немає ні відкинути, ні взяти}{Схема наша, не історична ілюстрація.}

\vsec{17.7}
\vdeva{वाञ्छा न विश्वविलये । न द्वेषस् तस्य च स्थितौ । यथा जीविकया तस्माद् । धन्य आस्ते यथा सुखं ॥}
\vtranslit{ваньчча на вішвавілає | на двешас тася ча стхітау | ятха джівікая тасмад | дганя асте ятха сукхам ||}
\vperek{Він не бажає, щоб світ загинув, і не має відрази до того, що світ триває. Тому блаженний сидить собі спокійно, як завгодно, — і живе з того, що має.}
\vcom{Вірш відкидає обидві есхатологічні пристрасті: і прагнення кінця світу, і відразу до його продовження. Це не байдужість до долі світу, а відмова вкладати в неї власне бажання. Друга половина — про спосіб життя: āste yathā sukham, «сидить як завгодно», поруч із yathā jīvikayā, «як дає прожиток». Разом це опис зовнішньої невимогливості, а не аскетичної програми; жодного слова про практику, обітницю чи стан життя (āśrama) тут немає.}
\vpar{Українська паралель (наша інтерпретація)}{Олександр Довженко, «Зачарована Десна» (1956).}{Наша паралель: у Довженка вміння жити з того, що маєш, теж подане як спокій, а не як бідність. Наше порівняння — за образом.}
\vfig{ilustratsii/17-7-space.pdf}{Небо без хмар: світ може йти й стояти}{Схема наша, не історична ілюстрація.}

\vsec{17.8}
\vdeva{कृतार्थो ऽनेन ज्ञानेनेत्य् । एवं गलितधीः कृती । पश्यन् शृण्वन् स्पृशन् जिघ्रन्न् । अश्नन्न् आस्ते यथा सुखं ॥}
\vtranslit{крітартхо ’нена джняненетй | евам ґалітадгіх кріті | пашян шрінван спрішан джіґгранн | ашнанн асте ятха сукхам ||}
\vperek{«Цим знанням я досяг мети» — і з цим, уже розтанувши розумом, він живе собі як завгодно: дивиться, слухає, торкається, нюхає, їсть.}
\vcom{Вірш цікавий тим, що єдине висловлювання, яке він цитує («цим знанням я досяг мети»), подане як думка, що вже відпрацювала себе: далі — не аргумент, а перелік найпростіших дій. Дієприкметники paśyan, śṛṇvan, spṛśan, jighran, aśnan ідуть без жодного оцінного слова; саме ця нейтральність і є змістом. galitadhī, «розталий розум», — образ, що не має філософського статусу терміна, і його не слід плутати з citta-vṛtti-nirodha інших текстів: тут нічого не спиняють, а лише констатують, що тримати вже нічого. Порядок дій (бачення, слух, дотик, нюх, їжа) повторений у 17.12, де до нього додано ще рух і мову.}
\vpar{Українська паралель (наша інтерпретація)}{Павло Тичина, «Сонячні кларнети» (1918).}{Наша паралель: у Тичини дивитися, слухати, торкатися теж стає способом жити, коли мету вже досягнуто. Наше порівняння — за образом.}
\vfig{ilustratsii/17-8-motion.pdf}{П'ять чуттів як спосіб жити}{Схема наша, не історична ілюстрація.}

\vsec{17.9}
\vdeva{शून्या दृष्टिर् वृथा चेष्टा । विकलानीन्द्रियाणि च । न स्पृहा न विरक्तिर् वा । क्षीणसंसारसागरे ॥}
\vtranslit{шуня дріштір врітха чешта | вікаланіндріяні ча | на спріха на вірактір ва | кшінасамсарасаґаре ||}
\vperek{Порожній погляд, марні рухи, чуття ніби понівечені; ні прагнення, ні відвернення — бо океан сансари в ньому висох.}
\vcom{Опис зовні непривабливий, і це навмисне: збоку такий стан виглядає як згасання. Три ознаки — порожній погляд, марна діяльність, немічні чуття — узяті з того, як людину бачать інші, а не з її самовідчуття. Локатив абсолютний kṣīṇasaṃsārasāgare ставить причину в кінці, і вона перевертає оцінку: це не втрата сил, а вичерпаність океану, який раніше треба було перепливати. Пара spṛhā / virakti знову знята — обидві названі як те, чого немає. Тут доречно сказати прямо: ані mukti, ані ātman, ані cit у вірші немає, і саме тому він не дає приводу назвати цей стан «звільненням» — це слово з'явиться в главі 20.}
\vpar{Українська паралель (ПОТРЕБУЄ ПЕРЕВІРКИ)}{українська поезія XX ст., не визначено.}{ПОТРЕБУЄ ПЕРЕВІРКИ: образ порожнього погляду й ніби понівечених чуттів близький до української поезії, але конкретного твору ми не звірили.}
\vfig{ilustratsii/17-9-neti.pdf}{Прагнення й відвернення відпали: океан висох}{Схема наша, не історична ілюстрація.}

\vsec{17.10}
\vdeva{न जागर्ति न निद्राति । नोन्मीलति न मीलति । अहो परदशा क्वापि । वर्तते मुक्तचेतसः ॥}
\vtranslit{на джаґарті на нідраті | нонмілаті на мілаті | ахо парадаша квапі | вартате муктачетасах ||}
\vperek{Він не пильнує і не спить, не розплющує й не заплющує очей. О диво: розум вільної людини перебуває в якомусь зовсім іншому стані.}
\vcom{Вірш будує заперечення парами — не-сон / не-пильність, не-розплющення / не-заплющення — так, що кожна пара замикає обидві можливості. Описати такий стан звичайними словами неможливо, і текст цього не приховує: він каже лише paradaśā, «інший стан», не називаючи його. Це місце легко перетворити на містичну загадку; варто пам'ятати, що йдеться про опис збоку, а не про техніку. Слова «свідомість» (cit) чи «свідок» (sākṣin) тут немає — є muktacetas, «вільний розумом», і це єдине, що сказано про того, хто так перебуває.}
\vpar{Українська паралель (ПОТРЕБУЄ ПЕРЕВІРКИ)}{українська містична проза, не визначено.}{ПОТРЕБУЄ ПЕРЕВІРКИ: опис стану, що не є ні сном, ні пильністю, найближчий до української містичної прози, але назвати твір без звірки ми не можемо.}
\vfig{ilustratsii/17-10-witness.pdf}{Стан поза сном і пильністю}{Схема наша, не історична ілюстрація.}

\vsec{17.11}
\vdeva{सर्वत्र दृश्यते स्वस्थः । सर्वत्र विमलाशयः । समस्तवासना मुक्तो । मुक्तः सर्वत्र राजते ॥}
\vtranslit{сарватра дрішяте свастхах | сарватра вімалашаях | самаставасана мукто | муктах сарватра раджате ||}
\vperek{Усюди він у собі, усюди його нутро чисте; вільний від усіх звичних слідів, звільнений сяє всюди.}
\vcom{Вірш побудований на повторі sarvatra і на подвійному mukta — «вільний [від васанових слідів]» і «звільнений». Це єдине місце глави, де сказано, від чого саме вільний: від vāsanā, звичних залишкових схильностей. Слово це важливе й легко перекладається неточно: vāsanā — не думка й не бажання, а слід, який минуле лишає в способі реагування. Порівняння з 16.7 показує, що там ішлося про той самий механізм з іншого боку: жадоба живиться станом без роздуму. rājate, «сяє», — дієслово панування, а не сяйва лампи; у перекладі збережено саме це.}
\vpar{Українська паралель (наша інтерпретація)}{Василь Стус, «Палімпсести».}{Наша паралель: у Стуса звільненість від звичних слідів теж описана як сяйво, а не як порожнеча. Наше порівняння — типологічне.}
\vfig{ilustratsii/17-11-flame.pdf}{Полумʼя без слідів: сяє всюди}{Схема наша, не історична ілюстрація.}

\vsec{17.12}
\vdeva{पश्यन् शृण्वन् स्पृशन् जिघ्रन्न् । अश्नन् गृह्णन् वदन्व्रजन् । ईहितानीहितैर् मुक्तो । मुक्त एव महाशयः ॥}
\vtranslit{пашян шрінван спрішан джіґгранн | ашнан ґріхнан ваданвраджан | іхітаніхітайр мукто | мукта ева махашаях ||}
\vperek{Дивиться, слухає, торкається, нюхає, їсть, бере, говорить, ходить — і вільний від «бажаного» й «небажаного». Він справді звільнений, цей великий душею.}
\vcom{Перелік дій тут довший, ніж у 17.8: додано брати, говорити, ходити. Це вже не лише сприйняття, а й активне втручання у світ — і воно нічим не відрізняється від решти. Повтор muktaḥ ... mukta eva звучить як поправка до читача, який готовий заперечити: мовляв, як це вільний, коли він бере й говорить. Відповідь — уточнення, від чого саме він вільний: від пари īhita / anīhita. Це не звільнення від дій, а звільнення від поділу дій на жадані й нежадані.}
\vpar{Українська паралель (наша інтерпретація)}{Михайло Коцюбинський, «Intermezzo» (1908).}{Наша паралель: у Коцюбинського звичайні дії теж не суперечать свободі від бажаного й небажаного. Наше порівняння — за будовою думки.}
\vfig{ilustratsii/17-12-motion.pdf}{Звичайні дії без бажаного й небажаного}{Схема наша, не історична ілюстрація.}

\vsec{17.13}
\vdeva{न निन्दति न च स्तौति । न हृष्यति न कुप्यति । न ददाति न गृह्णाति । मुक्तः सर्वत्र नीरसः ॥}
\vtranslit{на ніндаті на ча стауті | на хрішяті на куп'яті | на дадаті на ґріхнаті | муктах сарватра нірасах ||}
\vperek{Не гудить і не хвалить, не тішиться й не гнівається, не дає й не бере — звільнений усюди без смаку до всього.}
\vcom{Шість дієслів ідуть трьома парами, і кожна пара охоплює протилежності: осуд і хвала, радість і гнів, давання й брання. Схема та сама, що в 16.8 і 16.10: не вибір між протилежностями, а вихід із самої пари. Останнє слово, nīrasa, навмисно негарне — воно забирає в цього стану будь-яку привабливість. Це варто зберегти в перекладі: текст не обіцяє блаженства, він описує відсутність смаку. Слів mukti чи ātman тут немає — лише muktaḥ як означення.}
\vpar{Українська паралель (наша інтерпретація)}{Григорій Сковорода, «Вдячний Еродій» (трактат).}{Наша паралель: у Сковороди відсутність смаку до всього теж не є втратою смаку, а зміною його предмета. Наше порівняння — типологічне.}
\vfig{ilustratsii/17-13-balance.pdf}{Дві шальки: хвала й осуд — обидві порожні}{Схема наша, не історична ілюстрація.}

\vsec{17.14}
\vdeva{सानुरागां स्त्रियं दृष्ट्वा । मृत्युं वा समुपस्थितम् । अविह्वलमनाः स्वस्थो । मुक्त एव महाशयः ॥}
\vtranslit{санураґам стріям дріштва | мрітюм ва самупастхітам | авіхваламанах свастхо | мукта ева махашаях ||}
\vperek{Побачивши жінку, повну пристрасті, або смерть, що вже насувається, він не збентежиться: він у собі, він справді звільнений, цей великий душею.}
\vcom{Вірш випробовує рівність на двох найсильніших подразниках — бажанні й страхові — і ставить їх в один граматичний ряд. Це не твердження про безстатевість чи безстрашність як чесноти: сказано лише, що розум не хитається (avihvala). Зверніть увагу, що йдеться про зустріч, а не про відмову від неї: текст не радить уникати ні того, ні того. Традиційні коментарі читають це як доказ рівності до приємного й болісного; сам вірш називає лише два конкретні випадки, і узагальнення тут — робота читача.}
\vpar{Українська паралель (наша інтерпретація)}{Леся Українка, «Камінний господар» (1912).}{Наша паралель: у Лесі Українки зустріч зі смертю теж не збентежує того, хто вже в собі. Наше порівняння — за будовою думки.}
\vfig{ilustratsii/17-14-witness.pdf}{Пристрасть і смерть навколо — він у собі}{Схема наша, не історична ілюстрація.}

\vsec{17.15}
\vdeva{सुखे दुःखे नरे नार्यां । संपत्सु च विपत्सु च । विशेषो नैव धीरस्य । सर्वत्र समदर्शिनः ॥}
\vtranslit{сукхе духкхе наре нар'ям | сампатсу ча віпатсу ча | вішешо найва дгірася | сарватра самадаршінах ||}
\vperek{У щасті й у горі, в чоловікові й у жінці, в удачі й у біді — для стійкого, хто всюди бачить однаково, різниці немає.}
\vcom{Перелік пар тут не філософський, а побутовий: щастя й горе, чоловік і жінка, удача й біда. Пару «чоловік і жінка» не слід вичитувати як твердження про стать у метафізичному сенсі: вона стоїть у ряді з іншими парами й означає лише, що оцінка не змінюється залежно від того, з ким маєш справу. sama-darśin — «той, хто бачить рівно», і це єдиний пізнавальний термін у вірші. Слова ātman, cit, sākṣin тут немає; «рівність бачення» не тотожна ані «свідку», ані «свідомості».}
\vpar{Українська паралель (наша інтерпретація)}{Іван Франко, «Смерть Каїна» (1889).}{Наша паралель: у Франка різниця між щастям і горем теж зникає для того, хто бачить однаково. Наше порівняння — типологічне.}
\vfig{ilustratsii/17-15-balance.pdf}{Дві шальки, що стоять рівно: щастя й горе}{Схема наша, не історична ілюстрація.}

\vsec{17.16}
\vdeva{न हिंसा नैव कारुण्यं । नौद्धत्यं न च दीनता । नाश्चर्यं नैव च क्षोभः । क्षीणसंसरणे नरे ॥}
\vtranslit{на хімса найва каруням | науддгатям на ча діната | нашчар'ям найва ча кшобгах | кшінасамсаране наре ||}
\vperek{Ні насильства, ні жалю, ні зарозумілості, ні приниженості, ні здивування, ні збурення — немає цього в людині, в якої скінчилося блукання.}
\vcom{Найнезвичніше тут — заперечення kāruṇya, жалю, поруч із запереченням hiṃsā, насильства. Звична етична схема вимагає одного й забороняє друге; вірш знімає обидва як ознаки того, хто ще всередині оцінки. Заперечено також āścarya, подив: стан не містить нічого несподіваного, бо немає точки, з якої щось могло б здивувати. Це місце легко прочитати як виправдання байдужості до чужого болю — і такий докір до тексту висловлювали; сам вірш, однак, описує не поведінку, а відсутність реакцій у людини, в якої вже немає нічого попереду.}
\vpar{Українська паралель (наша інтерпретація)}{Ольга Кобилянська, «Земля» (1902).}{Наша паралель: у Кобилянської відсутність насильства й зарозумілості теж є ознакою людини, в якої скінчилося блукання. Наше порівняння — типологічне.}
\vfig{ilustratsii/17-16-neti.pdf}{Шість рухів, яких немає в звільненого}{Схема наша, не історична ілюстрація.}

\vsec{17.17}
\vdeva{न मुक्तो विषयद्वेष्टा । न वा विषयलोलुपः । असंसक्तमना नित्यं । प्राप्ताप्राप्तम् उपाश्नुते ॥}
\vtranslit{на мукто вішаядвешта | на ва вішаялолупах | асамсактамана нітям | праптапраптам упашнуте ||}
\vperek{Звільнений не ненавидить чуттєвих речей і не жадібний до них; його розум ні до чого не прив'язаний, і він бере те, що прийшло, і те, що не прийшло.}
\vcom{Вірш повторює схему 16.6, але додає дію: не лише немає ненависті й жадоби, а й є прийняття того, що трапляється. Це відповідь на очевидне заперечення — якщо немає потягу, чи не стане людина пасивною? Текст відповідає дієсловом upāśnute, «споживає, приймає»: дія не зникає, зникає прив'язаність до її результату. Тут знову немає ні vairāgya, ні mukti як термінів: muktaḥ — означення, а не назва стану.}
\vpar{Українська паралель (наша інтерпретація)}{Євген Плужник, «Рання осінь» (1927).}{Наша паралель: у Плужника розум, ні до чого не прив'язаний, теж бере те, що прийшло, і не сумує за тим, що не прийшло. Наше порівняння — за будовою думки.}
\vfig{ilustratsii/17-17-motion.pdf}{Бере й те, що прийшло, і те, що не прийшло}{Схема наша, не історична ілюстрація.}

\vsec{17.18}
\vdeva{समाधानासमाधान । हिताहितविकल्पनाः । शून्यचित्तो न जानाति । कैवल्यम् इव संस्थितः ॥}
\vtranslit{самадганасамадгана | хітахітавікалпанах | шунячітто на джанаті | кайвалям іва самстхітах ||}
\vperek{Розрізнень «зібраний я чи ні», «корисно це чи шкідливо» порожній розум уже не знає: він стоїть ніби в самотності.}
\vcom{Порожній розум (śūnyacitta) тут не мета, а причина незнання розрізнень: якщо немає з чим порівнювати, нема чого й розрізняти. Важливо, що серед переліку стоїть samādhāna — той самий корінь, що в samādhi, який у 16.2 був одним із трьох занять. Тут він опиняється в ряді vikalpanā, тобто серед уявних побудов. kaivalya вжите обережно, з iva, «наче»: текст не каже, що це і є звільнення, а каже, що стан схожий на нього. У главі 20 kaivalya з'явиться знову — уже в переліку того, чого немає.}
\vpar{Українська паралель (наша інтерпретація)}{Володимир Свідзінський, «Медобір» (1975, посмертно).}{Наша паралель: у Свідзінського самотність теж постає як стан, у якому розрізнення втрачають сенс. Наше порівняння — за образом.}
\vfig{ilustratsii/17-18-space.pdf}{Самотність порожнього розуму}{Схема наша, не історична ілюстрація.}

\vsec{17.19}
\vdeva{निर्ममो निरहंकारो । न किंचिद् इति निश्चितः । अन्तर्गलितसर्वाशः । कुर्वन्न् अपि करोति न ॥}
\vtranslit{нірмамо нірахамкаро | на кімчід іті нішчітах | антарґалітасарвашах | курванн апі кароті на ||}
\vperek{Без «моє» і без «я роблю», з певністю, що немає нічого, — усі його надії розтанули зсередини, і він навіть діючи не діє.}
\vcom{Вірш зводить докупи три заперечення — «моє», «я роблю», «щось є» — і робить висновок про дію. Остання пада, kurvann api karoti na, — класична формула недіяння в діянні; вона не означає, що тіло спинилося, а що дію не приписують собі. Саме тут свідки розходяться: SARIT має na lipyate, «не зазнає плям», і подає karoti na в дужках як варіант. Обидві формули відомі з інших текстів, і вибір між ними змінює акцент — онтологічний («не діє») чи етичний («не плямується»). Яке читання первісне, за двома свідками не встановлюється.}
\vpar{Українська паралель (наша інтерпретація)}{Богдан-Ігор Антонич, «Зелена Євангелія» (1938).}{Наша паралель: в Антонича дія без діяча теж описана як природний рух. Наше порівняння — типологічне.}
\vfig{ilustratsii/17-19-neti.pdf}{Моє, я роблю й надії розтанули}{Схема наша, не історична ілюстрація.}

\vsec{17.20}
\vdeva{मनःप्रकाशसम्मोह । स्वप्नजाड्यविवर्जितः । दशां काम् अपि संप्राप्तो । भवेद् गलितमानसः ॥}
\vtranslit{манахпракашасаммоха | свапнаджадявіварджітах | дашам кам апі сампрапто | бгавед ґалітаманасах ||}
\vperek{Він вільний і від ясності розуму, і від затьмарення, і від сну, і від заціпеніння — і досяг якогось [невідомого] стану, цей чоловік із розталим розумом.}
\vcom{Перелік станів розуму навмисно незвичний: до нього входять не лише затьмарення, сон і заціпеніння, а й prakāśa — ясність. Тобто заперечено не тільки погані стани, а й добрий: ясність розуму тут така сама ознака розуму, як і сон. Останній стан названо неозначено, daśāṃ kām api, і текст не намагається його описати. Форма bhavet (оптатив) додає обережності: «міг би стати», а не «став». Це останній вірш глави, і він лишає читача без назви того, про що йшлося двадцять віршів.}
\vpar{Українська паралель (ПОТРЕБУЄ ПЕРЕВІРКИ)}{українська філософська проза, не визначено.}{ПОТРЕБУЄ ПЕРЕВІРКИ: опис стану поза ясністю, затьмаренням, сном і заціпенінням найближчий до української філософської прози, але конкретного твору ми не перевірили.}
\vfig{ilustratsii/17-20-neti.pdf}{Чотири стани розуму відпали: лишається невідомий}{Схема наша, не історична ілюстрація.}


% ЗГЕНЕРОВАНО scripts/build_format_v2_chapter.py — не редагувати вручну.
% Джерело: editorial/cards/18.*.json, format-v2/rozdil-18.json, paraleli.json.

\naukchapter{18}{100}{Глава 18. Сто віршів про того, хто вільний}
\naukintro{Вісімнадцята глава — половина всієї книги за обсягом: сто віршів. Вона не будує вчення, а перебирає його з усіх боків: кожен вірш — це та сама думка, сказана інакше, і саме повтор і є формою. Читачеві варто знати, що тут немає розвитку від вірша до вірша; є один стан, описаний сто разів, і кілька віршів (18.31–18.39) прямо говорять про те, що вправляння й зусилля цього стану не дають.}
\vsec{18.1}
\vdeva{यस्य बोधोदये तावत् स्वप्नवद् भवति भ्रमः । तस्मै सुखैकरूपाय नमः शान्ताय तेजसे ॥}
\vtranslit{яся бодгодає тават свапнавад бгаваті бграмах | тасмай сукхайкарупая намах шантая теджасе ||}
\vperek{Той, з чиїм пробудженням усвідомлення омана стає наче сон, — тому, чия єдина природа щастя, уклін: втихлому, сяйному.}
\vcom{Глава 18 відкривається єдиним у книзі `namaḥ` — поклоном. До цього текст усі 17 глав звертався до читача на «ти», а тут раптом з'являється третя особа: уклін комусь іншому. `tāvat … svapnavat` — не «омана зникла», а «омана стала наче сон»: вона не анульована, а переведена в інший регістр. Переклад мусить тримати саме це, а не «звільнився від омани».}
\vpar{Українська паралель (наша інтерпретація)}{Григорій Сковорода, «Сад божественних пісень».}{Наша паралель: у Сковороди пробудження теж описується як те, з чим омана стає сном, а не як те, що додає нових знань. Наше порівняння — типологічне.}
\vfig{ilustratsii/18-1-flame.pdf}{Полумʼя пробудження: омана стає сном}{Схема наша, не історична ілюстрація.}

\vsec{18.2}
\vdeva{अर्जयित्वाखिलान् अर्थान् भोगान् आप्नोति पुष्कलान् । न हि सर्वपरित्यागम् अन्तरेण सुखी भवेत् ॥}
\vtranslit{арджаїтвакхілан артхан бгоґан апноті пушкалан | на хі сарвапарітяґам антарена сукхі бгавет ||}
\vperek{Здобувши всі блага, людина отримує рясні насолоди — але без повного зречення всього вона не стане щасливою.}
\vcom{Вірш прямо зіставляє накопичення й щастя і не заперечує першого: насолоди справді приходять (āpnoti, дійсний спосіб). Заперечується лише висновок — sukhī bhavet стоїть в optativus («став би»), тобто це не факт, а умова, яка не виконується. Найсильніше слово — sarvaparityāga («зречення ВСЬОГО»): не більшості, не зайвого, а всього без останку. ПРИМІТКА ПРО ОБСЯГ. Глава 18 — найдовша в книзі: 99 віршів, третина тексту. Вона опрацьовується БЛОКАМИ. Ця картка належить до блоку 18.2–18.10; вірш 18.1 опрацьовано окремо (PR АҐ-18.1), а 18.11–18.99 ще не мають карток.}
\vpar{Українська паралель (наша інтерпретація)}{Григорій Сковорода, «Убогий Жайворонок».}{Наша паралель: у Сковороди без зречення теж немає щастя, хоч би скільки благ зібралося. Наше порівняння — типологічне, не твердження про джерело.}
\vfig{ilustratsii/18-2-neti.pdf}{Блага й насолоди без зречення}{Схема наша, не історична ілюстрація.}

\vsec{18.3}
\vdeva{कर्तव्यदुःखमार्तण्डज्वालादग्धान्तरात्मनः । कुतः प्रशमपीयूषधारासारम् ऋते सुखम् ॥}
\vtranslit{картав'ядухкхамартандаджваладаґдгантаратманах | кутах прашамапіюшадгарасарам ріте сукхам ||}
\vperek{У того, чий внутрішній Атман спалений полум'ям сонця страждання від належного, — звідки візьметься щастя без суті потоку нектару вгамування?}
\vcom{mārtaṇḍa («сонце») буквально означає «той, що народився з мертвого яйця» — сонячне божество; тут воно вжите як звичайна назва сонця, але образ жорстокий: сонце не світить, а палить. Складене kartavyaduḥkha («страждання від належного») — точна діагностика: палить не біль як такий, а думка про те, що ще треба зробити. Друга половина вводить pīyūṣa («нектар») і praśama («вгамування») як єдині ліки, і це та сама логіка, що в 2.16 («немає інших ліків»).}
\vfig{ilustratsii/18-3-flame.pdf}{Полумʼя страждання в серці}{Схема наша, не історична ілюстрація.}

\vsec{18.4}
\vdeva{भवो ऽयं भावनामात्रो न किंचित् परमार्थतः । नास्त्य् अभावः स्वभावानां भावाभावविभाविनाम् ॥}
\vtranslit{бгаво ’ям бгаванаматро на кімчіт парамартхатах | настй абгавах свабгаванам бгавабгававібгавінам ||}
\vperek{Це буття — тільки уявлення, у вищому сенсі — ніщо. І немає небуття для тих речей, що розрізняють буття й небуття.}
\vcom{Вірш робить подвійний хід, і його легко прочитати наполовину. Спершу: буття — лише bhāvanā («уявлення»), і за остаточною істиною (paramārthataḥ) воно ніщо. Далі — несподіване: abhāvaḥ nāsti, «небуття немає» — тобто заперечується не лише буття, а й його відсутність. І заперечення стосується саме тих, хто розрізняє (vibhāvin): для того, хто ставить питання «є чи немає», обидві відповіді хибні. Це найточніше формулювання того, що в 2.14 було подано через athavā («або ж»).}
\vfig{ilustratsii/18-4-illusion.pdf}{Буття під серпанком: уявлення}{Схема наша, не історична ілюстрація.}

\vsec{18.5}
\vdeva{न दूरं न च सङ्कोचाल् लब्धम् एवात्मनः पदम् । निर्विकल्पं निरायासं निर्विकारं निरञ्जनम् ॥}
\vtranslit{на дурам на ча санкочал лабдгам еватманах падам | нірвікалпам ніраясам нірвікарам ніраньджанам ||}
\vperek{Стопа Атмана не далеко і не досягається стисненням — вона вже досягнута: без розрізнень, без зусилля, без змін, без плями.}
\vcom{Заперечуються два способи промахнутися: уявити мету далекою (dūram) або досягти її напруженням (saṅkoca — буквально «стиснення», звуження себе). Далі чотири nir- означення стоять без дієслова й без сполучників — це не опис шляху, а перелік того, чого в досягнутому немає. nirāyāsa («без зусилля») стоїть другим, одразу після nirvikalpa, і це відповідь на saṅkoca: те, що досягнуте без стиснення, і далі не потребує зусиль.}
\vpar{Українська паралель (наша інтерпретація)}{Тарас Шевченко, «Подражаніє Екклезіястові» (1845).}{Наша паралель: у Шевченка те, що не далеко й не близько, теж не досягається зусиллям. Наше порівняння — за будовою думки.}
\vfig{ilustratsii/18-5-space.pdf}{Стопа Атмана: не далеко й не близько}{Схема наша, не історична ілюстрація.}

\vsec{18.6}
\vdeva{व्यामोहमात्रविरतौ स्वरूपादानमात्रतः । वीतशोका विराजन्ते निरावरणदृष्टयः ॥}
\vtranslit{в'ямохаматравіратау сварупаданаматратах | віташока віраджанте ніраваранадріштаях ||}
\vperek{Коли припиняється сама лише завороженість, це вже й є прийняття власної природи: ті, чий зір без покрову, сяють, вільні від смутку.}
\vcom{Двічі повторене mātra («лише») робить основну роботу: припиняється ЛИШЕ завороженість — нічого не додається, і приймається ЛИШЕ власна природа — нічого не набувається. Тобто між двома половинами немає переходу, вони те саме, побачене з двох боків. nirāvaraṇa («без покрову») — точне слово: покров не усувається зусиллям, він просто не накидається. І зір (dṛṣṭi) тут не метафора знання, а саме зір.}
\vfig{ilustratsii/18-6-neti.pdf}{Завороженість припинилася — зір без покрову}{Схема наша, не історична ілюстрація.}

\vsec{18.7}
\vdeva{समस्तं कल्पनामात्रम् आत्मा मुक्तः सनातनः । इति विज्ञाय धीरो हि किम् अभ्यस्यति बालवत् ॥}
\vtranslit{самастам калпанаматрам атма муктах санатанах | іті віджняя дгіро хі кім абгясяті балават ||}
\vperek{Знаючи, що все разом — тільки побудова уяви, а Атман вільний і споконвічний, — чого ж мудрий вправляється, як дитина?}
\vcom{abhyasyati — дієслово повторювальної вправи (звідси «абг'яса» в йозі), і саме воно тут висміяне. Порівняння bālavat («як дитина») не образливе, а точне: дитина повторює те саме, поки не навчиться, — а тут навчатися вже нема кому й нема чого. sanātana («споконвічний») стоїть при ātmā і знімає час: те, що не почалося, не потребує вправи.}
\vfig{ilustratsii/18-7-illusion.pdf}{Побудова уяви й вільний Атман}{Схема наша, не історична ілюстрація.}

\vsec{18.8}
\vdeva{आत्मा ब्रह्मेति निश्चित्य भावाभावौ च कल्पितौ । निष्कामः किं विजानाति किं ब्रूते च करोति किम् ॥}
\vtranslit{атма брахметі нішчітя бгавабгавау ча калпітау | нішкамах кім віджанаті кім бруте ча кароті кім ||}
\vperek{Встановивши, що Атман — це Брахман, а буття й небуття — вигадані, — що безбажанний знає, що говорить, що робить?}
\vcom{Три питання в кінці — про знання, мову й дію — охоплюють усі три способи, якими людина взагалі себе виявляє. Кожне залишене без відповіді, і це не риторика заради ефекту: якщо буття й небуття вигадані, то питання «що є» втрачає ґрунт, а разом із ним і три відповіді. Формула ātmā brahma («Атман — це Брахман») тут не висновок, а niścaya — встановлене рішення, після якого питання стають безпредметними.}
\vfig{ilustratsii/18-8-neti.pdf}{Буття й небуття вигадані: мовчання безбажанного}{Схема наша, не історична ілюстрація.}

\vsec{18.9}
\vdeva{अयं सो ऽहम् अयं नाहम् इति क्षीणा विकल्पनाः । सर्वम् आत्मेति निश्चित्य तूष्णीभूतस्य योगिनः ॥}
\vtranslit{аям со ’хам аям нахам іті кшіна вікалпанах | сарвам атметі нішчітя тушнібгутася йоґінах ||}
\vperek{У йога, що став безмовним, вичерпані розрізнення «це він», «це я», «це не я», — бо він встановив: «все — Атман».}
\vcom{Вірш побудований як речення без головного дієслова: kṣīṇāḥ («вичерпані») — присудок, але вся конструкція стоїть у родовому (tūṣṇībhūtasya yoginaḥ), тобто описує стан, а не подію. tūṣṇībhūta («той, хто став безмовним») — не обітниця мовчання, а наслідок: розрізнення вичерпалися, і говорити нема про що. Три формули в лапках — «це він», «це я», «це не я» — це не філософські тези, а щоденна робота розуму, названа вголос.}
\vpar{Українська паралель (наша інтерпретація)}{Митрополит Іларіон Київський, «Слово про закон і благодать» (XI ст.).}{Наша паралель: у київському Слові встановлення теж знімає розрізнення, а не додає нового. Наше порівняння — типологічне; датування цієї книги науково не встановлене.}
\vfig{ilustratsii/18-9-neti.pdf}{Три розрізнення відкинуті: усе — Атман}{Схема наша, не історична ілюстрація.}

\vsec{18.10}
\vdeva{न विक्षेपो न चैकाग्र्यं नातिबोधो न मूढता । न सुखं न च वा दुःखम् उपशान्तस्य योगिनः ॥}
\vtranslit{на вікшепо на чайкаґр'ям натібодго на мудгата | на сукхам на ча ва духкхам упашантася йоґінах ||}
\vperek{У вгамованого йога немає ні розсіяності, ні зосередженості, ні надмірного знання, ні затьмареності, ні щастя, ні страждання.}
\vcom{Три пари, і кожна заперечує обидва боки. Найважливіша перша: vikṣepa і aikāgrya — розсіяність і зосередженість — це те, чим зазвичай міряють стан розуму, і тут обидва відсутні, тобто міряти нічим. Друга пара (atibodha / mūḍhatā) забирає різницю між мудрим і дурним, третя (sukha / duḥkha) — між приємним і болісним. Це та сама схема, що в 19.4–19.7, але тут вона подана як опис, а не як питання.}
\vfig{ilustratsii/18-10-balance.pdf}{Дві шальки, яких немає у вгамованого}{Схема наша, не історична ілюстрація.}

\vsec{18.11}
\vdeva{स्वाराज्ये भैक्ष्यवृत्तौ च लाभालाभे जने वने । निर्विकल्पस्वभावस्य न विशेषो ऽस्ति योगिनः ॥}
\vtranslit{свараджє бгайкшявріттау ча лабгалабге джане ване | нірвікалпасвабгавася на вішешо ’сті йоґінах ||}
\vperek{У власному царстві чи в житті з подаяння, у здобутку й нездобутку, серед людей чи в лісі — для йогина, чия природа без домислювання, немає різниці.}
\vcom{Чотири пари протилежностей названі поіменно — царство й милостиня, здобуток і втрата, люди й ліс — щоб жодну не лишити поза дією. `svārājya` — не «своє царство» як володіння, а само-володарювання, стан, у якому немає над собою нікого; у главі 18 це слово трапляється не раз, і переклад має бути сталим. `viśeṣa` — «різниця», не «перевага».}
\vpar{Українська паралель (наша інтерпретація)}{Олександр Довженко, «Зачарована Десна» (1956).}{Наша паралель: у Довженка багатство й убогість місця теж не міняють того, хто вміє жити. Наше порівняння — за образом.}
\vfig{ilustratsii/18-11-balance.pdf}{Дві шальки: царство й життя з подаяння}{Схема наша, не історична ілюстрація.}

\vsec{18.12}
\vdeva{क्व धर्मः क्व च वा कामः क्व चार्थः क्व विवेकता । इदं कृतम् इदं नेति द्वन्द्वैर् मुक्तस्य योगिनः ॥}
\vtranslit{ква дгармах ква ча ва камах ква чартхах ква вівеката | ідам крітам ідам неті двандвайр муктася йоґінах ||}
\vperek{Де дгарма, де бажання, де користь, де вміння розрізняти? У йога, вільного від пар протилежностей, — «це зроблено, це ні» — немає нічого з цього.}
\vcom{Той самий зачин, що в 19.2, але з іншим висновком: там було «у мене, що стоїть у власній величі», тут — «у вільного від пар». Пара «зроблено / не зроблено» — це dvandva у найпростішому вигляді: оцінка дії за наслідком. Вірш каже, що разом із парами зникає й сама можливість зарахувати щось як зроблене.}
\vfig{ilustratsii/18-12-neti.pdf}{Дгарма, бажання, користь і розрізнення відпали}{Схема наша, не історична ілюстрація.}

\vsec{18.13}
\vdeva{कृत्यं किम् अपि नैवास्ति न कापि हृदि रञ्जना । यथाजीवनम् एवेह जीवन्मुक्तस्य योगिनः ॥}
\vtranslit{крітям кім апі найвасті на капі хріді раньджана | ятхаджіванам евеха джіванмуктася йоґінах ||}
\vperek{Немає нічого, що слід зробити, і жодної вподоби в серці — у йога, звільненого ще за життя, все йде просто як живеться, тут-таки.}
\vcom{jīvanmukta («звільнений за життя») — технічний термін адвайти: форма jīvanmuktasya засвідчена у 18.13, а jīvanmuktaḥ — також у 18.26; споріднений іменник jīvanmukti («звільнення за життя») трапляється в 18.80 і 20.4 (за SARIT та GRETIL). Друга половина — найважча: yathājīvanam може читатися як «протягом усього життя» (прислівник) або як «як живеться» (порівняння). SARIT і GRETIL членують по-різному, і обидва прочитання дають сенс. rañjanā («вподоба, забавка») точніше за «бажання»: це те, що фарбує — те саме слово, від якого rāga («пристрасть») і «фарба».}
\vfig{ilustratsii/18-13-space.pdf}{Нічого не слід зробити: все йде як живеться}{Схема наша, не історична ілюстрація.}

\vsec{18.14}
\vdeva{क्व मोहः क्व च वा विश्वं क्व तद् ध्यानं क्व मुक्तता । सर्वसङ्कल्पसीमायां विश्रान्तस्य महात्मनः ॥}
\vtranslit{ква мохах ква ча ва вішвам ква тад дгянам ква муктата | сарвасанкалпасімаям вішрантася махатманах ||}
\vperek{Де омана, де світ, де те споглядання, де сам стан свободи — у великого духом, що спочив у межі всіх задумів?}
\vcom{Чотири питання, і останнє найсміливіше: питається й про muktatā — стан свободи. saṅkalpa («задум») тут ключове: межа всіх задумів — це не межа думок, а межа намірів, тобто те, де закінчується будь-яке «хай буде так». Саме там і спочиває mahātman, і тому навіть свобода не має місця як мета.}
\vfig{ilustratsii/18-14-neti.pdf}{Омана, світ, споглядання й свобода відпали}{Схема наша, не історична ілюстрація.}

\vsec{18.15}
\vdeva{येन विश्वम् इदं दृष्टं स नास्तीति करोतु वै । निर्वासनः किं कुरुते पश्यन्न् अपि न पश्यति ॥}
\vtranslit{єна вішвам ідам дріштам са настіті кароту вай | нірвасанах кім куруте пашянн апі на пашяті ||}
\vperek{Нехай той, ким побачений цей світ, каже, що його немає, — а що робить той, у кого не лишилося слідів? Він і бачить, і не бачить.}
\vcom{nāstīti karotu («нехай робить „нема“») — формула заперечення, подана як дія, яку можна виконувати. Протиставлений їй nirvāsanaḥ: не той, хто заперечує, а той, у кого не лишилося vāsanā — слідів, навичок, залишків минулого досвіду. Остання половина — найточніша: paśyann api na paśyati («хоч бачить — не бачить») не парадокс заради ефекту, а опис: бачення є, а впізнавання як «свого» немає.}
\vpar{Українська паралель (ПОТРЕБУЄ ПЕРЕВІРКИ)}{українська містична проза, не визначено.}{ПОТРЕБУЄ ПЕРЕВІРКИ: стан, у якому людина і бачить, і не бачить, найближчий до української містичної прози, але конкретного твору ми не звірили.}
\vfig{ilustratsii/18-15-illusion.pdf}{Світ, якого немає, і той, хто не має слідів}{Схема наша, не історична ілюстрація.}

\vsec{18.16}
\vdeva{येन दृष्टं परं ब्रह्म सो ऽहं ब्रह्मेति चिन्तयेत् । किं चिन्तयति निश्चिन्तो द्वितीयं यो न पश्यति ॥}
\vtranslit{єна дріштам парам брахма со ’хам брахметі чінтаєт | кім чінтаяті нішчінто двітіям йо на пашяті ||}
\vperek{Нехай той, ким побачено вище Брахмо, повторює: «я — Брахман». А що думає той, хто вільний від думки й не бачить нічого другого?}
\vcom{Формула so 'haṃ brahma («я — Брахман») подана як cintayet — те, що можна МІРКУВАТИ, тобто тримати в думці. Протиставлений їй niścintaḥ («без думки») — і питання kim cintayati показує, що для нього міркування вже не робота. dvitīya («друге») — точний термін: не «інше» взагалі, а те, що стоїть поряд із першим. Хто не бачить другого, тому нема з чим ототожнюватися.}
\vfig{ilustratsii/18-16-neti.pdf}{Я — Брахман сказано; вільного від думки не питають}{Схема наша, не історична ілюстрація.}

\vsec{18.17}
\vdeva{दृष्टो येनात्मविक्षेपो निरोधं कुरुते त्व् असौ । उदारस् तु न विक्षिप्तः साध्याभावात् करोति किम् ॥}
\vtranslit{дрішто єнатмавікшепо ніродгам куруте тв асау | ударас ту на вікшіптах садгябгават кароті кім ||}
\vperek{Хто побачив у собі розсіяння, той береться його спиняти. Але великодухий не розсіяний — і, коли немає чого досягати, що йому робити?}
\vcom{Вірш ділить людей на два типи за однією ознакою — чи бачать вони в собі розсіяння. Для першого nirodha («припинення») є справою; для другого sādhya («досяжне») просто немає, бо немає від чого відштовхуватися. Слово udāra («великодухий») ужите тут не як похвала, а як точна ознака: широкий — отже, не звужений до одного напрямку, а тому й не може бути «розсіяним» у поганому сенсі.}
\vfig{ilustratsii/18-17-knot.pdf}{Вузол розсіяння, якого немає в великодухого}{Схема наша, не історична ілюстрація.}

\vsec{18.18}
\vdeva{धीरो लोकविपर्यस्तो वर्तमानो ऽपि लोकवत् । न समाधिं न विक्षेपं न लेपं स्वस्य पश्यति ॥}
\vtranslit{дгіро локавіпар'ясто вартамано ’пі локават | на самадгім на вікшепам на лепам свася пашяті ||}
\vperek{Мудрий, збочений разом зі світом, поводиться як увесь світ — і не бачить у себе ні самадгі, ні розсіяння, ні налипання.}
\vcom{Найсміливіше місце блоку: dhīra описаний як lokaviparyasta — «збочений разом зі світом». Це не докір, а точність: він поводиться як усі, і ззовні його не відрізнити. Три заперечення (samādhi, vikṣepa, lepa) забирають і високу практику, і її відсутність, і забрудненість — тобто всю шкалу, якою зазвичай міряють. Порівняй 14.4, де така людина постає «як божевільна»: тут те саме сказано спокійніше й точніше.}
\vfig{ilustratsii/18-18-motion.pdf}{Мудрий, збочений разом зі світом}{Схема наша, не історична ілюстрація.}

\vsec{18.19}
\vdeva{भावाभावविहीनो यस् तृप्तो निर्वासनो बुधः । नैव किंचित् कृतं तेन लोकदृष्ट्या विकुर्वता ॥}
\vtranslit{бгавабгававіхіно яс тріпто нірвасано будгах | найва кімчіт крітам тена локадріштя вікурвата ||}
\vperek{Тим, хто позбавлений буття й небуття, хто вдоволений, без слідів і тверезий, — не зроблено нічого, хоч з погляду світу він діє.}
\vcom{Вірш розводить дві оцінки однієї дії: lokadṛṣṭyā («з погляду світу») дія є, а kṛtam («зроблене») не накопичується. Це та сама думка, що в 13.3, але тут вона подана як висновок, а не як питання. tṛpta («вдоволений») стоїть у переліку ознак не випадково: незадоволеність — те, що змушує зараховувати дії собі.}
\vfig{ilustratsii/18-19-neti.pdf}{Позбавлений буття й небуття — і все ж діє}{Схема наша, не історична ілюстрація.}

\vsec{18.20}
\vdeva{प्रवृत्तौ वा निवृत्तौ वा नैव धीरस्य दुर्ग्रहः । यदा यत् कर्तुम् आयाति तत् कृत्वा तिष्ठतः सुखम् ॥}
\vtranslit{правріттау ва нівріттау ва найва дгірася дурґрахах | яда ят картум аяті тат крітва тіштхатах сукхам ||}
\vperek{Ні в діянні, ні у відступі немає в мудрого зацикленості: що б не випало зробити, зробивши те, він перебуває зручно.}
\vcom{pravṛtti / nivṛtti («виступання» / «відступ») — класична пара індійської думки: діяти в світі чи відійти від нього. Вірш каже, що для мудрого вона не має ваги, і причина названа точно: durgraha («вперте хапання») — не вибір, а зацикленість на виборі. Друга половина повторює 13.3 майже дослівно (yadā yat kartum āyāti), але замість «сиджу, як завгодно» тут «перебуває зручно» — та сама формула, інший відтінок.}
\vpar{Українська паралель (наша інтерпретація)}{Іван Франко, «Мойсей» (1905).}{Наша паралель: у Франка відсутність зацикленості теж постає як ознака того, хто дійшов. Наше порівняння — типологічне.}
\vfig{ilustratsii/18-20-balance.pdf}{Дві шальки: діяння й відступ}{Схема наша, не історична ілюстрація.}

\vsec{18.21}
\vdeva{निर्वासनो निरालम्बः स्वच्छन्दो मुक्तबन्धनः । क्षिप्तः संस्कारवातेन चेष्टते शुष्कपर्णवत् ॥}
\vtranslit{нірвасано ніраламбах сваччандо муктабандганах | кшіптах самскараватена чештате шушкапарнават ||}
\vperek{Без слідів, без опори, за власною волею, з відпущеними путами — гнаний вітром самскар, він ворушиться, як сухий лист.}
\vcom{Образ сухого листа точний у двох відношеннях: він рухається, хоч сам не має сили руху, і він ні до чого не прикріплений. Тобто дія є, а діяча немає. saṃskāravāta («вітер самскар») — те саме, що cittavāta в главах 2 і 7, але тут вітер жене вже не човен і не хвилю, а лист: різниця мас. Чотири ознаки на початку (nirvāsana, nirālamba, svacchanda, muktabandhana) ідуть без дієслова — це не опис шляху, а перелік того, чого вже немає.}
\vfig{ilustratsii/18-21-thread.pdf}{Сухий лист на вітрі самскар}{Схема наша, не історична ілюстрація.}

\vsec{18.22}
\vdeva{असंसारस्य तु क्वापि न हर्षो न विषादिता । स शीतलमना नित्यं विदेह इव राजते ॥}
\vtranslit{асамсарася ту квапі на харшо на вішадіта | са шіталамана нітям відеха іва раджате ||}
\vperek{У того, хто поза сансарою, ніде немає ні захвату, ні смутку: прохолодний розумом, він завжди сяє, наче розвтілеснений.}
\vcom{śītamana («прохолодний розумом») — не «байдужий»: прохолода тут властивість, а не відсутність тепла. І пара harṣa / viṣāda («захват» / «смуток») та сама, що в 4.2 і 3.9, але там ішлося про те, чого мудрий не відчуває, а тут — про те, що його стан видно ззовні. videha («без тіла») з iva («наче») — не твердження, що тіла немає, а опис того, як він виглядає: тіло є, а прив'язки до нього не видно.}
\vfig{ilustratsii/18-22-flame.pdf}{Прохолодне полумʼя без тіла}{Схема наша, не історична ілюстрація.}

\vsec{18.23}
\vdeva{कुत्रापि न जिहासास्ति नाशो वापि न कुत्रचित् । आत्मारामस्य धीरस्य शीतलाच्छतरात्मनः ॥}
\vtranslit{кутрапі на джіхасасті нашо вапі на кутрачіт | атмарамася дгірася шіталаччатаратманах ||}
\vperek{Ніде немає бажання відкинути, і ніде немає загибелі — у мудрого, чия втіха — Атман, чий Атман прохолодний і прозорий.}
\vcom{jihāsā («бажання відкинути») — дезидератив, тобто «хотіння покинути»; разом із nāśa («загибель») він закриває обидва способи втрати: відмовитися самому або бути знищеним. ātmārāma («той, чия втіха — Атман») — слово, відоме з «Бгаґавад-ґіти»; тут воно вжите без пояснень, як звичний зворот. Форма śītalācchatarātmanaḥ важка й, найімовірніше, зіпсована; я розібрав її як складене «прохолодний-прозорий-Атман», але це тлумачення, а не встановлене членування.}
\vfig{ilustratsii/18-23-neti.pdf}{Немає що відкидати й немає загибелі}{Схема наша, не історична ілюстрація.}

\vsec{18.24}
\vdeva{प्रकृत्या शून्यचित्तस्य कुर्वतो ऽस्य यदृच्छया । प्राकृतस्येव धीरस्य न मानो नावमानता ॥}
\vtranslit{пракрітя шунячіттася курвато ’ся ядріччая | пракрітасєва дгірася на мано наваманата ||}
\vperek{У того, чий розум за природою порожній і хто діє, як випаде, — як у звичайної людини: немає ні гордості, ні приниженості.}
\vcom{prākṛtasya iva («як у звичайного») — найважливіше: мудрий не вирізняється, і саме тому в нього немає ні māna, ні avamānatā. Ці два слова — не просто «гордість» і «скромність», а стан у чиїхось очах: māna — коли тебе шанують, avamānatā — коли тебе принижують. Обидва потребують глядача, і вірш каже, що глядача немає. Порівняй 18.18 (lokaviparyasta): та сама думка, інший бік.}
\vfig{ilustratsii/18-24-balance.pdf}{Дві шальки: гордість і приниженість}{Схема наша, не історична ілюстрація.}

\vsec{18.25}
\vdeva{कृतं देहेन कर्मेदं न मया शुद्धरूपिणा । इति चिन्तानुरोधी यः कुर्वन्न् अपि करोति न ॥}
\vtranslit{крітам дехена кармедам на мая шуддгарупіна | іті чінтануродгі ях курванн апі кароті на ||}
\vperek{«Це діяння зробило тіло, не я, чистої природи» — хто тримається цієї думки, той і роблячи не робить.}
\vcom{Формула «зробило тіло, не я» — класичне виправдання, і вірш ставиться до неї стримано: cintānurodhī («той, хто йде за думкою») означає, що сама думка ще не є свободою, а лише тим, за чим ідуть. Висновок kurvann api karoti na («хоч робить — не робить») стосується не того, хто проголошує формулу, а того, хто нею живе. Різниця між 18.19 (там «не зроблено нічого» сказано про мудрого) і тут — у тому, що тут ідеться про думку, а не про стан.}
\vpar{Українська паралель (наша інтерпретація)}{Григорій Сковорода, «Наркіс. Розмова про те: Пізнай себе».}{Наша паралель: у Сковороди діяння теж приписується тілу, а не тому, хто в тілі бачить. Наше порівняння — типологічне.}
\vfig{ilustratsii/18-25-neti.pdf}{Діяння тіла, а не мене}{Схема наша, не історична ілюстрація.}

\vsec{18.26}
\vdeva{अतद्वादीव कुरुते न भवेद् अपि बालिशः । जीवन्मुक्तः सुखी श्रीमान् संसरन्न् अपि शोभते ॥}
\vtranslit{атадвадіва куруте на бгавед апі балішах | джіванмуктах сукхі шріман самсаранн апі шобгате ||}
\vperek{Навіть якщо він каже не те, він не стає від цього дурнем: звільнений за життя, щасливий, сяйний — він гарний навіть тоді, коли блукає сансарою.}
\vcom{Вірш захищає того, хто говорить неправильно: atadvādī («той, хто каже „не те“») не стає bāliśa («дурнем»), бо слова не міняють стану. Друга половина повертає jīvanmukta з 18.13 і додає saṃsaran api («навіть блукаючи»): свобода не скасовує блукання, а робить його гарним. śobhate («сяє, гарно виглядає») — слово про вигляд, не про суть, і це навмисно: йдеться про те, як це видно.}
\vfig{ilustratsii/18-26-flame.pdf}{Сяйво звільненого за життя}{Схема наша, не історична ілюстрація.}

\vsec{18.27}
\vdeva{नानाविचारसुश्रान्तो धीरो विश्रान्तिम् आगतः । न कल्पते न जानाति न शृणोति न पश्यति ॥}
\vtranslit{нанавічарасушранто дгіро вішрантім аґатах | на калпате на джанаті на шріноті на пашяті ||}
\vperek{Мудрий, стомлений різноманітними роздумами, дійшов до спочинку: він не утворюється, не знає, не чує, не бачить.}
\vcom{Чотири заперечення наприкінці — це не опис нечутливості, а скасування чотирьох способів бути при справах: kalpate («вважається, утворюється»), jānāti, śṛṇoti, paśyati. Найцікавіше перше: na kalpate означає, що він більше не «складається» в якусь роль. Стомленість названа suśrānta («вельми втомлений») і стоїть як причина, а не як вада: спочинок приходить саме тому, що роздумів було надто багато.}
\vfig{ilustratsii/18-27-neti.pdf}{Роздуми вичерпані: спочинок}{Схема наша, не історична ілюстрація.}

\vsec{18.28}
\vdeva{असमाधेर् अविक्षेपान् न मुमुक्षुर् न चेतरः । निश्चित्य कल्पितं पश्यन् ब्रह्मैवास्ते महाशयः ॥}
\vtranslit{асамадгер авікшепан на мумукшур на четарах | нішчітя калпітам пашян брахмайвасте махашаях ||}
\vperek{Ні через відсутність самадгі, ні через відсутність розсіяння він не є ні тим, хто прагне визволення, ні кимось іншим: вирішивши так і бачачи все вигаданим, великодухий перебуває як саме Брахмо.}
\vcom{Подвійний ablativus (asamādheḥ, avikṣepāt) закриває обидва боки: він не mumukṣu ні тому, що не має зосередження, ні тому, що не має розсіяння. Тобто ознака прагнення не залежить від стану розуму взагалі. brahma eva āste («перебуває як саме Брахмо») — найсильніше твердження блоку: не «усвідомлює себе Брахмом», а перебуває ним.}
\vfig{ilustratsii/18-28-witness.pdf}{Поза самадгі й розсіянням — як саме Брахмо}{Схема наша, не історична ілюстрація.}

\vsec{18.29}
\vdeva{यस्यान्तः स्याद् अहङ्कारो न करोति करोति सः । निरहङ्कारधीरेण न किंचिद् अकृतं कृतम् ॥}
\vtranslit{ясянтах сяд аханкаро на кароті кароті сах | ніраханкарадгірена на кімчід акрітам крітам ||}
\vperek{У кого всередині є «я-роблення», той не робить, хоч і робить. А в мудрого без «я-роблення» не лишається нічого незробленого.}
\vcom{ahaṅkāra — не «его» в психологічному сенсі, а буквально «роблення „я“»: акт, яким дія зараховується собі. Тому перша половина й парадоксальна: хто робить цей акт, той насправді не діє. Друга половина — дзеркальна: без цього акту немає незробленого, бо нема кому вести рахунок. Це найточніше формулювання тієї самої думки, що в 13.3 і 18.19.}
\vfig{ilustratsii/18-29-knot.pdf}{Вузол я-роблення}{Схема наша, не історична ілюстрація.}

\vsec{18.30}
\vdeva{नोद्विग्नं न च सन्तुष्टम् अकर्तृ स्पन्दवर्जितम् । निराशं गतसन्देहं चित्तं मुक्तस्य राजते ॥}
\vtranslit{нодвіґнам на ча сантуштам акартри спандаварджітам | нірашам ґатасандехам чіттам муктася раджате ||}
\vperek{Не схвильований і не вдоволений, без діяча, без пориву, без сподівань і без сумніву — розум вільного сяє.}
\vcom{Шість ознак, і всі — заперечні або позбавні, крім останнього дієслова rājate («сяє»). Це навмисна будова: описується не те, який розум є, а те, чого в ньому немає, і лише сяйво ствердне. akartṛ («без діяча») стоїть у середині переліку і є його віссю: прибравши діяча, решта ознак випливає сама. spanda («тремтіння») — те саме слово, що в назві школи кашмірського шіваїзму, але тут воно означає просто порив.}
\vfig{ilustratsii/18-30-flame.pdf}{Розум вільного сяє}{Схема наша, не історична ілюстрація.}

\vsec{18.31}
\vdeva{निर्ध्यातुं चेष्टितुं वापि यच् चित्तं न प्रवर्तते । निर्निमित्तम् इदं किन्तु निर्ध्यायति विचेष्टते ॥}
\vtranslit{нірдгятум чештітум вапі яч чіттам на правартате | нірніміттам ідам кінту нірдгяяті вічештате ||}
\vperek{Ум, що не береться ні замислитися, ні діяти, — але ж він, без причини, замислюється й діє.}
\vcom{Вірш описує парадокс, а не заклик: ум не має наміру ні думати, ні діяти — і все ж думає й діє. `nirnimitta` («без причини, без приводу») — ключове слово: дія відбувається, але без того, хто її задумав. Це та сама структура, що в 15.16, де немає ні того, хто в сансарі, ні того, хто поза нею. `kintu` («але ж») тримає здивування, і його не можна викидати.}
\vpar{Українська паралель (ПОТРЕБУЄ ПЕРЕВІРКИ)}{українська філософська лірика, не визначено.}{ПОТРЕБУЄ ПЕРЕВІРКИ: ум, що замислюється без причини, — думка, близька до української філософської лірики, але конкретного твору ми не перевірили.}
\vfig{ilustratsii/18-31-motion.pdf}{Ум, що замислюється без причини}{Схема наша, не історична ілюстрація.}

\vsec{18.32}
\vdeva{तत्त्वं यथार्थम् आकर्ण्य मन्दः प्राप्नोति मूढताम् । अथ वायाति सङ्कोचम् अमूढः को ऽपि मूढवत् ॥}
\vtranslit{таттвам ятхартхам акарня мандах прапноті мудгатам | атха ваяті санкочам амудгах ко ’пі мудгават ||}
\vperek{Почувши правильно істину, тупий стає ще затьмаренішим — або ж звужується; а буває й незатьмарений, що поводиться як затьмарений.}
\vcom{Вірш чесно називає, що буває з людиною від почутого: або гірше, або звуження. saṅkoca («стиснення») — те саме слово, що в 18.5, де воно було хибним способом досягнення; тут воно наслідок неправильно прийнятої істини. Остання половина додає третій випадок: amūḍha, який поводиться як mūḍha — і це вже не про розум, а про вигляд.}
\vfig{ilustratsii/18-32-illusion.pdf}{Той, хто чув істину: тупий і незатьмарений}{Схема наша, не історична ілюстрація.}

\vsec{18.33}
\vdeva{एकाग्रता निरोधो वा मूढैर् अभ्यस्यते भृशम् । धीराः कृत्यं न पश्यन्ति सुप्तवत् स्वपदे स्थिताः ॥}
\vtranslit{екаґрата ніродго ва мудгайр абгясяте бгрішам | дгірах крітям на пашянті суптават свападе стхітах ||}
\vperek{Зосередженість або припинення старанно вправляють дурні. Мудрі не бачать, що їм робити, — вони стоять на власному місці, як сплячі.}
\vcom{Дві вправи названі поіменно (ekāgratā, nirodha) і віддані дурням — це рідкісна в книзі пряма оцінка практики. suptavat («як сплячі») не принизливе: сплячий не робить нічого й не потребує робити. svapada («власне місце») — те саме слово, що в 18.5 (ātmanaḥ padam).}
\vfig{ilustratsii/18-33-neti.pdf}{Зосередженість і припинення відпадають}{Схема наша, не історична ілюстрація.}

\vsec{18.34}
\vdeva{अप्रयत्नात् प्रयत्नाद् वा मूढो नाप्नोति निर्वृतिम् । तत्त्वनिश्चयमात्रेण प्राज्ञो भवति निर्वृतः ॥}
\vtranslit{апраятнат праятнад ва мудго напноті нірврітім | таттванішчаяматрена праджньо бгаваті нірврітах ||}
\vperek{Ні без зусилля, ні із зусиллям дурний не досягає заспокоєння: розумний заспокоюється самим лише рішенням про істину.}
\vcom{Подвійний ablativus (aprayatnāt, prayatnāt) закриває обидва боки так само, як у 18.28, — і тут це відповідь на 18.33, де дурні старанно вправляються. tattvaniścayamātra («саме лише рішення про істину») — не «знання», а niścaya, тобто остаточне встановлення, після якого питання знімається.}
\vfig{ilustratsii/18-34-balance.pdf}{Дві шальки: зусилля й відсутність зусилля}{Схема наша, не історична ілюстрація.}

\vsec{18.35}
\vdeva{शुद्धं बुद्धं प्रियं पूर्णं निष्प्रपञ्चं निरामयम् । आत्मानं तं न जानन्ति तत्राभ्यासपरा जनाः ॥}
\vtranslit{шуддгам буддгам пріям пурнам нішпрапаньчам нірамаям | атманам там на джананті татрабгясапара джанах ||}
\vperek{Атман — чистий, пізнаний, любий, повний, без множинності, без недуги — такого не знають ті, для кого головне — вправа.}
\vcom{Шість ознак із шести — ствердні, і це єдиний такий перелік у главі: попередні будувалися на запереченнях. abhyāsaparāḥ («ті, чиє головне — вправа») — точна адреса докору: не «ліниві», а саме ті, хто зробив практику головним. priya («любий») у переліку ознак Атмана несподіване й навмисне: не «блаженний» як властивість, а той, кого люблять.}
\vfig{ilustratsii/18-35-neti.pdf}{Атман поза вправою й множинністю}{Схема наша, не історична ілюстрація.}

\vsec{18.36}
\vdeva{नाप्नोति कर्मणा मोक्षं विमूढो ऽभ्यासरूपिणा । धन्यो विज्ञानमात्रेण मुक्तस् तिष्ठत्य् अविक्रियः ॥}
\vtranslit{напноті кармана мокшам вімудго ’бгясарупіна | дганйо віджнянаматрена муктас тіштхатй авікріях ||}
\vperek{Зовсім затьмарений не досягає визволення ні діянням, ні тим, що має вигляд вправи; а блаженний, вільний самим лише розпізнанням, перебуває незмінним.}
\vcom{abhyāsarūpa («те, що має вигляд вправи») — формула точна: йдеться не про вправу, а про її видимість. Протиставлення vijñānamātra («саме лише розпізнання») ставить розпізнання на місце дії, і avikriya («незмінний») — наслідок: те, що не змінюється, не потребує ані дії, ані вправи.}
\vfig{ilustratsii/18-36-balance.pdf}{Дві шальки: діяння й вправа}{Схема наша, не історична ілюстрація.}

\vsec{18.37}
\vdeva{मूढो नाप्नोति तद् ब्रह्म यतो भवितुम् इच्छति । अनिच्छन्न् अपि धीरो हि परब्रह्मस्वरूपभाक् ॥}
\vtranslit{мудго напноті тад брахма ято бгавітум іччаті | аніччанн апі дгіро хі парабрахмасварупабгак ||}
\vperek{Дурний не досягає того Брахма, бо хоче ним стати. А мудрий, і не бажаючи, має власну природу вищого Брахма.}
\vcom{Причина названа точно: icchati («хоче стати») — саме бажання стати й не дає досягти, бо досягнення передбачає, що зараз не так. anicchan api («і не бажаючи») — не «без бажання» як заслуга, а просто відсутність цієї операції. bhāk («той, хто має частку») замість «той, хто є» — обережніше слово, ніж здається.}
\vfig{ilustratsii/18-37-knot.pdf}{Вузол бажання стати Брахмою}{Схема наша, не історична ілюстрація.}

\vsec{18.38}
\vdeva{निराधारा ग्रहव्यग्रा मूढाः संसारपोषकाः । एतस्यानर्थमूलस्य मूलच्छेदः कृतो बुधैः ॥}
\vtranslit{нірадгара ґрахав'яґра мудгах самсарапошаках | етасянартхамулася мулаччедах кріто будгайх ||}
\vperek{Без опори, заклопотані хапанням, дурні годують сансару. А розумні зробили відтин кореня цього кореня лиха.}
\vcom{saṃsārapoṣaka («той, що годує сансару») — сильне слово: сансара не стан, у якому перебувають, а те, що хтось годує. mūlaccheda («відтин кореня») — не «видалення», а саме підтин під основу, після якого нічого не відростає. graha тут двозначне: і «хапання», і «планета» — тобто заклопотаність тим, що хапає.}
\vfig{ilustratsii/18-38-thread.pdf}{Відтин кореня лиха}{Схема наша, не історична ілюстрація.}

\vsec{18.39}
\vdeva{न शान्तिं लभते मूढो यतः शमितुम् इच्छति । धीरस् तत्त्वं विनिश्चित्य सर्वदा शान्तमानसः ॥}
\vtranslit{на шантім лабгате мудго ятах шамітум іччаті | дгірас таттвам вінішчітя сарвада шантаманасах ||}
\vperek{Дурний не отримує спокою, бо хоче заспокоїтися. А мудрий, остаточно встановивши істину, завжди має спокійний розум.}
\vcom{Той самий хід, що в 18.37: бажання стану не дає стану. Різниця в тому, що тут ідеться про śānti, і паралель точна — icchati («хоче») у першій половині проти viniścitya («встановивши») у другій: бажання й встановлення стоять на одному місці й не можуть бути разом.}
\vfig{ilustratsii/18-39-knot.pdf}{Вузол бажання спокою}{Схема наша, не історична ілюстрація.}

\vsec{18.40}
\vdeva{क्वात्मनो दर्शनं तस्य यद् दृष्टम् अवलम्बते । धीरास् तं तं न पश्यन्ति पश्यन्त्य् आत्मानम् अव्ययम् ॥}
\vtranslit{кватмано даршанам тася яд дріштам аваламбате | дгірас там там на пашянті пашянтй атманам ав'яям ||}
\vperek{Де в того бачення Атмана, хто тримається побаченого? Мудрі такого не бачать — вони бачать Атман, незнищенний.}
\vcom{Вірш грає двома смислами paśyati: «бачити [людину]» і «бачити [Атман]». dhīrāḥ tam na paśyanti — вони не бачать того, хто спирається на побачене, тобто не визнають його за того, ким він себе вважає; і в тому ж руслі paśyanti ātmānam. avyaya («незнищенний») — точніше «невитратний»: те, що не зменшується від використання.}
\vfig{ilustratsii/18-40-witness.pdf}{Побачене й Атман, який бачать мудрі}{Схема наша, не історична ілюстрація.}

\vsec{18.41}
\vdeva{क्व निरोधो विमूढस्य यो निर्बन्धं करोति वै । स्वारामस्यैव धीरस्य सर्वदासाव् अकृत्रिमः ॥}
\vtranslit{ква ніродго вімудгася йо нірбандгам кароті вай | сварамасяйва дгірася сарвадасав акрітрімах ||}
\vperek{Де припинення в того, хто зовсім затьмарився і наполягає? У мудрого, чия втіха — в собі, воно завжди невдаване.}
\vcom{Протиставлення точне: nirodha можна робити (karoti), а akṛtrima — «невдаване» — робити не можна за визначенням. Тому питання kva («де?») стосується не місця, а самої можливості. svārāma («чия втіха — в собі») — те саме слово, що в 18.23; тут воно протиставлене nirbandha («наполегливості»).}
\vfig{ilustratsii/18-41-space.pdf}{Припинення, яке не треба вдавати}{Схема наша, не історична ілюстрація.}

\vsec{18.42}
\vdeva{भावस्य भावकः कश्चिन् न किंचिद् भावको ऽपरः । उभयाभावकः कश्चिद् एवम् एव निराकुलः ॥}
\vtranslit{бгавася бгаваках кашчін на кімчід бгавако ’парах | убгаябгаваках кашчід евам ева ніракулах ||}
\vperek{Один наділяє буттям, інший не наділяє нічим, дехто не наділяє ні тим, ні тим. Саме так він безтурботний.}
\vcom{Рідкісний у книзі перелік із трьох типів, поданий без оцінки: перший уявляє буття, другий — ніщо (заперечення), третій не робить ні того, ні того. Безтурботність (nirākula) приписана саме третьому, і це важливо: не заперечення рятує, а вихід із пари «уявляти / не уявляти». Форма bhāvaka (каузативна від bhū) означає активне наділяння буттям, а не пасивне сприйняття.}
\vfig{ilustratsii/18-42-motion.pdf}{Три способи наділяти буттям}{Схема наша, не історична ілюстрація.}

\vsec{18.43}
\vdeva{शुद्धम् अद्वयम् आत्मानं भावयन्ति कुबुद्धयः । न तु जानन्ति संमोहाद् यावज्जीवम् अनिर्वृताः ॥}
\vtranslit{шуддгам адваям атманам бгаваянті кубуддгаях | на ту джананті саммохад яваджджівам анірврітах ||}
\vperek{Дурні уявляють Атман чистим і недвоїстим, але через затьмарення не знають його — і до кінця життя лишаються незаспокоєними.}
\vcom{Розведено bhāvayanti («уявляють, плекають у думці») і jānanti («знають»): перше доступне й дурням, друге — ні. Тобто правильна формула ще не є знанням. yāvajjīvam («до кінця життя») ставить ціну помилки: не тимчасове непорозуміння, а ціле життя.}
\vfig{ilustratsii/18-43-illusion.pdf}{Уявлений Атман і незнаний Атман}{Схема наша, не історична ілюстрація.}

\vsec{18.44}
\vdeva{मुमुक्षोर् बुद्धिर् आलम्बम् अन्तरेण न विद्यते । निरालम्बैव निष्कामा बुद्धिर् मुक्तस्य सर्वदा ॥}
\vtranslit{мумукшор буддгір аламбам антарена на відяте | ніраламбайва нішкама буддгір муктася сарвада ||}
\vperek{У того, хто прагне визволення, розум не буває без опори. А у вільного розум без опори й без бажань — завжди.}
\vcom{Пара mumukṣu / mukta («той, хто прагне» / «вільний») і єдине слово ālamba («опора») розділяє їх: прагнення за визначенням потребує опори, бо мусить на щось спертися. Друга половина додає niṣkāmā («без бажань») — тобто відсутність опори й відсутність бажання тут не два факти, а один.}
\vfig{ilustratsii/18-44-balance.pdf}{Дві шальки: опора прагнення й порожнеча волі}{Схема наша, не історична ілюстрація.}

\vsec{18.45}
\vdeva{विषयद्वीपिनो वीक्ष्य चकिताः शरणार्थिनः । विशन्ति झटिति क्रोडं निरोधैकाग्र्यसिद्धये ॥}
\vtranslit{вішаядвіпіно вікшя чакітах шаранартхінах | вішанті джгатіті кродам ніродгайкаґр'ясіддгає ||}
\vperek{Побачивши пантер предметів, налякані й шукаючи притулку, вони стрімко кидаються в лоно — задля успіху в припиненні й зосередженості.}
\vcom{Тут і далі (18.46) тваринна метафора, і вона несподівано точна: у практику йдуть не з рішучості, а з переляку. `dvīpin` — саме пантера, не слон; слони з'являться в наступному вірші. Слово `kroḍa` («лоно») вжите без пояснення, що це за лоно, — і це навмисно: читач має сам упізнати притулок.}
\vfig{ilustratsii/18-45-motion.pdf}{Пантери предметів і втеча в припинення}{Схема наша, не історична ілюстрація.}

\vsec{18.46}
\vdeva{निर्वासनं हरिं दृष्ट्वा तूष्णीं विषयदन्तिनः । पलायन्ते न शक्तास् ते सेवन्ते कृतचाटवः ॥}
\vtranslit{нірвасанам харім дріштва тушнім вішаядантінах | палаянте на шактас те севанте крітачатавах ||}
\vperek{Побачивши лева, вільного від слідів, слони предметів мовчки тікають. Нездатні, вони більше не прислуговують, хоч раніше й підлещувалися.}
\vcom{Продовження 18.45, і метафора перевернута: там предмети — пантери, що лякають людину, тут вони слони, що тікають від лева. Те, що раніше гнало, тепер тікає. `kṛtacāṭava` («ті, що зробили підлабузництво») — незвичайно точне слово про предмети: вони не просто спокушали, а прислуговували з лестощами. `tūṣṇīm` («мовчки») підкреслює: втеча без опору.}
\vfig{ilustratsii/18-46-motion.pdf}{Слони предметів тікають від лева без слідів}{Схема наша, не історична ілюстрація.}

\vsec{18.47}
\vdeva{न मुक्तिकारिकां धत्ते निःशङ्को युक्तमानसः । पश्यञ् छृण्वन् स्पृशञ् जिघ्रन्न् अश्नन्न् आस्ते यथासुखम् ॥}
\vtranslit{на муктікарікам дгатте ніхшанко юктаманасах | пашянь ччрінван спрішань джіґгранн ашнанн асте ятхасукхам ||}
\vperek{Він не носить на собі справи визволення — без сумніву, зі з'єднаним розумом. Дивиться, слухає, торкається, нюхає, їсть — і перебуває, як завгодно.}
\vcom{П'ять дієприкметників перелічують чуттєві дії поспіль, без оцінки — і саме це тут твердження: чуття працюють, а справи визволення немає. `muktikārikā` («справа визволення») подана як ноша: її не тримають на собі. Формула `yathāsukham` повторює рефрен глави 13 — та сама поза, інший контекст.}
\vfig{ilustratsii/18-47-motion.pdf}{Звичайні дії без ноші визволення}{Схема наша, не історична ілюстрація.}

\vsec{18.48}
\vdeva{वस्तुश्रवणमात्रेण शुद्धबुद्धिर् निराकुलः । नैवाचारमनाचारम् औदास्यं वा प्रपश्यति ॥}
\vtranslit{вастушраванаматрена шуддгабуддгір ніракулах | найвачараманачарам аудасям ва прапашяті ||}
\vperek{Від самого лише слухання суті — чистий розумом, безтурботний — він не бачить ані різниці між належним і неналежним, ані байдужості.}
\vcom{Три речі названі й усі три знято: ācāra, anācāra і audāsya («байдужість»). Остання — найважливіша, бо це типова підміна: коли зникає різниця між належним і неналежним, на її місце стає байдужість, і вірш каже, що її теж немає. `vastuśravaṇamātra` («саме лише слухання суті») — єдина умова, названа тут.}
\vfig{ilustratsii/18-48-neti.pdf}{Різниця й байдужість відпали}{Схема наша, не історична ілюстрація.}

\vsec{18.49}
\vdeva{यदा यत् कर्तुम् आयाति तदा तत् कुरुते ऋजुः । शुभं वाप्य् अशुभं वापि तस्य चेष्टा हि बालवत् ॥}
\vtranslit{яда ят картум аяті тада тат куруте ріджух | шубгам вапй ашубгам вапі тася чешта хі балават ||}
\vperek{Коли що випадає зробити, тоді він те й робить — він прямий. Добре чи недобре — його діяння як у дитини.}
\vcom{Та сама формула, що в 13.3 і 18.20 (`yadā yat kartum āyāti`), і втретє в книзі — це вже сталий зворот. Нове тут — `bālavat` («як у дитини») і `śubha / aśubha`: дитина не розрізняє добра й недобра, і саме тому її діяння не зараховується. `ṛju` («прямий») пояснює, чому: він не вибирає, а робить те, що випало.}
\vpar{Українська паралель (наша інтерпретація)}{Василь Стефаник, «Камінний хрест» (1900).}{Наша паралель: у Стефаника простота дії теж не є недбальством, а способом бути прямим. Наше порівняння — типологічне.}
\vfig{ilustratsii/18-49-motion.pdf}{Діяння як у дитини: що випало, те й робить}{Схема наша, не історична ілюстрація.}

\vsec{18.50}
\vdeva{स्वातन्त्र्यात् सुखम् आप्नोति स्वातन्त्र्याल् लभते परम् । स्वातन्त्र्यान् निर्वृतिं गच्छेत् स्वातन्त्र्यात् परमं पदम् ॥}
\vtranslit{сватантр'ят сукхам апноті сватантр'ял лабгате парам | сватантр'ян нірврітім ґаччет сватантр'ят парамам падам ||}
\vperek{Від свободи отримує щастя, від свободи досягає вищого, від свободи приходить до заспокоєння, від свободи — до вищої стопи.}
\vcom{Чотириразова анафора `svātantryāt`, і це єдиний вірш книги, де одне слово відкриває всі чотири частини. `svātantrya` — не «свобода» взагалі, а буквально «правління собою»; тому воно й стоїть в ablativus: усе перелічене походить від нього, а не досягається заради нього. Порядок зростає: sukha → param → nirvṛti → paramaṃ pada, тобто від відчуття до останньої стопи.}
\vfig{ilustratsii/18-50-flame.pdf}{Полумʼя свободи: щастя й вища стопа}{Схема наша, не історична ілюстрація.}

\vsec{18.51}
\vdeva{अकर्तृत्वम् अभोक्तृत्वं स्वात्मनो मन्यते यदा । तदा क्षीणा भवन्त्य् एव समस्ताश् चित्तवृत्तयः ॥}
\vtranslit{акартрітвам абгоктрітвам сватмано маняте яда | тада кшіна бгавантй ева самасташ чіттавріттаях ||}
\vperek{Коли людина вважає свій Атман не-діячем і не-споживачем — тоді всі рухи розуму справді вичерпуються.}
\vcom{Дві ознаки названі разом і в правильному порядку: не-діяч (akartṛ) і не-споживач (abhoktṛ) — тобто знімається і дія, і той, хто привласнює її плід. `cittavṛtti` — технічний термін йоги (звідси «читта-врітті-ніродха»); тут сказано не «припиняються зусиллям», а `kṣīṇāḥ bhavanti` — «стають вичерпаними», самі.}
\vfig{ilustratsii/18-51-neti.pdf}{Не діяч і не споживач: рухи розуму вичерпані}{Схема наша, не історична ілюстрація.}

\vsec{18.52}
\vdeva{उच्छृङ्खलाप्य् अकृतिका स्थितिर् धीरस्य राजते । न तु सस्पृहचित्तस्य शान्तिर् मूढस्य कृत्रिमा ॥}
\vtranslit{уччрінкхалапй акрітіка стхітір дгірася раджате | на ту саспріхачіттася шантір мудгася крітріма ||}
\vperek{Навіть розкутий і безформний стан мудрого сяє. А спокій затьмареного з пожадливим розумом — удаваний.}
\vcom{Протиставлення `akṛtikā` / `kṛtrimā` — те саме слово в двох формах: стан мудрого без вигляду (не підробляється), а спокій дурного — зроблений. `ucchṛṅkhalā` («розкута») і `api` («навіть») разом означають: свобода видна навіть тоді, коли нічого не показує.}
\vfig{ilustratsii/18-52-illusion.pdf}{Удаваний спокій і справжній}{Схема наша, не історична ілюстрація.}

\vsec{18.53}
\vdeva{विलसन्ति महाभोगैर् विशन्ति गिरिगह्वरान् । निरस्तकल्पना धीरा अबद्धा मुक्तबुद्धयः ॥}
\vtranslit{віласанті махабгоґайр вішанті ґіріґахваран | нірастакалпана дгіра абаддга муктабуддгаях ||}
\vperek{Вони граються великими розкошами й входять у гірські нетрі — мудрі, що відкинули побудови, незв'язані, з визволеним розумом.}
\vcom{Обидві крайнощі названі поспіль, і жодна не оцінена: розкоші й печери. Це відповідь на всю книгу уявлень про аскетизм — `abaddha` («незв'язаний») означає, що місце не має значення. `nirastakalpanā` («ті, що відкинули побудови») — причина, чому обидві крайнощі можливі.}
\vfig{ilustratsii/18-53-motion.pdf}{Розкоші й нетрі — однаково для мудрих}{Схема наша, не історична ілюстрація.}

\vsec{18.54}
\vdeva{श्रोत्रियं देवतां तीर्थम् अङ्गनां भूपतिं प्रियम् । दृष्ट्वा सम्पूज्य धीरस्य न कापि हृदि वासना ॥}
\vtranslit{шротріям деватам тіртхам анґанам бгупатім пріям | дріштва сампуджя дгірася на капі хріді васана ||}
\vperek{Побачивши знавця писань, божество, святе місце, жінку, царя, любу людину — і вшанувавши їх, — мудрий не має в серці жодного сліду.}
\vcom{Перелік навмисно строкатий: те, чому поклоняються, і те, до чого тягнуться. `sampūjya` («вшанувавши») важливе — він не уникає нічого з переліченого, але після вшанування `vāsanā` не лишається. Тобто дія виконана повністю, а сліду немає.}
\vfig{ilustratsii/18-54-neti.pdf}{Шанування без сліду в серці}{Схема наша, не історична ілюстрація.}

\vsec{18.55}
\vdeva{भृत्यैः पुत्रैः कलत्रैश् च दौहित्रैश् चापि गोत्रजैः । विहस्य धिक्कृतो योगी न याति विकृतिं मनाक् ॥}
\vtranslit{бгрітяйх путрайх калатрайш ча даухітрайш чапі ґотраджайх | віхася дгіккріто йоґі на яті вікрітім манак ||}
\vperek{Висміяний і зневажений слугами, синами, жінкою, онуками й родичами, йог не змінюється анітрохи.}
\vcom{Найближче коло назване повністю — і саме тому воно найважче: не вороги, а ті, з ким живеш. `manāk` («анітрохи») — слово, що закриває питання про міру: не «мало змінюється», а зовсім. `vikṛti` («видозміна») — те саме слово, що в 18.5 (nirvikāra).}
\vpar{Українська паралель (наша інтерпретація)}{Василь Стус, «Веселий цвинтар» (1985).}{Наша паралель: у Стуса зневага близьких теж не змінює того, хто стоїть. Наше порівняння — за будовою думки.}
\vfig{ilustratsii/18-55-balance.pdf}{Дві шальки: висміяний і зневажений}{Схема наша, не історична ілюстрація.}

\vsec{18.56}
\vdeva{सन्तुष्टो ऽपि न सन्तुष्टः खिन्नो ऽपि न च खिद्यते । तस्याश्चर्यदशां तां तां तादृशा एव जानते ॥}
\vtranslit{сантушто ’пі на сантуштах кхінно ’пі на ча кхідяте | тасяшчар'ядашам там там тадріша ева джанате ||}
\vperek{І вдоволений — не вдоволений; і засмучений — не засмучується. Його дивовижні стани — такі-то й такі — знають лише такі ж, як він.}
\vcom{Закінчення дослівно повторює 14.4 (`tāṃ tāṃ tādṛśā eva jānate`) — тобто свідчення ззовні неможливе, і це вже вдруге в книзі. Перша половина — не парадокс, а точність: стан є, назви немає, бо `santuṣṭa` і `khinna` описують те, що всередині вже не відбувається.}
\vfig{ilustratsii/18-56-illusion.pdf}{Стан, що його знають лише такі ж}{Схема наша, не історична ілюстрація.}

\vsec{18.57}
\vdeva{कर्तव्यतैव संसारो न तां पश्यन्ति सूरयः । शून्याकारा निराकारा निर्विकारा निरामयाः ॥}
\vtranslit{картав'ятайва самсаро на там пашянті сураях | шунякара ніракара нірвікара нірамаях ||}
\vperek{Сансара — це саме «треба зробити». Мудрі її не бачать: вони з порожнім виглядом, безформні, без видозмін, без недуги.}
\vcom{Найкоротше й найсильніше визначення сансари в книзі: не світ, не страждання, а `kartavyatā` — сам стан «треба». Тому й не бачать: сансара не річ, яку можна побачити, а спосіб, у якому речі постають. Чотири nir- означення наприкінці — та сама низка, що в 18.5.}
\vfig{ilustratsii/18-57-neti.pdf}{Треба зробити — і сансари немає}{Схема наша, не історична ілюстрація.}

\vsec{18.58}
\vdeva{अकुर्वन्न् अपि संक्षोभाद् व्यग्रः सर्वत्र मूढधीः । कुर्वन्न् अपि तु कृत्यानि कुशलो हि निराकुलः ॥}
\vtranslit{акурванн апі самкшобгад в'яґрах сарватра мудгадгіх | курванн апі ту крітяні кушало хі ніракулах ||}
\vperek{Навіть нічого не роблячи, дурний розумом скрізь збентежений через хвилювання. А вмілий, хоч і робить справи, — безтурботний.}
\vcom{Точне протиставлення за однією ознакою: справа не визначає стан. `saṃkṣobha` («збурення») назване причиною — не бездіяльність дає спокій, а відсутність збурення. `kuśala` («вмілий») тут не похвала майстерності, а ознака: той, хто вміє, не витрачається на саму дію.}
\vfig{ilustratsii/18-58-balance.pdf}{Дві шальки: неробство й діяння}{Схема наша, не історична ілюстрація.}

\vsec{18.59}
\vdeva{सुखम् आस्ते सुखं शेते सुखम् आयाति याति च । सुखं वक्ति सुखं भुङ्क्ते व्यवहारे ऽपि शान्तधीः ॥}
\vtranslit{сукхам асте сукхам шете сукхам аяті яті ча | сукхам вакті сукхам бгункте в'явахаре ’пі шантадгіх ||}
\vperek{Зручно сидить, зручно лежить, зручно приходить і йде, зручно говорить, зручно їсть — навіть у справах світу той, чий розум спокійний.}
\vcom{Шестиразова анафора `sukham`, і вона охоплює все тіло й усе спілкування: сидіти, лежати, приходити, йти, говорити, їсти. Це не «насолода», а `sukha` в значенні зручності, легкості. `vyavahāre 'pi` («навіть у справах») закриває сумнів: ідеться не про самотність, а про життя серед людей.}
\vpar{Українська паралель (наша інтерпретація)}{Олександр Довженко, «Щоденник» (1941–1956).}{Наша паралель: у Довженка зручність у звичайних справах теж постає як ознака спокійного розуму, а не як байдужість. Наше порівняння — за образом.}
\vfig{ilustratsii/18-59-motion.pdf}{Зручно сидить, зручно їсть: розум спокійний}{Схема наша, не історична ілюстрація.}

\vsec{18.60}
\vdeva{स्वभावाद् यस्य नैवार्तिर् लोकवद् व्यवहारिणः । महाह्रद इवाक्षोभ्यो गतक्लेशः स शोभते ॥}
\vtranslit{свабгавад яся найвартір локавад в'явахарінах | махахрада івакшобгйо ґатаклешах са шобгате ||}
\vperek{У кого за природою немає страждання, хто поводиться, як усі в світі, — той сяє, як велике озеро: незбурюваний, з минулими стражданнями.}
\vcom{`lokavat` («як увесь світ») — те саме слово, що `prākṛta` у 18.24 і `lokaviparyasta` у 18.18: мудрого не відрізнити ззовні. `mahāhrada` («велике озеро») точніше за «океан»: озеро незбурюване, бо глибоке, а не безмежне. `kleśa` — технічний термін (занечищення в йозі), тому «страждання» тут у строгому сенсі, не побутовому.}
\vfig{ilustratsii/18-60-space.pdf}{Велике озеро без вітру}{Схема наша, не історична ілюстрація.}

\vsec{18.61}
\vdeva{निवृत्तेर् अपि मूढस्य प्रवृत्तिर् उपजायते । प्रवृत्तिर् अपि धीरस्य निवृत्तिफलदायिनी ॥}
\vtranslit{нівріттер апі мудгася правріттір упаджаяте | правріттір апі дгірася нівріттіпхаладаїні ||}
\vperek{Навіть від припинення в затьмареного народжується діяльність. А діяльність мудрого дає плід припинення.}
\vcom{Найточніше формулювання того, що в книзі повторюється багато разів: справа не в дії й не в бездіяльності, а в тому, що з них виходить. У дурного навіть `nivṛtti` обертається `pravṛtti`; у мудрого `pravṛtti` дає плід `nivṛtti`. Обидві половини будуються на одній парі слів, переставленій місцями.}
\vfig{ilustratsii/18-61-flame.pdf}{Діяльність, що дає плід припинення}{Схема наша, не історична ілюстрація.}

\vsec{18.62}
\vdeva{परिग्रहेषु वैराग्यं प्रायो मूढस्य दृश्यते । देहे विगलिताशस्य क्व रागः क्व विरागता ॥}
\vtranslit{паріґрахешу вайраґям прайо мудгася дрішяте | дехе віґаліташася ква раґах ква віраґата ||}
\vperek{Відречення від набутків здебільшого видно в затьмареного. А в того, у кого в тілі вислизнула надія, — де пристрасть і де безпристрасність?}
\vcom{`vairāgya` тут не чеснота, а симптом: воно видно в дурного, бо дурний відрікається від того, чого не має. Друга половина закриває питання — `vigalitāśa` («той, у кого вислизнула надія») не має ні пристрасті, ні її протилежності, бо обидві потребують надії. Це те саме місце, яке місія просить перевіряти окремо (`vairāgya`).}
\vfig{ilustratsii/18-62-balance.pdf}{Дві шальки: пристрасть і безпристрасність}{Схема наша, не історична ілюстрація.}

\vsec{18.63}
\vdeva{भावनाभावनासक्ता दृष्टिर् मूढस्य सर्वदा । भाव्यभावनया सा तु स्वस्थस्यादृष्टिरूपिणी ॥}
\vtranslit{бгаванабгаванасакта дріштір мудгася сарвада | бгав'ябгаваная са ту свастхасядріштірупіні ||}
\vperek{Погляд затьмареного завжди прив'язаний до уявлення й неуявлення. А в того, хто в собі, він має вигляд не-бачення.}
\vcom{Та сама пара, що в 18.42 (`bhāvaka` / abhāvaka): дурний застрягає на обох боках. `adṛṣṭirūpiṇī` («та, що має ВИГЛЯД не-бачення») — обережна формула: не «не бачить», а має такий вигляд. `svastha` («той, хто в собі») — те саме слово, що `svāsthya` у 13.1.}
\vfig{ilustratsii/18-63-illusion.pdf}{Погляд уявлення й не-бачення}{Схема наша, не історична ілюстрація.}

\vsec{18.64}
\vdeva{सर्वारम्भेषु निष्कामो यश् चरेद् बालवन् मुनिः । न लेपस् तस्य शुद्धस्य क्रियमाणे ऽपि कर्मणि ॥}
\vtranslit{сарварамбгешу нішкамо яш чаред балаван муніх | на лепас тася шуддгася кріямане ’пі кармані ||}
\vperek{Мудрець, що без бажань у всіх починаннях, поводиться як дитина. У чистого немає налипання, навіть коли діяння робиться.}
\vcom{`bālavat` («як дитина») — те саме слово, що в 18.7 і 18.49, і тут воно означає не наївність, а відсутність розрахунку. `lepa` («налипання») — технічне слово про те, як дія пристає до діяча; `kriyamāṇe 'pi karmaṇi` (locativus absolutus) ставить умову максимально невигідно: дія справді робиться, а налипання немає.}
\vfig{ilustratsii/18-64-motion.pdf}{Мудрець поводиться як дитина}{Схема наша, не історична ілюстрація.}

\vsec{18.65}
\vdeva{स एव धन्य आत्मज्ञः सर्वभावेषु यः समः । पश्यन् शृण्वन् स्पृशन् जिघ्रन्न् अश्नन् निस्तर्षमानसः ॥}
\vtranslit{са ева дганя атмаджнях сарвабгавешу ях самах | пашян шрінван спрішан джіґгранн ашнан ністаршаманасах ||}
\vperek{Саме він блаженний — той, хто знає Атман і рівний до всіх: дивиться, слухає, торкається, нюхає, їсть — а розум без спраги.}
\vcom{П'ять чуттєвих дієприкметників — той самий перелік, що в 18.47, і це вже вдруге: книга навмисно повторює його, щоб зняти підозру, ніби йдеться про відмову від чуттів. `nistarṣa` («без спраги») — точніше за «без бажань»: спрага (tarṣa) — те, що змушує пити знову.}
\vfig{ilustratsii/18-65-motion.pdf}{П'ять чуттів і розум без спраги}{Схема наша, не історична ілюстрація.}

\vsec{18.66}
\vdeva{क्व संसारः क्व चाभासः क्व साध्यं क्व च साधनम् । आकाशस्येव धीरस्य निर्विकल्पस्य सर्वदा ॥}
\vtranslit{ква самсарах ква чабгасах ква садгям ква ча садганам | акашасєва дгірася нірвікалпася сарвада ||}
\vperek{Де сансара, де навіть її видимість, де мета, де засіб — у мудрого, вільного від розрізнень, завжди, як у простору?}
\vcom{Чотири питання закривають усю схему практики: світ, його видимість, мета, засіб. `ābhāsa` («відблиск») доданий навмисно — після сансари питається й про її тінь, тобто навіть заперечення не має місця. Порівняння з простором (`ākāśa`) — те саме, що в 6.1, але тут воно в родовому: не «я — простір», а «в нього, як у простору».}
\vfig{ilustratsii/18-66-space.pdf}{Сансара й мета — як у простору}{Схема наша, не історична ілюстрація.}

\vsec{18.67}
\vdeva{स जयत्य् अर्थसंन्यासी पूर्णस्वरसविग्रहः । अकृत्रिमो ऽनवच्छिन्ने समाधिर् यस्य वर्तते ॥}
\vtranslit{са джаятй артхасамнясі пурнасварасавіґрахах | акрітрімо ’наваччінне самадгір яся вартате ||}
\vperek{Перемагає той, хто зрікся користі, чиє тіло — повний власний смак: у нього самадгі невдаване й триває в необмеженому.}
\vcom{`pūrṇasvarasavigraha` — одне з найщільніших складених у книзі: тіло (vigraha) тут не вмістилище, а сам смак (rasa) і власний тон (svara). `anavacchinna` («невідтинене») — точніше за «безмежне»: самадгі не має краю, бо його ніщо не відтяло. `akṛtrima` повторює 18.41 і 18.52: невдаваність — сталий критерій у цій главі.}
\vfig{ilustratsii/18-67-flame.pdf}{Невдаване самадгі}{Схема наша, не історична ілюстрація.}

\vsec{18.68}
\vdeva{बहुनात्र किम् उक्तेन ज्ञाततत्त्वो महाशयः । भोगमोक्षनिराकाङ्क्षी सदा सर्वत्र नीरसः ॥}
\vtranslit{бахунатра кім уктена джнятататтво махашаях | бгоґамокшаніраканкші сада сарватра нірасах ||}
\vperek{До чого тут багато слів: великодухий, що пізнав істину, не прагне ні насолод, ні визволення — завжди й скрізь він без смаку до них.}
\vcom{`bhogamokṣanirākāṅkṣī` — найсильніше складове блоку: насолода й визволення поставлені в один ряд як рівноцінні об'єкти прагнення, і обидва зняті. `nīrasa` («без соку») — не «байдужий», а саме позбавлений смаку до них: метафора їжі, не емоції. `bahunātra kim uktena` — формула, що закриває розмову, і вона стоїть передостаннім віршем глави.}
\vpar{Українська паралель (наша інтерпретація)}{Леся Українка, «Оргія» (1913).}{Наша паралель: у Лесі Українки відсутність смаку і до насолод, і до визволення теж є ознакою того, хто вже вільний. Наше порівняння — типологічне.}
\vfig{ilustratsii/18-68-neti.pdf}{Ні насолод, ні визволення}{Схема наша, не історична ілюстрація.}

\vsec{18.69}
\vdeva{महदादि जगद्द्वैतं नाममात्रविजृम्भितम् । विहाय शुद्धबोधस्य किं कृत्यम् अवशिष्यते ॥}
\vtranslit{махададі джаґаддвайтам намаматравіджрімбгітам | віхая шуддгабодгася кім крітям авашішяте ||}
\vperek{Двоїстість світу, починаючи з махат, розквітла лише з імені. Коли це відкинуто — яка справа лишається в чистого знання?}
\vcom{`nāmamātra` («лише ім'я») — найточніше формулювання: світ не ілюзія й не реальність, а розквітле ім'я. `mahadādi` відсилає до санкх'ї (mahat — перший принцип проявленого), тобто вірш прямо називає філософську схему й каже, що вона тримається на назвах. Питання `kim kṛtyam avaśiṣyate` («яка справа лишається») — останнє питання глави про дію.}
\vfig{ilustratsii/18-69-illusion.pdf}{Двоїстість, що розквітла з імені}{Схема наша, не історична ілюстрація.}

\vsec{18.70}
\vdeva{भ्रमभूतम् इदं सर्वं किंचिन् नास्तीति निश्चयी । अलक्ष्यस्फुरणः शुद्धः स्वभावेनैव शाम्यति ॥}
\vtranslit{бграмабгутам ідам сарвам кімчін настіті нішчаї | алакшяспхуранах шуддгах свабгавенайва шам'яті ||}
\vperek{Хто впевнений, що все це має природу омани й ніщо не є, — той, із сяянням без ознак, чистий, заспокоюється самою природою.}
\vcom{`alakṣyasphuraṇa` («сяяння без ознак») — точна формула на межу мови: те, що сяє, але не має жодної ознаки, за якою його можна впізнати. `svabhāvena eva` («самою природою») закриває главу без зусилля: заспокоєння не досягається, а є властивістю. Порядок важливий — спершу niścaya (рішення), потім śāmyati (заспокоєння), і між ними нічого.}
\vfig{ilustratsii/18-70-neti.pdf}{Усе має природу омани — і заспокоєння}{Схема наша, не історична ілюстрація.}

\vsec{18.71}
\vdeva{शुद्धस्फुरणरूपस्य दृश्यभावम् अपश्यतः । क्व विधिः क्व च वैराग्यं क्व त्यागः क्व शमो ऽपि वा ॥}
\vtranslit{шуддгаспхуранарупася дрішябгавам апашятах | ква відгіх ква ча вайраґям ква тяґах ква шамо ’пі ва ||}
\vperek{У того, чия природа — чисте сяяння, хто не бачить, що видиме існує: де припис, де відречення, де покидання, де й заспокоєння?}
\vcom{Чотири `kva` закривають усю практику: vidhi (припис), vairāgya (відречення), tyāga (покидання), śama (заспокоєння). Зверни увагу: `vairāgya` стоїть у переліку того, що зникає, — тобто в цій главі воно не чеснота, а те, що відпадає разом з іншим. `śuddhasphuraṇa` — те саме сяяння, що `alakṣyasphuraṇa` у 18.70.}
\vfig{ilustratsii/18-71-neti.pdf}{Припис і відречення відпали}{Схема наша, не історична ілюстрація.}

\vsec{18.72}
\vdeva{स्फुरतो ऽनन्तरूपेण प्रकृतिं च न पश्यतः । क्व बन्धः क्व च वा मोक्षः क्व हर्षः क्व विषादिता ॥}
\vtranslit{спхурато ’нантарупена пракрітім ча на пашятах | ква бандгах ква ча ва мокшах ква харшах ква вішадіта ||}
\vperek{У того, хто сяє безконечною формою і не бачить природи: де пута, де визволення, де захват, де смуток?}
\vcom{Та сама будова, що в 18.71, але пари інші: пута й визволення стоять поруч як рівноцінно відсутні, і так само захват із смутком. `viṣāditā` — те саме слово, що в 18.22, і повтор тут не випадковий: глава збирає власні пари.}
\vfig{ilustratsii/18-72-neti.pdf}{Пута й визволення відпали}{Схема наша, не історична ілюстрація.}

\vsec{18.73}
\vdeva{बुद्धिपर्यन्तसंसारे मायामात्रं विवर्तते । निर्ममो निरहङ्कारो निष्कामः शोभते बुधः ॥}
\vtranslit{буддгіпар'янтасамсаре маяматрам вівартате | нірмамо ніраханкаро нішкамах шобгате будгах ||}
\vperek{У сансарі, що сягає межі розуму, обертається лише майя. Вільний від «моє», без «я-роблення», без бажань — сяє той, хто знає.}
\vcom{`buddhiparyanta` («що закінчується на розумі») — точне обмеження: сансара не безмежна, вона сягає рівно до розуму й далі не йде. Три ознаки в другій половині — та сама трійця, що в 18.29 (`ahaṅkāra`) і 15.x (`mama`), і тут вона подана як причина сяйва, а не як наслідок.}
\vfig{ilustratsii/18-73-illusion.pdf}{Сансара, у якій обертається лише майя}{Схема наша, не історична ілюстрація.}

\vsec{18.74}
\vdeva{अक्षयं गतसन्तापम् आत्मानं पश्यतो मुनेः । क्व विद्या च क्व वा विश्वं क्व देहो ऽहं ममेति वा ॥}
\vtranslit{акшаям ґатасантапам атманам пашято мунех | ква відя ча ква ва вішвам ква дехо ’хам маметі ва ||}
\vperek{У мудреця, що бачить Атман — незнищенний, з минулим жаром: де знання, де всесвіт, де тіло, де «я» й «моє»?}
\vcom{У переліку стоїть і `vidyā` — знання, тобто те, чим цю книгу читають. Це найнезручніше місце блоку: знання зникає разом з тілом і світом. `akṣaya` («невитратний») — точніше за «вічний»: те, що не зменшується від ужитку.}
\vfig{ilustratsii/18-74-neti.pdf}{Знання, всесвіт і тіло відпали}{Схема наша, не історична ілюстрація.}

\vsec{18.75}
\vdeva{निरोधादीनि कर्माणि जहाति जडधीर् यदि । मनोरथान् प्रलापांश् च कर्तुम् आप्नोति तत्क्षणात् ॥}
\vtranslit{ніродгадіні кармані джахаті джададгір яді | маноратхан пралапамш ча картум апноті таткшанат ||}
\vperek{Якщо тупий розумом покидає діяння, починаючи з припинення, — він береться за забаганки й балачки, і саме з тієї ж миті.}
\vcom{Точний механізм підміни: покинувши практику, людина не заспокоюється, а починає `manoratha` (забаганки) і `pralāpa` (балачки). `manoratha` буквально «колісниця розуму» — те, що везе думку геть. `tatkṣaṇāt` («з тієї ж миті») робить це не поступовим, а негайним.}
\vfig{ilustratsii/18-75-knot.pdf}{Вузол: покинув діяння — і взявся за балачки}{Схема наша, не історична ілюстрація.}

\vsec{18.76}
\vdeva{मन्दः श्रुत्वापि तद्वस्तु न जहाति विमूढताम् । निर्विकल्पो बहिर्यत्नाद् अन्तर्विषयलालसः ॥}
\vtranslit{мандах шрутвапі тадвасту на джахаті вімудгатам | нірвікалпо бахір'ятнад антарвішаялаласах ||}
\vperek{Тупий, навіть почувши ту суть, не покидає цілковитої затьмареності: зовні він без розрізнень — із зовнішнього зусилля, — а всередині пожадливий до предметів.}
\vcom{`bahiryatnāt` / `antarviṣayalālasaḥ` — пара «зовні / усередині», і вона викриває механізм: `nirvikalpa` досягнуто зусиллям ззовні, тому всередині лишилося пожадання. Це пряме попередження проти наслідування: безрозрізнення можна зобразити, але не втримати. Слово `lālasa` («пожадливий») рідкісне й сильніше за звичайне kāma.}
\vfig{ilustratsii/18-76-illusion.pdf}{Зовні без розрізнень, усередині пожадливий}{Схема наша, не історична ілюстрація.}

\vsec{18.77}
\vdeva{ज्ञानाद् गलितकर्मा यो लोकदृष्ट्यापि कर्मकृत् । नाप्नोत्य् अवसरं कर्तुं वक्तुम् एव न किंचन ॥}
\vtranslit{джнянад ґалітакарма йо локадріштяпі кармакрит | напнотй авасарам картум вактум ева на кімчана ||}
\vperek{Хто через знання має відпалі діяння — хоч і діє з огляду на світ — не знаходить нагоди зробити чи сказати хоч що-небудь.}
\vcom{`galitakarman` («той, у кого діяння відпали») — не «той, хто їх покинув»: відпадання без дії з його боку, і це принципово інша модель, ніж у 18.75, де покидання веде до балачок. `lokadṛṣṭyā api` («хоч і з огляду на світ») зберігає видимість дії: він діє, але нагоди не має.}
\vfig{ilustratsii/18-77-neti.pdf}{Немає нагоди зробити чи сказати}{Схема наша, не історична ілюстрація.}

\vsec{18.78}
\vdeva{क्व तमः क्व प्रकाशो वा हानं क्व च न किञ्चन । निर्विकारस्य धीरस्य निरातङ्कस्य सर्वदा ॥}
\vtranslit{ква тамах ква пракашо ва ханам ква ча на кіньчана | нірвікарася дгірася ніратанкася сарвада ||}
\vperek{Де темрява, де світло, де відкидання, де «ніщо» — у незмінного, мудрого, безтривожного, завжди?}
\vcom{Пара `tamas` / `prakāśa` (темрява / світло) — та сама, що в 18.9, але тут до неї додано `hāna` (відкидання) і `na kiñcana` («ніщо»). Тобто знято не лише протилежності, а й сам вихід із них: навіть «ніщо» не лишається як позиція.}
\vfig{ilustratsii/18-78-neti.pdf}{Темрява й світло відпали}{Схема наша, не історична ілюстрація.}

\vsec{18.79}
\vdeva{क्व धैर्यं क्व विवेकित्वं क्व निरातङ्कतापि वा । अनिर्वाच्यस्वभावस्य निःस्वभावस्य योगिनः ॥}
\vtranslit{ква дгайр'ям ква вівекітвам ква ніратанкатапі ва | анірвачясвабгавася ніхсвабгавася йоґінах ||}
\vperek{Де стійкість, де розрізнення, де й безтривожність — у йога, чия природа невимовна, хто без природи?}
\vcom{`anirvācyasvabhāva` і `niḥsvabhāva` стоять поруч, і це не тавтологія: перше каже, що природу не можна назвати, друге — що її немає. Разом вони закривають обидва виходи. `vivekitva` («розрізнюваність») у переліку того, що зникає, — найнезручніше для читача, бо це чеснота розпізнавання істинного від хибного.}
\vfig{ilustratsii/18-79-neti.pdf}{Стійкість і розрізнення відпали}{Схема наша, не історична ілюстрація.}

\vsec{18.80}
\vdeva{न स्वर्गो नैव नरको जीवन्मुक्तिर् न चैव हि । बहुनात्र किम् उक्तेन योगदृष्ट्या न किंचन ॥}
\vtranslit{на сварґо найва нарако джіванмуктір на чайва хі | бахунатра кім уктена йоґадріштя на кімчана ||}
\vperek{Немає неба, немає пекла, немає й визволення за життя. До чого багато слів: з погляду йоги — ніщо.}
\vcom{`jīvanmukti` знято разом із небом і пеклом — тобто технічний термін книги поставлено в один ряд із релігійними уявленнями. Це найрадикальніше місце глави, і воно стоїть перед останніми віршами. `bahunātra kim uktena` повторює 18.68: та сама формула закриття, і вдруге вона закриває вже не розмову, а весь перелік.}
\vfig{ilustratsii/18-80-neti.pdf}{Небо, пекло й визволення — ніщо}{Схема наша, не історична ілюстрація.}

\vsec{18.81}
\vdeva{नैव प्रार्थयते लाभं नालाभेनानुशोचति । धीरस्य शीतलं चित्तम् अमृतेनैव पूरितम् ॥}
\vtranslit{найва прартхаяте лабгам налабгенанушочаті | дгірася шіталам чіттам амрітенайва пурітам ||}
\vperek{Він не прагне набутку й не сумує через ненабуток: розум мудрого прохолодний, сповнений амрітою.}
\vcom{Пара `lābha` / `alābha` (набуток / ненабуток) закриває обидва боки, як і в 18.28 і 18.34. `śītala` («прохолодний») повторює 18.22 і 18.30 — це стале слово глави для позначення стану, який не нагрівається ні втратою, ні придбанням. `amṛta` тут не метафора безсмертя, а те, чим розум сповнений замість очікування.}
\vfig{ilustratsii/18-81-balance.pdf}{Дві шальки: набуток і ненабуток}{Схема наша, не історична ілюстрація.}

\vsec{18.82}
\vdeva{न शान्तं स्तौति निष्कामो न दुष्टम् अपि निन्दति । समदुःखसुखस् तृप्तः किंचित् कृत्यं न पश्यति ॥}
\vtranslit{на шантам стауті нішкамо на душтам апі ніндаті | самадухкхасукхас тріптах кімчіт крітям на пашяті ||}
\vperek{Без бажань він не хвалить спокійного й не ганить лихого. Рівний у стражданні й щасті, насичений, він не бачить жодної справи.}
\vcom{Пара `śānta` / `duṣṭa` (спокійний / лихий) цікава тим, що обидві оцінки знято разом: похвала й осуд — один механізм. `tṛpta` («насичений») — не «задоволений», а саме той, хто вже наївся, тому `kṛtya` («справа») йому не видно. Це відповідь на 18.57, де сансара була `kartavyatā`.}
\vfig{ilustratsii/18-82-balance.pdf}{Дві шальки: хвала й осуд}{Схема наша, не історична ілюстрація.}

\vsec{18.83}
\vdeva{धीरो न द्वेष्टि संसारम् आत्मानं न दिदृक्षति । हर्षामर्षविनिर्मुक्तो न मृतो न च जीवति ॥}
\vtranslit{дгіро на двешті самсарам атманам на дідрікшаті | харшамаршавінірмукто на мріто на ча джіваті ||}
\vperek{Мудрий не ненавидить сансару й не прагне побачити Атман. Вільний від захвату й гніву, він ні мертвий, ні живий.}
\vcom{`didṛkṣati` (дезидератив «хоче побачити») — точне слово: знято не знання Атмана, а **бажання його побачити**, тобто саму позицію того, хто ще не бачив. `na mṛto na ca jīvati` — не парадокс заради ефекту: обидва слова описують стан того, хто має межі, а він їх не має.}
\vfig{ilustratsii/18-83-balance.pdf}{Дві шальки: сансара й Атман}{Схема наша, не історична ілюстрація.}

\vsec{18.84}
\vdeva{निःस्नेहः पुत्रदारादौ निष्कामो विषयेषु च । निश्चिन्तः स्वशरीरे ऽपि निराशः शोभते बुधः ॥}
\vtranslit{ніхснехах путрадарадау нішкамо вішаєшу ча | нішчінтах свашаріре ’пі нірашах шобгате будгах ||}
\vperek{Без прив'язаності до синів, жінки й інших, без бажань до предметів, без турботи навіть про власне тіло — без надій сяє той, хто знає.}
\vcom{`sneha` буквально «вологість» — те, чим речі прилипають; тому `niḥsneha` не «байдужий», а «незмочений». Перелік іде від найближчого до найдальшого й закінчується власним тілом — тобто остання прив'язка названа останньою. `nirāśa` («без надій») стоїть перед `śobhate`, і це єдина ствердна дія у вірші.}
\vfig{ilustratsii/18-84-neti.pdf}{Син, жінка, тіло — без надій}{Схема наша, не історична ілюстрація.}

\vsec{18.85}
\vdeva{तुष्टिः सर्वत्र धीरस्य यथापतितवर्तिनः । स्वच्छन्दं चरतो देशान् यत्रास्तमितशायिनः ॥}
\vtranslit{туштіх сарватра дгірася ятхапатітавартінах | сваччандам чарато дешан ятрастаміташаїнах ||}
\vperek{Вдоволення скрізь у мудрого, що живе як випаде, вільно ходить землями й лягає там, де сонце зайшло.}
\vcom{`yathāpatita` («як випаде») — те саме слово, що `yadṛcchā` у 18.24 і 15.x. Остання частина конкретна й несподівано побутова: `astamitaśāyin` — лягає спати там, де його застав захід сонця. Це не метафора, а опис вільного мандрівника без розкладу.}
\vfig{ilustratsii/18-85-motion.pdf}{Вдоволення скрізь, де зайшло сонце}{Схема наша, не історична ілюстрація.}

\vsec{18.86}
\vdeva{पततूदेतु वा देहो नास्य चिन्ता महात्मनः । स्वभावभूमिविश्रान्तिविस्मृताशेषसंसृतेः ॥}
\vtranslit{пататудету ва дехо нася чінта махатманах | свабгавабгумівішрантівісмріташешасамсрітех ||}
\vperek{Хай тіло падає або встає — у великого немає турботи: він спочив на ґрунті власної природи й забув усю сансару без останку.}
\vcom{`patatu ... udetu` — два імперативи, і вони стосуються тіла, не людини: тіло поставлене в граматичну позицію того, хто діє сам. `svabhāvabhūmi` («ґрунт власної природи») — точніше за «основа»: це те, на чому стоїть, а не те, чим є. `aśeṣa` («без останку») закриває: забуто не частину сансари, а всю.}
\vfig{ilustratsii/18-86-balance.pdf}{Дві шальки: тіло падає чи встає}{Схема наша, не історична ілюстрація.}

\vsec{18.87}
\vdeva{अकिंचनः कामचारो निर्द्वन्द्वश् छिन्नसंशयः । असक्तः सर्वभावेषु केवलो रमते बुधः ॥}
\vtranslit{акімчанах камачаро нірдвандваш ччіннасамшаях | асактах сарвабгавешу кевало рамате будгах ||}
\vperek{Без жодної речі, той, хто йде куди хоче, без двоїстості, з відтятим сумнівом, неприв'язаний до всіх станів — сам-один утішається той, хто знає.}
\vcom{`kevala` тут не «лише», а «сам-один, без домішку» — технічне слово, що пізніше стало назвою школи. `akiṃcana` («без хоч чогось») — найсильніша форма відсутності володіння в книзі. `ramate` («утішається») — єдине дієслово радості в цьому блоці, і воно стоїть останнім, після п'яти заперечень.}
\vfig{ilustratsii/18-87-space.pdf}{Без жодної речі — сам-один}{Схема наша, не історична ілюстрація.}

\vsec{18.88}
\vdeva{निर्ममः शोभते धीरः समलोष्टाश्मकाञ्चनः । सुभिन्नहृदयग्रन्थिर् विनिर्धूतरजस्तमः ॥}
\vtranslit{нірмамах шобгате дгірах самалошташмаканьчанах | субгіннахрідаяґрантхір вінірдгутараджастамах ||}
\vperek{Без «моє» сяє мудрий, для кого грудка, камінь і золото те саме; у кого добре розв'язаний вузол серця, хто струсив геть пил і темряву.}
\vcom{`loṣṭa` («грудка землі») навмисно нижче за камінь і золото: перелік іде від найдешевшого до найдорожчого, і всі три зрівняні. `hṛdayagranthi` («вузол серця») — усталений образ, але `su-bhinna` («добре розбитий») додає різкості: вузол не розв'язують, а розбивають. `rajas` і `tamas` — дві з трьох ґун, названі прямо.}
\vfig{ilustratsii/18-88-balance.pdf}{Дві шальки: грудка й золото}{Схема наша, не історична ілюстрація.}

\vsec{18.89}
\vdeva{सर्वत्रानवधानस्य न किञ्चिद् वासना हृदि । मुक्तात्मनो वितृप्तस्य तुलना केन जायते ॥}
\vtranslit{сарватранавадганася на кіньчід васана хріді | муктатмано вітріптася тулана кена джаяте ||}
\vperek{У того, хто ніде не має опори, немає жодного сліду в серці. З ким зрівняти визволеного Атмана, що наситився?}
\vcom{`anavadhāna` («без закріплення уваги») — не неуважність, а відсутність точки, за яку можна зачепитися; саме тому далі немає й `vāsanā`. Питання `tulanā kena jāyate` («з ким зрівняти») закриває останню можливість: порівняння потребує другого, а другого немає.}
\vpar{Українська паралель (ПОТРЕБУЄ ПЕРЕВІРКИ)}{українська релігійна поезія, не визначено.}{ПОТРЕБУЄ ПЕРЕВІРКИ: питання, з ким зрівняти визволеного, близьке до української релігійної поезії, але конкретного твору ми не звірили.}
\vfig{ilustratsii/18-89-witness.pdf}{Без опори — і немає сліду в серці}{Схема наша, не історична ілюстрація.}

\vsec{18.90}
\vdeva{जानन्न् अपि न जानाति पश्यन्न् अपि न पश्यति । ब्रुवन्न् अपि न च ब्रूते को ऽन्यो निर्वासनाद् ऋते ॥}
\vtranslit{джананн апі на джанаті пашянн апі на пашяті | бруванн апі на ча бруте ко ’нйо нірвасанад ріте ||}
\vperek{І знаючи — не знає, і бачачи — не бачить, і говорячи — не говорить. Хто інший, крім вільного від слідів?}
\vcom{Три дієприкметники з `api` («і... таки») — найчистіша форма цієї думки в книзі: дія відбувається, діяч відсутній. `nirvāsana` («без васан») повертає до 18.21, де він був першим словом блоку, — глава замикається на собі. Питання `kaḥ anyaḥ` замість твердження: відповіді не дано, бо й питання стосується того, кого нема з чим порівняти (пор. 18.89).}
\vfig{ilustratsii/18-90-neti.pdf}{Знає — не знає; бачить — не бачить}{Схема наша, не історична ілюстрація.}

\vsec{18.91}
\vdeva{भिक्षुर् वा भूपतिर् वापि यो निष्कामः स शोभते । भावेषु गलिता यस्य शोभनाशोभना मतिः ॥}
\vtranslit{бгікшур ва бгупатір вапі йо нішкамах са шобгате | бгавешу ґаліта яся шобганашобгана матіх ||}
\vperek{Жебрак чи володар — той, хто без бажань, сяє: у кого в усіх станах відпала думка про добре й недобре.}
\vcom{`bhikṣu` і `bhūpati` поставлені як рівноцінні, і це єдине місце в книзі, де соціальна крайність названа прямо. `galitā` («відпала, стікла») — те саме слово, що в 18.77; тут воно про `mati`, тобто про саму здатність оцінювати. `śobhana` / `aśobhana` — не «добре/погане» в етичному сенсі, а «гарне/негарне»: оцінка як смак.}
\vfig{ilustratsii/18-91-balance.pdf}{Дві шальки: жебрак і володар}{Схема наша, не історична ілюстрація.}

\vsec{18.92}
\vdeva{क्व स्वाच्छन्द्यं क्व सङ्कोचः क्व वा तत्त्वविनिश्चयः । निर्व्याजार्जवभूतस्य चरितार्थस्य योगिनः ॥}
\vtranslit{ква сваччандям ква санкочах ква ва таттвавінішчаях | нірв'яджарджавабгутася чарітартхася йоґінах ||}
\vperek{Де свобода, де стиснення, де рішення про істину — у йога, чия істота нелукаво-пряма, хто досяг мети?}
\vcom{`svācchandya` / `saṅkoca` — та сама пара, що в 18.32, але там `saṅkoca` був наслідком неправильно прийнятої істини, а тут знято обидві. `nirvyāja` («без обману») — рідкісне й точне слово: не «щирий» за наміром, а без підміни взагалі. `caritārtha` («чия справа зроблена») закриває тему дії, яку глава вела від 18.19.}
\vfig{ilustratsii/18-92-neti.pdf}{Свобода й стиснення відпали}{Схема наша, не історична ілюстрація.}

\vsec{18.93}
\vdeva{आत्मविश्रान्तितृप्तेन निराशेन गतार्तिना । अन्तर् यद् अनुभूयेत तत् कथं कस्य कथ्यते ॥}
\vtranslit{атмавішрантітріптена нірашена ґатартіна | антар яд анубгуєта тат катхам кася катхяте ||}
\vperek{Те, що переживається всередині в насиченого спочинком в Атмані, без надій, з минулим стражданням — як і кому це розповісти?}
\vcom{Подвійне питання `katham` / `kasya` («як?» і «кому?») закриває можливість передачі: не бракує слів, а немає ні способу, ні адресата. Форма `anubhūyeta` — оптатив пасиву, тобто «те, що могло б переживатися»: мова йде про можливість, а не про подію. Це пряме пояснення, чому глава весь час говорить запереченнями.}
\vfig{ilustratsii/18-93-flame.pdf}{Те, що переживається всередині, — не розповідається}{Схема наша, не історична ілюстрація.}

\vsec{18.94}
\vdeva{सुप्तो ऽपि न सुषुप्तौ च स्वप्ने ऽपि शयितो न च । जागरे ऽपि न जागर्ति धीरस् तृप्तः पदे पदे ॥}
\vtranslit{супто ’пі на сушуптау ча свапне ’пі шаїто на ча | джаґаре ’пі на джаґарті дгірас тріптах паде паде ||}
\vperek{І сплячи — не в глибокому сні; і в сні — не лежить; і в неспанні — не неспить. Мудрий, насичений на кожнім кроці.}
\vcom{Три стани свідомості (suṣupti, svapna, jāgara) названі поіменно й усі три знято — це відповідь усій традиції аналізу трьох станів. `pade pade` («на кожнім кроці») повертає до `svapada` з 18.33 і 18.5: те саме слово «крок/місце», і тут воно вжите як прислівник.}
\vfig{ilustratsii/18-94-balance.pdf}{Дві шальки: сон і неспання}{Схема наша, не історична ілюстрація.}

\vsec{18.95}
\vdeva{ज्ञः सचिन्तो ऽपि निश्चिन्तः सेन्द्रियो ऽपि निरिन्द्रियः । सबुद्धिर् अपि निर्बुद्धिः साहङ्कारो ऽनहंकृतिः ॥}
\vtranslit{джнях сачінто ’пі нішчінтах сендрійо ’пі ніріндріях | сабуддгір апі нірбуддгіх саханкаро ’нахамкрітіх ||}
\vperek{Знаючий, хоч і з думкою — без думки; хоч і з чуттями — без чуттів; хоч і з розумом — без розуму; хоч і з «я-робленням» — без нього.}
\vcom{Чотири пари з `api` («хоч і... таки») — найчистіша форма цього ходу в книзі: кожна ознака є й водночас не діє. `ahaṃkṛti` тут замість звичного `ahaṃkāra` — і це не варіант: `kṛti` («роблення») наголошує саме на акті, а не на результаті. Порівняй 18.29, де той самий механізм описано розгорнуто.}
\vfig{ilustratsii/18-95-neti.pdf}{З думкою — без думки; з розумом — без розуму}{Схема наша, не історична ілюстрація.}

\vsec{18.96}
\vdeva{न सुखी न च वा दुःखी न विरक्तो न सङ्गवान् । न मुमुक्षुर् न वा मुक्तो न किंचिन् न च किंचन ॥}
\vtranslit{на сукхі на ча ва духкхі на віракто на санґаван | на мумукшур на ва мукто на кімчін на ча кімчана ||}
\vperek{Не щасливий і не нещасний, не відречений і не прив'язаний, не прагне визволення й не визволений. Він не «щось» і не «ніщо».}
\vcom{Шість заперечень, і серед них `mumukṣu` / `mukta` — тобто знято обидва кінці шляху, який книга описує. Остання пара `kiṃcit` / `kiṃcana` тонка: перше — «щось», друге — «хоч щось, абищо»; разом вони закривають і позицію, і її відсутність. Це найповніший перелік заперечень у книзі.}
\vpar{Українська паралель (наша інтерпретація)}{Ліна Костенко, «Сад нетанучих скульптур» (1987).}{Наша паралель: у Костенко стан поза парою щось і ніщо теж постає як точність, а не як розмитість. Наше порівняння — за будовою думки.}
\vfig{ilustratsii/18-96-neti.pdf}{Ні щасливий, ні нещасний; ні щось, ні ніщо}{Схема наша, не історична ілюстрація.}

\vsec{18.97}
\vdeva{विक्षेपे ऽपि न विक्षिप्तः समाधौ न समाधिमान् । जाड्ये ऽपि न जडो धन्यः पाण्डित्ये ऽपि न पण्डितः ॥}
\vtranslit{вікшепе ’пі на вікшіптах самадгау на самадгіман | джадє ’пі на джадо дганях пандітє ’пі на пандітах ||}
\vperek{І в розсіянні — не розсіяний; і в самадгі — не той, хто має самадгі; і в тупості — не тупий. Блаженний, і у вченості — не вчений.}
\vcom{Пари навмисно принизливі: `jāḍya` («тупість») і `pāṇḍitya` («вченість») стоять поруч як рівноцінні — тобто вченість не краща за тупість, бо обидві є станом. `samādhimān` («той, хто має самадгі») — найтонше заперечення: не «не в самадгі», а не володіє ним. Порівняй 18.33, де зосередженість віддана дурням.}
\vfig{ilustratsii/18-97-balance.pdf}{Розсіяння й самадгі — не його}{Схема наша, не історична ілюстрація.}

\vsec{18.98}
\vdeva{मुक्तो यथास्थितिस्वस्थः कृतकर्तव्यनिर्वृतः । समः सर्वत्र वैतृष्ण्यान् न स्मरत्य् अकृतं कृतम् ॥}
\vtranslit{мукто ятхастхітісвастхах крітакартав'янірврітах | самах сарватра вайтрішнян на смаратй акрітам крітам ||}
\vperek{Вільний, у собі в тому, як є; заспокоєний у зробленому-що-треба; рівний скрізь. Через безспраглість він не пам'ятає ні зробленого, ні незробленого.}
\vcom{`kṛtakartavya` («зроблене-що-треба») — формула, відома з інших текстів, тут у складі: він заспокоєний саме в тому, що справу закрито. `vaitṛṣṇya` («безспраглість») названа причиною забуття, і це точніше, ніж «не пам'ятає»: пам'ять тримається спраги, без неї триматися нема за що.}
\vfig{ilustratsii/18-98-space.pdf}{Безспраглість: не памʼятає зробленого}{Схема наша, не історична ілюстрація.}

\vsec{18.99}
\vdeva{न प्रीयते वन्द्यमानो निन्द्यमानो न कुप्यति । नैवोद्विजति मरणे जीवने नाभिनन्दति ॥}
\vtranslit{на пріяте вандямано ніндямано на куп'яті | найводвіджаті маране джіване набгінандаті ||}
\vperek{Його не тішить, коли хвалять, і не гнівить, коли ганять. Він не лякається смерті й не радіє життю.}
\vcom{Останній вірш книги будується на двох парах, і обидві — про те, що приходить ззовні: похвала й осуд, смерть і життя. Дві форми пасиву (`vandyamāna`, `nindyamāna`) ставлять його в позицію того, з ким це роблять, — і саме тому він не відповідає. `udvijati` («здригається») — тілесне слово: не «боїться думки про смерть», а не здригається перед нею.}
\vpar{Українська паралель (наша інтерпретація)}{Тарас Шевченко, «Неофіти» (1857).}{Наша паралель: у Шевченка хвала й ганьба теж не міняють того, хто стоїть. Наше порівняння — типологічне, не про джерело.}
\vfig{ilustratsii/18-99-balance.pdf}{Дві шальки: хвала й ганьба}{Схема наша, не історична ілюстрація.}

\vsec{18.100}
\vdeva{न धावति जनाकीर्णं नारण्यम् उपशान्तधीः । यथातथा यत्रतत्र सम एवावतिष्ठते ॥}
\vtranslit{на дгаваті джанакірнам нараням упашантадгіх | ятхататха ятрататра сама еваватіштхате ||}
\vperek{Той, чия думка втихла, не біжить ні до переповненого людьми, ні до лісу; так чи інакше, там чи там він перебуває рівним.}
\vcom{Останній вірш книги. Дві відмови — не тікати до людей і не тікати до лісу — повторюють пару з 18.11, але тут вона про **рух**, а не про умови. `yathātathā yatratatra` — подвійна байдужість до способу й до місця; ця формула вже була в 9.4 і 15.18, і переклад мусить бути той самий. `sama eva avatiṣṭhate` — не «лишається спокійним», а «стоїть рівним»: `sama` — рівність, а не тиша.}
\vfig{ilustratsii/18-100-motion.pdf}{Ні до людей, ні до лісу — рівний скрізь}{Схема наша, не історична ілюстрація.}


% ЗГЕНЕРОВАНО scripts/build_format_v2_chapter.py — не редагувати вручну.
% Джерело: editorial/cards/19.*.json, format-v2/rozdil-19.json, paraleli.json.

\naukchapter{19}{8}{Глава 19. Де це все?}
\naukintro{Девʼятнадцята глава складається з восьми віршів, і сім із них мають одну форму: питання де те й де те — і відповідь у мене, що стоїть у власній величі. Це не заперечення перелічених речей: текст питає, де вони, у того, хто вже стоїть. Восьмий вірш обриває навіть цю форму: досить розмов.}
\vsec{19.1}
\vdeva{तत्त्वविज्ञानसन्दंशम् आदाय हृदयोदरात् । नानाविधपरामर्शशल्योद्धारः कृतो मया ॥}
\vtranslit{таттвавіджнянасандамшам адая хрідайодарат | нанавідгапарамаршашалйоддгарах кріто мая ||}
\vperek{Взявши кліщі знання істини, я витяг із нутра серця колючки різноманітних хибних уявлень.}
\vcom{Построкове зіставлення: редакторська чернетка, не рецензовано.}
\vcom{Рядок а}
\vdeva{तत्त्वविज्ञानसन्दंशम् आदाय हृदयोदरात्}
\vtranslit{таттвавіджнянасандамшам адая хрідайодарат}
\vperek{Взявши кліщі знання істини,}
\vcom{tattvavijñāna («знання істини») входить до складного слова із sandaṃśa («щипці, кліщі»); ādāya — дієприслівникове «взявши». hṛdayodarāt означає «з нутра серця», але в українському цілому реченні цей зворот стоїть у другій частині: поділ перекладу є смисловим, не послівним.}
\vcom{Рядок б}
\vdeva{नानाविधपरामर्शशल्योद्धारः कृतो मया}
\vtranslit{нанавідгапарамаршашалйоддгарах кріто мая}
\vperek{я витяг із нутра серця колючки різноманітних хибних уявлень.}
\vcom{nānāvidha — «різноманітний», parāmarśa тут стосується схоплювання або уявлень, śalya — «колючка». uddhāraḥ kṛto mayā передає виконане «мною» вилучення; український активний вираз «я витяг» зберігає чинну редакцію, але не відтворює буквально пасивну побудову санскриту.}
\vcom{Вірш — не декларація, а опис процедури, і метафора тут хірургічна: sandaṃśa — кліщі, якими виймають стороннє тіло з рани. Три слова працюють разом: śalya («колючка») — те, що застрягло; uddhāra («витягання») — дія; hṛdayodara («нутро серця») — місце, і воно назване нутром, а не поверхнею. parāmarśa тут не «думка», а хибне схоплювання — розумове прийняття чогось за інше. Тобто йдеться не про прибирання думок, а про виймання того, що застрягло й болить. РЕФРЕН ГЛАВИ. 19.2–19.7 закінчуються однаково: svamahimni sthitasya me («у мене, що стоїть у власній величі»). Кожен із цих шести віршів — низка питань kva («де»), і жодне не дістає відповіді. 19.1 і 19.8 стоять поза серією: перший описує, ЯК це сталося, останній — ЩО з цього випливає для мови.}
\vpar{Українська паралель (наша інтерпретація)}{Григорій Сковорода, «Асхань. Симфонія, наречена книга Асхань, про пізнання самого себе».}{Наша паралель: у Сковороди пізнання себе теж постає як виймання зайвого, а не як набирання нового. Наше порівняння — типологічне, не твердження про джерело.}
\vfig{ilustratsii/19-1-thread.pdf}{Кліщі знання витягають колючки з серця}{Схема наша, не історична ілюстрація.}

\vsec{19.2}
\vdeva{क्व धर्मः क्व च वा कामः क्व चार्थः क्व विवेकता । क्व द्वैतं क्व च वाद्वैतं स्वमहिम्नि स्थितस्य मे ॥}
\vtranslit{ква дгармах ква ча ва камах ква чартхах ква вівеката | ква двайтам ква ча вадвайтам свамахімні стхітася ме ||}
\vperek{Де дгарма, де бажання, де користь, де вміння розрізняти? Де двоїстість, де недвоїстість — у мене, що стоїть у власній величі?}
\vcom{Построкове зіставлення: редакторська чернетка, не рецензовано.}
\vcom{Рядок а}
\vdeva{क्व धर्मः क्व च वा कामः क्व चार्थः क्व विवेकता}
\vtranslit{ква дгармах ква ча ва камах ква чартхах ква вівеката}
\vperek{Де дгарма, де бажання, де користь, де вміння розрізняти?}
\vcom{Повтор kva («де?») об'єднує dharmaḥ, kāmaḥ, arthaḥ і термін розрізнення. Чинна картка має vivekatā, тоді як електронний GRETIL читає vivekitā: цю лексичну відмінність потрібно окремо перевірити; механічно виправляти авторський український переклад не можна.}
\vcom{Рядок б}
\vdeva{क्व द्वैतं क्व च वाद्वैतं स्वमहिम्नि स्थितस्य मे}
\vtranslit{ква двайтам ква ча вадвайтам свамахімні стхітася ме}
\vperek{Де двоїстість, де недвоїстість — у мене, що стоїть у власній величі?}
\vcom{Два питання kva dvaitaṃ і kva ... advaitaṃ охоплюють і двоїстість, і недвоїстість. svamahimni — місцевий відмінок «у власній величі», sthitasya me — родове «мене, що стоїть»; формула повторюється у віршах 19.2–19.7, і її український відповідник збережено.}
\vcom{Перший із шести віршів-переліків. Три з перших чотирьох понять — це три з чотирьох пурушартг (дгарма, користь, бажання); четверте, визволення, не назване, бо його місце займає сама позиція того, хто питає. Найважливіше — остання пара: питається не лише «де двоїстість», а й «де недвоїстість». Тобто заперечується не світ на користь єдності, а сама вісь, на якій обидва поняття мають сенс. svamahimni («у власній величі») — не гордість, а міра: те, що не має ні межі, ні місця. РЕФРЕН ГЛАВИ. 19.2–19.7 закінчуються однаково: svamahimni sthitasya me («у мене, що стоїть у власній величі»). Кожен із цих шести віршів — низка питань kva («де»), і жодне не дістає відповіді. 19.1 і 19.8 стоять поза серією: перший описує, ЯК це сталося, останній — ЩО з цього випливає для мови.}
\vpar{Українська паралель (наша інтерпретація)}{Митрополит Іларіон Київський, «Слово про закон і благодать» (XI ст.).}{Наша паралель: у київському Слові перелік того, що вже не важить, теж будується як питання, а не як твердження. Наше порівняння — типологічне; датування цієї книги науково не встановлене.}
\vfig{ilustratsii/19-2-neti.pdf}{Дгарма, бажання й двоїстість відпали}{Схема наша, не історична ілюстрація.}

\vsec{19.3}
\vdeva{क्व भूतं क्व भविष्यद् वा वर्तमानम् अपि क्व वा । क्व देशः क्व च वा नित्यं स्वमहिम्नि स्थितस्य मे ॥}
\vtranslit{ква бгутам ква бгавішяд ва вартаманам апі ква ва | ква дешах ква ча ва нітям свамахімні стхітася ме ||}
\vperek{Де минуле, де майбутнє, де навіть теперішнє? Де сторона світу, де вічність — у мене, що стоїть у власній величі?}
\vcom{Построкове зіставлення: редакторська чернетка, не рецензовано.}
\vcom{Рядок а}
\vdeva{क्व भूतं क्व भविष्यद् वा वर्तमानम् अपि क्व वा}
\vtranslit{ква бгутам ква бгавішяд ва вартаманам апі ква ва}
\vperek{Де минуле, де майбутнє, де навіть теперішнє?}
\vcom{bhūtaṃ, bhaviṣyad і vartamānam позначають минуле, майбутнє та теперішнє; api посилює згадку про теперішнє («навіть»). Питальне kva повторюється, а не називає фактичної відсутності календарного часу в житті читача.}
\vcom{Рядок б}
\vdeva{क्व देशः क्व च वा नित्यं स्वमहिम्नि स्थितस्य मे}
\vtranslit{ква дешах ква ча ва нітям свамахімні стхітася ме}
\vperek{Де сторона світу, де вічність — у мене, що стоїть у власній величі?}
\vcom{deśaḥ — місце або сторона, nityaṃ — постійне чи вічне. Рефрен svamahimni sthitasya me співвідносить простір і тривалість зі станом мовця; це філософське питання, не фізичне заперечення простору.}
\vcom{Три часи названі поспіль, і слово api («навіть») стоїть саме при теперішньому: минуле й майбутнє відсутні легко, а теперішнє — те, що здається найбільш незаперечним. Друга пара — deśa («місце») і nitya («вічне») — зіставляє простір із часом; разом обидві пари охоплюють усе, у чому можна було б розміститися. Вірш показує, що «стояти у власній величі» означає не займати іншого місця, а не мати місця взагалі. РЕФРЕН ГЛАВИ. 19.2–19.7 закінчуються однаково: svamahimni sthitasya me («у мене, що стоїть у власній величі»). Кожен із цих шести віршів — низка питань kva («де»), і жодне не дістає відповіді. 19.1 і 19.8 стоять поза серією: перший описує, ЯК це сталося, останній — ЩО з цього випливає для мови.}
\vpar{Українська паралель (ПОТРЕБУЄ ПЕРЕВІРКИ)}{українська філософська лірика, не визначено.}{ПОТРЕБУЄ ПЕРЕВІРКИ: питання, де минуле й майбутнє, близьке до української філософської лірики, але конкретного твору ми не звірили.}
\vfig{ilustratsii/19-3-neti.pdf}{Минуле, майбутнє й вічність відпали}{Схема наша, не історична ілюстрація.}

\vsec{19.4}
\vdeva{क्व चात्मा क्व च वानात्मा क्व शुभं क्वाशुभं तथा । क्व चिन्ता क्व च वाचिन्ता स्वमहिम्नि स्थितस्य मे ॥}
\vtranslit{ква чатма ква ча ванатма ква шубгам квашубгам татха | ква чінта ква ча вачінта свамахімні стхітася ме ||}
\vperek{Де Атман, де не-Атман, де добре, де недобре? Де думка, де відсутність думки — у мене, що стоїть у власній величі?}
\vcom{Построкове зіставлення: редакторська чернетка, не рецензовано.}
\vcom{Рядок а}
\vdeva{क्व चात्मा क्व च वानात्मा क्व शुभं क्वाशुभं तथा}
\vtranslit{ква чатма ква ча ванатма ква шубгам квашубгам татха}
\vperek{Де Атман, де не-Атман, де добре, де недобре?}
\vcom{Пари ātmā / anātmā та śubhaṃ / aśubhaṃ запитально зіставлені через kva. Чинна картка бере SARIT kvāśubhaṃ tathā; GRETIL має kvāsubhaṃ yathā — відмінність торкається і слова, і зв'язки. Жодного з цих електронних читань не оголошено незалежно підтвердженим рукописним оригіналом.}
\vcom{Рядок б}
\vdeva{क्व चिन्ता क्व च वाचिन्ता स्वमहिम्नि स्थितस्य मे}
\vtranslit{ква чінта ква ча вачінта свамахімні стхітася ме}
\vperek{Де думка, де відсутність думки — у мене, що стоїть у власній величі?}
\vcom{cintā означає думку або турботу, acintā — її відсутність; запитання поширюється на обидва боки пари. Рефрен svamahimni sthitasya me збережено однаково з сусідніми віршами. Цей рядок не дає підстав наказувати українському читачеві відмовитися від мислення.}
\vcom{Три пари, і в кожній заперечується не один бік, а обидва. Найгостріша — перша: питати «де Атман» і «де не-Атман» однаково означає не знайти місця ні для того, ні для того, тобто сама пара втрачає ґрунт. Друга пара (śubha / aśubha) — та сама, що в 13.7, де вона відкидалася цілою. Третя (cintā / acintā) важлива тим, що сюди потрапляє й «відсутність думки»: навіть спокій розуму не стає позицією, з якої можна стояти. РЕФРЕН ГЛАВИ. 19.2–19.7 закінчуються однаково: svamahimni sthitasya me («у мене, що стоїть у власній величі»). Кожен із цих шести віршів — низка питань kva («де»), і жодне не дістає відповіді. 19.1 і 19.8 стоять поза серією: перший описує, ЯК це сталося, останній — ЩО з цього випливає для мови.}
\vpar{Українська паралель (наша інтерпретація)}{Олександр Потебня, «Думка й мова» (1862).}{Наша паралель: Потебня теж показує, що пара думка й відсутність думки є справою самого мірила, а не речей. Наше порівняння — типологічне.}
\vfig{ilustratsii/19-4-neti.pdf}{Атман і не-Атман, думка й не-думка відпали}{Схема наша, не історична ілюстрація.}

\vsec{19.5}
\vdeva{क्व स्वप्नः क्व सुषुप्तिर् वा क्व च जागरणं तथा । क्व तुरीयं भयं वापि स्वमहिम्नि स्थितस्य मे ॥}
\vtranslit{ква свапнах ква сушуптір ва ква ча джаґаранам татха | ква туріям бгаям вапі свамахімні стхітася ме ||}
\vperek{Де сон зі сновидіннями, де глибокий сон, де неспання? Де четвертий стан, де навіть страх — у мене, що стоїть у власній величі?}
\vcom{Построкове зіставлення: редакторська чернетка, не рецензовано.}
\vcom{Рядок а}
\vdeva{क्व स्वप्नः क्व सुषुप्तिर् वा क्व च जागरणं तथा}
\vtranslit{ква свапнах ква сушуптір ва ква ча джаґаранам татха}
\vperek{Де сон зі сновидіннями, де глибокий сон, де неспання?}
\vcom{svapnaḥ — сон зі сновидіннями, suṣuptiḥ — глибокий сон, jāgaraṇaṃ — неспання. kva ставить під питання всі три стани; така трійця має філософський контекст, але тут не є описом медичного або психологічного дослідження сну.}
\vcom{Рядок б}
\vdeva{क्व तुरीयं भयं वापि स्वमहिम्नि स्थितस्य मे}
\vtranslit{ква туріям бгаям вапі свамахімні стхітася ме}
\vperek{Де четвертий стан, де навіть страх — у мене, що стоїть у власній величі?}
\vcom{turīyaṃ — «четвертий» стан, bhayaṃ — «страх», api посилює друге питання. Відмінність turīyaṃ / turiyaṃ між електронними текстами потребує орфографічного обліку; висновок про рукописне різночитання з неї не випливає.}
\vcom{Три стани — неспання, сон зі сновидіннями, глибокий сон — це стандартна схема упанішад, і четвертий (turīya) стоїть поза нею. Тут відкидається й він: навіть те, що мало б бути виходом, не має місця. Страх (bhaya) доданий останнім, і це не випадково: страх — те, що тримає в межах трьох станів, і якщо його немає, то й виходити нікуди. РЕФРЕН ГЛАВИ. 19.2–19.7 закінчуються однаково: svamahimni sthitasya me («у мене, що стоїть у власній величі»). Кожен із цих шести віршів — низка питань kva («де»), і жодне не дістає відповіді. 19.1 і 19.8 стоять поза серією: перший описує, ЯК це сталося, останній — ЩО з цього випливає для мови.}
\vpar{Українська паралель (ПОТРЕБУЄ ПЕРЕВІРКИ)}{українська релігійна проза, не визначено.}{ПОТРЕБУЄ ПЕРЕВІРКИ: перелік станів від сну до четвертого стану найближчий до української релігійної прози, але назвати твір без звірки ми не можемо.}
\vfig{ilustratsii/19-5-neti.pdf}{Сон, неспання й четвертий стан відпали}{Схема наша, не історична ілюстрація.}

\vsec{19.6}
\vdeva{क्व दूरं क्व समीपं वा बाह्यं क्वाभ्यन्तरं क्व वा । क्व स्थूलं क्व च वा सूक्ष्मं स्वमहिम्नि स्थितस्य मे ॥}
\vtranslit{ква дурам ква саміпам ва бахям квабгянтарам ква ва | ква стхулам ква ча ва сукшмам свамахімні стхітася ме ||}
\vperek{Де далеке, де близьке, де зовнішнє, де внутрішнє? Де грубе, де тонке — у мене, що стоїть у власній величі?}
\vcom{Построкове зіставлення: редакторська чернетка, не рецензовано.}
\vcom{Рядок а}
\vdeva{क्व दूरं क्व समीपं वा बाह्यं क्वाभ्यन्तरं क्व वा}
\vtranslit{ква дурам ква саміпам ва бахям квабгянтарам ква ва}
\vperek{Де далеке, де близьке, де зовнішнє, де внутрішнє?}
\vcom{dūraṃ / samīpaṃ позначають далеке й близьке; bāhyaṃ / abhyantaraṃ — зовнішнє й внутрішнє. Це триває низка питань kva, а не твердження про зникнення фізичних відстаней. У GRETIL можливе написання bahyaṃ, у SARIT bāhyaṃ: перевірити довготу голосної при фаховій колації.}
\vcom{Рядок б}
\vdeva{क्व स्थूलं क्व च वा सूक्ष्मं स्वमहिम्नि स्थितस्य मे}
\vtranslit{ква стхулам ква ча ва сукшмам свамахімні стхітася ме}
\vperek{Де грубе, де тонке — у мене, що стоїть у власній величі?}
\vcom{sthūlaṃ та sūkṣmaṃ зіставляють грубе або відчутне з тонким; обидва стоять у питальній конструкції. svamahimni sthitasya me — той самий рефрен, що в 19.2–19.5; переклад «у власній величі» не змінено.}
\vcom{Три пари, і кожна описує простір по-своєму: відстань (далеке / близьке), напрямок (зовнішнє / внутрішнє), щільність (грубе / тонке). Разом вони мали б вичерпати всі способи бути десь — і саме тому всі три відкидаються разом. Пара sthūla / sūkṣma особлива: це не лише «велике й мале», а й «тілесне й тонке», тобто остання межа між видимим і невидимим. РЕФРЕН ГЛАВИ. 19.2–19.7 закінчуються однаково: svamahimni sthitasya me («у мене, що стоїть у власній величі»). Кожен із цих шести віршів — низка питань kva («де»), і жодне не дістає відповіді. 19.1 і 19.8 стоять поза серією: перший описує, ЯК це сталося, останній — ЩО з цього випливає для мови.}
\vpar{Українська паралель (наша інтерпретація)}{Володимир Свідзінський, «Медобір» (1975, посмертно).}{Наша паралель: у Свідзінського далеке й близьке теж втрачають міру, коли зникає той, хто міряє. Наше порівняння — за образом.}
\vfig{ilustratsii/19-6-neti.pdf}{Далеке й близьке, грубе й тонке відпали}{Схема наша, не історична ілюстрація.}

\vsec{19.7}
\vdeva{क्व मृत्युर् जीवितं वा क्व लोकाः क्वास्य क्व लौकिकम् । क्व लयः क्व समाधिर् वा स्वमहिम्नि स्थितस्य मे ॥}
\vtranslit{ква мрітюр джівітам ва ква локах квася ква лаукікам | ква лаях ква самадгір ва свамахімні стхітася ме ||}
\vperek{Де смерть, де життя, де світи, де мирське? Де розчинення, де самадгі — у мене, що стоїть у власній величі?}
\vcom{Построкове зіставлення: редакторська чернетка, не рецензовано.}
\vcom{Рядок а}
\vdeva{क्व मृत्युर् जीवितं वा क्व लोकाः क्वास्य क्व लौकिकम्}
\vtranslit{ква мрітюр джівітам ва ква локах квася ква лаукікам}
\vperek{Де смерть, де життя, де світи, де мирське?}
\vcom{mṛtyuḥ — «смерть», jīvitaṃ — «життя», lokāḥ — «світи», laukikam — «мирське». Форма kvāsya у чинному санскритському рядку граматично неясна й не має однозначного українського відповідника в чинній картці; її **не** можна непомітно тлумачити як встановлене читання.}
\vcom{Рядок б}
\vdeva{क्व लयः क्व समाधिर् वा स्वमहिम्नि स्थितस्य मे}
\vtranslit{ква лаях ква самадгір ва свамахімні стхітася ме}
\vperek{Де розчинення, де самадгі — у мене, що стоїть у власній величі?}
\vcom{layaḥ означає розчинення, samādhiḥ — зосередження або самадгі; kva ставить обидві практичні цілі під питання. Рефрен svamahimni sthitasya me зберігається; вірш не є доказом безкорисності будь-якої медитації поза власним філософським контекстом.}
\vcom{Останній із шести переліків, і він забирає найвище: laya і samādhi — те, до чого глави 5 і 13 кликали, тут теж не мають місця. Це замикає книгу на собі: навіть мета попередніх глав не стає притулком. Форма kvāsya (kva + asya, «де цього») темна: asya — родового від займенника, але без опори в реченні; можливо, це зіпсоване читання. У перекладі цей член пропущено, і це позначено, а не замасковано. РЕФРЕН ГЛАВИ. 19.2–19.7 закінчуються однаково: svamahimni sthitasya me («у мене, що стоїть у власній величі»). Кожен із цих шести віршів — низка питань kva («де»), і жодне не дістає відповіді. 19.1 і 19.8 стоять поза серією: перший описує, ЯК це сталося, останній — ЩО з цього випливає для мови.}
\vpar{Українська паралель (наша інтерпретація)}{Леся Українка, «Кассандра» (1907).}{Наша паралель: у Лесі Українки життя й смерть теж не утворюють пари, вартої вибору, для того, хто вже стоїть поза нею. Наше порівняння — типологічне.}
\vfig{ilustratsii/19-7-neti.pdf}{Смерть, життя й самадгі відпали}{Схема наша, не історична ілюстрація.}

\vsec{19.8}
\vdeva{अलं त्रिवर्गकथया योगस्य कथयाप्य् अलम् । अलं विज्ञानकथया विश्रान्तस्य ममात्मनि ॥}
\vtranslit{алам тріварґакатхая йоґася катхаяпй алам | алам віджнянакатхая вішрантася маматмані ||}
\vperek{Досить розмов про трицінність, досить і розмов про йогу, досить розмов про знання — у мене, що спочив в Атмані.}
\vcom{Построкове зіставлення: редакторська чернетка, не рецензовано.}
\vcom{Рядок а}
\vdeva{अलं त्रिवर्गकथया योगस्य कथयाप्य् अलम्}
\vtranslit{алам тріварґакатхая йоґася катхаяпй алам}
\vperek{Досить розмов про трицінність, досить і розмов про йогу,}
\vcom{Повтор alam («досить») охоплює trivarga («три групи цілей» — дгарма, користь, бажання) і розмову про йогу; kathayā стоїть в орудному відмінку. Це підсумкова риторична формула, а не наказ заборонити українській спільноті обговорення знань.}
\vcom{Рядок б}
\vdeva{अलं विज्ञानकथया विश्रान्तस्य ममात्मनि}
\vtranslit{алам віджнянакатхая вішрантася маматмані}
\vperek{досить розмов про знання — у мене, що спочив в Атмані.}
\vcom{vijñānakathayā — «розмовою про знання». viśrāntasya mamātmani означає стан мовця, що спочив в Атмані; завершення відрізняється від рефрену svamahimni sthitasya me у віршах 19.2–19.7. Третє alam підкреслює завершальний характер вислову.}
\vcom{Вірш не має рефрену попередніх шести й замінює його іншим: не svamahimni sthitasya («що стоїть у власній величі»), а viśrāntasya mamātmani («що спочив в Атмані»). Різниця точна: стояння — це позиція, спочинок — її відсутність потреби. Тричі повторене alam («досить») знищує три теми поспіль: етику (трицінність), практику (йога) і теорію (знання). Це останній вірш книги перед главой 20, і він закриває розмову — буквально: kathā означає «розмова», і саме розмову тут оголошено зайвою.}
\vpar{Українська паралель (наша інтерпретація)}{Григорій Сковорода, «Розмова пʼяти подорожніх про істинне щастя в житті».}{Наша паралель: у Сковороди розмови про знання теж урешті виявляються зайвими перед тим, що вже є. Наше порівняння — типологічне.}
\vfig{ilustratsii/19-8-flame.pdf}{Полумʼя, що втихло: досить розмов}{Схема наша, не історична ілюстрація.}


% ЗГЕНЕРОВАНО scripts/build_format_v2_chapter.py — не редагувати вручну.
% Джерело: editorial/cards/20.*.json, format-v2/rozdil-20.json, paraleli.json.

\naukchapter{20}{14}{Глава 20. Ніщо не постає}
\naukintro{Остання глава повторює форму девʼятнадцятої й доводить її до кінця: чотирнадцять разів питає де те й де те, і щоразу відповідає не переліком, а собою. Останній вірш книги обриває й самий перелік: та й навіщо багато слів — у мене не постає ніщо. Книга, що почалася питанням учня, закінчується тим, що питати більше нічого.}
\vsec{20.1}
\vdeva{क्व भूतानि क्व देहो वा । क्वेन्द्रियाणि क्व वा मनः । क्व शून्यं क्व च नैराश्यं । मत्स्वरूपे निरञ्जने ॥}
\vtranslit{ква бгутані ква дехо ва | квендріяні ква ва манах | ква шуням ква ча найрашям | матсварупе ніраньджане ||}
\vperek{Де першоелементи, де тіло, де чуття, де розум? Де порожнеча, де без-надійність — у моїй власній формі, що не знає плями?}
\vcom{Глава 20 складається майже цілком із питань kva, і це не пошук, а спосіб сказати, що переліченого немає де шукати. Перший вірш задає рамку: усі пари й категорії поставлено до однієї locative-форми — matsvarūpe, «у моїй власній формі». Зверніть увагу, що поряд із тілом і чуттями стоїть nairāśya, без-надійність: це не чеснота, а один із станів, який теж зникає. Слово ātman тут не вжите — є svarūpa, «власна форма»; ототожнювати їх можна лише як тлумачення.}
\vpar{Українська паралель (наша інтерпретація)}{Григорій Сковорода, «Наркіс. Розмова про те: Пізнай себе».}{Наша паралель: у Сковороди питання де тіло й де чуття теж обертається на питання про того, хто питає. Наше порівняння — типологічне.}
\vfig{ilustratsii/20-1-neti.pdf}{Першоелементи й розум відпали: лишається власна форма}{Схема наша, не історична ілюстрація.}

\vsec{20.2}
\vdeva{क्व शास्त्रं क्वात्मविज्ञानं । क्व वा निर्विषयं मनः । क्व तृप्तिः क्व वितृष्णत्वं । गतद्वन्द्वस्य मे सदा ॥}
\vtranslit{ква шастрам кватмавіджнянам | ква ва нірвішаям манах | ква тріптіх ква вітрішнатвам | ґатадвандвася ме сада ||}
\vperek{Де писання, де пізнання Я, де розум без об'єктів? Де вдоволеність, де свобода від спраги — у мене, в кого пари вже відпали назавжди?}
\vcom{Другий вірш ставить під питання й засоби пізнання: писання, пізнання Я, безоб'єктний розум. Це важливо для розуміння всієї глави — заперечення не обминає духовних досягнень, воно починає саме з них. tṛpti і vitṛṣṇatva названо поруч і досить несподівано: вдоволеність теж виявляється тим, чого немає де шукати. Термін toṣa, з яким tṛpti часто ототожнюють, тут не вжитий; це різні слова, і переклад їх не змішує.}
\vpar{Українська паралель (наша інтерпретація)}{Іван Вишенський, «Книжка» (кінець XVI ст.).}{Наша паралель: у Вишенського писання й пізнання теж втрачають вагу перед тим, у кого пари вже відпали. Наше порівняння — типологічне.}
\vfig{ilustratsii/20-2-neti.pdf}{Писання й пізнання Я відпали}{Схема наша, не історична ілюстрація.}

\vsec{20.3}
\vdeva{क्व विद्या क्व च वाविद्या । क्वाहं क्वेदं मम क्व वा । क्व बन्धः क्व च वा मोक्षः । स्वरूपस्य क्व रूपिता ॥}
\vtranslit{ква відя ква ча вавідя | квахам кведам мама ква ва | ква бандгах ква ча ва мокшах | сварупася ква рупіта ||}
\vperek{Де знання й де незнання, де «я», де «це», де «моє»? Де пута, де визволення — і де взагалі якась оформленість у того, хто сам є своєю формою?}
\vcom{Вірш перебирає основні опозиції будь-якого вчення — знання й незнання, я й це, моє, пута й визволення — і всі ставить під те саме питання. Порядок показовий: mokṣa з'являється не як відповідь, а як останній член пари з bandha. Остання пада замикає коло: якщо він є власна форма (svarūpa), то приписувати йому ще й оформленість (rūpitā) немає де. Саме тут свідки розходяться в одному слові — rūpitā / rūpatā, — і вибір між ними змінює відтінок, але не напрям думки.}
\vpar{Українська паралель (наша інтерпретація)}{Памфіл Юркевич, «Серце і його значення в духовному житті людини» (1860).}{Наша паралель: коли вірш питає, де я й де моє, він розводить того, хто говорить, і те, що він собі приписує. Юркевич розводить їх так само рішуче. Наше порівняння — за способом розрізняти.}
\vfig{ilustratsii/20-3-neti.pdf}{Я, моє, пута й визволення відпали}{Схема наша, не історична ілюстрація.}

\vsec{20.4}
\vdeva{क्व प्रारब्धानि कर्माणि । जीवन्मुक्तिर् अपि क्व वा । क्व तद् विदेहकैवल्यं । निर्विशेषस्य सर्वदा ॥}
\vtranslit{ква прарабдгані кармані | джіванмуктір апі ква ва | ква тад відехакайвалям | нірвішешася сарвада ||}
\vperek{Де вже запущені в дію діяння, де звільнення ще за життя? Де те безтілесне кайвальям — у того, хто завжди без ознак?}
\vcom{Тут заперечено не лише пута, а й обидві форми звільнення, відомі традиції: прижиттєве (jīvanmukti) і посмертне (videhakaivalya). Разом із ними зникає й prārabdha-karman — залишкова карма, якою зазвичай пояснюють, чому звільнений ще живе в тілі. Логіка вірша: усе це ознаки, а того, хто без ознак (nirviśeṣa), означити нічим. Зауважмо, що слово mukti в главах 16, 17 і 20 ужите лише двічі — тут і в 20.12, — і обидва рази в переліку того, чого немає.}
\vpar{Українська паралель (ПОТРЕБУЄ ПЕРЕВІРКИ)}{українська філософська проза, не визначено.}{ПОТРЕБУЄ ПЕРЕВІРКИ: питання про безтілесне звільнення найближче до української філософської прози, але конкретного твору ми не перевірили.}
\vfig{ilustratsii/20-4-neti.pdf}{Діяння й безтілесне звільнення відпали}{Схема наша, не історична ілюстрація.}

\vsec{20.5}
\vdeva{क्व कर्ता क्व च वा भोक्ता । निष्क्रियं स्फुरणं क्व वा । क्वापरोक्षं फलं वा क्व । निःस्वभावस्य मे सदा ॥}
\vtranslit{ква карта ква ча ва бгокта | нішкріям спхуранам ква ва | квапарокшам пхалам ва ква | ніхсвабгавася ме сада ||}
\vperek{Де діяч, де той, хто зазнає? Де бездіяльне сяйво? Де безпосереднє знання, де його плід — у мене, в кого немає власної природи?}
\vcom{Вірш забирає найстійкішу пару будь-якої системи — діяча й споживача — і додає до неї niṣkriya sphuraṇa, бездіяльне сяйво: те, що виявляється, не діючи. Це єдина в главі спроба назвати щось позитивне, і вона одразу поставлена під питання kva. Останнє слово — niḥsvabhāva, «без власної природи»; воно не означає порожнечі, а вказує, що жодна природа не може бути йому приписана. Термінів cit чи sākṣin тут немає, хоч sphuraṇa за змістом найближче до них.}
\vpar{Українська паралель (наша інтерпретація)}{Микола Зеров, «Камена» (1924).}{Наша паралель: у Зерова бездіяльне сяйво теж постає як повнота, а не як відсутність. Наше порівняння — типологічне.}
\vfig{ilustratsii/20-5-flame.pdf}{Бездіяльне сяйво без власної природи}{Схема наша, не історична ілюстрація.}

\vsec{20.6}
\vdeva{क्व लोकः क्व मुमुक्षुर् वा । क्व योगी ज्ञानवान् क्व वा । क्व बद्धः क्व च वा मुक्तः । स्वस्वरूपे ऽहम् अद्वये ॥}
\vtranslit{ква локах ква мумукшур ва | ква йоґі джнянаван ква ва | ква баддгах ква ча ва муктах | свасварупе ’хам адває ||}
\vperek{Де світ, де той, хто прагне звільнення, де йог чи знавець? Де зв'язаний, де звільнений — я у своїй власній формі, яка не знає двоїстості.}
\vcom{Уперше в главі з'являється ствердне речення: aham ... svasvarūpe advaye, «я — у своїй недвоїстій формі». Усе інше лишається питанням. Показово, що серед зниклих постатей є й muktaḥ, звільнений: у недвоїстості немає ні того, хто зв'язаний, ні того, хто звільнений. Слово advaya, «недвоїстий», — єдиний явний адвайтичний термін у цих трьох главах; усе решта сказано описово. Це місце варто читати уважно: «я» тут не протиставлене світові, воно вказане як те, у чому світу не знаходиться місця.}
\vpar{Українська паралель (наша інтерпретація)}{Тарас Шевченко, «Кобзар» (1840).}{Наша паралель: у Шевченка питання де звʼязаний, де звільнений теж звучить як докір мірі, а не як опис двох людей. Наше порівняння — типологічне, не про джерело.}
\vfig{ilustratsii/20-6-balance.pdf}{Дві шальки: світ і визволення — обидві порожні}{Схема наша, не історична ілюстрація.}

\vsec{20.7}
\vdeva{क्व सृष्टिः क्व च संहारः । क्व साध्यं क्व च साधनम् । क्व साधकः क्व सिद्धिर् वा । स्वस्वरूपे ऽहम् अद्वये ॥}
\vtranslit{ква сріштіх ква ча самхарах | ква садгям ква ча садганам | ква садгаках ква сіддгір ва | свасварупе ’хам адває ||}
\vperek{Де творення, де згортання, де мета, де засіб? Де той, хто практикує, де здобуття — я у своїй власній формі, яка не знає двоїстості.}
\vcom{Вірш знімає всю структуру практики: мету, засіб, виконавця й результат. Це найпряміше місце глави, де сказано, що шлях як такий не має опори. Друга пада повторює 20.6 дослівно — прийом навмисний: після кожного переліку читача повертають до того самого місця. Слово siddhi тут ужите в звичайному значенні «здобуття», а не як назва надприродних сил у йозі.}
\vpar{Українська паралель (наша інтерпретація)}{Богдан-Ігор Антонич, «Ротації» (1938).}{Наша паралель: в Антонича творення й згортання теж описані як рух самої природи, а не чиясь справа. Наше порівняння — за образом.}
\vfig{ilustratsii/20-7-motion.pdf}{Творення й згортання відпали}{Схема наша, не історична ілюстрація.}

\vsec{20.8}
\vdeva{क्व प्रमाता प्रमाणं वा । क्व प्रमेयं क्व च प्रमा । क्व किंचित् क्व न किंचिद् वा । सर्वदा विमलस्य मे ॥}
\vtranslit{ква прамата праманам ва | ква прамеям ква ча прама | ква кімчіт ква на кімчід ва | сарвада вімалася ме ||}
\vperek{Де той, хто пізнає, де засіб пізнання, де предмет, де саме пізнання? Де «щось», де «ніщо» — у мене, хто завжди без плями?}
\vcom{Тут розібрано повний апарат пізнання — чотири члени, які індійські школи сперечалися уточнювати. Текст не заперечує жодного з них окремо: він питає, де вони, коли немає місця, у якому вони могли б стояти. Пада kva kiṃcit kva na kiṃcid vā забирає й останню опору — саму пару «є / немає». Це важливо для розуміння всієї глави: заперечення не є твердженням порожнечі, воно знімає обидві можливості.}
\vpar{Українська паралель (ПОТРЕБУЄ ПЕРЕВІРКИ)}{українська релігійна філософія, не визначено.}{ПОТРЕБУЄ ПЕРЕВІРКИ: перелік пізнавача, засобу й предмета, після якого питати нічого, близький до української релігійної філософії, але конкретного місця ми не звірили.}
\vfig{ilustratsii/20-8-neti.pdf}{Пізнавач, предмет і пізнання відпали}{Схема наша, не історична ілюстрація.}

\vsec{20.9}
\vdeva{क्व विक्षेपः क्व चैकाग्र्यं । क्व निर्बोधः क्व मूढता । क्व हर्षः क्व विषादो वा । सर्वदा निष्क्रियस्य मे ॥}
\vtranslit{ква вікшепах ква чайкаґр'ям | ква нірбодгах ква мудгата | ква харшах ква вішадо ва | сарвада нішкріяся ме ||}
\vperek{Де розсіяність, де зосередженість, де несвідомість, де затьмареність? Де радість, де відчай — у мене, хто завжди бездіяльний?}
\vcom{Вірш повторює схему 18.10 — три пари, кожна заперечує обидва боки, — але з іншим закінченням: замість «у вгамованого йога» тут «у мене, хто завжди бездіяльний» (niṣkriya). niṣkriya не означає «лінивий»: це відсутність діяння як діяча, а не відсутність руху. Пари підібрані так, щоб охопити і стан уваги (розсіяність / зосередженість), і стан розуму (несвідомість / затьмареність), і стан почуття (радість / відчай). Те, що лишається поза цими парами, у вірші не названо.}
\vpar{Українська паралель (наша інтерпретація)}{Василь Стус, «Палімпсести».}{Наша паралель: у Стуса розсіяність і зосередженість теж не творять пари для того, хто вже стоїть. Наше порівняння — типологічне.}
\vfig{ilustratsii/20-9-balance.pdf}{Дві шальки: розсіяність і зосередженість}{Схема наша, не історична ілюстрація.}

\vsec{20.10}
\vdeva{क्व चैष व्यवहारो वा । क्व च सा परमार्थता । क्व सुखं क्व च वा दुःखं । निर्विमर्शस्य मे सदा ॥}
\vtranslit{ква чайша в'явахаро ва | ква ча са парамартхата | ква сукхам ква ча ва духкхам | нірвімаршася ме сада ||}
\vperek{Де ця буденна практика, де та вища істина? Де щастя, де страждання — у мене, хто вже завжди нічого не розглядає?}
\vcom{Найнебезпечніше місце для читача: тут знято пару vyavahāra / paramārtha, тобто буденний і вищий рівні істини. Зазвичай цю пару використовують, щоб пояснити, чому світ хибний, а Брахмо істинне; вірш не робить і цього кроку — обидва рівні однаково не мають де стояти. nirvimarśa, «без розгляду», — не заклик до нерозуміння, а опис: розглядати нічого, бо немає двох речей, які можна було б зіставити.}
\vpar{Українська паралель (наша інтерпретація)}{Ліна Костенко, «Берестечко» (1999).}{Наша паралель: у Костенко буденна практика й вища істина теж втрачають різницю перед обличчям того, що вже є. Наше порівняння — за будовою думки.}
\vfig{ilustratsii/20-10-balance.pdf}{Дві шальки: практика й вища істина}{Схема наша, не історична ілюстрація.}

\vsec{20.11}
\vdeva{क्व माया क्व च संसारः । क्व प्रीतिर् विरतिः क्व वा । क्व जीवः क्व च तद्ब्रह्म । सर्वदा विमलस्य मे ॥}
\vtranslit{ква мая ква ча самсарах | ква прітір віратіх ква ва | ква джівах ква ча тадбрахма | сарвада вімалася ме ||}
\vperek{Де майя, де сансара, де прихильність і де відвернення? Де жива істота, де те Брахмо — у мене, хто завжди без плями?}
\vcom{Вірш називає основні поняття будь-якого вчення про визволення — майю, сансару, прихильність і відвернення, живу істоту й Брахмо — і ставить їх в один ряд із питанням kva. Найрадикальніше тут те, що Брахмо стоїть у тому самому ряду, що й сансара: текст не протиставляє їх як хибне й істинне. Це місце легко прочитати як спростування адвайти; точніше сказати, що вірш відмовляється від її словника, не пропонуючи іншого.}
\vpar{Українська паралель (наша інтерпретація)}{Михайло Коцюбинський, «Тіні забутих предків» (1911).}{Наша паралель: у Коцюбинського мая й сансара теж постають як те, що показується й зникає, не питаючи людину. Наше порівняння — за образом.}
\vfig{ilustratsii/20-11-illusion.pdf}{Мая й сансара під серпанком}{Схема наша, не історична ілюстрація.}

\vsec{20.12}
\vdeva{क्व प्रवृत्तिर् निवृत्तिर् वा । क्व मुक्तिः क्व च बन्धनम् । कूटस्थनिर्विभागस्य । स्वस्थस्य मम सर्वदा ॥}
\vtranslit{ква правріттір нівріттір ва | ква муктіх ква ча бандганам | кутастханірвібгаґася | свастхася мама сарвада ||}
\vperek{Де діяння і де припинення, де звільнення, де пута — у мене, незмінного й неділимого, що завжди в собі?}
\vcom{Пару pravṛtti / nivṛtti вже було знято в 16.8, і тут вона повертається разом із парою mukti / bandhana. Друга половина вірша — єдина в главі, де подано ствердні означення: kūṭastha, nirvibhāga, svastha. Усі три описують не властивість, а відсутність поділу: незмінний, неділимий, у собі. Саме тому питання kva не має відповіді-місця — ділити нема на що.}
\vpar{Українська паралель (наша інтерпретація)}{Іван Франко, «Мойсей» (1905).}{Наша паралель: у Франка діяння й припинення теж не утворюють вибору для того, хто незмінний. Наше порівняння — типологічне.}
\vfig{ilustratsii/20-12-balance.pdf}{Дві шальки: діяння й припинення}{Схема наша, не історична ілюстрація.}

\vsec{20.13}
\vdeva{क्वोपदेशः क्व वा शास्त्रं । क्व शिष्यः क्व च वा गुरुः । क्व चास्ति पुरुषार्थो वा । निरुपाधेः शिवस्य मे ॥}
\vtranslit{квопадешах ква ва шастрам | ква шішях ква ча ва ґурух | ква часті пурушартхо ва | нірупадгех шівася ме ||}
\vperek{Де настановлення, де писання, де учень, де вчитель? Де й мета людського життя — у мене, хто без жодної опори, у мене, блаженного?}
\vcom{Після того як у 20.2 знято писання, тут знято всю установу передачі знання: настановлення, текст, учня, вчителя. Це послідовно: якщо немає ні знання, ні незнання, то немає й того, кому й від кого його передавати. puruṣārtha зникає останнім — разом із самою ідеєю мети. śiva тут майже напевно прикметник «блаженний, ласкавий», а не ім'я бога; вірш не дає приводу читати його як теологічне твердження, і переклад цього не робить.}
\vpar{Українська паралель (наша інтерпретація)}{Григорій Сковорода, «Алфавіт, чи Буквар миру» (діалог).}{Наша паралель: у Сковороди настановлення й писання теж відступають перед тим, хто вже не має опори. Наше порівняння — типологічне.}
\vfig{ilustratsii/20-13-neti.pdf}{Настановлення, учень і вчитель відпали}{Схема наша, не історична ілюстрація.}

\vsec{20.14}
\vdeva{क्व चास्ति क्व च वा नास्ति । क्वास्ति चैकं क्व च द्वयम् । बहुनात्र किम् उक्तेन । किंचिन् नोत्तिष्ठते मम ॥}
\vtranslit{ква часті ква ча ва насті | квасті чайкам ква ча дваям | бахунатра кім уктена | кімчін ноттіштхате мама ||}
\vperek{Де «є» і де «немає», де «одне» і де «двоє»? Та й навіщо багато слів: у мене не постає ніщо.}
\vcom{Останній вірш закриває главу відмовою від подальшої мови: після того як знято пару «є / немає» і пару «одне / двоє», говорити більше нема про що. Формула bahunātra kim uktena — усталений спосіб обірвати перелік, і вона тут доречна подвійно: як риторика і як зміст. Слово kiṃcit у заперечному реченні означає «ніщо»; це не твердження порожнечі, а констатація, що ніщо не підводиться як предмет. Цим і завершується глава: не відповіддю, а припиненням питання.}
\vpar{Українська паралель (наша інтерпретація)}{Олександр Довженко, «Зачарована Десна» (1956).}{Наша паралель: у Довженка останнє слово теж не додає нічого до сказаного, а лише закриває розмову. Наше порівняння — за будовою думки.}
\vfig{ilustratsii/20-14-space.pdf}{Ніщо не постає: останнє слово}{Схема наша, не історична ілюстрація.}



% =====================================================================
%  ПОЕТИЧНИЙ ПЕРЕКЛАД
% =====================================================================
\backmatter
\chapter*{Поетичний переклад. Наш, не канонічний}
\addcontentsline{toc}{chapter}{Поетичний переклад (наш, не канонічний)}
\markboth{Поетичний переклад}{Поетичний переклад}

\noindent{\sffamily\small
\textbf{Це наш поетичний переклад, а не канонічний і не науковий.}
Він не замінює переклад у науковій частині: обидва йдуть у книзі, і читач
бачить обидва. Тут збережено зміст кожного вірша, але змінено порядок слів,
довжину рядків і ритм. Нічого, чого немає в санскритському тексті, не додано.}

% ЗГЕНЕРОВАНО scripts/build_format_v2_chapter.py — не редагувати вручну.
% Джерело: books/ashtavakra-gita/poetychnyi/rozdil-1.md

\poemverse{1.1}{Як знання до людини приходять? \\ Як свобода в ній настає? \\ Як безсторонність вона пізнає? \\ Скажи мені, пане, — усе.}
\poemverse{1.2}{Якщо волі, сину, ти хочеш — \\ відкинь чуттєве, мов отруту. \\ Терпимість, щирість, співчуття, \\ вдоволеність, правдивість \\ шануй, мов нектар, повсякчас.}
\poemverse{1.3}{Ти — не земля, не вода, не вогонь, \\ не вітер, не небо, і не з них. \\ Заради волі пізнай Атман — \\ свідка стихій, чия природа — cit.}
\poemverse{1.4}{Якщо, від тіла відірвавшись, \\ у cit спочинеш і стоїш, — \\ уже тепер ти будеш щасливий, \\ умиротворений, без пут.}
\poemverse{1.5}{Ти не брахман і не інший стан, \\ не учень і не відречений, \\ і око тебе не бачить. \\ Ти без прив'язності, без форми, \\ свідок усього. Будь щасливий.}
\poemverse{1.6}{Дхарма й адхарма, щастя й мука — \\ усе це в умі, не твоє, всепроникний. \\ Ти не дієць і не споживач — \\ ти вільний завжди.}
\poemverse{1.7}{Ти єдиний бачитель усього \\ і завжди майже вільний. \\ Ось у чому твоя уза: \\ ти бачиш іншого бачителя, крім себе.}
\poemverse{1.8}{«Я діяч» — це самомніння, \\ укус великого чорного змія. \\ Випий нектар упевненості: \\ «я не діяч» — і будь щасливий.}
\poemverse{1.9}{«Я — єдине чисте прозріння»: \\ цим вогнем упевненості \\ спали хащу незнання — \\ будь без жалю, будь щасливий.}
\poemverse{1.10}{Там, де все це сяє — \\ вигадане, як змія в мотузці, — \\ ти є те прозріння, \\ що є радість і найвища радість. \\ Ходи щасливо.}
\poemverse{1.11}{Хто вважає себе вільним — вільний; \\ хто вважає себе зв'язаним — зв'язаний. \\ Що тут іще сказати? \\ Яке переконання — така й доля.}
\poemverse{1.12}{Без прив'язності, без жадання, втихлий — \\ через оману він видається тим, хто в сансарі.}
\poemverse{1.13}{Осягай Атман — незмінне, \\ недвоїсте усвідомлення. \\ Відкинувши оману «я — лише відблиск», \\ відкинь і стан зовнішній, і внутрішній.}
\poemverse{1.14}{Ти довго зв'язаний, синку, \\ мотузкою приписування тіла. \\ Скажи: «я — усвідомлення» — \\ і відсічи те мечем знання. \\ Будь щасливий.}
\poemverse{1.15}{Ти без прив'язності, без діяння, \\ самосяйний, незаплямлений. \\ Ось саме це твої пута: \\ те, що ти вправляєшся в самадгі.}
\poemverse{1.16}{Цей всесвіт просякнутий тобою \\ і справді нанизаний на тебе. \\ Ти маєш природу чистого й просвітленого — \\ не йди до дріб'язковості ума.}
\poemverse{1.17}{Без очікування, без видозмін, \\ без тягаря, з прохолодним нутром, \\ із незміренною думкою, незбурений — \\ будь тим, чия схильність є сама лише свідомість.}
\poemverse{1.18}{Знай: те, що має образ, — несправжнє; \\ а те, що без образу — непохитне. \\ Цим настановленням про суть \\ немає народження знову.}
\poemverse{1.19}{Як образ у дзеркалі — \\ він усередині й навколо нього; \\ так само і в цьому тілі — \\ всередині й навколо верховний володар.}
\poemverse{1.20}{Як один усюдисущий простір — \\ зовні й усередині глечика, \\ так повсякчас і без перерви \\ брахман у всій сукупності істот.}

% ЗГЕНЕРОВАНО scripts/build_format_v2_chapter.py — не редагувати вручну.
% Джерело: books/ashtavakra-gita/poetychnyi/rozdil-2.md

\poemverse{2.1}{Я — усвідомлення без плями, \\ спокій, що вище за природу. \\ І весь цей час \\ дурила мене тільки вона — \\ омана, і більше ніщо.}
\poemverse{2.2}{Я освітлюю це тіло — \\ і тим освітлюю світ. \\ Тож світ — мій. \\ Або ж — аніщо.}
\poemverse{2.3}{Скинув світ разом із тілом — \\ і, самим лише вмінням, \\ невідомо звідки, \\ бачу те, що вище за мене.}
\poemverse{2.4}{Хвиля, піна, бульбашка — \\ усе це вода, і тільки. \\ Так і світ, що вийшов із мене: \\ не віддільний, як хвиля від моря.}
\poemverse{2.5}{Тканину дослідиш — \\ лишиться нитка. \\ Дослідиш світ — \\ лишиться Я. \\ І глибше немає нічого.}
\poemverse{2.6}{Цукор уже в тростині, \\ уже в її соку. \\ Так і світ — уже в мені, \\ просякнутий мною до дна.}
\poemverse{2.7}{Поки не знаєш Я — світ сяє. \\ Знаєш — і сяйво гасне. \\ Поки не знаєш мотузки — сяє змія. \\ Знаєш — і змія зникає.}
\poemverse{2.8}{Світло — не позичене: \\ світло і є я. \\ Коли сяє світ, \\ сяє тільки моє сяйво.}
\poemverse{2.9}{Світ, вигаданий незнанням, \\ сяє в мені — \\ срібло в перламутрі, \\ змій у мотузці, \\ вода в мареві променя.}
\poemverse{2.10}{Світ вийшов із мене — \\ у мені й розчиниться: \\ горщик у глині, \\ хвиля у воді, \\ браслет у золоті.}
\poemverse{2.11}{Поклін мені самому — \\ тому, кого не вбити, \\ хто стоїть, \\ коли гине світ \\ від Брахми до стеблини трави.}
\poemverse{2.12}{Поклін мені самому! \\ Я один. \\ Маю тіло — і не йду, не приходжу: \\ охопив усе — і стою.}
\poemverse{2.13}{Поклін мені самому — \\ тому, хто не торкається тілом, \\ а все тримає довго. \\ І немає нікого, \\ хто зрівнявся б зі мною в цьому.}
\poemverse{2.14}{Поклін мені самому! \\ У мене немає нічого — \\ або ж у мене є все, \\ до чого дістає слово й розум.}
\poemverse{2.15}{Пізнавач, пізнаване, знання — \\ трійця, якої нема. \\ Те, у чому вона сяє з незнання, \\ і є я — \\ незаплямований.}
\poemverse{2.16}{Корінь болю — двоїстість. \\ Інших ліків немає. \\ Все, що видно, — хибне. \\ Я один: смак свідомості, без плями.}
\poemverse{2.17}{Я — саме усвідомлення. \\ Незнанням я домалював собі межі. \\ Хто так думає безупинно, \\ стоїть у не-розрізненні. \\ Це моє стояння.}
\poemverse{2.18}{Світ наче в мені — \\ а в мені не вміщується. \\ Ні пут у мене, ні визволення: \\ омана вщухла, \\ бо втратила опору.}
\poemverse{2.19}{Світ разом із тілом — ніщо, \\ так установлено. \\ А я — тільки чиста свідомість. \\ Що ж іще домислювати?}
\poemverse{2.20}{Тіло, небо, пекло, \\ пута, визволення, страх — \\ усе це домалюване. \\ Що мені робити, \\ коли я — свідомість?}
\poemverse{2.21}{Навіть у натовпі, \\ де не бачу двоїстості, \\ все стає лісом. \\ То де ж мені шукати втіхи?}
\poemverse{2.22}{Я не тіло, і тіло не моє. \\ Я не живчик — я свідомість. \\ Ось де був вузол: \\ спрага жити.}
\poemverse{2.23}{Різнобарвні хвилі світів \\ раптом устали — \\ то в мені, безмежному океані, \\ здіймається вітер розуму.}
\poemverse{2.24}{Коли в безмежному океані \\ вітер розуму вщухає — \\ торговцеві-живчику лихо: \\ човен світу тоне.}
\poemverse{2.25}{У безмежному океані — диво: \\ хвилі живих істот \\ устають, бʼються, граються \\ й входять самі собою, \\ за власною природою.}

% ЗГЕНЕРОВАНО scripts/build_format_v2_chapter.py — не редагувати вручну.
% Джерело: books/ashtavakra-gita/poetychnyi/rozdil-3.md

\poemverse{3.1}{Пізнав ти Атман — незнищенний, єдиний. \\ То де тобі втіха з набутків, \\ атмазнавцю, мудрий?}
\poemverse{3.2}{Не знаєш Атмана — \\ і прикипаєш до речей у полі омани. \\ Не знаєш перламутру — \\ і жадібно хапаєш срібло, якого нема.}
\poemverse{3.3}{Пізнав: я — те, \\ у чому сяє весь цей світ, \\ як хвилі в морі. \\ Чого ж ти біжиш, наче вбогий?}
\poemverse{3.4}{Ти чув про Атман — \\ чисту свідомість, надзвичайно прекрасну. \\ Але хто надто привʼязаний до того, що під рукою, \\ той набирає бруду.}
\poemverse{3.5}{Мудрець знає Атман у всіх істотах \\ і всі істоти — в Атмані. \\ А своє досі тримається.}
\poemverse{3.6}{Став у вищій недвоїстості, \\ стоїш навіть у самій справі звільнення — \\ а бажання все ж веде тебе. \\ Безпорадний, мов у дитячій грі.}
\poemverse{3.7}{Зрозумів: знання було поганим другом. \\ І все ж, надзвичайно слабкий, \\ він прагнув би бажання \\ аж до кінця часу.}
\poemverse{3.8}{Хто відвернувся від цього світу й того, \\ хто розрізняє вічне й невічне, \\ хто хоче звільнення, — \\ той боїться саме звільнення.}
\poemverse{3.9}{Голублять його чи мучать — \\ він бачить тільки Атман. \\ Тому не радіє \\ і не гнівається.}
\poemverse{3.10}{Своє тіло, що рухається, \\ він бачить як тіло чуже. \\ То як його збентежить хвала \\ або ганьба?}
\poemverse{3.11}{Бачить: весь цей світ — тільки мая. \\ Охота до нього згасла. \\ То як злякається він, \\ коли смерть уже поруч, — \\ той, чий розум стійкий?}
\poemverse{3.12}{Розум без прагнень. \\ Великий дух — навіть у безнадії. \\ Вдоволений знанням Атмана. \\ З ким його й порівнювати?}
\poemverse{3.13}{Знає самою природою: \\ все, що видно, — аніщо. \\ Взяти це, відкинути те — \\ то що ж він бачить, \\ чий розум стійкий?}
\poemverse{3.14}{Каламуть усередині відкинуто. \\ Немає двоїстості, немає сподівань. \\ Все, що випадає само собою, \\ не стає ні болем, ні втіхою.}

% ЗГЕНЕРОВАНО scripts/build_format_v2_chapter.py — не редагувати вручну.
% Джерело: books/ashtavakra-gita/poetychnyi/rozdil-4.md

\poemverse{4.1}{О! Мудрому, що знає Атман \\ і грається грою досвідів, \\ немає рівні з дурнями, \\ що тягнуть за собою сансару.}
\poemverse{4.2}{О диво: \\ до того, чого прагнуть усі вбогі боги, \\ починаючи з Шакри, \\ йог, що вже стоїть у цьому, \\ не відчуває захоплення.}
\poemverse{4.3}{Того, хто осягнув Істину, \\ не зачіпають ні заслуги, ні провини. \\ Небо не зливається з димом, \\ хоч на око вони ніби поєднані.}
\poemverse{4.4}{Хто стримає того, \\ хто пізнав увесь світ як єдиний Атман \\ і живе так, \\ як складаються обставини?}
\poemverse{4.5}{Серед чотирьох родів істот, \\ від Брахми до стеблини трави, \\ лише той, хто пізнав, \\ спроможний полишити \\ і бажання, і небажання.}
\poemverse{4.6}{Хтось знає Атман — \\ недвоїстий, володар світу. \\ Що знає, те й робить. \\ І страху в нього немає ніде.}

% ЗГЕНЕРОВАНО scripts/build_format_v2_chapter.py — не редагувати вручну.
% Джерело: books/ashtavakra-gita/poetychnyi/rozdil-5.md

\poemverse{5.1}{Нема в тебе нічого, до чого б ти приріс, — \\ чого ж ти, чистий, прагнеш відкидати? \\ Розчини скупчення — і так само йди до розчинення.}
\poemverse{5.2}{Світ устає з тебе, мов бульбашка з води. \\ Знай: Атман один — і так само йди до розчинення.}
\poemverse{5.3}{Хоч світ і перед очима — він не річ у тобі, незаплямований: \\ змій у мотузці. І так само йди до розчинення.}
\poemverse{5.4}{Рівний у смутку й радості, повний, рівний у надії і в безнадії, \\ рівний у житті й смерті — так само йди до розчинення.}

% ЗГЕНЕРОВАНО scripts/build_format_v2_chapter.py — не редагувати вручну.
% Джерело: books/ashtavakra-gita/poetychnyi/rozdil-6.md

\poemverse{6.1}{Я — простір безкраїй, а світ — як глечик глиняний. \\ У цьому знанні — ні відкидання, ні взяття, ні розчинення.}
\poemverse{6.2}{Я — океан великий, а світ — лиш хвиля на воді. \\ У цьому знанні — ні відкидання, ні взяття, ні розчинення.}
\poemverse{6.3}{Я — перламутр, а уявний світ — мов срібло. \\ У цьому знанні — ні відкидання, ні взяття, ні розчинення.}
\poemverse{6.4}{Я — в кожній істоті, а кожна істота — в мені. \\ У цьому знанні — ні відкидання, ні взяття, ні розчинення.}

% ЗГЕНЕРОВАНО scripts/build_format_v2_chapter.py — не редагувати вручну.
% Джерело: books/ashtavakra-gita/poetychnyi/rozdil-7.md

\poemverse{7.1}{У мені, безмежному океані, \\ човен світу блукає \\ туди й сюди — \\ вітром власного нутра. \\ І в мене немає нетерпимості.}
\poemverse{7.2}{У мені, безмежному океані \\ хвиля світу — за власною природою. \\ Хай піднімається, хай спадає: \\ у мене немає \\ ні приросту, ні втрати.}
\poemverse{7.3}{У мені, безмежному океані \\ те, що зветься світом, — домислювання. \\ Вельми втихлий, \\ без форми — \\ саме в цьому я стою.}
\poemverse{7.4}{Я не в станах буття, \\ і стан — не там, \\ у безмежному, незаплямленому. \\ Так: непривʼязаний, без спраги, втихлий — \\ саме в цьому я стою.}
\poemverse{7.5}{О диво: \\ я — саме лише свідомість, \\ а світ подібний до чару. \\ Тож звідки й де в мені \\ домислювання про те, \\ що відкинути і що прийняти?}

% ЗГЕНЕРОВАНО scripts/build_format_v2_chapter.py — не редагувати вручну.
% Джерело: books/ashtavakra-gita/poetychnyi/rozdil-8.md

\poemverse{8.1}{Полон — це коли розум \\ чогось бажає, \\ сумує, \\ щось відпускає або бере, \\ чомусь радіє \\ чи гнівається.}
\poemverse{8.2}{Визволення — це коли розум \\ не бажає, \\ не сумує, \\ не відпускає й не бере, \\ не радіє \\ й не гнівається.}
\poemverse{8.3}{Полон — коли розум привʼязаний \\ до якихось поглядів. \\ Визволення — коли він \\ не привʼязаний до жодних.}
\poemverse{8.4}{Немає я — визволення. \\ Є я — полон. \\ Знаючи це, легко: \\ нічого не бери \\ й нічого не відпускай.}

% ЗГЕНЕРОВАНО scripts/build_format_v2_chapter.py — не редагувати вручну.
% Джерело: books/ashtavakra-gita/poetychnyi/rozdil-9.md

\poemverse{9.1}{Зроблене й незроблене, \\ пари протилежностей — \\ коли вони втихли, \\ то чиї вони? \\ Знаючи це, будь тут, у відреченні, \\ схильним до відкидання — \\ і без обітниць.}
\poemverse{9.2}{У якогось щасливого, любий, \\ від споглядання того, \\ як поводиться світ, \\ бажання жити, \\ голод і жага пізнавати \\ дійшли до втишення.}
\poemverse{9.3}{Усе це нестале, \\ зіпсоване троїстою мукою, \\ без суті, \\ осуджене, гідне відкинення. \\ Так вирішивши, людина втихає.}
\poemverse{9.4}{Який же той час, вік або що, \\ коли в людей немає пар протилежностей? \\ Знехтувавши їх, \\ хто живе як випало — \\ той досяг би здійснення.}
\poemverse{9.5}{Побачивши, скільки різних поглядів \\ мають великі мудреці, \\ праведні йоги, — \\ хто, дійшовши відречення, \\ не втихне?}
\poemverse{9.6}{Хто, пізнавши образ свідомості, — \\ чи не вчитель? — \\ той через відречення, рівність і зʼєднаність \\ переправляє з сансари.}
\poemverse{9.7}{Дивись на видозміни сущого \\ як на самі лише елементи, \\ як є насправді, — \\ і тієї ж миті, вільний від пут, \\ станеш тим, хто стоїть у власній формі.}
\poemverse{9.8}{Сансара — це саме усталені схильності, так. \\ Відпусти їх усі. \\ Їх відкидання — через відкидання схильностей; \\ стояння в собі — нині, \\ так чи інакше.}

% ЗГЕНЕРОВАНО scripts/build_format_v2_chapter.py — не редагувати вручну.
% Джерело: books/ashtavakra-gita/poetychnyi/rozdil-10.md

\poemverse{10.1}{Покинув вороже бажання, \\ покинув користь, сповнену лиха, — \\ навіть дхарму, що є причиною обох. \\ Усюди зневажай.}
\poemverse{10.2}{Дивись на друзів, поле, гроші, \\ дім, дружину, спадок \\ і решту багатств — \\ як на сон і чари, \\ на три дні або пʼять.}
\poemverse{10.3}{Де б не виникла спрага — \\ знай: саме там і сансара. \\ Опираючись на зрілу безсторонність, \\ стань тим, у кого спрага минула, — \\ і будь щасливий.}
\poemverse{10.4}{Полон має природу самої лише спраги; \\ її знищення називають визволенням. \\ Від самої лише непривʼязаності до буття \\ приходить вдоволеність досягнутим — \\ знову й знову.}
\poemverse{10.5}{Ти один — свідомий, чистий; \\ світ же — неживий і несправжній. \\ Навіть незнання — аніщо. \\ То яке в тебе, при всьому тому, \\ бажання пізнавати?}
\poemverse{10.6}{Царство, сини, дружини, тіла й утіхи — \\ усе це вже гинуло для тебе, \\ привʼязаного, — \\ з народження в народження.}
\poemverse{10.7}{Досить уже й багатства, і бажання, \\ і навіть добрих учинків: \\ у цьому лісі сансари \\ твій розум так і не знав відпочинку.}
\poemverse{10.8}{Скільки ж за ці народження \\ зроблено тілом, розумом і мовою \\ діл, що дають біль і втому! \\ Тож нехай це припиниться хоч тепер.}

% ЗГЕНЕРОВАНО scripts/build_format_v2_chapter.py — не редагувати вручну.
% Джерело: books/ashtavakra-gita/poetychnyi/rozdil-11.md

\poemverse{11.1}{Видозміна буття й небуття — \\ за власною природою. \\ Хто має таке переконання, \\ той без видозмін, \\ з минулими турбаціями, \\ втихає легко.}
\poemverse{11.2}{Творець усього — Ішвара, \\ і немає тут іншого. \\ Хто має таке переконання, \\ той із розчиненими всередині сподіваннями, \\ втихлий, \\ ніде не застрягає.}
\poemverse{11.3}{Нещастя й достаток у свій час — \\ саме від долі. \\ Хто має таке переконання, \\ той вдоволений, \\ із чуттями в порядку, \\ повсякчас не бажає й не сумує.}
\poemverse{11.4}{Щастя й страждання, \\ народження й смерть — \\ саме від долі. \\ Хто має таке переконання, \\ той, не бачачи нічого досяжного, \\ без натуги, навіть діючи, \\ не забруднюється.}
\poemverse{11.5}{Страждання народжується з думання, \\ не інакше. \\ Хто має таке переконання, \\ той, позбавлений його, \\ щасливий, втихлий, \\ скрізь із опалою спрагою.}
\poemverse{11.6}{Я не тіло, і тіло не моє; \\ я — усвідомлення. \\ Хто має таке переконання, \\ той, наче досягши окремішності, \\ не памʼятає незробленого як зробленого.}
\poemverse{11.7}{Від Брахми до стеблини — \\ це саме я. \\ Хто має таке переконання, \\ той без домислювання, \\ чистий, втихлий, \\ вдоволений і досягнутим, і недосягнутим.}
\poemverse{11.8}{Цей світ, \\ різноманітно дивовижний, — \\ аніщо. \\ Хто має таке переконання, \\ той без усталених схильностей, \\ саме лише сяйнування, \\ втихає, наче аніщо.}

% ЗГЕНЕРОВАНО scripts/build_format_v2_chapter.py — не редагувати вручну.
% Джерело: books/ashtavakra-gita/poetychnyi/rozdil-12.md

\poemverse{12.1}{Спершу нездатний до тілесних справ, \\ потім нездатний до поширення мови, \\ далі нездатний до думання — \\ тому саме в цьому я стою.}
\poemverse{12.2}{Через відсутність прихильності \\ до звуку й подібного \\ та через невидимість його для самого себе — \\ серце, для якого розсіяність і зосередженість одне: \\ саме в цьому я стою.}
\poemverse{12.3}{Коли розсіяність походить \\ від прагнення самадгі та подібного, \\ то й звичайна діяльність — задля самадгі. \\ Так побачивши це правило, \\ саме в цьому я стою.}
\poemverse{12.4}{Через відсутність того, що відкидати й приймати, \\ а отже — через відсутність радості й смутку, \\ нині, о Брахмане, \\ саме в цьому я стою.}
\poemverse{12.5}{Медитація зі стадією життя і без неї, \\ прийняття й відкидання, \\ зроблені умом, — \\ усе це домислювання. \\ Побачивши це в собі, \\ через саме це я так стою.}
\poemverse{12.6}{Виконання обрядів через незнання \\ і припинення їх — однаково. \\ Зрозумівши правильно цю суть, \\ я так стою.}
\poemverse{12.7}{Незбагненне, навіть коли його осмислюють, \\ набирає форми думки. \\ Тому, відкинувши те плекання, \\ я так стою.}
\poemverse{12.8}{У кого це так зроблено — \\ той стає тим, чию мету здійснено; \\ у кого така природа — \\ той стає тим, чию мету здійснено.}

% ЗГЕНЕРОВАНО scripts/build_format_v2_chapter.py — не редагувати вручну.
% Джерело: books/ashtavakra-gita/poetychnyi/rozdil-13.md

\poemverse{13.1}{Стан неволодіння — це бути в собі; \\ навіть у жебрацькому поясі він рідкісний. \\ Відкинувши віддавання й приймання, \\ я сиджу, як мені завгодно.}
\poemverse{13.2}{Десь втомлюється тіло, \\ десь втомлюється язик, \\ десь розум. \\ Відкинувши це, \\ я стою в меті людини, щасливо.}
\poemverse{13.3}{Подумавши по-справжньому: \\ зробленого не буде нічого — \\ що б не випало зробити, \\ зробивши те, \\ я сиджу, як мені завгодно.}
\poemverse{13.4}{У йога, що перебуває в тілі, \\ є стани, до яких його змушують \\ дія й недіяння. \\ Та через відсутність зʼєднання й розʼєднання \\ я сиджу, як мені завгодно.}
\poemverse{13.5}{Ні користі, ні шкоди немає мені \\ від стояння, ходіння чи лежання. \\ Тому, стоячи, ходячи чи сплячи, \\ я сиджу, як мені завгодно.}
\poemverse{13.6}{У сплячого немає мені втрати, \\ як і немає досягнення \\ в того, хто докладає зусиль. \\ Відкинувши загибель і піднесення, \\ я сиджу, як мені завгодно.}
\poemverse{13.7}{Побачивши не раз у всіх станах \\ те правило, що має вигляд задоволення, \\ відкинувши добре й погане, \\ я сиджу, як мені завгодно.}

% ЗГЕНЕРОВАНО scripts/build_format_v2_chapter.py — не редагувати вручну.
% Джерело: books/ashtavakra-gita/poetychnyi/rozdil-14.md

\poemverse{14.1}{Хто за природою має порожній розум — \\ навіть якщо з недбальства плекає в собі стани, — \\ той наче заснув і прокинувся: \\ блукання його вичерпане.}
\poemverse{14.2}{Де багатства, де друзі, \\ де мої слуги — чуттєві предмети? \\ Де писання, де знання — \\ коли в мене розтанула спрага?}
\poemverse{14.3}{Коли пізнано пурушу-свідка, \\ вищого Атмана й Ішвару — \\ а також безнадію, узу й визволення — \\ у мене немає турботи про визволення.}
\poemverse{14.4}{У того, хто всередині порожній від розрізнень, \\ а зовні ходить, як сам захоче, \\ стани бувають усякі, наче в божевільного; \\ і знають його тільки такі самі, як він.}

% ЗГЕНЕРОВАНО scripts/build_format_v2_chapter.py — не редагувати вручну.
% Джерело: books/ashtavakra-gita/poetychnyi/rozdil-15.md

\poemverse{15.1}{Хто має ясний розум, \\ той здійснює мету настановою як є; \\ а допитливий, хоч би все життя шукав, \\ інший у цьому збивається з пуття.}
\poemverse{15.2}{Визволення — коли чуттєве втрачає свій смак; \\ неволя — коли воно й далі вабить. \\ Ось і все розуміння. \\ Чини, як бажаєш.}
\poemverse{15.3}{Це пізнання суті \\ робить красномовних, тямущих і вельми завзятих \\ німими, тупими й лінивими — \\ тому голодні його відкидають.}
\poemverse{15.4}{Ти не тіло, і тіло не твоє; \\ і не ти — той, хто споживає чи чинить. \\ Ти маєш природу свідомості, \\ повсякчас свідок, без залежності — \\ живи легко.}
\poemverse{15.5}{Пристрасть і відраза — властивості ума, \\ і ум ніколи не твій. \\ Ти без домислювання, твоя суть — усвідомлення, \\ ти без видозмін — \\ живи легко.}
\poemverse{15.6}{Пізнавши себе у всіх істотах \\ і всі істоти — у собі, \\ ти без я-роблю, без моє — \\ будь щасливий.}
\poemverse{15.7}{Те, в чому цей всесвіт постає, \\ як хвилі в морі, — \\ це саме ти, немає сумніву, \\ о втілення свідомості; \\ будь без трепету.}
\poemverse{15.8}{Повір, любий, повір; \\ не роби тут затьмарення, друже. \\ Твоя суть — знання; \\ бгаґаван ātman — ти, понад природою.}
\poemverse{15.9}{Тіло, обгорнуте гунами, \\ стоїть, приходить і йде; \\ ātman не йде і не приходить — \\ чого ти за ним жалкуєш?}
\poemverse{15.10}{Нехай тіло стоїть до кінця світового періоду — \\ або нехай іде хоч сьогодні. \\ Де приріст і де втрата для тебе, \\ чия природа — сама лише свідомість?}
\poemverse{15.11}{У тобі, безмежному великому океані, \\ хвиля всесвіту за власною природою \\ нехай піднімається або спадає — \\ немає в тебе ні приросту, ні втрати.}
\poemverse{15.12}{Любий, твоя природа — сама лише свідомість; \\ цей світ не відмінний від тебе. \\ Тому — для кого, як і де \\ вигадувати те, що відкидати й приймати?}
\poemverse{15.13}{У тобі — одному, незмінному, втихлому, \\ незаплямованому просторі свідомості — \\ звідки народження, звідки діяння, \\ звідки ж я-роблю?}
\poemverse{15.14}{Те, що ти бачиш, — \\ там сяєш один саме ти. \\ Чи сяє окремо від золота \\ браслет, наруч або ножна прикраса?}
\poemverse{15.15}{Це — він, це я; це — не я — \\ відкинь такий поділ. \\ Вирішивши, що все — ātman, \\ будь без побудов і щасливий.}
\poemverse{15.16}{Всесвіт — саме через твоє незнання; \\ у справжньому сенсі ти один. \\ Немає іншого, ніж ти — \\ ні того, хто в сансарі, ні того, хто поза нею.}
\poemverse{15.17}{Цей всесвіт — саме лише блукання, аніщо — \\ хто має таке переконання, \\ той без усталених схильностей, \\ саме лише сяйнування, \\ втихає, наче аніщо.}
\poemverse{15.18}{Один-єдиний в океані буття — \\ був, є, буде. \\ Немає в тебе ні пут, ні визволення; \\ той, чию справу зроблено, — \\ живи легко.}
\poemverse{15.19}{Не збурюй ум наміром і домислюванням, \\ о сповнений свідомості. \\ Втихнувши, стій легко у собі — \\ втіленні блаженства.}
\poemverse{15.20}{Відкинь саму медитацію, \\ скрізь не тримай нічого в серці. \\ Ти — ātman, ти вже вільний: \\ що ти робитимеш, далі розмірковуючи?}

% ЗГЕНЕРОВАНО scripts/build_format_v2_chapter.py — не редагувати вручну.
% Джерело: books/ashtavakra-gita/poetychnyi/rozdil-16.md

\poemverse{16.1}{Розповідай чи слухай, любий, \\ хоч би скільки різних наук, — \\ спокою тобі не буде, \\ поки не забудеш усе.}
\poemverse{16.2}{Насолоджуйся, дій або зосереджуйся, \\ обізнаний, — \\ та по-справжньому смакуватиме тобі \\ лише тоді, \\ коли розум відкине всі надії.}
\poemverse{16.3}{Від натуги всі нещасні, \\ і ніхто з них цього не бачить. \\ Блаженний же досягає вгамування \\ самим лише цим настановленням.}
\poemverse{16.4}{Кому важко навіть повіки звести — \\ той перший серед лінивих: \\ немає щастя ні йому, \\ ні комусь іншому такому.}
\poemverse{16.5}{Коли розум звільниться від пари \\ це я зробив — а це ні, \\ він стане байдужим \\ і до дгарми, і до артги, \\ і до ками, і до мокші.}
\poemverse{16.6}{Відвернений — це той, хто ненавидить чуттєве; \\ пристрасний — той, хто до нього жадібний. \\ А хто не хапає і не відкидає — \\ той уже ні відвернений, ні пристрасний.}
\poemverse{16.7}{Розрізнення відкинути це, взяти те — \\ лише пагін на дереві сансари; \\ він росте, доки жива жадоба, \\ а жадоба — осідок життя без роздуму.}
\poemverse{16.8}{Коли дієш — народжується потяг, \\ коли спиняєшся — відраза. \\ Мудрий, вільний від цієї пари, \\ стоїть так, як стоїть дитина.}
\poemverse{16.9}{Пристрасний хоче втекти з сансари, \\ бо хоче позбутися страждання. \\ А в кого пристрасті вже немає, \\ той не страждає — \\ і навіть у сансарі не мучиться.}
\poemverse{16.10}{Якщо людина претендує навіть на звільнення \\ і водночас відчуває тіло як своє, \\ вона не знавець і не йог — \\ вона просто має свою частку страждання.}
\poemverse{16.11}{Навіть якби твоїм наставником був сам Шіва, \\ Вішну чи Брахма, — \\ спокою тобі не буде, \\ поки не забудеш усе.}

% ЗГЕНЕРОВАНО scripts/build_format_v2_chapter.py — не редагувати вручну.
% Джерело: books/ashtavakra-gita/poetychnyi/rozdil-17.md

\poemverse{17.1}{Хто вдоволений, \\ чиї чуття прозорі \\ і хто завжди сам — \\ той і справді радіє: \\ він уже має плід знання \\ і плід вправ у йозі.}
\poemverse{17.2}{Ніколи в цьому світі не мучиться \\ той, хто знає істину: \\ о, тим одним сповнене все — \\ усе коло світів.}
\poemverse{17.3}{Жодні чуттєві речі не тішать того, \\ кому втіха — у ньому самому, \\ як не тішить слона, що любить пагони салакі, \\ гірке листя німби.}
\poemverse{17.4}{Хто вже зазнаних утіх не тягне за собою, \\ як просякнутий їхнім запахом, \\ і не прагне незазнаних — \\ така людина в цьому світі рідкість.}
\poemverse{17.5}{У цьому світі видно і тих, хто голодний до нього, \\ і тих, хто прагне звільнення. \\ Але хто не прагне ні втіх, ні звільнення — \\ такий великий і рідкісний.}
\poemverse{17.6}{Для великодушної людини \\ немає відкинути й узяти — \\ ні щодо дгарми, артги, ками й мокші, \\ ні щодо життя й смерті.}
\poemverse{17.7}{Він не бажає, щоб світ загинув, \\ і не має відрази до того, що світ триває. \\ Тому блаженний сидить собі спокійно, \\ як завгодно, — \\ і живе з того, що має.}
\poemverse{17.8}{Цим знанням я досяг мети — \\ і з цим, уже розтанувши розумом, \\ він живе собі як завгодно: \\ дивиться, слухає, торкається, нюхає, їсть.}
\poemverse{17.9}{Порожній погляд, \\ марні рухи, \\ чуття ніби понівечені; \\ ні прагнення, ні відвернення — \\ бо океан сансари в ньому висох.}
\poemverse{17.10}{Він не пильнує і не спить, \\ не розплющує й не заплющує очей. \\ О диво: розум вільної людини \\ перебуває в якомусь зовсім іншому стані.}
\poemverse{17.11}{Усюди він у собі, \\ усюди його нутро чисте; \\ вільний від усіх звичних слідів, \\ звільнений сяє всюди.}
\poemverse{17.12}{Дивиться, слухає, торкається, \\ нюхає, їсть, бере, \\ говорить, ходить — \\ і вільний від бажаного й небажаного. \\ Він справді звільнений, цей великий душею.}
\poemverse{17.13}{Не гудить і не хвалить, \\ не тішиться й не гнівається, \\ не дає й не бере — \\ звільнений усюди без смаку до всього.}
\poemverse{17.14}{Побачивши жінку, повну пристрасті, \\ або смерть, що вже насувається, \\ він не збентежиться: \\ він у собі, він справді звільнений, \\ цей великий душею.}
\poemverse{17.15}{У щасті й у горі, \\ в чоловікові й у жінці, \\ в удачі й у біді — \\ для стійкого, хто всюди бачить однаково, \\ різниці немає.}
\poemverse{17.16}{Ні насильства, ні жалю, \\ ні зарозумілості, ні приниженості, \\ ні здивування, ні збурення — \\ немає цього в людині, \\ в якої скінчилося блукання.}
\poemverse{17.17}{Звільнений не ненавидить чуттєвих речей \\ і не жадібний до них; \\ його розум ні до чого не привʼязаний, \\ і він бере те, що прийшло, \\ і те, що не прийшло.}
\poemverse{17.18}{Розрізнень зібраний я чи ні, \\ корисно це чи шкідливо \\ порожній розум уже не знає: \\ він стоїть ніби в самотності.}
\poemverse{17.19}{Без моє і без я роблю, \\ з певністю, що немає нічого, — \\ усі його надії розтанули зсередини, \\ і він навіть діючи не діє.}
\poemverse{17.20}{Він вільний і від ясності розуму, \\ і від затьмарення, \\ і від сну, і від заціпеніння — \\ і досяг якогось невідомого стану, \\ цей чоловік із розталим розумом.}

% ЗГЕНЕРОВАНО scripts/build_format_v2_chapter.py — не редагувати вручну.
% Джерело: books/ashtavakra-gita/poetychnyi/rozdil-18.md

\poemverse{18.1}{Той, з чиїм пробудженням \\ усвідомлення омана стає наче сон, — \\ тому, чия єдина природа — щастя, \\ уклін: втихлому, сяйному.}
\poemverse{18.2}{Здобувши всі блага, \\ людина отримує рясні насолоди — \\ але без повного зречення всього \\ вона не стане щасливою.}
\poemverse{18.3}{У того, чий внутрішній Атман спалений \\ полумʼям сонця страждання від належного, — \\ звідки візьметься щастя \\ без суті потоку нектару вгамування?}
\poemverse{18.4}{Це буття — тільки уявлення, \\ у вищому сенсі — ніщо. \\ І немає небуття для тих речей, \\ що розрізняють буття й небуття.}
\poemverse{18.5}{Стопа Атмана не далеко \\ і не досягається стисненням — \\ вона вже досягнута: \\ без розрізнень, без зусилля, \\ без змін, без плями.}
\poemverse{18.6}{Коли припиняється сама лише завороженість, \\ це вже й є прийняття власної природи: \\ ті, чий зір без покрову, \\ сяють, вільні від смутку.}
\poemverse{18.7}{Знаючи, що все разом — тільки побудова уяви, \\ а Атман вільний і споконвічний, — \\ чого ж мудрий вправляється, як дитина?}
\poemverse{18.8}{Встановивши, що Атман — це Брахман, \\ а буття й небуття — вигадані, — \\ що безбажанний знає, \\ що говорить, що робить?}
\poemverse{18.9}{У йога, що став безмовним, \\ вичерпані розрізнення це він, це я, це не я, — \\ бо він встановив: усе — Атман.}
\poemverse{18.10}{У вгамованого йога немає \\ ні розсіяності, ні зосередженості, \\ ні надмірного знання, ні затьмареності, \\ ні щастя, ні страждання.}
\poemverse{18.11}{У власному царстві чи в житті з подаяння, \\ у здобутку й нездобутку, \\ серед людей чи в лісі — \\ для йогина, чия природа без домислювання, \\ немає різниці.}
\poemverse{18.12}{Де дгарма, де бажання, \\ де користь, де вміння розрізняти? \\ У йога, вільного від пар протилежностей, — \\ це зроблено, це ні — \\ немає нічого з цього.}
\poemverse{18.13}{Немає нічого, що слід зробити, \\ і жодної вподоби в серці — \\ у йога, звільненого ще за життя, \\ все йде просто як живеться, тут-таки.}
\poemverse{18.14}{Де омана, де світ, \\ де те споглядання, де сам стан свободи — \\ у великого духом, \\ що спочив у межі всіх задумів?}
\poemverse{18.15}{Нехай той, ким побачений цей світ, \\ каже, що його немає, — \\ а що робить той, у кого не лишилося слідів? \\ Він і бачить, і не бачить.}
\poemverse{18.16}{Нехай той, ким побачено вище Брахмо, \\ повторює: я — Брахман. \\ А що думає той, хто вільний від думки \\ й не бачить нічого другого?}
\poemverse{18.17}{Хто побачив у собі розсіяння, \\ той береться його спиняти. \\ Але великодухий не розсіяний — \\ і, коли немає чого досягати, \\ що йому робити?}
\poemverse{18.18}{Мудрий, збочений разом зі світом, \\ поводиться як увесь світ — \\ і не бачить у себе \\ ні самадгі, ні розсіяння, ні налипання.}
\poemverse{18.19}{Тим, хто позбавлений буття й небуття, \\ хто вдоволений, без слідів і тверезий, — \\ не зроблено нічого, \\ хоч з погляду світу він діє.}
\poemverse{18.20}{Ні в діянні, ні у відступі \\ немає в мудрого зацикленості: \\ що б не випало зробити, зробивши те, \\ він перебуває зручно.}
\poemverse{18.21}{Без слідів, без опори, \\ за власною волею, з відпущеними путами — \\ гнаний вітром самскар, \\ він ворушиться, як сухий лист.}
\poemverse{18.22}{У того, хто поза сансарою, \\ ніде немає ні захвату, ні смутку: \\ прохолодний розумом, \\ він завжди сяє, наче розвтілеснений.}
\poemverse{18.23}{Ніде немає бажання відкинути, \\ і ніде немає загибелі — \\ у мудрого, чия втіха — Атман, \\ чий Атман прохолодний і прозорий.}
\poemverse{18.24}{У того, чий розум за природою порожній \\ і хто діє, як випаде, — \\ як у звичайної людини: \\ немає ні гордості, ні приниженості.}
\poemverse{18.25}{Це діяння зробило тіло, \\ не я, чистої природи — \\ хто тримається цієї думки, \\ той і роблячи не робить.}
\poemverse{18.26}{Навіть якщо він каже не те, \\ він не стає від цього дурнем: \\ звільнений за життя, щасливий, сяйний — \\ він гарний навіть тоді, \\ коли блукає сансарою.}
\poemverse{18.27}{Мудрий, стомлений різноманітними роздумами, \\ дійшов до спочинку: \\ він не утворюється, \\ не знає, не чує, не бачить.}
\poemverse{18.28}{Ні через відсутність самадгі, \\ ні через відсутність розсіяння \\ він не є ні тим, хто прагне визволення, \\ ні кимось іншим: \\ вирішивши так і бачачи все вигаданим, \\ великодухий перебуває як саме Брахмо.}
\poemverse{18.29}{У кого всередині є я-роблення, \\ той не робить, хоч і робить. \\ А в мудрого без я-роблення \\ не лишається нічого незробленого.}
\poemverse{18.30}{Не схвильований і не вдоволений, \\ без діяча, без пориву, \\ без сподівань і без сумніву — \\ розум вільного сяє.}
\poemverse{18.31}{Ум, що не береться ні замислитися, ні діяти, — \\ але ж він, без причини, \\ замислюється й діє.}
\poemverse{18.32}{Почувши правильно істину, \\ тупий стає ще затьмаренішим — \\ або ж звужується; \\ а буває й незатьмарений, \\ що поводиться як затьмарений.}
\poemverse{18.33}{Зосередженість або припинення \\ старанно вправляють дурні. \\ Мудрі не бачать, що їм робити, — \\ вони стоять на власному місці, як сплячі.}
\poemverse{18.34}{Ні без зусилля, ні із зусиллям \\ дурний не досягає заспокоєння: \\ розумний заспокоюється \\ самим лише рішенням про істину.}
\poemverse{18.35}{Атман — чистий, пізнаний, любий, \\ повний, без множинності, без недуги — \\ такого не знають ті, \\ для кого головне — вправа.}
\poemverse{18.36}{Зовсім затьмарений не досягає визволення \\ ні діянням, ні тим, що має вигляд вправи; \\ а блаженний, вільний самим лише розпізнанням, \\ перебуває незмінним.}
\poemverse{18.37}{Дурний не досягає того Брахма, \\ бо хоче ним стати. \\ А мудрий, і не бажаючи, \\ має власну природу вищого Брахма.}
\poemverse{18.38}{Без опори, заклопотані хапанням, \\ дурні годують сансару. \\ А розумні зробили відтин \\ кореня цього кореня лиха.}
\poemverse{18.39}{Дурний не отримує спокою, \\ бо хоче заспокоїтися. \\ А мудрий, остаточно встановивши істину, \\ завжди має спокійний розум.}
\poemverse{18.40}{Де в того бачення Атмана, \\ хто тримається побаченого? \\ Мудрі такого не бачать — \\ вони бачать Атман, незнищенний.}
\poemverse{18.41}{Де припинення в того, \\ хто зовсім затьмарився і наполягає? \\ У мудрого, чия втіха — в собі, \\ воно завжди невдаване.}
\poemverse{18.42}{Один наділяє буттям, \\ інший не наділяє нічим, \\ дехто не наділяє ні тим, ні тим. \\ Саме так він безтурботний.}
\poemverse{18.43}{Дурні уявляють Атман чистим і недвоїстим, \\ але через затьмарення не знають його — \\ і до кінця життя лишаються незаспокоєними.}
\poemverse{18.44}{У того, хто прагне визволення, \\ розум не буває без опори. \\ А у вільного розум без опори \\ й без бажань — завжди.}
\poemverse{18.45}{Побачивши пантер предметів, \\ налякані й шукаючи притулку, \\ вони стрімко кидаються в лоно — \\ задля успіху в припиненні й зосередженості.}
\poemverse{18.46}{Побачивши лева, вільного від слідів, \\ слони предметів мовчки тікають. \\ Нездатні, вони більше не прислуговують, \\ хоч раніше й підлещувалися.}
\poemverse{18.47}{Він не носить на собі справи визволення — \\ без сумніву, зі зʼєднаним розумом. \\ Дивиться, слухає, торкається, нюхає, їсть — \\ і перебуває, як завгодно.}
\poemverse{18.48}{Від самого лише слухання суті — \\ чистий розумом, безтурботний — \\ він не бачить ані різниці між належним і неналежним, \\ ані байдужості.}
\poemverse{18.49}{Коли що випадає зробити, \\ тоді він те й робить — \\ він прямий. \\ Добре чи недобре — \\ його діяння як у дитини.}
\poemverse{18.50}{Від свободи отримує щастя, \\ від свободи досягає вищого, \\ від свободи приходить до заспокоєння, \\ від свободи — до вищої стопи.}
\poemverse{18.51}{Коли людина вважає свій Атман \\ не-діячем і не-споживачем — \\ тоді всі рухи розуму \\ справді вичерпуються.}
\poemverse{18.52}{Навіть розкутий і безформний стан мудрого сяє. \\ А спокій затьмареного \\ з пожадливим розумом — удаваний.}
\poemverse{18.53}{Вони граються великими розкошами \\ й входять у гірські нетрі — \\ мудрі, що відкинули побудови, \\ незвʼязані, з визволеним розумом.}
\poemverse{18.54}{Побачивши знавця писань, божество, \\ святе місце, жінку, царя, любу людину — \\ і вшанувавши їх, — \\ мудрий не має в серці жодного сліду.}
\poemverse{18.55}{Висміяний і зневажений слугами, \\ синами, жінкою, онуками й родичами, \\ йог не змінюється анітрохи.}
\poemverse{18.56}{І вдоволений — не вдоволений; \\ і засмучений — не засмучується. \\ Його дивовижні стани — такі-то й такі — \\ знають лише такі ж, як він.}
\poemverse{18.57}{Сансара — це саме треба зробити. \\ Мудрі її не бачать: \\ вони з порожнім виглядом, \\ безформні, без видозмін, без недуги.}
\poemverse{18.58}{Навіть нічого не роблячи, \\ дурний розумом скрізь збентежений \\ через хвилювання. \\ А вмілий, хоч і робить справи, — безтурботний.}
\poemverse{18.59}{Зручно сидить, зручно лежить, \\ зручно приходить і йде, \\ зручно говорить, зручно їсть — \\ навіть у справах світу той, \\ чий розум спокійний.}
\poemverse{18.60}{У кого за природою немає страждання, \\ хто поводиться, як усі в світі, — \\ той сяє, як велике озеро: \\ незбурюваний, з минулими стражданнями.}
\poemverse{18.61}{Навіть від припинення в затьмареного \\ народжується діяльність. \\ А діяльність мудрого \\ дає плід припинення.}
\poemverse{18.62}{Відречення від набутків \\ здебільшого видно в затьмареного. \\ А в того, у кого в тілі вислизнула надія, — \\ де пристрасть і де безпристрасність?}
\poemverse{18.63}{Погляд затьмареного завжди привʼязаний \\ до уявлення й неуявлення. \\ А в того, хто в собі, \\ він має вигляд не-бачення.}
\poemverse{18.64}{Мудрець, що без бажань у всіх починаннях, \\ поводиться як дитина. \\ У чистого немає налипання, \\ навіть коли діяння робиться.}
\poemverse{18.65}{Саме він блаженний — \\ той, хто знає Атман і рівний до всіх: \\ дивиться, слухає, торкається, нюхає, їсть — \\ а розум без спраги.}
\poemverse{18.66}{Де сансара, де навіть її видимість, \\ де мета, де засіб — \\ у мудрого, вільного від розрізнень, \\ завжди, як у простору?}
\poemverse{18.67}{Перемагає той, хто зрікся користі, \\ чиє тіло — повний власний смак: \\ у нього самадгі невдаване \\ й триває в необмеженому.}
\poemverse{18.68}{До чого тут багато слів: \\ великодухий, що пізнав істину, \\ не прагне ні насолод, ні визволення — \\ завжди й скрізь він без смаку до них.}
\poemverse{18.69}{Двоїстість світу, починаючи з махат, \\ розквітла лише з імені. \\ Коли це відкинуто — \\ яка справа лишається в чистого знання?}
\poemverse{18.70}{Хто впевнений, що все це має природу омани \\ й ніщо не є, — \\ той, із сяянням без ознак, чистий, \\ заспокоюється самою природою.}
\poemverse{18.71}{У того, чия природа — чисте сяяння, \\ хто не бачить, що видиме існує: \\ де припис, де відречення, \\ де покидання, де й заспокоєння?}
\poemverse{18.72}{У того, хто сяє безконечною формою \\ і не бачить природи: \\ де пута, де визволення, \\ де захват, де смуток?}
\poemverse{18.73}{У сансарі, що сягає межі розуму, \\ обертається лише майя. \\ Вільний від моє, без я-роблення, без бажань — \\ сяє той, хто знає.}
\poemverse{18.74}{У мудреця, що бачить Атман — \\ незнищенний, з минулим жаром: \\ де знання, де всесвіт, де тіло, \\ де я й моє?}
\poemverse{18.75}{Якщо тупий розумом покидає діяння, \\ починаючи з припинення, — \\ він береться за забаганки й балачки, \\ і саме з тієї ж миті.}
\poemverse{18.76}{Тупий, навіть почувши ту суть, \\ не покидає цілковитої затьмареності: \\ зовні він без розрізнень — \\ із зовнішнього зусилля, — \\ а всередині пожадливий до предметів.}
\poemverse{18.77}{Хто через знання має відпалі діяння — \\ хоч і діє з огляду на світ — \\ не знаходить нагоди зробити \\ чи сказати хоч що-небудь.}
\poemverse{18.78}{Де темрява, де світло, \\ де відкидання, де ніщо — \\ у незмінного, мудрого, безтривожного, завжди?}
\poemverse{18.79}{Де стійкість, де розрізнення, \\ де й безтривожність — \\ у йога, чия природа невимовна, \\ хто без природи?}
\poemverse{18.80}{Немає неба, немає пекла, \\ немає й визволення за життя. \\ До чого багато слів: \\ з погляду йоги — ніщо.}
\poemverse{18.81}{Він не прагне набутку \\ й не сумує через ненабуток: \\ розум мудрого прохолодний, \\ сповнений амрітою.}
\poemverse{18.82}{Без бажань він не хвалить спокійного \\ й не ганить лихого. \\ Рівний у стражданні й щасті, насичений, \\ він не бачить жодної справи.}
\poemverse{18.83}{Мудрий не ненавидить сансару \\ й не прагне побачити Атман. \\ Вільний від захвату й гніву, \\ він ні мертвий, ні живий.}
\poemverse{18.84}{Без привʼязаності до синів, жінки й інших, \\ без бажань до предметів, \\ без турботи навіть про власне тіло — \\ без надій сяє той, хто знає.}
\poemverse{18.85}{Вдоволення скрізь у мудрого, \\ що живе як випаде, \\ вільно ходить землями \\ й лягає там, де сонце зайшло.}
\poemverse{18.86}{Хай тіло падає або встає — \\ у великого немає турботи: \\ він спочив на ґрунті власної природи \\ й забув усю сансару без останку.}
\poemverse{18.87}{Без жодної речі, \\ той, хто йде куди хоче, \\ без двоїстості, з відтятим сумнівом, \\ непривʼязаний до всіх станів — \\ сам-один утішається той, хто знає.}
\poemverse{18.88}{Без моє сяє мудрий, \\ для кого грудка, камінь і золото те саме; \\ у кого добре розвʼязаний вузол серця, \\ хто струсив геть пил і темряву.}
\poemverse{18.89}{У того, хто ніде не має опори, \\ немає жодного сліду в серці. \\ З ким зрівняти визволеного Атмана, \\ що наситився?}
\poemverse{18.90}{І знаючи — не знає, \\ і бачачи — не бачить, \\ і говорячи — не говорить. \\ Хто інший, крім вільного від слідів?}
\poemverse{18.91}{Жебрак чи володар — \\ той, хто без бажань, сяє: \\ у кого в усіх станах \\ відпала думка про добре й недобре.}
\poemverse{18.92}{Де свобода, де стиснення, \\ де рішення про істину — \\ у йога, чия істота нелукаво-пряма, \\ хто досяг мети?}
\poemverse{18.93}{Те, що переживається всередині \\ в насиченого спочинком в Атмані, \\ без надій, з минулим стражданням — \\ як і кому це розповісти?}
\poemverse{18.94}{І сплячи — не в глибокому сні; \\ і в сні — не лежить; \\ і в неспанні — не неспить. \\ Мудрий, насичений на кожнім кроці.}
\poemverse{18.95}{Знаючий, хоч і з думкою — без думки; \\ хоч і з чуттями — без чуттів; \\ хоч і з розумом — без розуму; \\ хоч і з я-робленням — без нього.}
\poemverse{18.96}{Не щасливий і не нещасний, \\ не відречений і не привʼязаний, \\ не прагне визволення й не визволений. \\ Він не щось і не ніщо.}
\poemverse{18.97}{І в розсіянні — не розсіяний; \\ і в самадгі — не той, хто має самадгі; \\ і в тупості — не тупий. \\ Блаженний, і у вченості — не вчений.}
\poemverse{18.98}{Вільний, у собі в тому, як є; \\ заспокоєний у зробленому-що-треба; \\ рівний скрізь. \\ Через безспраглість він не памʼятає \\ ні зробленого, ні незробленого.}
\poemverse{18.99}{Його не тішить, коли хвалять, \\ і не гнівить, коли ганять. \\ Він не лякається смерті \\ й не радіє життю.}
\poemverse{18.100}{Той, чия думка втихла, \\ не біжить ні до переповненого людьми, \\ ні до лісу; \\ так чи інакше, там чи там \\ він перебуває рівним.}

% ЗГЕНЕРОВАНО scripts/build_format_v2_chapter.py — не редагувати вручну.
% Джерело: books/ashtavakra-gita/poetychnyi/rozdil-19.md

\poemverse{19.1}{Взявши кліщі знання істини, \\ я витяг із нутра серця \\ колючки різноманітних хибних уявлень.}
\poemverse{19.2}{Де дгарма, де бажання, \\ де користь, де вміння розрізняти? \\ Де двоїстість, де недвоїстість — \\ у мене, що стоїть у власній величі?}
\poemverse{19.3}{Де минуле, де майбутнє, \\ де навіть теперішнє? \\ Де сторона світу, де вічність — \\ у мене, що стоїть у власній величі?}
\poemverse{19.4}{Де Атман, де не-Атман, \\ де добре, де недобре? \\ Де думка, де відсутність думки — \\ у мене, що стоїть у власній величі?}
\poemverse{19.5}{Де сон зі сновидіннями, \\ де глибокий сон, де неспання? \\ Де четвертий стан, де навіть страх — \\ у мене, що стоїть у власній величі?}
\poemverse{19.6}{Де далеке, де близьке, \\ де зовнішнє, де внутрішнє? \\ Де грубе, де тонке — \\ у мене, що стоїть у власній величі?}
\poemverse{19.7}{Де смерть, де життя, \\ де світи, де мирське? \\ Де розчинення, де самадгі — \\ у мене, що стоїть у власній величі?}
\poemverse{19.8}{Досить розмов про трицінність, \\ досить і розмов про йогу, \\ досить розмов про знання — \\ у мене, що спочив в Атмані.}

% ЗГЕНЕРОВАНО scripts/build_format_v2_chapter.py — не редагувати вручну.
% Джерело: books/ashtavakra-gita/poetychnyi/rozdil-20.md

\poemverse{20.1}{Де першоелементи, де тіло, \\ де чуття, де розум? \\ Де порожнеча, де без-надійність — \\ у моїй власній формі, \\ що не знає плями?}
\poemverse{20.2}{Де писання, де пізнання Я, \\ де розум без обʼєктів? \\ Де вдоволеність, де свобода від спраги — \\ у мене, в кого пари вже відпали назавжди?}
\poemverse{20.3}{Де знання й де незнання, \\ де я, де це, де моє? \\ Де пута, де визволення — \\ і де взагалі якась оформленість \\ у того, хто сам є своєю формою?}
\poemverse{20.4}{Де вже запущені в дію діяння, \\ де звільнення ще за життя? \\ Де те безтілесне кайвальям — \\ у того, хто завжди без ознак?}
\poemverse{20.5}{Де діяч, де той, хто зазнає? \\ Де бездіяльне сяйво? \\ Де безпосереднє знання, де його плід — \\ у мене, в кого немає власної природи?}
\poemverse{20.6}{Де світ, де той, хто прагне звільнення, \\ де йог чи знавець? \\ Де звʼязаний, де звільнений — \\ я у своїй власній формі, \\ яка не знає двоїстості.}
\poemverse{20.7}{Де творення, де згортання, \\ де мета, де засіб? \\ Де той, хто практикує, де здобуття — \\ я у своїй власній формі, \\ яка не знає двоїстості.}
\poemverse{20.8}{Де той, хто пізнає, де засіб пізнання, \\ де предмет, де саме пізнання? \\ Де щось, де ніщо — \\ у мене, хто завжди без плями?}
\poemverse{20.9}{Де розсіяність, де зосередженість, \\ де несвідомість, де затьмареність? \\ Де радість, де відчай — \\ у мене, хто завжди бездіяльний?}
\poemverse{20.10}{Де ця буденна практика, \\ де та вища істина? \\ Де щастя, де страждання — \\ у мене, хто вже завжди нічого не розглядає?}
\poemverse{20.11}{Де мая, де сансара, \\ де прихильність і де відвернення? \\ Де жива істота, де те Брахмо — \\ у мене, хто завжди без плями?}
\poemverse{20.12}{Де діяння і де припинення, \\ де звільнення, де пута — \\ у мене, незмінного й неділимого, \\ що завжди в собі?}
\poemverse{20.13}{Де настановлення, де писання, \\ де учень, де вчитель? \\ Де й мета людського життя — \\ у мене, хто без жодної опори, \\ у мене, блаженного?}
\poemverse{20.14}{Де є і де немає, \\ де одне і де двоє? \\ Та й навіщо багато слів: \\ у мене не постає ніщо.}


\chapter*{Джерела й межі цього видання}
\addcontentsline{toc}{chapter}{Джерела й межі}
\markboth{Джерела й межі}{Джерела й межі}

\noindent\textbf{\sffamily Що є джерелом нашого тексту.}
Власні редакторські картки \texttt{editorial/cards/*.json}: поля
\texttt{ukrainian\_translation} і \texttt{commentary}. Деванагарі й українська
транслітерація — наше механічне перетворення поля \texttt{iast} тих самих
карток. Коментарі, паралелі, схеми й поетичний переклад — наш авторський
внесок під ліцензією «ВОЛЬНІСТЬ».

\vspace{0.6em}
\noindent\textbf{\sffamily Що ми не відтворюємо.}
Транскрипції GRETIL і SARIT, їхню розмітку та редакторський апарат. Локатори
(\texttt{Avg\_1.1}, \texttt{verse\_1.1}) подано в картках як покажчик для звірки.

\vspace{0.6em}
\noindent\textbf{\sffamily Чого ми не стверджуємо.}
Що це критичне видання. Що переклад вивірений. Що коментар рецензований. Що
права на всі використані матеріали остаточно зʼясовані. Що наш поетичний
переклад є канонічним.

\end{document}
